% $Header: /cvsroot/latex-beamer/latex-beamer/solutions/generic-talks/generic-ornate-15min-45min.en.tex,v 1.5 2007/01/28 20:48:23 tantau Exp $

\documentclass{beamer}
%\documentclass[handout]{beamer}
\usepackage{pgfpages}
\pgfpagesuselayout{2 on 1}[a4paper,border shrink=5mm]

% This file is a solution template for:

% - Giving a talk on some subject.
% - The talk is between 15min and 45min long.
% - Style is ornate.



% Copyright 2004 by Till Tantau <tantau@users.sourceforge.net>.
%
% In principle, this file can be redistributed and/or modified under
% the terms of the GNU Public License, version 2.
%
% However, this file is supposed to be a template to be modified
% for your own needs. For this reason, if you use this file as a
% template and not specifically distribute it as part of a another
% package/program, I grant the extra permission to freely copy and
% modify this file as you see fit and even to delete this copyright
% notice. 


\mode<presentation>
{
  \usetheme{Warsaw}
  % or ...

  \setbeamercovered{transparent}
  % or whatever (possibly just delete it)
}
\usepackage{url}
\usepackage[normalem]{ulem}
\usepackage{listings}
\usepackage{pstricks,pst-node,pst-tree,graphics}
\newpsobject{showgrid}{psgrid}{subgriddiv=1,griddots=10,gridlabels=6pt}
\usepackage{amssymb,marvosym}
\usepackage{colortbl}
\usepackage{array}
\newcolumntype{I}{!{\vrule width 1pt}}
\newlength\savedwidth
\newcommand\whline{\noalign{\global\savedwidth\arrayrulewidth
                            \global\arrayrulewidth 1pt}%
                   \hline
                   \noalign{\global\arrayrulewidth\savedwidth}}
\usepackage{textpos}
\usepackage{textcomp,pifont}
\newcommand{\CC}[1]{\Tcircle[linestyle=solid]{#1}}
\newcommand{\REM}[2]{\only<#1>{\textcolor{blue}{\ding{233} \textit{#2}}}}
\newcommand{\CURR}{\Tcircle[fillstyle=solid,fillcolor=blue,linestyle=solid]{.}}
\newcommand{\FAILTREE}[2]{\Ttri[fillstyle=solid,fillcolor=red,linestyle=solid]{\hyperlink{#1}{#2}}}
\newcommand{\FAIL}[1]{\Tcircle[fillstyle=solid,fillcolor=red,linestyle=solid]{\hyperlink{#1}{\ding{54}}}}
\newcommand{\SUCC}[1]{\Tcircle[fillstyle=solid,fillcolor=green,linestyle=solid]{\hyperlink{#1}{\checkmark}}}
\newcommand{\G}{\cellcolor[gray]{0.7}}
\newcommand{\R}{\cellcolor[gray]{0.5}\ding{54}}
\newcommand{\X}{\cellcolor[gray]{0.5}-}
\newcommand{\D}{\cellcolor[gray]{0.5}|}
\newcommand{\A}{\cellcolor[rgb]{0.5,0,0}}
\newcommand{\Y}{\cellcolor[rgb]{1.0,1.0,0}|}
\newcommand{\F}{\cellcolor[rgb]{1.0,0,0}\ding{89}}


\usepackage[english]{babel}
% or whatever

\usepackage[latin1]{inputenc}
% or whatever

\usepackage{times}
\usepackage[T1]{fontenc}
% Or whatever. Note that the encoding and the font should match. If T1
% does not look nice, try deleting the line with the fontenc.


\title[Global Constraints] % (optional, use only with long paper titles)
{Chapter 5: The Finite Domain Solver - Global Constraints}

\subtitle
{} % (optional)

\author % (optional, use only with lots of authors)
{Helmut Simonis}
% - Use the \inst{?} command only if the authors have different
%   affiliation.

\institute[4C] % (optional, but mostly needed)
{Cork Constraint Computation Centre\\Computer Science Department\\University College Cork\\Ireland}
% - Use the \inst command only if there are several affiliations.
% - Keep it simple, no one is interested in your street address.
\logo{\includegraphics[width=15mm]{4c}}
%\logo{\includegraphics[width=40mm]{ucclogo}}
\date[\today] % (optional)
{ECLiPSe ELearning}

\subject{Talks}
% This is only inserted into the PDF information catalog. Can be left
% out. 



% If you have a file called "university-logo-filename.xxx", where xxx
% is a graphic format that can be processed by latex or pdflatex,
% resp., then you can add a logo as follows:

% \pgfdeclareimage[height=0.5cm]{university-logo}{university-logo-filename}
% \logo{\pgfuseimage{university-logo}}



% Delete this, if you do not want the table of contents to pop up at
% the beginning of each subsection:
\AtBeginSection[]
{
  \begin{frame}<beamer>
    \frametitle{Outline}
    \tableofcontents[sectionstyle=show/shaded,subsectionstyle=show/show/hide]
  \end{frame}
}


% If you wish to uncover everything in a step-wise fashion, uncomment
% the following command: 

%\beamerdefaultoverlayspecification{<+->}


\begin{document}
%\textblockcolour{red}

\begin{frame}
  \titlepage
\end{frame}

\begin{frame}
\frametitle{Outline}
  \tableofcontents[hideallsubsections]
  % You might wish to add the option [pausesections]
\end{frame}


% Since this a solution template for a generic talk, very little can
% be said about how it should be structured. However, the talk length
% of between 15min and 45min and the theme suggest that you stick to
% the following rules:  

% - Exactly two or three sections (other than the summary).
% - At *most* three subsections per section.
% - Talk about 30s to 2min per frame. So there should be between about
%   15 and 30 frames, all told.

\begin{frame}
\frametitle{What we want to introduce}
\begin{itemize}
\item Global Constraints
\begin{itemize}
\item Powerful modelling abstractions
\item Non-trivial propagation
\end{itemize}
\item Consistency Levels
\begin{itemize}
\item Tradeoff between speed and propagation
\item Characterisation of reasoning power
\end{itemize}
\item Example: Alldifferent
\begin{itemize}
\item 3 variants shown
\end{itemize}
\end{itemize}
\end{frame}

\begin{frame}
\frametitle{Methodology}
\begin{itemize}
\item Evaluation on Sudoku puzzle
\item Comparing
\begin{itemize}
\item Initial setup
\item Search
\item Performance
\end{itemize}
\item Explaining reasoning inside constraint
\item Link to general classification of global constraints
\end{itemize}
\end{frame}

\section{Problem}
\begin{frame}
\frametitle{Problem Definition}
\begin{block}{Sudoku}
Fill in numbers from 1 to 9 so that each row, column and block contain each number exactly once
\end{block}
\begin{tabular}{cc}
\scalebox{0.7}{\begin{tabular}{rrrrr}
  & S & E & N & D \\
+ & M & O & R & E \\ \hline
M & O & N & E & Y
\end{tabular}
}
&
\only<2>{\scalebox{0.7}{\input{dump/dump_gac_31}}}
\end{tabular}
\end{frame}

\begin{frame}
\frametitle{Model}
\begin{itemize}
\item A variable for each cell, ranging from 1 to 9
\item A 9x9 matrix of variables describing the problem
\item Preassigned integers for the given hints
\item alldifferent constraints for each row, column and 3x3 block
\end{itemize}
\end{frame}

\begin{frame}
\frametitle{Reminder: alldifferent}
\begin{itemize}
\item Argument: list of variables
\item Meaning: variables are pairwise different
\item Reasoning: Forward Checking (FC)
\begin{itemize}
\item When variable is assigned to value, remove the value from all other variables
\item If a variable has only one possible value, then it is assigned
\item If a variable has no possible values, then the constraint fails
\item Constraint is checked whenever one of its variables is assigned
\item Equivalent to decomposition into binary disequality constraints
\end{itemize}
\end{itemize}
\end{frame}

\section{Program}

\begin{frame}[fragile]
\frametitle{Declarations}
\begin{semiverbatim}
:-{\bf{}module}(sudoku).
:-{\bf{}export}(top/0).
:-{\bf{}lib}(ic).

top:-
    problem(Matrix),
    model(Matrix),
    {\bf{}writeln}(Matrix).
\end{semiverbatim}
\end{frame}

\begin{frame}[fragile]
\frametitle{Data}        
\begin{semiverbatim}
problem([]([](4, _, 8, _, _, _, _, _, _), 
           [](_, _, _, 1, 7, _, _, _, _), 
           [](_, _, _, _, 8, _, _, 3, 2), 
           [](_, _, 6, _, _, 8, 2, 5, _), 
           [](_, 9, _, _, _, _, _, 8, _), 
           [](_, 3, 7, 6, _, _, 9, _, _), 
           [](2, 7, _, _, 5, _, _, _, _), 
           [](_, _, _, _, 1, 4, _, _, _), 
           [](_, _, _, _, _, _, 6, _, 4))).
\end{semiverbatim}
\end{frame}

\begin{frame}[fragile]
\frametitle{Main Program}
\begin{semiverbatim}
model(Matrix):-
    Matrix[1..9,1..9] :: 1..9,
    (for(I,1,9),
     param(Matrix) do
        {\bf{}alldifferent}(Matrix[I,1..9]),
        {\bf{}alldifferent}(Matrix[1..9,I])
    ),
    (multifor([I,J],[1,1],[7,7],[3,3]),
     param(Matrix) do
        {\bf{}alldifferent}(flatten(Matrix[I..I+2,J..J+2]))
    ),
    {\bf{}flatten_array}(Matrix,List),
    {\bf{}labeling}(List).
\end{semiverbatim}
\end{frame}

\begin{frame}
\frametitle{Domain Visualizer}
\begin{textblock}{5.5}(10,-2)
\scalebox{0.68}{\input{dump/dump_fc_1}}
\end{textblock}
\begin{itemize}
\item Problem shown as matrix
\item Each cell corresponds to a variable
\item Instantiated: Shows integer value (large)
\item Uninstantiated: Shows values in domain
\end{itemize}
\end{frame}

\begin{frame}
\frametitle{Constraint Visualizer}
\begin{textblock}{5.5}(10,-2)
\scalebox{0.68}{\input{dump/dump_fc_25}}
\end{textblock}
\begin{itemize}
\item Problem shown as matrix
\item Currently active constraint highlighted
\item Values removed at this step shown in red
\item Values assigned at this step shown in blue
\end{itemize}
\end{frame}

\section{Initial Propagation (Forward Checking)}

\begin{frame}
\frametitle{Initial State (Forward Checking)}
\input{dump/dump_fc_1}
\end{frame}

\begin{frame}
\frametitle{Propagation Steps (Forward Checking)}
\only<1>{\input{dump/dump_fc_1}}
\only<2>{\input{dump/dump_fc_2}}
\only<3>{\input{dump/dump_fc_3}}
\only<4>{\input{dump/dump_fc_4}}
\only<5>{\input{dump/dump_fc_5}}
\only<6>{\input{dump/dump_fc_6}}
\only<7>{\input{dump/dump_fc_7}}
\only<8>{\input{dump/dump_fc_8}}
\only<9>{\input{dump/dump_fc_9}}
\only<10>{\input{dump/dump_fc_10}}
\only<11>{\input{dump/dump_fc_11}}
\only<12>{\input{dump/dump_fc_12}}
\only<13>{\input{dump/dump_fc_13}}
\only<14>{\input{dump/dump_fc_14}}
\only<15>{\input{dump/dump_fc_15}}
\only<16>{\input{dump/dump_fc_16}}
\only<17>{\input{dump/dump_fc_17}}
\only<18>{\input{dump/dump_fc_18}}
\only<19>{\input{dump/dump_fc_19}}
\only<20>{\input{dump/dump_fc_20}}
\only<21>{\input{dump/dump_fc_21}}
\only<22>{\input{dump/dump_fc_22}}
\only<23>{\input{dump/dump_fc_23}}
\only<24>{\input{dump/dump_fc_24}}
\only<25>{\input{dump/dump_fc_25}}
\only<26>{\input{dump/dump_fc_26}}
\only<27>{\input{dump/dump_fc_27}}
\only<28>{\input{dump/dump_fc_28}}
\vfill
\hyperlinkframestart<2-28>{\beamerreturnbutton{Back to Start}}
\hyperlinkframeend<1-27>{\beamerskipbutton{Skip Animation}}
\end{frame}

\begin{frame}
\frametitle{After Setup (Forward Checking)}
\input{dump/dump_fc_29}
\end{frame}

\section{Improved Reasoning}

\begin{frame}
\frametitle{Can we do better?}
\begin{itemize}
\item The alldifferent constraint is missing propagation
\begin{itemize}
\item How can we do more propagation?
\item Do we know when we derive all possible information from the constraint?
\end{itemize}
\item Constraints only interact by changing domains of variables
\end{itemize}
\end{frame}

\begin{frame}[fragile]
\frametitle{A Simpler Example}
\begin{semiverbatim}
:-{\bf{}lib}(ic).

top:-
    X :: 1..2,
    Y :: 1..2,
    Z :: 1..3,
    {\bf{}alldifferent}([X,Y,Z]),
    {\bf{}writeln}([X,Y,Z]).
\end{semiverbatim}
\end{frame}

\begin{frame}
\frametitle{Using Forward Checking}
\begin{itemize}
\item No variable is assigned
\item No reduction of domains
\item But, values 1 and 2 can be removed from Z
\item This means that Z is assigned to 3
\end{itemize}
\end{frame}

\begin{frame}
\frametitle{Visualization of alldifferent as Graph}
\begin{textblock}{5.5}(12,-2)
\scalebox{0.6}{\begin{psmatrix}
X & & 1 \\
Y & & 2 \\
Z & & 3
\ncline{1,1}{1,3}
\ncline{1,1}{2,3}
\ncline{2,1}{1,3}
\ncline{2,1}{2,3}
\ncline{3,1}{1,3}
\ncline{3,1}{2,3}
\ncline{3,1}{3,3}
\end{psmatrix}
}
\end{textblock}
\begin{itemize}
\item Show problem as graph with two types of nodes
\begin{itemize}
\item Variables on the left
\item Values on the right
\end{itemize}
\item If value is in domain of variable, show link between them
\item This is called a {\em bipartite} graph
\end{itemize}
\end{frame}

\begin{frame}[fragile]
\frametitle{A Simpler Example}
\begin{textblock}{8}(6,0)
\only<1>{Value Graph for 
\begin{semiverbatim}
X :: 1..2, 
Y :: 1..2, 
Z :: 1..3
\end{semiverbatim}}
\only<2>{Check interval [1,2]}
\only<3>{
\begin{itemize}
\item Find variables completely contained in interval
\item There are two: X and Y
\item This uses up the capacity of the interval
\end{itemize}
}
\only<4>{No other variable can use that interval}
\only<5>{Only one value left in domain of Z, this can be assigned}
\end{textblock}
\only<1>{
\begin{psmatrix}
X & & 1 \\
Y & & 2 \\
Z & & 3
\ncline{1,1}{1,3}
\ncline{1,1}{2,3}
\ncline{2,1}{1,3}
\ncline{2,1}{2,3}
\ncline{3,1}{1,3}
\ncline{3,1}{2,3}
\ncline{3,1}{3,3}
\end{psmatrix}
}
\only<2>{
\begin{psmatrix}
X & & \textcolor{insight-blue}{1} \\
Y & & \textcolor{insight-blue}{2} \\
Z & & 3
\ncline{1,1}{1,3}
\ncline{1,1}{2,3}
\ncline{2,1}{1,3}
\ncline{2,1}{2,3}
\ncline{3,1}{1,3}
\ncline{3,1}{2,3}
\ncline{3,1}{3,3}
\end{psmatrix}
}
\only<3>{
\begin{psmatrix}
\textcolor{insight-blue}{X} & & \textcolor{insight-blue}{1} \\
\textcolor{insight-blue}{Y} & & \textcolor{insight-blue}{2} \\
Z & & 3
\ncline[linecolor=insight-blue]{1,1}{1,3}
\ncline[linecolor=insight-blue]{1,1}{2,3}
\ncline[linecolor=insight-blue]{2,1}{1,3}
\ncline[linecolor=insight-blue]{2,1}{2,3}
\ncline{3,1}{1,3}
\ncline{3,1}{2,3}
\ncline{3,1}{3,3}
\end{psmatrix}
}
\only<4>{
\begin{psmatrix}
\textcolor{insight-blue}{X} & & \textcolor{insight-blue}{1} \\
\textcolor{insight-blue}{Y} & & \textcolor{insight-blue}{2} \\
Z & & 3
\ncline[linecolor=insight-blue]{1,1}{1,3}
\ncline[linecolor=insight-blue]{1,1}{2,3}
\ncline[linecolor=insight-blue]{2,1}{1,3}
\ncline[linecolor=insight-blue]{2,1}{2,3}
\ncline[linestyle=dotted,linecolor=insight-pink]{3,1}{1,3}
\ncline[linestyle=dotted,linecolor=insight-pink]{3,1}{2,3}
\ncline{3,1}{3,3}
\end{psmatrix}
}
\only<5>{
\begin{psmatrix}
\textcolor{insight-blue}{X} & & \textcolor{insight-blue}{1} \\
\textcolor{insight-blue}{Y} & & \textcolor{insight-blue}{2} \\
Z & & 3
\ncline[linecolor=insight-blue]{1,1}{1,3}
\ncline[linecolor=insight-blue]{1,1}{2,3}
\ncline[linecolor=insight-blue]{2,1}{1,3}
\ncline[linecolor=insight-blue]{2,1}{2,3}
\ncline[linecolor=insight-pink]{3,1}{3,3}
\end{psmatrix}
}

\end{frame}

\begin{frame}
\frametitle{Idea (Hall Intervals)}
\begin{itemize}
\item Take each interval of possible values, say size $N$
\item Find all $K$ variables whose domain is completely contained in interval
\item If $K>N$ then the constraint is infeasible
\item If $K=N$ then no other variable can use that interval
\item Remove values from such variables if their bounds change
\item If $K<N$ do nothing
\item Re-check whenever domain bounds change
\end{itemize}
\end{frame}

\begin{frame}
\frametitle{Implementation}
\begin{itemize}
\item Problem: Too many intervals ($O(n^2)$) to consider
\item Solution:
\begin{itemize}
\item Check only those intervals which update bounds
\item Enumerate intervals incrementally
\item Starting from lowest(highest) value
\item Using sorted list of variables
\end{itemize}
\item Complexity: $O(n\log(n))$ in standard implementations
\item Important: Only looks at min/max bounds of variables
\end{itemize}
\end{frame}

\begin{frame}
\frametitle{Bounds Consistency}
\begin{block}{Definition}
A constraint achieves {\em bounds consistency}, if for the lower and upper bound of every variable, it is possible to find values for all other variables between their lower and upper bounds which satisfy the constraint.
\end{block}
\end{frame}

\begin{frame}
\frametitle{Can we do better?}
\begin{itemize}
\item Bounds consistency only considers min/max bounds
\item Ignores ``holes'' in domain
\item Sometimes we can improve propagation looking at those holes
\end{itemize}
\end{frame}

\begin{frame}[fragile]
\frametitle{Another Simple Example}
\begin{semiverbatim}
:-{\bf{}lib}(ic).

top:-
    X :: [1,3],
    Y :: [1,3],
    Z :: 1..3,
    {\bf{}alldifferent}([X,Y,Z]),
    {\bf{}writeln}([X,Y,Z]).
\end{semiverbatim}
\end{frame}


\begin{frame}[fragile]
\frametitle{Another Simple Example}
\begin{textblock}{8}(6,0)
\only<1>{Value Graph for 
\begin{semiverbatim}
X :: [1,3], 
Y :: [1,3], 
Z :: 1..3
\end{semiverbatim}}
\only<2>{\begin{itemize}
\item Check interval [1,2]
\item No domain of a variable completely contained in interval
\item No propagation
\end{itemize}}
\only<3>{\begin{itemize}
\item Check interval [2,3]
\item No domain of a variable completely contained in interval
\item No propagation
\end{itemize}}
\only<4>{But, more propagation is possible, there are only two solutions}
\only<5>{Solution 1: assignment in blue}
\only<6>{Solution 2: assignment in green}
\end{textblock}
\only<1>{
\begin{psmatrix}
X & & 1 \\
Y & & 2 \\
Z & & 3
\ncline{1,1}{1,3}
\ncline{1,1}{3,3}
\ncline{2,1}{1,3}
\ncline{2,1}{3,3}
\ncline{3,1}{1,3}
\ncline{3,1}{2,3}
\ncline{3,1}{3,3}
\end{psmatrix}
}
\only<2>{
\begin{psmatrix}
X & & \textcolor{insight-blue}{1} \\
Y & & \textcolor{insight-blue}{2} \\
Z & & 3
\ncline{1,1}{1,3}
\ncline{1,1}{3,3}
\ncline{2,1}{1,3}
\ncline{2,1}{3,3}
\ncline{3,1}{1,3}
\ncline{3,1}{2,3}
\ncline{3,1}{3,3}
\end{psmatrix}
}
\only<3>{
\begin{psmatrix}
X & & 1 \\
Y & & \textcolor{insight-blue}{2} \\
Z & & \textcolor{insight-blue}{3}
\ncline{1,1}{1,3}
\ncline{1,1}{3,3}
\ncline{2,1}{1,3}
\ncline{2,1}{3,3}
\ncline{3,1}{1,3}
\ncline{3,1}{2,3}
\ncline{3,1}{3,3}
\end{psmatrix}
}
\only<4>{
\begin{psmatrix}
X & & 1 \\
Y & & 2 \\
Z & & 3
\ncline{1,1}{1,3}
\ncline{1,1}{3,3}
\ncline{2,1}{1,3}
\ncline{2,1}{3,3}
\ncline{3,1}{1,3}
\ncline{3,1}{2,3}
\ncline{3,1}{3,3}
\end{psmatrix}
}
\only<5>{
\begin{psmatrix}
X & & 1 \\
Y & & 2 \\
Z & & 3
\ncline[linecolor=insight-blue]{1,1}{1,3}
\ncline[linestyle=dotted]{1,1}{3,3}
\ncline[linestyle=dotted]{2,1}{1,3}
\ncline[linecolor=insight-blue]{2,1}{3,3}
\ncline[linestyle=dotted]{3,1}{1,3}
\ncline[linecolor=insight-blue]{3,1}{2,3}
\ncline[linestyle=dotted]{3,1}{3,3}
\end{psmatrix}
}
\only<6>{
\begin{psmatrix}
X & & 1 \\
Y & & 2 \\
Z & & 3
\ncline[linestyle=dotted]{1,1}{1,3}
\ncline[linecolor=insight-green]{1,1}{3,3}
\ncline[linecolor=insight-green]{2,1}{1,3}
\ncline[linestyle=dotted]{2,1}{3,3}
\ncline[linestyle=dotted]{3,1}{1,3}
\ncline[linecolor=insight-green]{3,1}{2,3}
\ncline[linestyle=dotted]{3,1}{3,3}
\end{psmatrix}
}
\only<7>{
\begin{psmatrix}
X & & 1 \\
Y & & 2 \\
Z & & 3
\ncline[linecolor=insight-blue]{1,1}{1,3}
\ncline[linestyle=dotted]{1,1}{3,3}
\ncline[linestyle=dotted]{2,1}{1,3}
\ncline[linecolor=insight-blue]{2,1}{3,3}
\ncline[linestyle=dotted]{3,1}{1,3}
\ncline[linecolor=insight-blue]{3,1}{2,3}
\ncline[linestyle=dotted]{3,1}{3,3}
\end{psmatrix}
\begin{psmatrix}
X & & 1 \\
Y & & 2 \\
Z & & 3
\ncline[linestyle=dotted]{1,1}{1,3}
\ncline[linecolor=insight-green]{1,1}{3,3}
\ncline[linecolor=insight-green]{2,1}{1,3}
\ncline[linestyle=dotted]{2,1}{3,3}
\ncline[linestyle=dotted]{3,1}{1,3}
\ncline[linecolor=insight-green]{3,1}{2,3}
\ncline[linestyle=dotted]{3,1}{3,3}
\end{psmatrix}
\begin{psmatrix}
X & & 1 \\
Y & & 2 \\
Z & & 3
\ncline[linecolor=insight-blue]{1,1}{1,3}
\ncline[linecolor=insight-green]{1,1}{3,3}
\ncline[linecolor=insight-green]{2,1}{1,3}
\ncline[linecolor=insight-blue]{2,1}{3,3}
\ncline[linecolor=insight-pink,linestyle=dotted]{3,1}{1,3}
\ncline[linecolor=insight-green]{3,1}{2,3}
\ncline[linecolor=insight-pink,linestyle=dotted]{3,1}{3,3}
\end{psmatrix}
}

\begin{textblock}{8}(2,0)
\only<7>{Combining solutions shows that Z=1 and Z=3 are not possible}
\only<8>{Can we deduce this without enumerating solutions?}
\end{textblock}
\end{frame}

\begin{frame}
\frametitle{Solutions and maximal matchings}
\begin{itemize}
\item A {\em Matching} is subset of edges which do not coincide in any node
\item No matching can have more edges than number of variables
\item Every solution corresponds to a {\em maximal matching} and vice versa
\item If a link does not belong to some maximal matching, then it can be removed 
\end{itemize}
\end{frame}

\begin{frame}
\frametitle{Implementation}
\begin{itemize}
\item Possible to compute all links which belong to some matching
\begin{itemize}
\item Without enumerating all of them!
\end{itemize}
\item Enough to compute {\bf one} maximal matching
\item Requires algorithm for {\em strongly connected components}
\item Extra work required if more values than variables
\item All links (values in domains) which are not supported can be removed
\item Complexity: $O(n^{1.5}d)$
\end{itemize}
\end{frame}

\begin{frame}
\frametitle{Domain Consistency}
\begin{block}{Definition}
A constraint achieves {\em domain consistency}, if for every variable and for every value in its domain, it is possible to find values in the domains of all other variables which satisfy the constraint.
\end{block}
\begin{itemize}
\item Also called {\em generalized arc consistency (GAC)}
\item or {\em hyper arc consistency}
\end{itemize}
\end{frame}

\begin{frame}
\frametitle{Can we still do better?}
\begin{itemize}
\item NO! This extracts all information from this one constraint
\item We could perhaps improve speed, but not propagation
\item But possible to use different model
\item Or model interaction of multiple constraints
\end{itemize}
\end{frame}

\begin{frame}
\frametitle{Should all constraints achieve domain consistency?}
\begin{itemize}
\item Domain consistency is usually more expensive than bounds consistency
\begin{itemize}
\item Overkill for simple problems
\item Nice to have choices
\end{itemize}
\item For some constraints achieving domain consistency is NP-hard
\begin{itemize}
\item We have to live with more restricted propagation
\end{itemize}
\end{itemize}
\end{frame}

\begin{frame}
\frametitle{Improved Propagation in ECLiPSe}
\begin{itemize}
\item \texttt{ic\_global} library bounds consistent version
\item \texttt{ic\_global\_gac} library domain consistent version 
\item Choose which version to use by using module annotation
\item Choice can be passed as parameter
\end{itemize}
\end{frame}

\begin{frame}[fragile]
\frametitle{Declarations}
\begin{semiverbatim}
:-{\bf{}module}(sudoku).
:-{\bf{}export}(top/0).
:-{\bf{}lib}(ic).
:-{\bf{}lib}(ic_global).
:-{\bf{}lib}(ic_global_gac).

top:-
    problem(Matrix),
    model(ic_global,Matrix),
    {\bf{}writeln}(Matrix).
\end{semiverbatim}
\end{frame}


\begin{frame}[fragile]
\frametitle{Main Program}
\begin{semiverbatim}
model(Method,Matrix):-
    Matrix[1..9,1..9] :: 1..9,
    (for(I,1,9),
     param(Method,Matrix) do
        Method:{\bf{}alldifferent}(Matrix[I,1..9]),
        Method:{\bf{}alldifferent}(Matrix[1..9,I])
    ),
    (multifor([I,J],[1,1],[7,7],[3,3]),
     param(Method,Matrix) do
        Method:{\bf{}alldifferent}(flatten(Matrix[I..I+2,
                                       J..J+2]))
    ),
    {\bf{}flatten_array}(Matrix,List),{\bf{}labeling}(List).
\end{semiverbatim}
\end{frame}

\subsection{Bounds Consistency}

\begin{frame}
\frametitle{Initial State (Bounds Consistency)}
\input{dump/dump_bc_1}
\end{frame}

\begin{frame}
\frametitle{Propagation Steps (Bounds Consistency)}
\only<1>{\input{dump/dump_bc_1}}
\only<2>{\input{dump/dump_bc_2}}
\only<3>{\input{dump/dump_bc_3}}
\only<4>{\input{dump/dump_bc_4}}
\only<5>{\input{dump/dump_bc_5}}
\only<6>{\input{dump/dump_bc_6}}
\only<7>{\input{dump/dump_bc_7}}
\only<8>{\input{dump/dump_bc_8}}
\only<9>{\input{dump/dump_bc_9}}
\only<10>{\input{dump/dump_bc_10}}
\only<11>{\input{dump/dump_bc_11}}
\only<12>{\input{dump/dump_bc_12}}
\only<13>{\input{dump/dump_bc_13}}
\only<14>{\input{dump/dump_bc_14}}
\only<15>{\input{dump/dump_bc_15}}
\only<16>{\input{dump/dump_bc_16}}
\only<17>{\input{dump/dump_bc_17}}
\only<18>{\input{dump/dump_bc_18}}
\only<19>{\input{dump/dump_bc_19}}
\only<20>{\input{dump/dump_bc_20}}
\only<21>{\input{dump/dump_bc_21}}
\only<22>{\input{dump/dump_bc_22}}
\only<23>{\input{dump/dump_bc_23}}
\only<24>{\input{dump/dump_bc_24}}
\only<25>{\input{dump/dump_bc_25}}
\only<26>{\input{dump/dump_bc_26}}
\only<27>{\input{dump/dump_bc_27}}
\only<28>{\input{dump/dump_bc_28}}
\vfill
\hyperlinkframestart<2-28>{\beamerreturnbutton{Back to Start}}
\hyperlinkframeend<1-27>{\beamerskipbutton{Skip Animation}}
\end{frame}

\begin{frame}
\frametitle{After Setup (Bounds Consistency)}
\input{dump/dump_bc_29}
\end{frame}

\subsection{Domain Consistency}

\begin{frame}
\frametitle{Initial State (Domain Consistency)}
\input{dump/dump_gac_1}
\end{frame}

\begin{frame}
\frametitle{Propagation Steps (Domain Consistency)}
\only<1>{\input{dump/dump_gac_1}}
\only<2>{\input{dump/dump_gac_2}}
\only<3>{\input{dump/dump_gac_3}}
\only<4>{\input{dump/dump_gac_4}}
\only<5>{\input{dump/dump_gac_5}}
\only<6>{\input{dump/dump_gac_6}}
\only<7>{\input{dump/dump_gac_7}}
\only<8>{\input{dump/dump_gac_8}}
\only<9>{\input{dump/dump_gac_9}}
\only<10>{\input{dump/dump_gac_10}}
\only<11>{\input{dump/dump_gac_11}}
\only<12>{\input{dump/dump_gac_12}}
\only<13>{\input{dump/dump_gac_13}}
\only<14>{\input{dump/dump_gac_14}}
\only<15>{\input{dump/dump_gac_15}}
\only<16>{\input{dump/dump_gac_16}}
\only<17>{\input{dump/dump_gac_17}}
\only<18>{\input{dump/dump_gac_18}}
\only<19>{\input{dump/dump_gac_19}}
\only<20>{\input{dump/dump_gac_20}}
\only<21>{\input{dump/dump_gac_21}}
\only<22>{\input{dump/dump_gac_22}}
\only<23>{\input{dump/dump_gac_23}}
\only<24>{\input{dump/dump_gac_24}}
\only<25>{\input{dump/dump_gac_25}}
\only<26>{\input{dump/dump_gac_26}}
\only<27>{\input{dump/dump_gac_27}}
\only<28>{\input{dump/dump_gac_28}}
\vfill
\hyperlinkframestart<2-28>{\beamerreturnbutton{Back to Start}}
\hyperlinkframeend<1-27>{\beamerskipbutton{Skip Animation}}
\end{frame}

\begin{frame}
\frametitle{After Setup (Domain Consistency)}
\input{dump/dump_gac_29}
\end{frame}

\subsection{Comparison}

\begin{frame}
\frametitle{Comparison}
\begin{tabular}{ccc}
Forward Checking & Bounds Consistency & Domain Consistency \\
\scalebox{0.55}{
\input{dump/dump_fc_29}
}&
\scalebox{0.55}{
\input{dump/dump_bc_29}
}&
\scalebox{0.55}{
\input{dump/dump_gac_29}
}
\end{tabular}
\end{frame}

\begin{frame}
\frametitle{Typical?}
\begin{itemize}
\item This does not always happen
\item Sometimes, two methods produce same amount of propagation
\item Possible to predict in certain special cases
\item In general, tradeoff between speed and propagation
\item Not always fastest to remove inconsistent values early
\item But often required to find a solution at all
\end{itemize}
\end{frame}

\section{Search}

\begin{frame}
\frametitle{Simple search routine}
\begin{itemize}
\item Enumerate variables in given order
\item Try values starting from smallest one in domain
\item Complete, chronological backtracking
\end{itemize}
\end{frame}

\begin{frame}[label=searchfc]
\frametitle{Search Tree (Forward Checking)}
\begin{textblock}{5.5}(9,0)
\only<1>{\scalebox{0.7}{\input{dump/dump_fc_29}}}
\only<2>{\scalebox{0.7}{\input{dump/dump_fc_30}}}
\only<3>{\scalebox{0.7}{\input{dump/dump_fc_31}}}
\only<7>{\scalebox{0.7}{\input{dump/dump_fc_32}}}
\only<12>{\scalebox{0.7}{\input{dump/dump_fc_33}}}
\only<13>{\scalebox{0.7}{\input{dump/dump_fc_34}}}
\only<16>{\scalebox{0.7}{\input{dump/dump_fc_35}}}
\only<18>{\scalebox{0.7}{\input{dump/dump_fc_36}}}
\only<19>{\scalebox{0.7}{\input{dump/dump_fc_37}}}
\only<20>{\scalebox{0.7}{\input{dump/dump_fc_38}}}
\end{textblock}
\includegraphics[width=8cm]{sudoku/BC/tree}

\vfill
\hyperlinkframestart<2-20>{\beamerreturnbutton{Back to Start}}
\hyperlinkframeend<1-19>{\beamerskipbutton{Skip Animation}}
\end{frame}


\begin{frame}[label=seachbc]
\frametitle{Search Tree (Bounds Consistency)}
\begin{textblock}{5.5}(9,0)
\only<1>{\scalebox{0.7}{\input{dump/dump_bc_29}}}
\only<3>{\scalebox{0.7}{\input{dump/dump_bc_30}}}
\only<4>{\scalebox{0.7}{\input{dump/dump_bc_31}}}
\end{textblock}
\only<1>{
\scalebox{ 0.50}{\psset{nodesep=0pt,levelsep=36pt}
\pstree{\CC{\hyperlink{searchbc<1>}{2}}}{
}
}
}
\only<2>{
\scalebox{ 0.50}{\psset{nodesep=0pt,levelsep=36pt}
\pstree{\CC{\hyperlink{searchbc<1>}{2}}}{
 \psset{linestyle=solid}\FAIL{searchbc<2>}\trput{1}
}
}
}
\only<3>{
\scalebox{ 0.50}{\psset{nodesep=0pt,levelsep=36pt}
\pstree{\CC{\hyperlink{searchbc<1>}{2}}}{
 \psset{linestyle=solid}\FAIL{searchbc<2>}\trput{1}
 \psset{linestyle=solid}\pstree{\CC{\hyperlink{searchbc<3>}{8}}\trput{2}}{
 }
}
}
}
\only<4>{
\scalebox{ 0.50}{\psset{nodesep=0pt,levelsep=36pt}
\pstree{\CC{\hyperlink{searchbc<1>}{2}}}{
 \psset{linestyle=solid}\FAIL{searchbc<2>}\trput{1}
 \psset{linestyle=solid}\pstree{\CC{\hyperlink{searchbc<3>}{8}}\trput{2}}{
  \psset{linestyle=solid}\SUCC{searchbc<4>}\trput{7}
 }
}
}
}

\vfill
\hyperlinkframestart<2-4>{\beamerreturnbutton{Back to Start}}
\hyperlinkframeend<1-3>{\beamerskipbutton{Skip Animation}}
\end{frame}

\begin{frame}[label=searchgac]
\frametitle{Search Tree (Domain Consistency)}
\begin{textblock}{5.5}(9,0)
\only<1>{\scalebox{0.7}{\input{dump/dump_gac_29}}}
\only<2>{\scalebox{0.7}{\input{dump/dump_gac_30}}}
\end{textblock}
\input{tree2}
\vfill
\hyperlinkframestart<2>{\beamerreturnbutton{Back to Start}}
\hyperlinkframeend<1>{\beamerskipbutton{Skip Animation}}
\end{frame}

\begin{frame}
\frametitle{Observations}
\begin{itemize}
\item Search tree much smaller for bounds/domain consistency
\item Does not always happen like this
\item Smaller tree = Less execution time
\item Less reasoning = Less execution time
\item Problem: Finding best balance
\item For Sudoku: not good enough, should not require any search!
\end{itemize}
\end{frame}

\subsection{Solution}
\begin{frame}
\frametitle{Solution}
\input{dump/dump_gac_31}
\end{frame}


\section{Lessons Learned}

\begin{frame}
\frametitle{Global Constraints}
\begin{itemize}
\item Powerful modelling abstractions
\item Efficient reasoning
\end{itemize}
\end{frame}


\begin{frame}
\frametitle{Consistency Levels}
\begin{itemize}
\item Defined levels of propagation
\item Tradeoff speed/reasoning
\item Characterisation of power of constraint
\end{itemize}
\end{frame}

\begin{frame}
\frametitle{Alldifferent Variants}
\begin{itemize}
\item Forward Checking
\begin{itemize}
\item Only reacts when variables are assigned
\item Equivalent to decomposition into binary constraints
\end{itemize}
\item Bounds Consistency
\begin{itemize}
\item Typical best compomise speed/reasoning
\item Works well if no holes in domain
\end{itemize}
\item Domain Consistency
\begin{itemize}
\item Extracts all information from single constraint
\item Cost only justified for very hard problems
\end{itemize}
\end{itemize}
\end{frame}

\begin{frame}
\frametitle{End of Chapter 5}
\begin{block}{Thank you!}
Some optional material follows
\end{block}
\end{frame}

\appendix
\section{Optional Material}

\subsection{Complete Example: Domain Consistent Alldifferent}

\begin{frame}[fragile]
\frametitle{Bigger Example}
\begin{semiverbatim}
:-lib(ic).
:-lib(ic_global_gac).

top:-
    [X,Y] :: 1..2,
    Z :: 2..5,
    [T,U] :: 3..5,
    V :: [2,4,6,7],
    ic_global_gac:alldifferent([X,Y,Z,T,U,V]).
\end{semiverbatim}
\end{frame}

\begin{frame}
\frametitle{Making constraint domain consistent}
\begin{textblock}{8}(6,0)
\only<1>{Problem shown as bipartite graph}
\only<2>{Find maximal matching (in blue)}
\only<3>{Orient graph (edges in matching from variables to values, all others from values to variables), mark edges in matching}
\only<4>{Find strongly connected components (green and brown), mark their edges}
\only<5>{Find unmatched value nodes (here node 7, magenta)}
\only<6>{Find alternating paths from such nodes (in magenta), mark their edges}
\only<7>{All unmarked edges can be removed}
\only<8>{Resulting graph, constraint is domain consistent}
\end{textblock}
\input{bigexample}
\end{frame}

\begin{frame}[fragile]
\frametitle{Extended Example}
\begin{semiverbatim}
:-lib(ic).
:-lib(ic_global_gac).

top:-
    \textcolor{red}{X :: 1..2},
    \textcolor{red}{Y :: [1,2,7]},
    Z :: 2..5,
    [T,U] :: 3..5,
    V :: [2,4,6,7],
    ic_global_gac:alldifferent([X,Y,Z,T,U,V]).
\end{semiverbatim}
\end{frame}

\begin{frame}
\frametitle{No propagation in expanded example}
\begin{textblock}{8}(6,0)
\only<1>{Problem shown as bipartite graph}
\only<2>{Find maximal matching (in blue)}
\only<3>{Orient graph (edges in matching from variables to values, all others from values to variables), mark edges in matching}
\only<4>{Find strongly connected components (green and brown), mark their edges}
\only<5>{Find unmatched value nodes (here node 7, magenta)}
\only<6>{Find alternating paths from such nodes (in magenta), mark their edges}
\only<7>{Continue with alternating paths}
\only<8>{Continue with alternating paths, all edges marked, no propagation, constraint is domain consistent}
\end{textblock}
\only<1>{
\begin{psmatrix}[rowsep=12pt]
X & & 1 \\
Y & & 2 \\
Z & & 3 \\
T & & 4 \\
U & & 5 \\
V & & 6 \\
  & & 7 
\ncline{1,1}{1,3} 
\ncline{1,1}{2,3} 
\ncline{2,1}{1,3} 
\ncline{2,1}{2,3} 
\ncline{2,1}{7,3} 
\ncline{3,1}{2,3} 
\ncline{3,1}{3,3} 
\ncline{3,1}{4,3} 
\ncline{3,1}{5,3} 
\ncline{4,1}{3,3} 
\ncline{4,1}{4,3} 
\ncline{4,1}{5,3} 
\ncline{5,1}{3,3} 
\ncline{5,1}{4,3} 
\ncline{5,1}{5,3} 
\ncline{6,1}{2,3} 
\ncline{6,1}{4,3} 
\ncline{6,1}{6,3} 
\ncline{6,1}{7,3} 
\end{psmatrix}
}
\only<2>{
\begin{psmatrix}[rowsep=12pt]
X & & 1 \\
Y & & 2 \\
Z & & 3 \\
T & & 4 \\
U & & 5 \\
V & & 6 \\
  & & 7 
\ncline[linecolor=insight-blue]{1,1}{1,3} 
\ncline{1,1}{2,3} 
\ncline{2,1}{1,3} 
\ncline[linecolor=insight-blue]{2,1}{2,3} 
\ncline{2,1}{7,3} 
\ncline{3,1}{2,3} 
\ncline{3,1}{3,3} 
\ncline[linecolor=insight-blue]{3,1}{4,3} 
\ncline{3,1}{5,3} 
\ncline{4,1}{3,3} 
\ncline{4,1}{4,3} 
\ncline[linecolor=insight-blue]{4,1}{5,3} 
\ncline[linecolor=insight-blue]{5,1}{3,3} 
\ncline{5,1}{4,3} 
\ncline{5,1}{5,3} 
\ncline{6,1}{2,3} 
\ncline{6,1}{4,3} 
\ncline[linecolor=insight-blue]{6,1}{6,3} 
\ncline{6,1}{7,3} 
\end{psmatrix}
}
\only<3>{
\begin{psmatrix}[rowsep=12pt]
X & & 1 \\
Y & & 2 \\
Z & & 3 \\
T & & 4 \\
U & & 5 \\
V & & 6 \\
  & & 7 
\ncline[arrows=->,linecolor=insight-blue]{1,1}{1,3} 
\ncline[linestyle=dashed,arrows=<-]{1,1}{2,3} 
\ncline[linestyle=dashed,arrows=<-]{2,1}{1,3} 
\ncline[arrows=->,linecolor=insight-blue]{2,1}{2,3} 
\ncline[linestyle=dashed,arrows=<-]{2,1}{7,3} 
\ncline[linestyle=dashed,arrows=<-]{3,1}{2,3} 
\ncline[linestyle=dashed,arrows=<-]{3,1}{3,3} 
\ncline[arrows=->,linecolor=insight-blue]{3,1}{4,3} 
\ncline[linestyle=dashed,arrows=<-]{3,1}{5,3} 
\ncline[linestyle=dashed,arrows=<-]{4,1}{3,3} 
\ncline[linestyle=dashed,arrows=<-]{4,1}{4,3} 
\ncline[arrows=->,linecolor=insight-blue]{4,1}{5,3} 
\ncline[arrows=->,linecolor=insight-blue]{5,1}{3,3} 
\ncline[linestyle=dashed,arrows=<-]{5,1}{4,3} 
\ncline[linestyle=dashed,arrows=<-]{5,1}{5,3} 
\ncline[linestyle=dashed,arrows=<-]{6,1}{2,3} 
\ncline[linestyle=dashed,arrows=<-]{6,1}{4,3} 
\ncline[arrows=->,linecolor=insight-blue]{6,1}{6,3} 
\ncline[linestyle=dashed,arrows=<-]{6,1}{7,3} 
\end{psmatrix}
}
\only<4>{
\begin{psmatrix}[rowsep=12pt]
X & & 1 \\
Y & & 2 \\
Z & & 3 \\
T & & 4 \\
U & & 5 \\
V & & 6 \\
  & & 7 
\ncline[arrows=->,linecolor=insight-green]{1,1}{1,3} 
\ncline[arrows=<-,linecolor=insight-green]{1,1}{2,3} 
\ncline[arrows=<-,linecolor=insight-green]{2,1}{1,3} 
\ncline[arrows=->,linecolor=insight-green]{2,1}{2,3} 
\ncline[linestyle=dashed,arrows=<-]{2,1}{7,3} 
\ncline[linestyle=dashed,arrows=<-]{3,1}{2,3} 
\ncline[arrows=<-,linecolor=brown]{3,1}{3,3} 
\ncline[arrows=->,linecolor=brown]{3,1}{4,3} 
\ncline[arrows=<-,linecolor=brown]{3,1}{5,3} 
\ncline[arrows=<-,linecolor=brown]{4,1}{3,3} 
\ncline[arrows=<-,linecolor=brown]{4,1}{4,3} 
\ncline[arrows=->,linecolor=brown]{4,1}{5,3} 
\ncline[arrows=->,linecolor=brown]{5,1}{3,3} 
\ncline[arrows=<-,linecolor=brown]{5,1}{4,3} 
\ncline[arrows=<-,linecolor=brown]{5,1}{5,3} 
\ncline[linestyle=dashed,arrows=<-]{6,1}{2,3} 
\ncline[linestyle=dashed,arrows=<-]{6,1}{4,3} 
\ncline[arrows=->]{6,1}{6,3} 
\ncline[linestyle=dashed,arrows=<-]{6,1}{7,3} 
\end{psmatrix}
}
\only<5>{
\begin{psmatrix}[rowsep=12pt]
X & & 1 \\
Y & & 2 \\
Z & & 3 \\
T & & 4 \\
U & & 5 \\
V & & 6 \\
  & & \textcolor{magenta}{7} 
\ncline[arrows=->]{1,1}{1,3} 
\ncline[arrows=<-]{1,1}{2,3} 
\ncline[arrows=<-]{2,1}{1,3} 
\ncline[arrows=->]{2,1}{2,3} 
\ncline[linestyle=dashed,arrows=<-]{2,1}{7,3} 
\ncline[linestyle=dashed,arrows=<-]{3,1}{2,3} 
\ncline[arrows=<-]{3,1}{3,3} 
\ncline[arrows=->]{3,1}{4,3} 
\ncline[arrows=<-]{3,1}{5,3} 
\ncline[arrows=<-]{4,1}{3,3} 
\ncline[arrows=<-]{4,1}{4,3} 
\ncline[arrows=->]{4,1}{5,3} 
\ncline[arrows=->]{5,1}{3,3} 
\ncline[arrows=<-]{5,1}{4,3} 
\ncline[arrows=<-]{5,1}{5,3} 
\ncline[linestyle=dashed,arrows=<-]{6,1}{2,3} 
\ncline[linestyle=dashed,arrows=<-]{6,1}{4,3} 
\ncline[arrows=->]{6,1}{6,3} 
\ncline[linestyle=dashed,arrows=<-]{6,1}{7,3} 
\end{psmatrix}
}
\only<6>{
\begin{psmatrix}[rowsep=12pt]
X & & 1 \\
Y & & 2 \\
Z & & 3 \\
T & & 4 \\
U & & 5 \\
V & & 6 \\
  & & \textcolor{magenta}{7} 
\ncline[arrows=->]{1,1}{1,3} 
\ncline[arrows=<-]{1,1}{2,3} 
\ncline[arrows=<-]{2,1}{1,3} 
\ncline[linecolor=magenta,arrows=->]{2,1}{2,3} 
\ncline[linecolor=magenta,arrows=<-]{2,1}{7,3} 
\ncline[linestyle=dashed,arrows=<-]{3,1}{2,3} 
\ncline[arrows=<-]{3,1}{3,3} 
\ncline[arrows=->]{3,1}{4,3} 
\ncline[arrows=<-]{3,1}{5,3} 
\ncline[arrows=<-]{4,1}{3,3} 
\ncline[arrows=<-]{4,1}{4,3} 
\ncline[arrows=->]{4,1}{5,3} 
\ncline[arrows=->]{5,1}{3,3} 
\ncline[arrows=<-]{5,1}{4,3} 
\ncline[arrows=<-]{5,1}{5,3} 
\ncline[linestyle=dashed,arrows=<-]{6,1}{2,3} 
\ncline[linestyle=dashed,arrows=<-]{6,1}{4,3} 
\ncline[linecolor=magenta,arrows=->]{6,1}{6,3} 
\ncline[linecolor=magenta,arrows=<-]{6,1}{7,3} 
\end{psmatrix}
}
\only<7>{
\begin{psmatrix}[rowsep=12pt]
X & & 1 \\
Y & & 2 \\
Z & & 3 \\
T & & 4 \\
U & & 5 \\
V & & 6 \\
  & & \textcolor{magenta}{7} 
\ncline[linecolor=magenta,arrows=->]{1,1}{1,3} 
\ncline[linecolor=magenta,arrows=<-]{1,1}{2,3} 
\ncline[arrows=<-]{2,1}{1,3} 
\ncline[linecolor=magenta,arrows=->]{2,1}{2,3} 
\ncline[linecolor=magenta,arrows=<-]{2,1}{7,3} 
\ncline[linecolor=magenta,arrows=<-]{3,1}{2,3} 
\ncline[arrows=<-]{3,1}{3,3} 
\ncline[linecolor=magenta,arrows=->]{3,1}{4,3} 
\ncline[arrows=<-]{3,1}{5,3} 
\ncline[arrows=<-]{4,1}{3,3} 
\ncline[arrows=<-]{4,1}{4,3} 
\ncline[arrows=->]{4,1}{5,3} 
\ncline[arrows=->]{5,1}{3,3} 
\ncline[arrows=<-]{5,1}{4,3} 
\ncline[arrows=<-]{5,1}{5,3} 
\ncline[linecolor=magenta,arrows=<-]{6,1}{2,3} 
\ncline[linestyle=dashed,arrows=<-]{6,1}{4,3} 
\ncline[linecolor=magenta,arrows=->]{6,1}{6,3} 
\ncline[linecolor=magenta,arrows=<-]{6,1}{7,3} 
\end{psmatrix}
}
\only<8>{
\begin{psmatrix}[rowsep=12pt]
X & & 1 \\
Y & & 2 \\
Z & & 3 \\
T & & 4 \\
U & & 5 \\
V & & 6 \\
  & & \textcolor{magenta}{7} 
\ncline[linecolor=magenta,arrows=->]{1,1}{1,3} 
\ncline[linecolor=magenta,arrows=<-]{1,1}{2,3} 
\ncline[linecolor=magenta,arrows=<-]{2,1}{1,3} 
\ncline[linecolor=magenta,arrows=->]{2,1}{2,3} 
\ncline[linecolor=magenta,arrows=<-]{2,1}{7,3} 
\ncline[linecolor=magenta,arrows=<-]{3,1}{2,3} 
\ncline[arrows=<-]{3,1}{3,3} 
\ncline[linecolor=magenta,arrows=->]{3,1}{4,3} 
\ncline[arrows=<-]{3,1}{5,3} 
\ncline[arrows=<-]{4,1}{3,3} 
\ncline[linecolor=magenta,arrows=<-]{4,1}{4,3} 
\ncline[arrows=->]{4,1}{5,3} 
\ncline[arrows=->]{5,1}{3,3} 
\ncline[linecolor=magenta,arrows=<-]{5,1}{4,3} 
\ncline[arrows=<-]{5,1}{5,3} 
\ncline[linecolor=magenta,arrows=<-]{6,1}{2,3} 
\ncline[linecolor=magenta,arrows=<-]{6,1}{4,3} 
\ncline[linecolor=magenta,arrows=->]{6,1}{6,3} 
\ncline[linecolor=magenta,arrows=<-]{6,1}{7,3} 
\end{psmatrix}
}

\end{frame}

\begin{frame}
\frametitle{Observation}
\begin{itemize}
\item A lot of effort for no propagation
\item Problem: Slows down search without any upside
\item Constraint is woken every time any domain is changed
\item How often does the constraint do actual pruning?
\end{itemize}
\end{frame}

\subsection{Generic Model}

\begin{frame}
\frametitle{Generalize Program for different sizes}
\begin{itemize}
\item How to generalize program for different sizes (4,9,16,25,36...)
\item Add parameter R (Order, number of blocks in a row/column)
\item Size N is square of R
\item Remove explicit integer bounds by expressions 
\item Useful to do this change as rewriting of working program
\end{itemize}

\end{frame}
\begin{frame}[fragile]
\frametitle{Main Program}
\begin{semiverbatim}
model(R,M,Matrix):-
    N is R*R,Matrix[1..N,1..N] :: 1..N,
    (for(I,1,N),
     param(N,M,Matrix) do
        M:{\bf{}alldifferent}(Matrix[I,1..N]),
        M:{\bf{}alldifferent}(Matrix[1..N,I])
    ),
    (multifor([I,J],[1,1],[N-R+1,N-R+1],[R,R]),
     param(R,M,Matrix) do
        M:{\bf{}alldifferent}(flatten(Matrix[I..I+R-1,
                                 J..J+R-1]))
    ),
    {\bf{}flatten_array}(Matrix,List),{\bf{}labeling}(List).
\end{semiverbatim}
\end{frame}

\end{document}
------------------------------------------------------------------------------

\subsection{Problem Variant}
\begin{frame}
\frametitle{Problem Variant}
\begin{block}{9 Spot }
Fill in numbers from 1 to 9 so that each row, column and block contain each number exactly once. Blocks contain 9 cells, but have irregular shape.
\end{block}
Problem and Instance provided by Raphael Finkel
\end{frame}

\begin{frame}
\frametitle{Model}
\begin{itemize}
\item Same as before
\item Blocks have irregular shape
\item Therefore, blocks must be given as part of data
\item Each cell belongs to exactly one block
\item Each block contains 9 cells
\end{itemize}
\end{frame}

\begin{frame}
\frametitle{Main difference}
\begin{itemize}
\item Interaction of rows/columns and blocks
\item Row and block share between 1 and 5 cells
\item Sudoku: exactly three cells
\item This seems to make problem harder
\end{itemize}
\end{frame}

\begin{frame}
\frametitle{Initial State (Domain Consistency)}
\input{spot/dump_gac_1}
\end{frame}

\begin{frame}
\frametitle{Propagation Steps (Domain Consistency)}
\only<1>{\input{spot/dump_gac_1}}
\only<2>{\input{spot/dump_gac_2}}
\only<3>{\input{spot/dump_gac_3}}
\only<4>{\input{spot/dump_gac_4}}
\only<5>{\input{spot/dump_gac_5}}
\only<6>{\input{spot/dump_gac_6}}
\only<7>{\input{spot/dump_gac_7}}
\only<8>{\input{spot/dump_gac_8}}
\only<9>{\input{spot/dump_gac_9}}
\only<10>{\input{spot/dump_gac_10}}
\only<11>{\input{spot/dump_gac_11}}
\only<12>{\input{spot/dump_gac_12}}
\only<13>{\input{spot/dump_gac_13}}
\only<14>{\input{spot/dump_gac_14}}
\only<15>{\input{spot/dump_gac_15}}
\only<16>{\input{spot/dump_gac_16}}
\only<17>{\input{spot/dump_gac_17}}
\only<18>{\input{spot/dump_gac_18}}
\only<19>{\input{spot/dump_gac_19}}
\only<20>{\input{spot/dump_gac_20}}
\only<21>{\input{spot/dump_gac_21}}
\only<22>{\input{spot/dump_gac_22}}
\only<23>{\input{spot/dump_gac_23}}
\only<24>{\input{spot/dump_gac_24}}
\only<25>{\input{spot/dump_gac_25}}
\only<26>{\input{spot/dump_gac_26}}
\only<27>{\input{spot/dump_gac_27}}
\only<28>{\input{spot/dump_gac_28}}
\vfill
\hyperlinkframestart<2-28>{\beamerreturnbutton{Back to Start}}
\hyperlinkframeend<1-27>{\beamerskipbutton{Skip Animation}}
\end{frame}

\begin{frame}
\frametitle{After Setup (Domain Consistency)}
\input{spot/dump_gac_29}
\end{frame}

\begin{frame}
\frametitle{Shaving Iterations (Domain Consistency)}
\only<1>{\input{spot/dump_gac_29}}
\only<2>{\input{spot/dump_gac_30}}
\only<3>{\input{spot/dump_gac_31}}
\only<4>{\input{spot/dump_gac_32}}
\only<5>{\input{spot/dump_gac_33}}
\only<6>{\input{spot/dump_gac_34}}
\vfill
\hyperlinkframestart<2-6>{\beamerreturnbutton{Back to Start}}
\hyperlinkframeend<1-5>{\beamerskipbutton{Skip Animation}}
\end{frame}

\begin{frame}
\frametitle{Overlap (Domain Consistency)}
\input{spot/dump_gac_35}
\end{frame}



\begin{frame}[label=spotgac]
\frametitle{9Spot Search (Domain Consistency)}
\only<1>{
\scalebox{ 0.40}{\psset{nodesep=0pt,levelsep=36pt}
\pstree{\CC{\hyperlink{spotgac<1>}{4}}}{
}
}
}
\only<2>{
\scalebox{ 0.40}{\psset{nodesep=0pt,levelsep=36pt}
\pstree{\CC{\hyperlink{spotgac<1>}{4}}}{
 \psset{linestyle=solid}\pstree{\CC{\hyperlink{spotgac<2>}{5}}\trput{2}}{
 }
}
}
}
\only<3>{
\scalebox{ 0.40}{\psset{nodesep=0pt,levelsep=36pt}
\pstree{\CC{\hyperlink{spotgac<1>}{4}}}{
 \psset{linestyle=solid}\pstree{\CC{\hyperlink{spotgac<2>}{5}}\trput{2}}{
  \psset{linestyle=solid}\pstree{\CC{\hyperlink{spotgac<3>}{6}}\trput{4}}{
  }
 }
}
}
}
\only<4>{
\scalebox{ 0.40}{\psset{nodesep=0pt,levelsep=36pt}
\pstree{\CC{\hyperlink{spotgac<1>}{4}}}{
 \psset{linestyle=solid}\pstree{\CC{\hyperlink{spotgac<2>}{5}}\trput{2}}{
  \psset{linestyle=solid}\pstree{\CC{\hyperlink{spotgac<3>}{6}}\trput{4}}{
   \psset{linestyle=solid}\pstree{\CC{\hyperlink{spotgac<4>}{7}}\trput{5}}{
   }
  }
 }
}
}
}
\only<5>{
\scalebox{ 0.40}{\psset{nodesep=0pt,levelsep=36pt}
\pstree{\CC{\hyperlink{spotgac<1>}{4}}}{
 \psset{linestyle=solid}\pstree{\CC{\hyperlink{spotgac<2>}{5}}\trput{2}}{
  \psset{linestyle=solid}\pstree{\CC{\hyperlink{spotgac<3>}{6}}\trput{4}}{
   \psset{linestyle=solid}\pstree{\CC{\hyperlink{spotgac<4>}{7}}\trput{5}}{
    \psset{linestyle=solid}\FAIL{spotgac<5>}\trput{6}
   }
  }
 }
}
}
}
\only<6>{
\scalebox{ 0.40}{\psset{nodesep=0pt,levelsep=36pt}
\pstree{\CC{\hyperlink{spotgac<1>}{4}}}{
 \psset{linestyle=solid}\pstree{\CC{\hyperlink{spotgac<2>}{5}}\trput{2}}{
  \psset{linestyle=solid}\pstree{\CC{\hyperlink{spotgac<3>}{6}}\trput{4}}{
   \psset{linestyle=solid}\pstree{\CC{\hyperlink{spotgac<4>}{7}}\trput{5}}{
    \psset{linestyle=solid}\FAIL{spotgac<5>}\trput{6}
    \psset{linestyle=solid}\pstree{\CC{\hyperlink{spotgac<6>}{12}}\trput{9}}{
    }
   }
  }
 }
}
}
}
\only<7>{
\scalebox{ 0.40}{\psset{nodesep=0pt,levelsep=36pt}
\pstree{\CC{\hyperlink{spotgac<1>}{4}}}{
 \psset{linestyle=solid}\pstree{\CC{\hyperlink{spotgac<2>}{5}}\trput{2}}{
  \psset{linestyle=solid}\pstree{\CC{\hyperlink{spotgac<3>}{6}}\trput{4}}{
   \psset{linestyle=solid}\pstree{\CC{\hyperlink{spotgac<4>}{7}}\trput{5}}{
    \psset{linestyle=solid}\FAIL{spotgac<5>}\trput{6}
    \psset{linestyle=solid}\pstree{\CC{\hyperlink{spotgac<6>}{12}}\trput{9}}{
     \psset{linestyle=solid}\pstree{\CC{\hyperlink{spotgac<7>}{19}}\trput{6}}{
     }
    }
   }
  }
 }
}
}
}
\only<8>{
\scalebox{ 0.40}{\psset{nodesep=0pt,levelsep=36pt}
\pstree{\CC{\hyperlink{spotgac<1>}{4}}}{
 \psset{linestyle=solid}\pstree{\CC{\hyperlink{spotgac<2>}{5}}\trput{2}}{
  \psset{linestyle=solid}\pstree{\CC{\hyperlink{spotgac<3>}{6}}\trput{4}}{
   \psset{linestyle=solid}\pstree{\CC{\hyperlink{spotgac<4>}{7}}\trput{5}}{
    \psset{linestyle=solid}\FAIL{spotgac<5>}\trput{6}
    \psset{linestyle=solid}\pstree{\CC{\hyperlink{spotgac<6>}{12}}\trput{9}}{
     \psset{linestyle=solid}\pstree{\CC{\hyperlink{spotgac<7>}{19}}\trput{6}}{
      \psset{linestyle=solid}\FAIL{spotgac<8>}\trput{6}
     }
    }
   }
  }
 }
}
}
}
\only<9>{
\scalebox{ 0.40}{\psset{nodesep=0pt,levelsep=36pt}
\pstree{\CC{\hyperlink{spotgac<1>}{4}}}{
 \psset{linestyle=solid}\pstree{\CC{\hyperlink{spotgac<2>}{5}}\trput{2}}{
  \psset{linestyle=solid}\pstree{\CC{\hyperlink{spotgac<3>}{6}}\trput{4}}{
   \psset{linestyle=solid}\pstree{\CC{\hyperlink{spotgac<4>}{7}}\trput{5}}{
    \psset{linestyle=solid}\FAIL{spotgac<5>}\trput{6}
    \psset{linestyle=solid}\pstree{\CC{\hyperlink{spotgac<6>}{12}}\trput{9}}{
     \psset{linestyle=solid}\pstree{\CC{\hyperlink{spotgac<7>}{19}}\trput{6}}{
      \psset{linestyle=solid}\FAIL{spotgac<8>}\trput{6}
      \psset{linestyle=solid}\pstree{\CC{\hyperlink{spotgac<9>}{20}}\trput{9}}{
      }
     }
    }
   }
  }
 }
}
}
}
\only<10>{
\scalebox{ 0.40}{\psset{nodesep=0pt,levelsep=36pt}
\pstree{\CC{\hyperlink{spotgac<1>}{4}}}{
 \psset{linestyle=solid}\pstree{\CC{\hyperlink{spotgac<2>}{5}}\trput{2}}{
  \psset{linestyle=solid}\pstree{\CC{\hyperlink{spotgac<3>}{6}}\trput{4}}{
   \psset{linestyle=solid}\pstree{\CC{\hyperlink{spotgac<4>}{7}}\trput{5}}{
    \psset{linestyle=solid}\FAIL{spotgac<5>}\trput{6}
    \psset{linestyle=solid}\pstree{\CC{\hyperlink{spotgac<6>}{12}}\trput{9}}{
     \psset{linestyle=solid}\pstree{\CC{\hyperlink{spotgac<7>}{19}}\trput{6}}{
      \psset{linestyle=solid}\FAIL{spotgac<8>}\trput{6}
      \psset{linestyle=solid}\pstree{\CC{\hyperlink{spotgac<9>}{20}}\trput{9}}{
       \psset{linestyle=solid}\pstree{\CC{\hyperlink{spotgac<10>}{23}}\trput{6}}{
       }
      }
     }
    }
   }
  }
 }
}
}
}
\only<11>{
\scalebox{ 0.40}{\psset{nodesep=0pt,levelsep=36pt}
\pstree{\CC{\hyperlink{spotgac<1>}{4}}}{
 \psset{linestyle=solid}\pstree{\CC{\hyperlink{spotgac<2>}{5}}\trput{2}}{
  \psset{linestyle=solid}\pstree{\CC{\hyperlink{spotgac<3>}{6}}\trput{4}}{
   \psset{linestyle=solid}\pstree{\CC{\hyperlink{spotgac<4>}{7}}\trput{5}}{
    \psset{linestyle=solid}\FAIL{spotgac<5>}\trput{6}
    \psset{linestyle=solid}\pstree{\CC{\hyperlink{spotgac<6>}{12}}\trput{9}}{
     \psset{linestyle=solid}\pstree{\CC{\hyperlink{spotgac<7>}{19}}\trput{6}}{
      \psset{linestyle=solid}\FAIL{spotgac<8>}\trput{6}
      \psset{linestyle=solid}\pstree{\CC{\hyperlink{spotgac<9>}{20}}\trput{9}}{
       \psset{linestyle=solid}\pstree{\CC{\hyperlink{spotgac<10>}{23}}\trput{6}}{
        \psset{linestyle=solid}\FAIL{spotgac<11>}\trput{2}
       }
      }
     }
    }
   }
  }
 }
}
}
}
\only<12>{
\scalebox{ 0.40}{\psset{nodesep=0pt,levelsep=36pt}
\pstree{\CC{\hyperlink{spotgac<1>}{4}}}{
 \psset{linestyle=solid}\pstree{\CC{\hyperlink{spotgac<2>}{5}}\trput{2}}{
  \psset{linestyle=solid}\pstree{\CC{\hyperlink{spotgac<3>}{6}}\trput{4}}{
   \psset{linestyle=solid}\pstree{\CC{\hyperlink{spotgac<4>}{7}}\trput{5}}{
    \psset{linestyle=solid}\FAIL{spotgac<5>}\trput{6}
    \psset{linestyle=solid}\pstree{\CC{\hyperlink{spotgac<6>}{12}}\trput{9}}{
     \psset{linestyle=solid}\pstree{\CC{\hyperlink{spotgac<7>}{19}}\trput{6}}{
      \psset{linestyle=solid}\FAIL{spotgac<8>}\trput{6}
      \psset{linestyle=solid}\pstree{\CC{\hyperlink{spotgac<9>}{20}}\trput{9}}{
       \psset{linestyle=solid}\pstree{\CC{\hyperlink{spotgac<10>}{23}}\trput{6}}{
        \psset{linestyle=solid}\FAIL{spotgac<11>}\trput{2}
        \psset{linestyle=solid}\FAIL{spotgac<12>}\trput{8}
       }
      }
     }
    }
   }
  }
 }
}
}
}
\only<13>{
\scalebox{ 0.40}{\psset{nodesep=0pt,levelsep=36pt}
\pstree{\CC{\hyperlink{spotgac<1>}{4}}}{
 \psset{linestyle=solid}\pstree{\CC{\hyperlink{spotgac<2>}{5}}\trput{2}}{
  \psset{linestyle=solid}\pstree{\CC{\hyperlink{spotgac<3>}{6}}\trput{4}}{
   \psset{linestyle=solid}\pstree{\CC{\hyperlink{spotgac<4>}{7}}\trput{5}}{
    \psset{linestyle=solid}\FAIL{spotgac<5>}\trput{6}
    \psset{linestyle=solid}\pstree{\CC{\hyperlink{spotgac<6>}{12}}\trput{9}}{
     \psset{linestyle=solid}\pstree{\CC{\hyperlink{spotgac<7>}{19}}\trput{6}}{
      \psset{linestyle=solid}\FAIL{spotgac<8>}\trput{6}
      \psset{linestyle=solid}\pstree{\CC{\hyperlink{spotgac<9>}{20}}\trput{9}}{
       \psset{linestyle=solid}\pstree{\CC{\hyperlink{spotgac<10>}{23}}\trput{6}}{
        \psset{linestyle=solid}\FAIL{spotgac<11>}\trput{2}
        \psset{linestyle=solid}\FAIL{spotgac<12>}\trput{8}
       }
       \psset{linestyle=solid}\FAIL{spotgac<13>}\trput{8}
      }
     }
    }
   }
  }
 }
}
}
}
\only<14>{
\scalebox{ 0.40}{\psset{nodesep=0pt,levelsep=36pt}
\pstree{\CC{\hyperlink{spotgac<1>}{4}}}{
 \psset{linestyle=solid}\pstree{\CC{\hyperlink{spotgac<2>}{5}}\trput{2}}{
  \psset{linestyle=solid}\pstree{\CC{\hyperlink{spotgac<3>}{6}}\trput{4}}{
   \psset{linestyle=solid}\pstree{\CC{\hyperlink{spotgac<4>}{7}}\trput{5}}{
    \psset{linestyle=solid}\FAIL{spotgac<5>}\trput{6}
    \psset{linestyle=solid}\pstree{\CC{\hyperlink{spotgac<6>}{12}}\trput{9}}{
     \psset{linestyle=solid}\pstree{\CC{\hyperlink{spotgac<7>}{19}}\trput{6}}{
      \psset{linestyle=solid}\FAIL{spotgac<8>}\trput{6}
      \psset{linestyle=solid}\pstree{\CC{\hyperlink{spotgac<9>}{20}}\trput{9}}{
       \psset{linestyle=solid}\pstree{\CC{\hyperlink{spotgac<10>}{23}}\trput{6}}{
        \psset{linestyle=solid}\FAIL{spotgac<11>}\trput{2}
        \psset{linestyle=solid}\FAIL{spotgac<12>}\trput{8}
       }
       \psset{linestyle=solid}\FAIL{spotgac<13>}\trput{8}
      }
     }
     \psset{linestyle=solid}\FAIL{spotgac<14>}\trput{8}
    }
   }
  }
 }
}
}
}
\only<15>{
\scalebox{ 0.40}{\psset{nodesep=0pt,levelsep=36pt}
\pstree{\CC{\hyperlink{spotgac<1>}{4}}}{
 \psset{linestyle=solid}\pstree{\CC{\hyperlink{spotgac<2>}{5}}\trput{2}}{
  \psset{linestyle=solid}\pstree{\CC{\hyperlink{spotgac<3>}{6}}\trput{4}}{
   \psset{linestyle=solid}\pstree{\CC{\hyperlink{spotgac<4>}{7}}\trput{5}}{
    \psset{linestyle=solid}\FAIL{spotgac<5>}\trput{6}
    \psset{linestyle=solid}\pstree{\CC{\hyperlink{spotgac<6>}{12}}\trput{9}}{
     \psset{linestyle=solid}\pstree{\CC{\hyperlink{spotgac<7>}{19}}\trput{6}}{
      \psset{linestyle=solid}\FAIL{spotgac<8>}\trput{6}
      \psset{linestyle=solid}\pstree{\CC{\hyperlink{spotgac<9>}{20}}\trput{9}}{
       \psset{linestyle=solid}\pstree{\CC{\hyperlink{spotgac<10>}{23}}\trput{6}}{
        \psset{linestyle=solid}\FAIL{spotgac<11>}\trput{2}
        \psset{linestyle=solid}\FAIL{spotgac<12>}\trput{8}
       }
       \psset{linestyle=solid}\FAIL{spotgac<13>}\trput{8}
      }
     }
     \psset{linestyle=solid}\FAIL{spotgac<14>}\trput{8}
    }
   }
   \psset{linestyle=solid}\FAIL{spotgac<15>}\trput{6}
  }
 }
}
}
}
\only<16>{
\scalebox{ 0.40}{\psset{nodesep=0pt,levelsep=36pt}
\pstree{\CC{\hyperlink{spotgac<1>}{4}}}{
 \psset{linestyle=solid}\pstree{\CC{\hyperlink{spotgac<2>}{5}}\trput{2}}{
  \psset{linestyle=solid}\pstree{\CC{\hyperlink{spotgac<3>}{6}}\trput{4}}{
   \psset{linestyle=solid}\pstree{\CC{\hyperlink{spotgac<4>}{7}}\trput{5}}{
    \psset{linestyle=solid}\FAIL{spotgac<5>}\trput{6}
    \psset{linestyle=solid}\pstree{\CC{\hyperlink{spotgac<6>}{12}}\trput{9}}{
     \psset{linestyle=solid}\pstree{\CC{\hyperlink{spotgac<7>}{19}}\trput{6}}{
      \psset{linestyle=solid}\FAIL{spotgac<8>}\trput{6}
      \psset{linestyle=solid}\pstree{\CC{\hyperlink{spotgac<9>}{20}}\trput{9}}{
       \psset{linestyle=solid}\pstree{\CC{\hyperlink{spotgac<10>}{23}}\trput{6}}{
        \psset{linestyle=solid}\FAIL{spotgac<11>}\trput{2}
        \psset{linestyle=solid}\FAIL{spotgac<12>}\trput{8}
       }
       \psset{linestyle=solid}\FAIL{spotgac<13>}\trput{8}
      }
     }
     \psset{linestyle=solid}\FAIL{spotgac<14>}\trput{8}
    }
   }
   \psset{linestyle=solid}\FAIL{spotgac<15>}\trput{6}
  }
  \psset{linestyle=solid}\pstree{\CC{\hyperlink{spotgac<16>}{6}}\trput{6}}{
  }
 }
}
}
}
\only<17>{
\scalebox{ 0.40}{\psset{nodesep=0pt,levelsep=36pt}
\pstree{\CC{\hyperlink{spotgac<1>}{4}}}{
 \psset{linestyle=solid}\pstree{\CC{\hyperlink{spotgac<2>}{5}}\trput{2}}{
  \psset{linestyle=solid}\pstree{\CC{\hyperlink{spotgac<3>}{6}}\trput{4}}{
   \psset{linestyle=solid}\pstree{\CC{\hyperlink{spotgac<4>}{7}}\trput{5}}{
    \psset{linestyle=solid}\FAIL{spotgac<5>}\trput{6}
    \psset{linestyle=solid}\pstree{\CC{\hyperlink{spotgac<6>}{12}}\trput{9}}{
     \psset{linestyle=solid}\pstree{\CC{\hyperlink{spotgac<7>}{19}}\trput{6}}{
      \psset{linestyle=solid}\FAIL{spotgac<8>}\trput{6}
      \psset{linestyle=solid}\pstree{\CC{\hyperlink{spotgac<9>}{20}}\trput{9}}{
       \psset{linestyle=solid}\pstree{\CC{\hyperlink{spotgac<10>}{23}}\trput{6}}{
        \psset{linestyle=solid}\FAIL{spotgac<11>}\trput{2}
        \psset{linestyle=solid}\FAIL{spotgac<12>}\trput{8}
       }
       \psset{linestyle=solid}\FAIL{spotgac<13>}\trput{8}
      }
     }
     \psset{linestyle=solid}\FAIL{spotgac<14>}\trput{8}
    }
   }
   \psset{linestyle=solid}\FAIL{spotgac<15>}\trput{6}
  }
  \psset{linestyle=solid}\pstree{\CC{\hyperlink{spotgac<16>}{6}}\trput{6}}{
   \psset{linestyle=solid}\pstree{\CC{\hyperlink{spotgac<17>}{10}}\trput{4}}{
   }
  }
 }
}
}
}
\only<18>{
\scalebox{ 0.40}{\psset{nodesep=0pt,levelsep=36pt}
\pstree{\CC{\hyperlink{spotgac<1>}{4}}}{
 \psset{linestyle=solid}\pstree{\CC{\hyperlink{spotgac<2>}{5}}\trput{2}}{
  \psset{linestyle=solid}\pstree{\CC{\hyperlink{spotgac<3>}{6}}\trput{4}}{
   \psset{linestyle=solid}\pstree{\CC{\hyperlink{spotgac<4>}{7}}\trput{5}}{
    \psset{linestyle=solid}\FAIL{spotgac<5>}\trput{6}
    \psset{linestyle=solid}\pstree{\CC{\hyperlink{spotgac<6>}{12}}\trput{9}}{
     \psset{linestyle=solid}\pstree{\CC{\hyperlink{spotgac<7>}{19}}\trput{6}}{
      \psset{linestyle=solid}\FAIL{spotgac<8>}\trput{6}
      \psset{linestyle=solid}\pstree{\CC{\hyperlink{spotgac<9>}{20}}\trput{9}}{
       \psset{linestyle=solid}\pstree{\CC{\hyperlink{spotgac<10>}{23}}\trput{6}}{
        \psset{linestyle=solid}\FAIL{spotgac<11>}\trput{2}
        \psset{linestyle=solid}\FAIL{spotgac<12>}\trput{8}
       }
       \psset{linestyle=solid}\FAIL{spotgac<13>}\trput{8}
      }
     }
     \psset{linestyle=solid}\FAIL{spotgac<14>}\trput{8}
    }
   }
   \psset{linestyle=solid}\FAIL{spotgac<15>}\trput{6}
  }
  \psset{linestyle=solid}\pstree{\CC{\hyperlink{spotgac<16>}{6}}\trput{6}}{
   \psset{linestyle=solid}\pstree{\CC{\hyperlink{spotgac<17>}{10}}\trput{4}}{
    \psset{linestyle=solid}\FAIL{spotgac<18>}\trput{4}
   }
  }
 }
}
}
}
\only<19>{
\scalebox{ 0.40}{\psset{nodesep=0pt,levelsep=36pt}
\pstree{\CC{\hyperlink{spotgac<1>}{4}}}{
 \psset{linestyle=solid}\pstree{\CC{\hyperlink{spotgac<2>}{5}}\trput{2}}{
  \psset{linestyle=solid}\pstree{\CC{\hyperlink{spotgac<3>}{6}}\trput{4}}{
   \psset{linestyle=solid}\pstree{\CC{\hyperlink{spotgac<4>}{7}}\trput{5}}{
    \psset{linestyle=solid}\FAIL{spotgac<5>}\trput{6}
    \psset{linestyle=solid}\pstree{\CC{\hyperlink{spotgac<6>}{12}}\trput{9}}{
     \psset{linestyle=solid}\pstree{\CC{\hyperlink{spotgac<7>}{19}}\trput{6}}{
      \psset{linestyle=solid}\FAIL{spotgac<8>}\trput{6}
      \psset{linestyle=solid}\pstree{\CC{\hyperlink{spotgac<9>}{20}}\trput{9}}{
       \psset{linestyle=solid}\pstree{\CC{\hyperlink{spotgac<10>}{23}}\trput{6}}{
        \psset{linestyle=solid}\FAIL{spotgac<11>}\trput{2}
        \psset{linestyle=solid}\FAIL{spotgac<12>}\trput{8}
       }
       \psset{linestyle=solid}\FAIL{spotgac<13>}\trput{8}
      }
     }
     \psset{linestyle=solid}\FAIL{spotgac<14>}\trput{8}
    }
   }
   \psset{linestyle=solid}\FAIL{spotgac<15>}\trput{6}
  }
  \psset{linestyle=solid}\pstree{\CC{\hyperlink{spotgac<16>}{6}}\trput{6}}{
   \psset{linestyle=solid}\pstree{\CC{\hyperlink{spotgac<17>}{10}}\trput{4}}{
    \psset{linestyle=solid}\FAIL{spotgac<18>}\trput{4}
    \psset{linestyle=solid}\FAIL{spotgac<19>}\trput{6}
   }
  }
 }
}
}
}
\only<20>{
\scalebox{ 0.40}{\psset{nodesep=0pt,levelsep=36pt}
\pstree{\CC{\hyperlink{spotgac<1>}{4}}}{
 \psset{linestyle=solid}\pstree{\CC{\hyperlink{spotgac<2>}{5}}\trput{2}}{
  \psset{linestyle=solid}\pstree{\CC{\hyperlink{spotgac<3>}{6}}\trput{4}}{
   \psset{linestyle=solid}\pstree{\CC{\hyperlink{spotgac<4>}{7}}\trput{5}}{
    \psset{linestyle=solid}\FAIL{spotgac<5>}\trput{6}
    \psset{linestyle=solid}\pstree{\CC{\hyperlink{spotgac<6>}{12}}\trput{9}}{
     \psset{linestyle=solid}\pstree{\CC{\hyperlink{spotgac<7>}{19}}\trput{6}}{
      \psset{linestyle=solid}\FAIL{spotgac<8>}\trput{6}
      \psset{linestyle=solid}\pstree{\CC{\hyperlink{spotgac<9>}{20}}\trput{9}}{
       \psset{linestyle=solid}\pstree{\CC{\hyperlink{spotgac<10>}{23}}\trput{6}}{
        \psset{linestyle=solid}\FAIL{spotgac<11>}\trput{2}
        \psset{linestyle=solid}\FAIL{spotgac<12>}\trput{8}
       }
       \psset{linestyle=solid}\FAIL{spotgac<13>}\trput{8}
      }
     }
     \psset{linestyle=solid}\FAIL{spotgac<14>}\trput{8}
    }
   }
   \psset{linestyle=solid}\FAIL{spotgac<15>}\trput{6}
  }
  \psset{linestyle=solid}\pstree{\CC{\hyperlink{spotgac<16>}{6}}\trput{6}}{
   \psset{linestyle=solid}\pstree{\CC{\hyperlink{spotgac<17>}{10}}\trput{4}}{
    \psset{linestyle=solid}\FAIL{spotgac<18>}\trput{4}
    \psset{linestyle=solid}\FAIL{spotgac<19>}\trput{6}
    \psset{linestyle=solid}\FAIL{spotgac<20>}\trput{8}
   }
  }
 }
}
}
}
\only<21>{
\scalebox{ 0.40}{\psset{nodesep=0pt,levelsep=36pt}
\pstree{\CC{\hyperlink{spotgac<1>}{4}}}{
 \psset{linestyle=solid}\pstree{\CC{\hyperlink{spotgac<2>}{5}}\trput{2}}{
  \psset{linestyle=solid}\pstree{\CC{\hyperlink{spotgac<3>}{6}}\trput{4}}{
   \psset{linestyle=solid}\pstree{\CC{\hyperlink{spotgac<4>}{7}}\trput{5}}{
    \psset{linestyle=solid}\FAIL{spotgac<5>}\trput{6}
    \psset{linestyle=solid}\pstree{\CC{\hyperlink{spotgac<6>}{12}}\trput{9}}{
     \psset{linestyle=solid}\pstree{\CC{\hyperlink{spotgac<7>}{19}}\trput{6}}{
      \psset{linestyle=solid}\FAIL{spotgac<8>}\trput{6}
      \psset{linestyle=solid}\pstree{\CC{\hyperlink{spotgac<9>}{20}}\trput{9}}{
       \psset{linestyle=solid}\pstree{\CC{\hyperlink{spotgac<10>}{23}}\trput{6}}{
        \psset{linestyle=solid}\FAIL{spotgac<11>}\trput{2}
        \psset{linestyle=solid}\FAIL{spotgac<12>}\trput{8}
       }
       \psset{linestyle=solid}\FAIL{spotgac<13>}\trput{8}
      }
     }
     \psset{linestyle=solid}\FAIL{spotgac<14>}\trput{8}
    }
   }
   \psset{linestyle=solid}\FAIL{spotgac<15>}\trput{6}
  }
  \psset{linestyle=solid}\pstree{\CC{\hyperlink{spotgac<16>}{6}}\trput{6}}{
   \psset{linestyle=solid}\pstree{\CC{\hyperlink{spotgac<17>}{10}}\trput{4}}{
    \psset{linestyle=solid}\FAIL{spotgac<18>}\trput{4}
    \psset{linestyle=solid}\FAIL{spotgac<19>}\trput{6}
    \psset{linestyle=solid}\FAIL{spotgac<20>}\trput{8}
   }
   \psset{linestyle=solid}\pstree{\CC{\hyperlink{spotgac<21>}{7}}\trput{5}}{
   }
  }
 }
}
}
}
\only<22>{
\scalebox{ 0.40}{\psset{nodesep=0pt,levelsep=36pt}
\pstree{\CC{\hyperlink{spotgac<1>}{4}}}{
 \psset{linestyle=solid}\pstree{\CC{\hyperlink{spotgac<2>}{5}}\trput{2}}{
  \psset{linestyle=solid}\pstree{\CC{\hyperlink{spotgac<3>}{6}}\trput{4}}{
   \psset{linestyle=solid}\pstree{\CC{\hyperlink{spotgac<4>}{7}}\trput{5}}{
    \psset{linestyle=solid}\FAIL{spotgac<5>}\trput{6}
    \psset{linestyle=solid}\pstree{\CC{\hyperlink{spotgac<6>}{12}}\trput{9}}{
     \psset{linestyle=solid}\pstree{\CC{\hyperlink{spotgac<7>}{19}}\trput{6}}{
      \psset{linestyle=solid}\FAIL{spotgac<8>}\trput{6}
      \psset{linestyle=solid}\pstree{\CC{\hyperlink{spotgac<9>}{20}}\trput{9}}{
       \psset{linestyle=solid}\pstree{\CC{\hyperlink{spotgac<10>}{23}}\trput{6}}{
        \psset{linestyle=solid}\FAIL{spotgac<11>}\trput{2}
        \psset{linestyle=solid}\FAIL{spotgac<12>}\trput{8}
       }
       \psset{linestyle=solid}\FAIL{spotgac<13>}\trput{8}
      }
     }
     \psset{linestyle=solid}\FAIL{spotgac<14>}\trput{8}
    }
   }
   \psset{linestyle=solid}\FAIL{spotgac<15>}\trput{6}
  }
  \psset{linestyle=solid}\pstree{\CC{\hyperlink{spotgac<16>}{6}}\trput{6}}{
   \psset{linestyle=solid}\pstree{\CC{\hyperlink{spotgac<17>}{10}}\trput{4}}{
    \psset{linestyle=solid}\FAIL{spotgac<18>}\trput{4}
    \psset{linestyle=solid}\FAIL{spotgac<19>}\trput{6}
    \psset{linestyle=solid}\FAIL{spotgac<20>}\trput{8}
   }
   \psset{linestyle=solid}\pstree{\CC{\hyperlink{spotgac<21>}{7}}\trput{5}}{
    \psset{linestyle=solid}\pstree{\CC{\hyperlink{spotgac<22>}{10}}\trput{4}}{
    }
   }
  }
 }
}
}
}
\only<23>{
\scalebox{ 0.40}{\psset{nodesep=0pt,levelsep=36pt}
\pstree{\CC{\hyperlink{spotgac<1>}{4}}}{
 \psset{linestyle=solid}\pstree{\CC{\hyperlink{spotgac<2>}{5}}\trput{2}}{
  \psset{linestyle=solid}\pstree{\CC{\hyperlink{spotgac<3>}{6}}\trput{4}}{
   \psset{linestyle=solid}\pstree{\CC{\hyperlink{spotgac<4>}{7}}\trput{5}}{
    \psset{linestyle=solid}\FAIL{spotgac<5>}\trput{6}
    \psset{linestyle=solid}\pstree{\CC{\hyperlink{spotgac<6>}{12}}\trput{9}}{
     \psset{linestyle=solid}\pstree{\CC{\hyperlink{spotgac<7>}{19}}\trput{6}}{
      \psset{linestyle=solid}\FAIL{spotgac<8>}\trput{6}
      \psset{linestyle=solid}\pstree{\CC{\hyperlink{spotgac<9>}{20}}\trput{9}}{
       \psset{linestyle=solid}\pstree{\CC{\hyperlink{spotgac<10>}{23}}\trput{6}}{
        \psset{linestyle=solid}\FAIL{spotgac<11>}\trput{2}
        \psset{linestyle=solid}\FAIL{spotgac<12>}\trput{8}
       }
       \psset{linestyle=solid}\FAIL{spotgac<13>}\trput{8}
      }
     }
     \psset{linestyle=solid}\FAIL{spotgac<14>}\trput{8}
    }
   }
   \psset{linestyle=solid}\FAIL{spotgac<15>}\trput{6}
  }
  \psset{linestyle=solid}\pstree{\CC{\hyperlink{spotgac<16>}{6}}\trput{6}}{
   \psset{linestyle=solid}\pstree{\CC{\hyperlink{spotgac<17>}{10}}\trput{4}}{
    \psset{linestyle=solid}\FAIL{spotgac<18>}\trput{4}
    \psset{linestyle=solid}\FAIL{spotgac<19>}\trput{6}
    \psset{linestyle=solid}\FAIL{spotgac<20>}\trput{8}
   }
   \psset{linestyle=solid}\pstree{\CC{\hyperlink{spotgac<21>}{7}}\trput{5}}{
    \psset{linestyle=solid}\pstree{\CC{\hyperlink{spotgac<22>}{10}}\trput{4}}{
     \psset{linestyle=solid}\FAIL{spotgac<23>}\trput{5}
    }
   }
  }
 }
}
}
}
\only<24>{
\scalebox{ 0.40}{\psset{nodesep=0pt,levelsep=36pt}
\pstree{\CC{\hyperlink{spotgac<1>}{4}}}{
 \psset{linestyle=solid}\pstree{\CC{\hyperlink{spotgac<2>}{5}}\trput{2}}{
  \psset{linestyle=solid}\pstree{\CC{\hyperlink{spotgac<3>}{6}}\trput{4}}{
   \psset{linestyle=solid}\pstree{\CC{\hyperlink{spotgac<4>}{7}}\trput{5}}{
    \psset{linestyle=solid}\FAIL{spotgac<5>}\trput{6}
    \psset{linestyle=solid}\pstree{\CC{\hyperlink{spotgac<6>}{12}}\trput{9}}{
     \psset{linestyle=solid}\pstree{\CC{\hyperlink{spotgac<7>}{19}}\trput{6}}{
      \psset{linestyle=solid}\FAIL{spotgac<8>}\trput{6}
      \psset{linestyle=solid}\pstree{\CC{\hyperlink{spotgac<9>}{20}}\trput{9}}{
       \psset{linestyle=solid}\pstree{\CC{\hyperlink{spotgac<10>}{23}}\trput{6}}{
        \psset{linestyle=solid}\FAIL{spotgac<11>}\trput{2}
        \psset{linestyle=solid}\FAIL{spotgac<12>}\trput{8}
       }
       \psset{linestyle=solid}\FAIL{spotgac<13>}\trput{8}
      }
     }
     \psset{linestyle=solid}\FAIL{spotgac<14>}\trput{8}
    }
   }
   \psset{linestyle=solid}\FAIL{spotgac<15>}\trput{6}
  }
  \psset{linestyle=solid}\pstree{\CC{\hyperlink{spotgac<16>}{6}}\trput{6}}{
   \psset{linestyle=solid}\pstree{\CC{\hyperlink{spotgac<17>}{10}}\trput{4}}{
    \psset{linestyle=solid}\FAIL{spotgac<18>}\trput{4}
    \psset{linestyle=solid}\FAIL{spotgac<19>}\trput{6}
    \psset{linestyle=solid}\FAIL{spotgac<20>}\trput{8}
   }
   \psset{linestyle=solid}\pstree{\CC{\hyperlink{spotgac<21>}{7}}\trput{5}}{
    \psset{linestyle=solid}\pstree{\CC{\hyperlink{spotgac<22>}{10}}\trput{4}}{
     \psset{linestyle=solid}\FAIL{spotgac<23>}\trput{5}
     \psset{linestyle=solid}\FAIL{spotgac<24>}\trput{6}
    }
   }
  }
 }
}
}
}
\only<25>{
\scalebox{ 0.40}{\psset{nodesep=0pt,levelsep=36pt}
\pstree{\CC{\hyperlink{spotgac<1>}{4}}}{
 \psset{linestyle=solid}\pstree{\CC{\hyperlink{spotgac<2>}{5}}\trput{2}}{
  \psset{linestyle=solid}\pstree{\CC{\hyperlink{spotgac<3>}{6}}\trput{4}}{
   \psset{linestyle=solid}\pstree{\CC{\hyperlink{spotgac<4>}{7}}\trput{5}}{
    \psset{linestyle=solid}\FAIL{spotgac<5>}\trput{6}
    \psset{linestyle=solid}\pstree{\CC{\hyperlink{spotgac<6>}{12}}\trput{9}}{
     \psset{linestyle=solid}\pstree{\CC{\hyperlink{spotgac<7>}{19}}\trput{6}}{
      \psset{linestyle=solid}\FAIL{spotgac<8>}\trput{6}
      \psset{linestyle=solid}\pstree{\CC{\hyperlink{spotgac<9>}{20}}\trput{9}}{
       \psset{linestyle=solid}\pstree{\CC{\hyperlink{spotgac<10>}{23}}\trput{6}}{
        \psset{linestyle=solid}\FAIL{spotgac<11>}\trput{2}
        \psset{linestyle=solid}\FAIL{spotgac<12>}\trput{8}
       }
       \psset{linestyle=solid}\FAIL{spotgac<13>}\trput{8}
      }
     }
     \psset{linestyle=solid}\FAIL{spotgac<14>}\trput{8}
    }
   }
   \psset{linestyle=solid}\FAIL{spotgac<15>}\trput{6}
  }
  \psset{linestyle=solid}\pstree{\CC{\hyperlink{spotgac<16>}{6}}\trput{6}}{
   \psset{linestyle=solid}\pstree{\CC{\hyperlink{spotgac<17>}{10}}\trput{4}}{
    \psset{linestyle=solid}\FAIL{spotgac<18>}\trput{4}
    \psset{linestyle=solid}\FAIL{spotgac<19>}\trput{6}
    \psset{linestyle=solid}\FAIL{spotgac<20>}\trput{8}
   }
   \psset{linestyle=solid}\pstree{\CC{\hyperlink{spotgac<21>}{7}}\trput{5}}{
    \psset{linestyle=solid}\pstree{\CC{\hyperlink{spotgac<22>}{10}}\trput{4}}{
     \psset{linestyle=solid}\FAIL{spotgac<23>}\trput{5}
     \psset{linestyle=solid}\FAIL{spotgac<24>}\trput{6}
     \psset{linestyle=solid}\FAIL{spotgac<25>}\trput{8}
    }
   }
  }
 }
}
}
}
\only<26>{
\scalebox{ 0.40}{\psset{nodesep=0pt,levelsep=36pt}
\pstree{\CC{\hyperlink{spotgac<1>}{4}}}{
 \psset{linestyle=solid}\pstree{\CC{\hyperlink{spotgac<2>}{5}}\trput{2}}{
  \psset{linestyle=solid}\pstree{\CC{\hyperlink{spotgac<3>}{6}}\trput{4}}{
   \psset{linestyle=solid}\pstree{\CC{\hyperlink{spotgac<4>}{7}}\trput{5}}{
    \psset{linestyle=solid}\FAIL{spotgac<5>}\trput{6}
    \psset{linestyle=solid}\pstree{\CC{\hyperlink{spotgac<6>}{12}}\trput{9}}{
     \psset{linestyle=solid}\pstree{\CC{\hyperlink{spotgac<7>}{19}}\trput{6}}{
      \psset{linestyle=solid}\FAIL{spotgac<8>}\trput{6}
      \psset{linestyle=solid}\pstree{\CC{\hyperlink{spotgac<9>}{20}}\trput{9}}{
       \psset{linestyle=solid}\pstree{\CC{\hyperlink{spotgac<10>}{23}}\trput{6}}{
        \psset{linestyle=solid}\FAIL{spotgac<11>}\trput{2}
        \psset{linestyle=solid}\FAIL{spotgac<12>}\trput{8}
       }
       \psset{linestyle=solid}\FAIL{spotgac<13>}\trput{8}
      }
     }
     \psset{linestyle=solid}\FAIL{spotgac<14>}\trput{8}
    }
   }
   \psset{linestyle=solid}\FAIL{spotgac<15>}\trput{6}
  }
  \psset{linestyle=solid}\pstree{\CC{\hyperlink{spotgac<16>}{6}}\trput{6}}{
   \psset{linestyle=solid}\pstree{\CC{\hyperlink{spotgac<17>}{10}}\trput{4}}{
    \psset{linestyle=solid}\FAIL{spotgac<18>}\trput{4}
    \psset{linestyle=solid}\FAIL{spotgac<19>}\trput{6}
    \psset{linestyle=solid}\FAIL{spotgac<20>}\trput{8}
   }
   \psset{linestyle=solid}\pstree{\CC{\hyperlink{spotgac<21>}{7}}\trput{5}}{
    \psset{linestyle=solid}\pstree{\CC{\hyperlink{spotgac<22>}{10}}\trput{4}}{
     \psset{linestyle=solid}\FAIL{spotgac<23>}\trput{5}
     \psset{linestyle=solid}\FAIL{spotgac<24>}\trput{6}
     \psset{linestyle=solid}\FAIL{spotgac<25>}\trput{8}
    }
    \psset{linestyle=solid}\pstree{\CC{\hyperlink{spotgac<26>}{10}}\trput{9}}{
    }
   }
  }
 }
}
}
}
\only<27>{
\scalebox{ 0.40}{\psset{nodesep=0pt,levelsep=36pt}
\pstree{\CC{\hyperlink{spotgac<1>}{4}}}{
 \psset{linestyle=solid}\pstree{\CC{\hyperlink{spotgac<2>}{5}}\trput{2}}{
  \psset{linestyle=solid}\pstree{\CC{\hyperlink{spotgac<3>}{6}}\trput{4}}{
   \psset{linestyle=solid}\pstree{\CC{\hyperlink{spotgac<4>}{7}}\trput{5}}{
    \psset{linestyle=solid}\FAIL{spotgac<5>}\trput{6}
    \psset{linestyle=solid}\pstree{\CC{\hyperlink{spotgac<6>}{12}}\trput{9}}{
     \psset{linestyle=solid}\pstree{\CC{\hyperlink{spotgac<7>}{19}}\trput{6}}{
      \psset{linestyle=solid}\FAIL{spotgac<8>}\trput{6}
      \psset{linestyle=solid}\pstree{\CC{\hyperlink{spotgac<9>}{20}}\trput{9}}{
       \psset{linestyle=solid}\pstree{\CC{\hyperlink{spotgac<10>}{23}}\trput{6}}{
        \psset{linestyle=solid}\FAIL{spotgac<11>}\trput{2}
        \psset{linestyle=solid}\FAIL{spotgac<12>}\trput{8}
       }
       \psset{linestyle=solid}\FAIL{spotgac<13>}\trput{8}
      }
     }
     \psset{linestyle=solid}\FAIL{spotgac<14>}\trput{8}
    }
   }
   \psset{linestyle=solid}\FAIL{spotgac<15>}\trput{6}
  }
  \psset{linestyle=solid}\pstree{\CC{\hyperlink{spotgac<16>}{6}}\trput{6}}{
   \psset{linestyle=solid}\pstree{\CC{\hyperlink{spotgac<17>}{10}}\trput{4}}{
    \psset{linestyle=solid}\FAIL{spotgac<18>}\trput{4}
    \psset{linestyle=solid}\FAIL{spotgac<19>}\trput{6}
    \psset{linestyle=solid}\FAIL{spotgac<20>}\trput{8}
   }
   \psset{linestyle=solid}\pstree{\CC{\hyperlink{spotgac<21>}{7}}\trput{5}}{
    \psset{linestyle=solid}\pstree{\CC{\hyperlink{spotgac<22>}{10}}\trput{4}}{
     \psset{linestyle=solid}\FAIL{spotgac<23>}\trput{5}
     \psset{linestyle=solid}\FAIL{spotgac<24>}\trput{6}
     \psset{linestyle=solid}\FAIL{spotgac<25>}\trput{8}
    }
    \psset{linestyle=solid}\pstree{\CC{\hyperlink{spotgac<26>}{10}}\trput{9}}{
     \psset{linestyle=solid}\pstree{\CC{\hyperlink{spotgac<27>}{11}}\trput{5}}{
     }
    }
   }
  }
 }
}
}
}
\only<28>{
\scalebox{ 0.40}{\psset{nodesep=0pt,levelsep=36pt}
\pstree{\CC{\hyperlink{spotgac<1>}{4}}}{
 \psset{linestyle=solid}\pstree{\CC{\hyperlink{spotgac<2>}{5}}\trput{2}}{
  \psset{linestyle=solid}\pstree{\CC{\hyperlink{spotgac<3>}{6}}\trput{4}}{
   \psset{linestyle=solid}\pstree{\CC{\hyperlink{spotgac<4>}{7}}\trput{5}}{
    \psset{linestyle=solid}\FAIL{spotgac<5>}\trput{6}
    \psset{linestyle=solid}\pstree{\CC{\hyperlink{spotgac<6>}{12}}\trput{9}}{
     \psset{linestyle=solid}\pstree{\CC{\hyperlink{spotgac<7>}{19}}\trput{6}}{
      \psset{linestyle=solid}\FAIL{spotgac<8>}\trput{6}
      \psset{linestyle=solid}\pstree{\CC{\hyperlink{spotgac<9>}{20}}\trput{9}}{
       \psset{linestyle=solid}\pstree{\CC{\hyperlink{spotgac<10>}{23}}\trput{6}}{
        \psset{linestyle=solid}\FAIL{spotgac<11>}\trput{2}
        \psset{linestyle=solid}\FAIL{spotgac<12>}\trput{8}
       }
       \psset{linestyle=solid}\FAIL{spotgac<13>}\trput{8}
      }
     }
     \psset{linestyle=solid}\FAIL{spotgac<14>}\trput{8}
    }
   }
   \psset{linestyle=solid}\FAIL{spotgac<15>}\trput{6}
  }
  \psset{linestyle=solid}\pstree{\CC{\hyperlink{spotgac<16>}{6}}\trput{6}}{
   \psset{linestyle=solid}\pstree{\CC{\hyperlink{spotgac<17>}{10}}\trput{4}}{
    \psset{linestyle=solid}\FAIL{spotgac<18>}\trput{4}
    \psset{linestyle=solid}\FAIL{spotgac<19>}\trput{6}
    \psset{linestyle=solid}\FAIL{spotgac<20>}\trput{8}
   }
   \psset{linestyle=solid}\pstree{\CC{\hyperlink{spotgac<21>}{7}}\trput{5}}{
    \psset{linestyle=solid}\pstree{\CC{\hyperlink{spotgac<22>}{10}}\trput{4}}{
     \psset{linestyle=solid}\FAIL{spotgac<23>}\trput{5}
     \psset{linestyle=solid}\FAIL{spotgac<24>}\trput{6}
     \psset{linestyle=solid}\FAIL{spotgac<25>}\trput{8}
    }
    \psset{linestyle=solid}\pstree{\CC{\hyperlink{spotgac<26>}{10}}\trput{9}}{
     \psset{linestyle=solid}\pstree{\CC{\hyperlink{spotgac<27>}{11}}\trput{5}}{
      \psset{linestyle=solid}\FAIL{spotgac<28>}\trput{6}
     }
    }
   }
  }
 }
}
}
}
\only<29>{
\scalebox{ 0.40}{\psset{nodesep=0pt,levelsep=36pt}
\pstree{\CC{\hyperlink{spotgac<1>}{4}}}{
 \psset{linestyle=solid}\pstree{\CC{\hyperlink{spotgac<2>}{5}}\trput{2}}{
  \psset{linestyle=solid}\pstree{\CC{\hyperlink{spotgac<3>}{6}}\trput{4}}{
   \psset{linestyle=solid}\pstree{\CC{\hyperlink{spotgac<4>}{7}}\trput{5}}{
    \psset{linestyle=solid}\FAIL{spotgac<5>}\trput{6}
    \psset{linestyle=solid}\pstree{\CC{\hyperlink{spotgac<6>}{12}}\trput{9}}{
     \psset{linestyle=solid}\pstree{\CC{\hyperlink{spotgac<7>}{19}}\trput{6}}{
      \psset{linestyle=solid}\FAIL{spotgac<8>}\trput{6}
      \psset{linestyle=solid}\pstree{\CC{\hyperlink{spotgac<9>}{20}}\trput{9}}{
       \psset{linestyle=solid}\pstree{\CC{\hyperlink{spotgac<10>}{23}}\trput{6}}{
        \psset{linestyle=solid}\FAIL{spotgac<11>}\trput{2}
        \psset{linestyle=solid}\FAIL{spotgac<12>}\trput{8}
       }
       \psset{linestyle=solid}\FAIL{spotgac<13>}\trput{8}
      }
     }
     \psset{linestyle=solid}\FAIL{spotgac<14>}\trput{8}
    }
   }
   \psset{linestyle=solid}\FAIL{spotgac<15>}\trput{6}
  }
  \psset{linestyle=solid}\pstree{\CC{\hyperlink{spotgac<16>}{6}}\trput{6}}{
   \psset{linestyle=solid}\pstree{\CC{\hyperlink{spotgac<17>}{10}}\trput{4}}{
    \psset{linestyle=solid}\FAIL{spotgac<18>}\trput{4}
    \psset{linestyle=solid}\FAIL{spotgac<19>}\trput{6}
    \psset{linestyle=solid}\FAIL{spotgac<20>}\trput{8}
   }
   \psset{linestyle=solid}\pstree{\CC{\hyperlink{spotgac<21>}{7}}\trput{5}}{
    \psset{linestyle=solid}\pstree{\CC{\hyperlink{spotgac<22>}{10}}\trput{4}}{
     \psset{linestyle=solid}\FAIL{spotgac<23>}\trput{5}
     \psset{linestyle=solid}\FAIL{spotgac<24>}\trput{6}
     \psset{linestyle=solid}\FAIL{spotgac<25>}\trput{8}
    }
    \psset{linestyle=solid}\pstree{\CC{\hyperlink{spotgac<26>}{10}}\trput{9}}{
     \psset{linestyle=solid}\pstree{\CC{\hyperlink{spotgac<27>}{11}}\trput{5}}{
      \psset{linestyle=solid}\FAIL{spotgac<28>}\trput{6}
      \psset{linestyle=solid}\FAIL{spotgac<29>}\trput{8}
     }
    }
   }
  }
 }
}
}
}
\only<30>{
\scalebox{ 0.40}{\psset{nodesep=0pt,levelsep=36pt}
\pstree{\CC{\hyperlink{spotgac<1>}{4}}}{
 \psset{linestyle=solid}\pstree{\CC{\hyperlink{spotgac<2>}{5}}\trput{2}}{
  \psset{linestyle=solid}\pstree{\CC{\hyperlink{spotgac<3>}{6}}\trput{4}}{
   \psset{linestyle=solid}\pstree{\CC{\hyperlink{spotgac<4>}{7}}\trput{5}}{
    \psset{linestyle=solid}\FAIL{spotgac<5>}\trput{6}
    \psset{linestyle=solid}\pstree{\CC{\hyperlink{spotgac<6>}{12}}\trput{9}}{
     \psset{linestyle=solid}\pstree{\CC{\hyperlink{spotgac<7>}{19}}\trput{6}}{
      \psset{linestyle=solid}\FAIL{spotgac<8>}\trput{6}
      \psset{linestyle=solid}\pstree{\CC{\hyperlink{spotgac<9>}{20}}\trput{9}}{
       \psset{linestyle=solid}\pstree{\CC{\hyperlink{spotgac<10>}{23}}\trput{6}}{
        \psset{linestyle=solid}\FAIL{spotgac<11>}\trput{2}
        \psset{linestyle=solid}\FAIL{spotgac<12>}\trput{8}
       }
       \psset{linestyle=solid}\FAIL{spotgac<13>}\trput{8}
      }
     }
     \psset{linestyle=solid}\FAIL{spotgac<14>}\trput{8}
    }
   }
   \psset{linestyle=solid}\FAIL{spotgac<15>}\trput{6}
  }
  \psset{linestyle=solid}\pstree{\CC{\hyperlink{spotgac<16>}{6}}\trput{6}}{
   \psset{linestyle=solid}\pstree{\CC{\hyperlink{spotgac<17>}{10}}\trput{4}}{
    \psset{linestyle=solid}\FAIL{spotgac<18>}\trput{4}
    \psset{linestyle=solid}\FAIL{spotgac<19>}\trput{6}
    \psset{linestyle=solid}\FAIL{spotgac<20>}\trput{8}
   }
   \psset{linestyle=solid}\pstree{\CC{\hyperlink{spotgac<21>}{7}}\trput{5}}{
    \psset{linestyle=solid}\pstree{\CC{\hyperlink{spotgac<22>}{10}}\trput{4}}{
     \psset{linestyle=solid}\FAIL{spotgac<23>}\trput{5}
     \psset{linestyle=solid}\FAIL{spotgac<24>}\trput{6}
     \psset{linestyle=solid}\FAIL{spotgac<25>}\trput{8}
    }
    \psset{linestyle=solid}\pstree{\CC{\hyperlink{spotgac<26>}{10}}\trput{9}}{
     \psset{linestyle=solid}\pstree{\CC{\hyperlink{spotgac<27>}{11}}\trput{5}}{
      \psset{linestyle=solid}\FAIL{spotgac<28>}\trput{6}
      \psset{linestyle=solid}\FAIL{spotgac<29>}\trput{8}
     }
     \psset{linestyle=solid}\FAIL{spotgac<30>}\trput{6}
    }
   }
  }
 }
}
}
}
\only<31>{
\scalebox{ 0.40}{\psset{nodesep=0pt,levelsep=36pt}
\pstree{\CC{\hyperlink{spotgac<1>}{4}}}{
 \psset{linestyle=solid}\pstree{\CC{\hyperlink{spotgac<2>}{5}}\trput{2}}{
  \psset{linestyle=solid}\pstree{\CC{\hyperlink{spotgac<3>}{6}}\trput{4}}{
   \psset{linestyle=solid}\pstree{\CC{\hyperlink{spotgac<4>}{7}}\trput{5}}{
    \psset{linestyle=solid}\FAIL{spotgac<5>}\trput{6}
    \psset{linestyle=solid}\pstree{\CC{\hyperlink{spotgac<6>}{12}}\trput{9}}{
     \psset{linestyle=solid}\pstree{\CC{\hyperlink{spotgac<7>}{19}}\trput{6}}{
      \psset{linestyle=solid}\FAIL{spotgac<8>}\trput{6}
      \psset{linestyle=solid}\pstree{\CC{\hyperlink{spotgac<9>}{20}}\trput{9}}{
       \psset{linestyle=solid}\pstree{\CC{\hyperlink{spotgac<10>}{23}}\trput{6}}{
        \psset{linestyle=solid}\FAIL{spotgac<11>}\trput{2}
        \psset{linestyle=solid}\FAIL{spotgac<12>}\trput{8}
       }
       \psset{linestyle=solid}\FAIL{spotgac<13>}\trput{8}
      }
     }
     \psset{linestyle=solid}\FAIL{spotgac<14>}\trput{8}
    }
   }
   \psset{linestyle=solid}\FAIL{spotgac<15>}\trput{6}
  }
  \psset{linestyle=solid}\pstree{\CC{\hyperlink{spotgac<16>}{6}}\trput{6}}{
   \psset{linestyle=solid}\pstree{\CC{\hyperlink{spotgac<17>}{10}}\trput{4}}{
    \psset{linestyle=solid}\FAIL{spotgac<18>}\trput{4}
    \psset{linestyle=solid}\FAIL{spotgac<19>}\trput{6}
    \psset{linestyle=solid}\FAIL{spotgac<20>}\trput{8}
   }
   \psset{linestyle=solid}\pstree{\CC{\hyperlink{spotgac<21>}{7}}\trput{5}}{
    \psset{linestyle=solid}\pstree{\CC{\hyperlink{spotgac<22>}{10}}\trput{4}}{
     \psset{linestyle=solid}\FAIL{spotgac<23>}\trput{5}
     \psset{linestyle=solid}\FAIL{spotgac<24>}\trput{6}
     \psset{linestyle=solid}\FAIL{spotgac<25>}\trput{8}
    }
    \psset{linestyle=solid}\pstree{\CC{\hyperlink{spotgac<26>}{10}}\trput{9}}{
     \psset{linestyle=solid}\pstree{\CC{\hyperlink{spotgac<27>}{11}}\trput{5}}{
      \psset{linestyle=solid}\FAIL{spotgac<28>}\trput{6}
      \psset{linestyle=solid}\FAIL{spotgac<29>}\trput{8}
     }
     \psset{linestyle=solid}\FAIL{spotgac<30>}\trput{6}
     \psset{linestyle=solid}\FAIL{spotgac<31>}\trput{8}
    }
   }
  }
 }
}
}
}
\only<32>{
\scalebox{ 0.40}{\psset{nodesep=0pt,levelsep=36pt}
\pstree{\CC{\hyperlink{spotgac<1>}{4}}}{
 \psset{linestyle=solid}\pstree{\CC{\hyperlink{spotgac<2>}{5}}\trput{2}}{
  \psset{linestyle=solid}\pstree{\CC{\hyperlink{spotgac<3>}{6}}\trput{4}}{
   \psset{linestyle=solid}\pstree{\CC{\hyperlink{spotgac<4>}{7}}\trput{5}}{
    \psset{linestyle=solid}\FAIL{spotgac<5>}\trput{6}
    \psset{linestyle=solid}\pstree{\CC{\hyperlink{spotgac<6>}{12}}\trput{9}}{
     \psset{linestyle=solid}\pstree{\CC{\hyperlink{spotgac<7>}{19}}\trput{6}}{
      \psset{linestyle=solid}\FAIL{spotgac<8>}\trput{6}
      \psset{linestyle=solid}\pstree{\CC{\hyperlink{spotgac<9>}{20}}\trput{9}}{
       \psset{linestyle=solid}\pstree{\CC{\hyperlink{spotgac<10>}{23}}\trput{6}}{
        \psset{linestyle=solid}\FAIL{spotgac<11>}\trput{2}
        \psset{linestyle=solid}\FAIL{spotgac<12>}\trput{8}
       }
       \psset{linestyle=solid}\FAIL{spotgac<13>}\trput{8}
      }
     }
     \psset{linestyle=solid}\FAIL{spotgac<14>}\trput{8}
    }
   }
   \psset{linestyle=solid}\FAIL{spotgac<15>}\trput{6}
  }
  \psset{linestyle=solid}\pstree{\CC{\hyperlink{spotgac<16>}{6}}\trput{6}}{
   \psset{linestyle=solid}\pstree{\CC{\hyperlink{spotgac<17>}{10}}\trput{4}}{
    \psset{linestyle=solid}\FAIL{spotgac<18>}\trput{4}
    \psset{linestyle=solid}\FAIL{spotgac<19>}\trput{6}
    \psset{linestyle=solid}\FAIL{spotgac<20>}\trput{8}
   }
   \psset{linestyle=solid}\pstree{\CC{\hyperlink{spotgac<21>}{7}}\trput{5}}{
    \psset{linestyle=solid}\pstree{\CC{\hyperlink{spotgac<22>}{10}}\trput{4}}{
     \psset{linestyle=solid}\FAIL{spotgac<23>}\trput{5}
     \psset{linestyle=solid}\FAIL{spotgac<24>}\trput{6}
     \psset{linestyle=solid}\FAIL{spotgac<25>}\trput{8}
    }
    \psset{linestyle=solid}\pstree{\CC{\hyperlink{spotgac<26>}{10}}\trput{9}}{
     \psset{linestyle=solid}\pstree{\CC{\hyperlink{spotgac<27>}{11}}\trput{5}}{
      \psset{linestyle=solid}\FAIL{spotgac<28>}\trput{6}
      \psset{linestyle=solid}\FAIL{spotgac<29>}\trput{8}
     }
     \psset{linestyle=solid}\FAIL{spotgac<30>}\trput{6}
     \psset{linestyle=solid}\FAIL{spotgac<31>}\trput{8}
    }
   }
  }
 }
 \psset{linestyle=solid}\pstree{\CC{\hyperlink{spotgac<32>}{5}}\trput{4}}{
 }
}
}
}
\only<33>{
\scalebox{ 0.40}{\psset{nodesep=0pt,levelsep=36pt}
\pstree{\CC{\hyperlink{spotgac<1>}{4}}}{
 \psset{linestyle=solid}\pstree{\CC{\hyperlink{spotgac<2>}{5}}\trput{2}}{
  \psset{linestyle=solid}\pstree{\CC{\hyperlink{spotgac<3>}{6}}\trput{4}}{
   \psset{linestyle=solid}\pstree{\CC{\hyperlink{spotgac<4>}{7}}\trput{5}}{
    \psset{linestyle=solid}\FAIL{spotgac<5>}\trput{6}
    \psset{linestyle=solid}\pstree{\CC{\hyperlink{spotgac<6>}{12}}\trput{9}}{
     \psset{linestyle=solid}\pstree{\CC{\hyperlink{spotgac<7>}{19}}\trput{6}}{
      \psset{linestyle=solid}\FAIL{spotgac<8>}\trput{6}
      \psset{linestyle=solid}\pstree{\CC{\hyperlink{spotgac<9>}{20}}\trput{9}}{
       \psset{linestyle=solid}\pstree{\CC{\hyperlink{spotgac<10>}{23}}\trput{6}}{
        \psset{linestyle=solid}\FAIL{spotgac<11>}\trput{2}
        \psset{linestyle=solid}\FAIL{spotgac<12>}\trput{8}
       }
       \psset{linestyle=solid}\FAIL{spotgac<13>}\trput{8}
      }
     }
     \psset{linestyle=solid}\FAIL{spotgac<14>}\trput{8}
    }
   }
   \psset{linestyle=solid}\FAIL{spotgac<15>}\trput{6}
  }
  \psset{linestyle=solid}\pstree{\CC{\hyperlink{spotgac<16>}{6}}\trput{6}}{
   \psset{linestyle=solid}\pstree{\CC{\hyperlink{spotgac<17>}{10}}\trput{4}}{
    \psset{linestyle=solid}\FAIL{spotgac<18>}\trput{4}
    \psset{linestyle=solid}\FAIL{spotgac<19>}\trput{6}
    \psset{linestyle=solid}\FAIL{spotgac<20>}\trput{8}
   }
   \psset{linestyle=solid}\pstree{\CC{\hyperlink{spotgac<21>}{7}}\trput{5}}{
    \psset{linestyle=solid}\pstree{\CC{\hyperlink{spotgac<22>}{10}}\trput{4}}{
     \psset{linestyle=solid}\FAIL{spotgac<23>}\trput{5}
     \psset{linestyle=solid}\FAIL{spotgac<24>}\trput{6}
     \psset{linestyle=solid}\FAIL{spotgac<25>}\trput{8}
    }
    \psset{linestyle=solid}\pstree{\CC{\hyperlink{spotgac<26>}{10}}\trput{9}}{
     \psset{linestyle=solid}\pstree{\CC{\hyperlink{spotgac<27>}{11}}\trput{5}}{
      \psset{linestyle=solid}\FAIL{spotgac<28>}\trput{6}
      \psset{linestyle=solid}\FAIL{spotgac<29>}\trput{8}
     }
     \psset{linestyle=solid}\FAIL{spotgac<30>}\trput{6}
     \psset{linestyle=solid}\FAIL{spotgac<31>}\trput{8}
    }
   }
  }
 }
 \psset{linestyle=solid}\pstree{\CC{\hyperlink{spotgac<32>}{5}}\trput{4}}{
  \psset{linestyle=solid}\pstree{\CC{\hyperlink{spotgac<33>}{6}}\trput{2}}{
  }
 }
}
}
}
\only<34>{
\scalebox{ 0.40}{\psset{nodesep=0pt,levelsep=36pt}
\pstree{\CC{\hyperlink{spotgac<1>}{4}}}{
 \psset{linestyle=solid}\pstree{\CC{\hyperlink{spotgac<2>}{5}}\trput{2}}{
  \psset{linestyle=solid}\pstree{\CC{\hyperlink{spotgac<3>}{6}}\trput{4}}{
   \psset{linestyle=solid}\pstree{\CC{\hyperlink{spotgac<4>}{7}}\trput{5}}{
    \psset{linestyle=solid}\FAIL{spotgac<5>}\trput{6}
    \psset{linestyle=solid}\pstree{\CC{\hyperlink{spotgac<6>}{12}}\trput{9}}{
     \psset{linestyle=solid}\pstree{\CC{\hyperlink{spotgac<7>}{19}}\trput{6}}{
      \psset{linestyle=solid}\FAIL{spotgac<8>}\trput{6}
      \psset{linestyle=solid}\pstree{\CC{\hyperlink{spotgac<9>}{20}}\trput{9}}{
       \psset{linestyle=solid}\pstree{\CC{\hyperlink{spotgac<10>}{23}}\trput{6}}{
        \psset{linestyle=solid}\FAIL{spotgac<11>}\trput{2}
        \psset{linestyle=solid}\FAIL{spotgac<12>}\trput{8}
       }
       \psset{linestyle=solid}\FAIL{spotgac<13>}\trput{8}
      }
     }
     \psset{linestyle=solid}\FAIL{spotgac<14>}\trput{8}
    }
   }
   \psset{linestyle=solid}\FAIL{spotgac<15>}\trput{6}
  }
  \psset{linestyle=solid}\pstree{\CC{\hyperlink{spotgac<16>}{6}}\trput{6}}{
   \psset{linestyle=solid}\pstree{\CC{\hyperlink{spotgac<17>}{10}}\trput{4}}{
    \psset{linestyle=solid}\FAIL{spotgac<18>}\trput{4}
    \psset{linestyle=solid}\FAIL{spotgac<19>}\trput{6}
    \psset{linestyle=solid}\FAIL{spotgac<20>}\trput{8}
   }
   \psset{linestyle=solid}\pstree{\CC{\hyperlink{spotgac<21>}{7}}\trput{5}}{
    \psset{linestyle=solid}\pstree{\CC{\hyperlink{spotgac<22>}{10}}\trput{4}}{
     \psset{linestyle=solid}\FAIL{spotgac<23>}\trput{5}
     \psset{linestyle=solid}\FAIL{spotgac<24>}\trput{6}
     \psset{linestyle=solid}\FAIL{spotgac<25>}\trput{8}
    }
    \psset{linestyle=solid}\pstree{\CC{\hyperlink{spotgac<26>}{10}}\trput{9}}{
     \psset{linestyle=solid}\pstree{\CC{\hyperlink{spotgac<27>}{11}}\trput{5}}{
      \psset{linestyle=solid}\FAIL{spotgac<28>}\trput{6}
      \psset{linestyle=solid}\FAIL{spotgac<29>}\trput{8}
     }
     \psset{linestyle=solid}\FAIL{spotgac<30>}\trput{6}
     \psset{linestyle=solid}\FAIL{spotgac<31>}\trput{8}
    }
   }
  }
 }
 \psset{linestyle=solid}\pstree{\CC{\hyperlink{spotgac<32>}{5}}\trput{4}}{
  \psset{linestyle=solid}\pstree{\CC{\hyperlink{spotgac<33>}{6}}\trput{2}}{
   \psset{linestyle=solid}\FAIL{spotgac<34>}\trput{5}
  }
 }
}
}
}
\only<35>{
\scalebox{ 0.40}{\psset{nodesep=0pt,levelsep=36pt}
\pstree{\CC{\hyperlink{spotgac<1>}{4}}}{
 \psset{linestyle=solid}\pstree{\CC{\hyperlink{spotgac<2>}{5}}\trput{2}}{
  \psset{linestyle=solid}\pstree{\CC{\hyperlink{spotgac<3>}{6}}\trput{4}}{
   \psset{linestyle=solid}\pstree{\CC{\hyperlink{spotgac<4>}{7}}\trput{5}}{
    \psset{linestyle=solid}\FAIL{spotgac<5>}\trput{6}
    \psset{linestyle=solid}\pstree{\CC{\hyperlink{spotgac<6>}{12}}\trput{9}}{
     \psset{linestyle=solid}\pstree{\CC{\hyperlink{spotgac<7>}{19}}\trput{6}}{
      \psset{linestyle=solid}\FAIL{spotgac<8>}\trput{6}
      \psset{linestyle=solid}\pstree{\CC{\hyperlink{spotgac<9>}{20}}\trput{9}}{
       \psset{linestyle=solid}\pstree{\CC{\hyperlink{spotgac<10>}{23}}\trput{6}}{
        \psset{linestyle=solid}\FAIL{spotgac<11>}\trput{2}
        \psset{linestyle=solid}\FAIL{spotgac<12>}\trput{8}
       }
       \psset{linestyle=solid}\FAIL{spotgac<13>}\trput{8}
      }
     }
     \psset{linestyle=solid}\FAIL{spotgac<14>}\trput{8}
    }
   }
   \psset{linestyle=solid}\FAIL{spotgac<15>}\trput{6}
  }
  \psset{linestyle=solid}\pstree{\CC{\hyperlink{spotgac<16>}{6}}\trput{6}}{
   \psset{linestyle=solid}\pstree{\CC{\hyperlink{spotgac<17>}{10}}\trput{4}}{
    \psset{linestyle=solid}\FAIL{spotgac<18>}\trput{4}
    \psset{linestyle=solid}\FAIL{spotgac<19>}\trput{6}
    \psset{linestyle=solid}\FAIL{spotgac<20>}\trput{8}
   }
   \psset{linestyle=solid}\pstree{\CC{\hyperlink{spotgac<21>}{7}}\trput{5}}{
    \psset{linestyle=solid}\pstree{\CC{\hyperlink{spotgac<22>}{10}}\trput{4}}{
     \psset{linestyle=solid}\FAIL{spotgac<23>}\trput{5}
     \psset{linestyle=solid}\FAIL{spotgac<24>}\trput{6}
     \psset{linestyle=solid}\FAIL{spotgac<25>}\trput{8}
    }
    \psset{linestyle=solid}\pstree{\CC{\hyperlink{spotgac<26>}{10}}\trput{9}}{
     \psset{linestyle=solid}\pstree{\CC{\hyperlink{spotgac<27>}{11}}\trput{5}}{
      \psset{linestyle=solid}\FAIL{spotgac<28>}\trput{6}
      \psset{linestyle=solid}\FAIL{spotgac<29>}\trput{8}
     }
     \psset{linestyle=solid}\FAIL{spotgac<30>}\trput{6}
     \psset{linestyle=solid}\FAIL{spotgac<31>}\trput{8}
    }
   }
  }
 }
 \psset{linestyle=solid}\pstree{\CC{\hyperlink{spotgac<32>}{5}}\trput{4}}{
  \psset{linestyle=solid}\pstree{\CC{\hyperlink{spotgac<33>}{6}}\trput{2}}{
   \psset{linestyle=solid}\FAIL{spotgac<34>}\trput{5}
   \psset{linestyle=solid}\FAIL{spotgac<35>}\trput{6}
  }
 }
}
}
}
\only<36>{
\scalebox{ 0.40}{\psset{nodesep=0pt,levelsep=36pt}
\pstree{\CC{\hyperlink{spotgac<1>}{4}}}{
 \psset{linestyle=solid}\pstree{\CC{\hyperlink{spotgac<2>}{5}}\trput{2}}{
  \psset{linestyle=solid}\pstree{\CC{\hyperlink{spotgac<3>}{6}}\trput{4}}{
   \psset{linestyle=solid}\pstree{\CC{\hyperlink{spotgac<4>}{7}}\trput{5}}{
    \psset{linestyle=solid}\FAIL{spotgac<5>}\trput{6}
    \psset{linestyle=solid}\pstree{\CC{\hyperlink{spotgac<6>}{12}}\trput{9}}{
     \psset{linestyle=solid}\pstree{\CC{\hyperlink{spotgac<7>}{19}}\trput{6}}{
      \psset{linestyle=solid}\FAIL{spotgac<8>}\trput{6}
      \psset{linestyle=solid}\pstree{\CC{\hyperlink{spotgac<9>}{20}}\trput{9}}{
       \psset{linestyle=solid}\pstree{\CC{\hyperlink{spotgac<10>}{23}}\trput{6}}{
        \psset{linestyle=solid}\FAIL{spotgac<11>}\trput{2}
        \psset{linestyle=solid}\FAIL{spotgac<12>}\trput{8}
       }
       \psset{linestyle=solid}\FAIL{spotgac<13>}\trput{8}
      }
     }
     \psset{linestyle=solid}\FAIL{spotgac<14>}\trput{8}
    }
   }
   \psset{linestyle=solid}\FAIL{spotgac<15>}\trput{6}
  }
  \psset{linestyle=solid}\pstree{\CC{\hyperlink{spotgac<16>}{6}}\trput{6}}{
   \psset{linestyle=solid}\pstree{\CC{\hyperlink{spotgac<17>}{10}}\trput{4}}{
    \psset{linestyle=solid}\FAIL{spotgac<18>}\trput{4}
    \psset{linestyle=solid}\FAIL{spotgac<19>}\trput{6}
    \psset{linestyle=solid}\FAIL{spotgac<20>}\trput{8}
   }
   \psset{linestyle=solid}\pstree{\CC{\hyperlink{spotgac<21>}{7}}\trput{5}}{
    \psset{linestyle=solid}\pstree{\CC{\hyperlink{spotgac<22>}{10}}\trput{4}}{
     \psset{linestyle=solid}\FAIL{spotgac<23>}\trput{5}
     \psset{linestyle=solid}\FAIL{spotgac<24>}\trput{6}
     \psset{linestyle=solid}\FAIL{spotgac<25>}\trput{8}
    }
    \psset{linestyle=solid}\pstree{\CC{\hyperlink{spotgac<26>}{10}}\trput{9}}{
     \psset{linestyle=solid}\pstree{\CC{\hyperlink{spotgac<27>}{11}}\trput{5}}{
      \psset{linestyle=solid}\FAIL{spotgac<28>}\trput{6}
      \psset{linestyle=solid}\FAIL{spotgac<29>}\trput{8}
     }
     \psset{linestyle=solid}\FAIL{spotgac<30>}\trput{6}
     \psset{linestyle=solid}\FAIL{spotgac<31>}\trput{8}
    }
   }
  }
 }
 \psset{linestyle=solid}\pstree{\CC{\hyperlink{spotgac<32>}{5}}\trput{4}}{
  \psset{linestyle=solid}\pstree{\CC{\hyperlink{spotgac<33>}{6}}\trput{2}}{
   \psset{linestyle=solid}\FAIL{spotgac<34>}\trput{5}
   \psset{linestyle=solid}\FAIL{spotgac<35>}\trput{6}
  }
  \psset{linestyle=solid}\pstree{\CC{\hyperlink{spotgac<36>}{6}}\trput{6}}{
  }
 }
}
}
}
\only<37>{
\scalebox{ 0.40}{\psset{nodesep=0pt,levelsep=36pt}
\pstree{\CC{\hyperlink{spotgac<1>}{4}}}{
 \psset{linestyle=solid}\pstree{\CC{\hyperlink{spotgac<2>}{5}}\trput{2}}{
  \psset{linestyle=solid}\pstree{\CC{\hyperlink{spotgac<3>}{6}}\trput{4}}{
   \psset{linestyle=solid}\pstree{\CC{\hyperlink{spotgac<4>}{7}}\trput{5}}{
    \psset{linestyle=solid}\FAIL{spotgac<5>}\trput{6}
    \psset{linestyle=solid}\pstree{\CC{\hyperlink{spotgac<6>}{12}}\trput{9}}{
     \psset{linestyle=solid}\pstree{\CC{\hyperlink{spotgac<7>}{19}}\trput{6}}{
      \psset{linestyle=solid}\FAIL{spotgac<8>}\trput{6}
      \psset{linestyle=solid}\pstree{\CC{\hyperlink{spotgac<9>}{20}}\trput{9}}{
       \psset{linestyle=solid}\pstree{\CC{\hyperlink{spotgac<10>}{23}}\trput{6}}{
        \psset{linestyle=solid}\FAIL{spotgac<11>}\trput{2}
        \psset{linestyle=solid}\FAIL{spotgac<12>}\trput{8}
       }
       \psset{linestyle=solid}\FAIL{spotgac<13>}\trput{8}
      }
     }
     \psset{linestyle=solid}\FAIL{spotgac<14>}\trput{8}
    }
   }
   \psset{linestyle=solid}\FAIL{spotgac<15>}\trput{6}
  }
  \psset{linestyle=solid}\pstree{\CC{\hyperlink{spotgac<16>}{6}}\trput{6}}{
   \psset{linestyle=solid}\pstree{\CC{\hyperlink{spotgac<17>}{10}}\trput{4}}{
    \psset{linestyle=solid}\FAIL{spotgac<18>}\trput{4}
    \psset{linestyle=solid}\FAIL{spotgac<19>}\trput{6}
    \psset{linestyle=solid}\FAIL{spotgac<20>}\trput{8}
   }
   \psset{linestyle=solid}\pstree{\CC{\hyperlink{spotgac<21>}{7}}\trput{5}}{
    \psset{linestyle=solid}\pstree{\CC{\hyperlink{spotgac<22>}{10}}\trput{4}}{
     \psset{linestyle=solid}\FAIL{spotgac<23>}\trput{5}
     \psset{linestyle=solid}\FAIL{spotgac<24>}\trput{6}
     \psset{linestyle=solid}\FAIL{spotgac<25>}\trput{8}
    }
    \psset{linestyle=solid}\pstree{\CC{\hyperlink{spotgac<26>}{10}}\trput{9}}{
     \psset{linestyle=solid}\pstree{\CC{\hyperlink{spotgac<27>}{11}}\trput{5}}{
      \psset{linestyle=solid}\FAIL{spotgac<28>}\trput{6}
      \psset{linestyle=solid}\FAIL{spotgac<29>}\trput{8}
     }
     \psset{linestyle=solid}\FAIL{spotgac<30>}\trput{6}
     \psset{linestyle=solid}\FAIL{spotgac<31>}\trput{8}
    }
   }
  }
 }
 \psset{linestyle=solid}\pstree{\CC{\hyperlink{spotgac<32>}{5}}\trput{4}}{
  \psset{linestyle=solid}\pstree{\CC{\hyperlink{spotgac<33>}{6}}\trput{2}}{
   \psset{linestyle=solid}\FAIL{spotgac<34>}\trput{5}
   \psset{linestyle=solid}\FAIL{spotgac<35>}\trput{6}
  }
  \psset{linestyle=solid}\pstree{\CC{\hyperlink{spotgac<36>}{6}}\trput{6}}{
   \psset{linestyle=solid}\FAIL{spotgac<37>}\trput{2}
  }
 }
}
}
}
\only<38>{
\scalebox{ 0.40}{\psset{nodesep=0pt,levelsep=36pt}
\pstree{\CC{\hyperlink{spotgac<1>}{4}}}{
 \psset{linestyle=solid}\pstree{\CC{\hyperlink{spotgac<2>}{5}}\trput{2}}{
  \psset{linestyle=solid}\pstree{\CC{\hyperlink{spotgac<3>}{6}}\trput{4}}{
   \psset{linestyle=solid}\pstree{\CC{\hyperlink{spotgac<4>}{7}}\trput{5}}{
    \psset{linestyle=solid}\FAIL{spotgac<5>}\trput{6}
    \psset{linestyle=solid}\pstree{\CC{\hyperlink{spotgac<6>}{12}}\trput{9}}{
     \psset{linestyle=solid}\pstree{\CC{\hyperlink{spotgac<7>}{19}}\trput{6}}{
      \psset{linestyle=solid}\FAIL{spotgac<8>}\trput{6}
      \psset{linestyle=solid}\pstree{\CC{\hyperlink{spotgac<9>}{20}}\trput{9}}{
       \psset{linestyle=solid}\pstree{\CC{\hyperlink{spotgac<10>}{23}}\trput{6}}{
        \psset{linestyle=solid}\FAIL{spotgac<11>}\trput{2}
        \psset{linestyle=solid}\FAIL{spotgac<12>}\trput{8}
       }
       \psset{linestyle=solid}\FAIL{spotgac<13>}\trput{8}
      }
     }
     \psset{linestyle=solid}\FAIL{spotgac<14>}\trput{8}
    }
   }
   \psset{linestyle=solid}\FAIL{spotgac<15>}\trput{6}
  }
  \psset{linestyle=solid}\pstree{\CC{\hyperlink{spotgac<16>}{6}}\trput{6}}{
   \psset{linestyle=solid}\pstree{\CC{\hyperlink{spotgac<17>}{10}}\trput{4}}{
    \psset{linestyle=solid}\FAIL{spotgac<18>}\trput{4}
    \psset{linestyle=solid}\FAIL{spotgac<19>}\trput{6}
    \psset{linestyle=solid}\FAIL{spotgac<20>}\trput{8}
   }
   \psset{linestyle=solid}\pstree{\CC{\hyperlink{spotgac<21>}{7}}\trput{5}}{
    \psset{linestyle=solid}\pstree{\CC{\hyperlink{spotgac<22>}{10}}\trput{4}}{
     \psset{linestyle=solid}\FAIL{spotgac<23>}\trput{5}
     \psset{linestyle=solid}\FAIL{spotgac<24>}\trput{6}
     \psset{linestyle=solid}\FAIL{spotgac<25>}\trput{8}
    }
    \psset{linestyle=solid}\pstree{\CC{\hyperlink{spotgac<26>}{10}}\trput{9}}{
     \psset{linestyle=solid}\pstree{\CC{\hyperlink{spotgac<27>}{11}}\trput{5}}{
      \psset{linestyle=solid}\FAIL{spotgac<28>}\trput{6}
      \psset{linestyle=solid}\FAIL{spotgac<29>}\trput{8}
     }
     \psset{linestyle=solid}\FAIL{spotgac<30>}\trput{6}
     \psset{linestyle=solid}\FAIL{spotgac<31>}\trput{8}
    }
   }
  }
 }
 \psset{linestyle=solid}\pstree{\CC{\hyperlink{spotgac<32>}{5}}\trput{4}}{
  \psset{linestyle=solid}\pstree{\CC{\hyperlink{spotgac<33>}{6}}\trput{2}}{
   \psset{linestyle=solid}\FAIL{spotgac<34>}\trput{5}
   \psset{linestyle=solid}\FAIL{spotgac<35>}\trput{6}
  }
  \psset{linestyle=solid}\pstree{\CC{\hyperlink{spotgac<36>}{6}}\trput{6}}{
   \psset{linestyle=solid}\FAIL{spotgac<37>}\trput{2}
   \psset{linestyle=solid}\FAIL{spotgac<38>}\trput{5}
  }
 }
}
}
}
\only<39>{
\scalebox{ 0.40}{\psset{nodesep=0pt,levelsep=36pt}
\pstree{\CC{\hyperlink{spotgac<1>}{4}}}{
 \psset{linestyle=solid}\pstree{\CC{\hyperlink{spotgac<2>}{5}}\trput{2}}{
  \psset{linestyle=solid}\pstree{\CC{\hyperlink{spotgac<3>}{6}}\trput{4}}{
   \psset{linestyle=solid}\pstree{\CC{\hyperlink{spotgac<4>}{7}}\trput{5}}{
    \psset{linestyle=solid}\FAIL{spotgac<5>}\trput{6}
    \psset{linestyle=solid}\pstree{\CC{\hyperlink{spotgac<6>}{12}}\trput{9}}{
     \psset{linestyle=solid}\pstree{\CC{\hyperlink{spotgac<7>}{19}}\trput{6}}{
      \psset{linestyle=solid}\FAIL{spotgac<8>}\trput{6}
      \psset{linestyle=solid}\pstree{\CC{\hyperlink{spotgac<9>}{20}}\trput{9}}{
       \psset{linestyle=solid}\pstree{\CC{\hyperlink{spotgac<10>}{23}}\trput{6}}{
        \psset{linestyle=solid}\FAIL{spotgac<11>}\trput{2}
        \psset{linestyle=solid}\FAIL{spotgac<12>}\trput{8}
       }
       \psset{linestyle=solid}\FAIL{spotgac<13>}\trput{8}
      }
     }
     \psset{linestyle=solid}\FAIL{spotgac<14>}\trput{8}
    }
   }
   \psset{linestyle=solid}\FAIL{spotgac<15>}\trput{6}
  }
  \psset{linestyle=solid}\pstree{\CC{\hyperlink{spotgac<16>}{6}}\trput{6}}{
   \psset{linestyle=solid}\pstree{\CC{\hyperlink{spotgac<17>}{10}}\trput{4}}{
    \psset{linestyle=solid}\FAIL{spotgac<18>}\trput{4}
    \psset{linestyle=solid}\FAIL{spotgac<19>}\trput{6}
    \psset{linestyle=solid}\FAIL{spotgac<20>}\trput{8}
   }
   \psset{linestyle=solid}\pstree{\CC{\hyperlink{spotgac<21>}{7}}\trput{5}}{
    \psset{linestyle=solid}\pstree{\CC{\hyperlink{spotgac<22>}{10}}\trput{4}}{
     \psset{linestyle=solid}\FAIL{spotgac<23>}\trput{5}
     \psset{linestyle=solid}\FAIL{spotgac<24>}\trput{6}
     \psset{linestyle=solid}\FAIL{spotgac<25>}\trput{8}
    }
    \psset{linestyle=solid}\pstree{\CC{\hyperlink{spotgac<26>}{10}}\trput{9}}{
     \psset{linestyle=solid}\pstree{\CC{\hyperlink{spotgac<27>}{11}}\trput{5}}{
      \psset{linestyle=solid}\FAIL{spotgac<28>}\trput{6}
      \psset{linestyle=solid}\FAIL{spotgac<29>}\trput{8}
     }
     \psset{linestyle=solid}\FAIL{spotgac<30>}\trput{6}
     \psset{linestyle=solid}\FAIL{spotgac<31>}\trput{8}
    }
   }
  }
 }
 \psset{linestyle=solid}\pstree{\CC{\hyperlink{spotgac<32>}{5}}\trput{4}}{
  \psset{linestyle=solid}\pstree{\CC{\hyperlink{spotgac<33>}{6}}\trput{2}}{
   \psset{linestyle=solid}\FAIL{spotgac<34>}\trput{5}
   \psset{linestyle=solid}\FAIL{spotgac<35>}\trput{6}
  }
  \psset{linestyle=solid}\pstree{\CC{\hyperlink{spotgac<36>}{6}}\trput{6}}{
   \psset{linestyle=solid}\FAIL{spotgac<37>}\trput{2}
   \psset{linestyle=solid}\FAIL{spotgac<38>}\trput{5}
  }
 }
 \psset{linestyle=solid}\pstree{\CC{\hyperlink{spotgac<39>}{5}}\trput{5}}{
 }
}
}
}
\only<40>{
\scalebox{ 0.40}{\psset{nodesep=0pt,levelsep=36pt}
\pstree{\CC{\hyperlink{spotgac<1>}{4}}}{
 \psset{linestyle=solid}\pstree{\CC{\hyperlink{spotgac<2>}{5}}\trput{2}}{
  \psset{linestyle=solid}\pstree{\CC{\hyperlink{spotgac<3>}{6}}\trput{4}}{
   \psset{linestyle=solid}\pstree{\CC{\hyperlink{spotgac<4>}{7}}\trput{5}}{
    \psset{linestyle=solid}\FAIL{spotgac<5>}\trput{6}
    \psset{linestyle=solid}\pstree{\CC{\hyperlink{spotgac<6>}{12}}\trput{9}}{
     \psset{linestyle=solid}\pstree{\CC{\hyperlink{spotgac<7>}{19}}\trput{6}}{
      \psset{linestyle=solid}\FAIL{spotgac<8>}\trput{6}
      \psset{linestyle=solid}\pstree{\CC{\hyperlink{spotgac<9>}{20}}\trput{9}}{
       \psset{linestyle=solid}\pstree{\CC{\hyperlink{spotgac<10>}{23}}\trput{6}}{
        \psset{linestyle=solid}\FAIL{spotgac<11>}\trput{2}
        \psset{linestyle=solid}\FAIL{spotgac<12>}\trput{8}
       }
       \psset{linestyle=solid}\FAIL{spotgac<13>}\trput{8}
      }
     }
     \psset{linestyle=solid}\FAIL{spotgac<14>}\trput{8}
    }
   }
   \psset{linestyle=solid}\FAIL{spotgac<15>}\trput{6}
  }
  \psset{linestyle=solid}\pstree{\CC{\hyperlink{spotgac<16>}{6}}\trput{6}}{
   \psset{linestyle=solid}\pstree{\CC{\hyperlink{spotgac<17>}{10}}\trput{4}}{
    \psset{linestyle=solid}\FAIL{spotgac<18>}\trput{4}
    \psset{linestyle=solid}\FAIL{spotgac<19>}\trput{6}
    \psset{linestyle=solid}\FAIL{spotgac<20>}\trput{8}
   }
   \psset{linestyle=solid}\pstree{\CC{\hyperlink{spotgac<21>}{7}}\trput{5}}{
    \psset{linestyle=solid}\pstree{\CC{\hyperlink{spotgac<22>}{10}}\trput{4}}{
     \psset{linestyle=solid}\FAIL{spotgac<23>}\trput{5}
     \psset{linestyle=solid}\FAIL{spotgac<24>}\trput{6}
     \psset{linestyle=solid}\FAIL{spotgac<25>}\trput{8}
    }
    \psset{linestyle=solid}\pstree{\CC{\hyperlink{spotgac<26>}{10}}\trput{9}}{
     \psset{linestyle=solid}\pstree{\CC{\hyperlink{spotgac<27>}{11}}\trput{5}}{
      \psset{linestyle=solid}\FAIL{spotgac<28>}\trput{6}
      \psset{linestyle=solid}\FAIL{spotgac<29>}\trput{8}
     }
     \psset{linestyle=solid}\FAIL{spotgac<30>}\trput{6}
     \psset{linestyle=solid}\FAIL{spotgac<31>}\trput{8}
    }
   }
  }
 }
 \psset{linestyle=solid}\pstree{\CC{\hyperlink{spotgac<32>}{5}}\trput{4}}{
  \psset{linestyle=solid}\pstree{\CC{\hyperlink{spotgac<33>}{6}}\trput{2}}{
   \psset{linestyle=solid}\FAIL{spotgac<34>}\trput{5}
   \psset{linestyle=solid}\FAIL{spotgac<35>}\trput{6}
  }
  \psset{linestyle=solid}\pstree{\CC{\hyperlink{spotgac<36>}{6}}\trput{6}}{
   \psset{linestyle=solid}\FAIL{spotgac<37>}\trput{2}
   \psset{linestyle=solid}\FAIL{spotgac<38>}\trput{5}
  }
 }
 \psset{linestyle=solid}\pstree{\CC{\hyperlink{spotgac<39>}{5}}\trput{5}}{
  \psset{linestyle=solid}\pstree{\CC{\hyperlink{spotgac<40>}{6}}\trput{2}}{
  }
 }
}
}
}
\only<41>{
\scalebox{ 0.40}{\psset{nodesep=0pt,levelsep=36pt}
\pstree{\CC{\hyperlink{spotgac<1>}{4}}}{
 \psset{linestyle=solid}\pstree{\CC{\hyperlink{spotgac<2>}{5}}\trput{2}}{
  \psset{linestyle=solid}\pstree{\CC{\hyperlink{spotgac<3>}{6}}\trput{4}}{
   \psset{linestyle=solid}\pstree{\CC{\hyperlink{spotgac<4>}{7}}\trput{5}}{
    \psset{linestyle=solid}\FAIL{spotgac<5>}\trput{6}
    \psset{linestyle=solid}\pstree{\CC{\hyperlink{spotgac<6>}{12}}\trput{9}}{
     \psset{linestyle=solid}\pstree{\CC{\hyperlink{spotgac<7>}{19}}\trput{6}}{
      \psset{linestyle=solid}\FAIL{spotgac<8>}\trput{6}
      \psset{linestyle=solid}\pstree{\CC{\hyperlink{spotgac<9>}{20}}\trput{9}}{
       \psset{linestyle=solid}\pstree{\CC{\hyperlink{spotgac<10>}{23}}\trput{6}}{
        \psset{linestyle=solid}\FAIL{spotgac<11>}\trput{2}
        \psset{linestyle=solid}\FAIL{spotgac<12>}\trput{8}
       }
       \psset{linestyle=solid}\FAIL{spotgac<13>}\trput{8}
      }
     }
     \psset{linestyle=solid}\FAIL{spotgac<14>}\trput{8}
    }
   }
   \psset{linestyle=solid}\FAIL{spotgac<15>}\trput{6}
  }
  \psset{linestyle=solid}\pstree{\CC{\hyperlink{spotgac<16>}{6}}\trput{6}}{
   \psset{linestyle=solid}\pstree{\CC{\hyperlink{spotgac<17>}{10}}\trput{4}}{
    \psset{linestyle=solid}\FAIL{spotgac<18>}\trput{4}
    \psset{linestyle=solid}\FAIL{spotgac<19>}\trput{6}
    \psset{linestyle=solid}\FAIL{spotgac<20>}\trput{8}
   }
   \psset{linestyle=solid}\pstree{\CC{\hyperlink{spotgac<21>}{7}}\trput{5}}{
    \psset{linestyle=solid}\pstree{\CC{\hyperlink{spotgac<22>}{10}}\trput{4}}{
     \psset{linestyle=solid}\FAIL{spotgac<23>}\trput{5}
     \psset{linestyle=solid}\FAIL{spotgac<24>}\trput{6}
     \psset{linestyle=solid}\FAIL{spotgac<25>}\trput{8}
    }
    \psset{linestyle=solid}\pstree{\CC{\hyperlink{spotgac<26>}{10}}\trput{9}}{
     \psset{linestyle=solid}\pstree{\CC{\hyperlink{spotgac<27>}{11}}\trput{5}}{
      \psset{linestyle=solid}\FAIL{spotgac<28>}\trput{6}
      \psset{linestyle=solid}\FAIL{spotgac<29>}\trput{8}
     }
     \psset{linestyle=solid}\FAIL{spotgac<30>}\trput{6}
     \psset{linestyle=solid}\FAIL{spotgac<31>}\trput{8}
    }
   }
  }
 }
 \psset{linestyle=solid}\pstree{\CC{\hyperlink{spotgac<32>}{5}}\trput{4}}{
  \psset{linestyle=solid}\pstree{\CC{\hyperlink{spotgac<33>}{6}}\trput{2}}{
   \psset{linestyle=solid}\FAIL{spotgac<34>}\trput{5}
   \psset{linestyle=solid}\FAIL{spotgac<35>}\trput{6}
  }
  \psset{linestyle=solid}\pstree{\CC{\hyperlink{spotgac<36>}{6}}\trput{6}}{
   \psset{linestyle=solid}\FAIL{spotgac<37>}\trput{2}
   \psset{linestyle=solid}\FAIL{spotgac<38>}\trput{5}
  }
 }
 \psset{linestyle=solid}\pstree{\CC{\hyperlink{spotgac<39>}{5}}\trput{5}}{
  \psset{linestyle=solid}\pstree{\CC{\hyperlink{spotgac<40>}{6}}\trput{2}}{
   \psset{linestyle=solid}\pstree{\CC{\hyperlink{spotgac<41>}{7}}\trput{4}}{
   }
  }
 }
}
}
}
\only<42>{
\scalebox{ 0.40}{\psset{nodesep=0pt,levelsep=36pt}
\pstree{\CC{\hyperlink{spotgac<1>}{4}}}{
 \psset{linestyle=solid}\pstree{\CC{\hyperlink{spotgac<2>}{5}}\trput{2}}{
  \psset{linestyle=solid}\pstree{\CC{\hyperlink{spotgac<3>}{6}}\trput{4}}{
   \psset{linestyle=solid}\pstree{\CC{\hyperlink{spotgac<4>}{7}}\trput{5}}{
    \psset{linestyle=solid}\FAIL{spotgac<5>}\trput{6}
    \psset{linestyle=solid}\pstree{\CC{\hyperlink{spotgac<6>}{12}}\trput{9}}{
     \psset{linestyle=solid}\pstree{\CC{\hyperlink{spotgac<7>}{19}}\trput{6}}{
      \psset{linestyle=solid}\FAIL{spotgac<8>}\trput{6}
      \psset{linestyle=solid}\pstree{\CC{\hyperlink{spotgac<9>}{20}}\trput{9}}{
       \psset{linestyle=solid}\pstree{\CC{\hyperlink{spotgac<10>}{23}}\trput{6}}{
        \psset{linestyle=solid}\FAIL{spotgac<11>}\trput{2}
        \psset{linestyle=solid}\FAIL{spotgac<12>}\trput{8}
       }
       \psset{linestyle=solid}\FAIL{spotgac<13>}\trput{8}
      }
     }
     \psset{linestyle=solid}\FAIL{spotgac<14>}\trput{8}
    }
   }
   \psset{linestyle=solid}\FAIL{spotgac<15>}\trput{6}
  }
  \psset{linestyle=solid}\pstree{\CC{\hyperlink{spotgac<16>}{6}}\trput{6}}{
   \psset{linestyle=solid}\pstree{\CC{\hyperlink{spotgac<17>}{10}}\trput{4}}{
    \psset{linestyle=solid}\FAIL{spotgac<18>}\trput{4}
    \psset{linestyle=solid}\FAIL{spotgac<19>}\trput{6}
    \psset{linestyle=solid}\FAIL{spotgac<20>}\trput{8}
   }
   \psset{linestyle=solid}\pstree{\CC{\hyperlink{spotgac<21>}{7}}\trput{5}}{
    \psset{linestyle=solid}\pstree{\CC{\hyperlink{spotgac<22>}{10}}\trput{4}}{
     \psset{linestyle=solid}\FAIL{spotgac<23>}\trput{5}
     \psset{linestyle=solid}\FAIL{spotgac<24>}\trput{6}
     \psset{linestyle=solid}\FAIL{spotgac<25>}\trput{8}
    }
    \psset{linestyle=solid}\pstree{\CC{\hyperlink{spotgac<26>}{10}}\trput{9}}{
     \psset{linestyle=solid}\pstree{\CC{\hyperlink{spotgac<27>}{11}}\trput{5}}{
      \psset{linestyle=solid}\FAIL{spotgac<28>}\trput{6}
      \psset{linestyle=solid}\FAIL{spotgac<29>}\trput{8}
     }
     \psset{linestyle=solid}\FAIL{spotgac<30>}\trput{6}
     \psset{linestyle=solid}\FAIL{spotgac<31>}\trput{8}
    }
   }
  }
 }
 \psset{linestyle=solid}\pstree{\CC{\hyperlink{spotgac<32>}{5}}\trput{4}}{
  \psset{linestyle=solid}\pstree{\CC{\hyperlink{spotgac<33>}{6}}\trput{2}}{
   \psset{linestyle=solid}\FAIL{spotgac<34>}\trput{5}
   \psset{linestyle=solid}\FAIL{spotgac<35>}\trput{6}
  }
  \psset{linestyle=solid}\pstree{\CC{\hyperlink{spotgac<36>}{6}}\trput{6}}{
   \psset{linestyle=solid}\FAIL{spotgac<37>}\trput{2}
   \psset{linestyle=solid}\FAIL{spotgac<38>}\trput{5}
  }
 }
 \psset{linestyle=solid}\pstree{\CC{\hyperlink{spotgac<39>}{5}}\trput{5}}{
  \psset{linestyle=solid}\pstree{\CC{\hyperlink{spotgac<40>}{6}}\trput{2}}{
   \psset{linestyle=solid}\pstree{\CC{\hyperlink{spotgac<41>}{7}}\trput{4}}{
    \psset{linestyle=solid}\FAIL{spotgac<42>}\trput{6}
   }
  }
 }
}
}
}
\only<43>{
\scalebox{ 0.40}{\psset{nodesep=0pt,levelsep=36pt}
\pstree{\CC{\hyperlink{spotgac<1>}{4}}}{
 \psset{linestyle=solid}\pstree{\CC{\hyperlink{spotgac<2>}{5}}\trput{2}}{
  \psset{linestyle=solid}\pstree{\CC{\hyperlink{spotgac<3>}{6}}\trput{4}}{
   \psset{linestyle=solid}\pstree{\CC{\hyperlink{spotgac<4>}{7}}\trput{5}}{
    \psset{linestyle=solid}\FAIL{spotgac<5>}\trput{6}
    \psset{linestyle=solid}\pstree{\CC{\hyperlink{spotgac<6>}{12}}\trput{9}}{
     \psset{linestyle=solid}\pstree{\CC{\hyperlink{spotgac<7>}{19}}\trput{6}}{
      \psset{linestyle=solid}\FAIL{spotgac<8>}\trput{6}
      \psset{linestyle=solid}\pstree{\CC{\hyperlink{spotgac<9>}{20}}\trput{9}}{
       \psset{linestyle=solid}\pstree{\CC{\hyperlink{spotgac<10>}{23}}\trput{6}}{
        \psset{linestyle=solid}\FAIL{spotgac<11>}\trput{2}
        \psset{linestyle=solid}\FAIL{spotgac<12>}\trput{8}
       }
       \psset{linestyle=solid}\FAIL{spotgac<13>}\trput{8}
      }
     }
     \psset{linestyle=solid}\FAIL{spotgac<14>}\trput{8}
    }
   }
   \psset{linestyle=solid}\FAIL{spotgac<15>}\trput{6}
  }
  \psset{linestyle=solid}\pstree{\CC{\hyperlink{spotgac<16>}{6}}\trput{6}}{
   \psset{linestyle=solid}\pstree{\CC{\hyperlink{spotgac<17>}{10}}\trput{4}}{
    \psset{linestyle=solid}\FAIL{spotgac<18>}\trput{4}
    \psset{linestyle=solid}\FAIL{spotgac<19>}\trput{6}
    \psset{linestyle=solid}\FAIL{spotgac<20>}\trput{8}
   }
   \psset{linestyle=solid}\pstree{\CC{\hyperlink{spotgac<21>}{7}}\trput{5}}{
    \psset{linestyle=solid}\pstree{\CC{\hyperlink{spotgac<22>}{10}}\trput{4}}{
     \psset{linestyle=solid}\FAIL{spotgac<23>}\trput{5}
     \psset{linestyle=solid}\FAIL{spotgac<24>}\trput{6}
     \psset{linestyle=solid}\FAIL{spotgac<25>}\trput{8}
    }
    \psset{linestyle=solid}\pstree{\CC{\hyperlink{spotgac<26>}{10}}\trput{9}}{
     \psset{linestyle=solid}\pstree{\CC{\hyperlink{spotgac<27>}{11}}\trput{5}}{
      \psset{linestyle=solid}\FAIL{spotgac<28>}\trput{6}
      \psset{linestyle=solid}\FAIL{spotgac<29>}\trput{8}
     }
     \psset{linestyle=solid}\FAIL{spotgac<30>}\trput{6}
     \psset{linestyle=solid}\FAIL{spotgac<31>}\trput{8}
    }
   }
  }
 }
 \psset{linestyle=solid}\pstree{\CC{\hyperlink{spotgac<32>}{5}}\trput{4}}{
  \psset{linestyle=solid}\pstree{\CC{\hyperlink{spotgac<33>}{6}}\trput{2}}{
   \psset{linestyle=solid}\FAIL{spotgac<34>}\trput{5}
   \psset{linestyle=solid}\FAIL{spotgac<35>}\trput{6}
  }
  \psset{linestyle=solid}\pstree{\CC{\hyperlink{spotgac<36>}{6}}\trput{6}}{
   \psset{linestyle=solid}\FAIL{spotgac<37>}\trput{2}
   \psset{linestyle=solid}\FAIL{spotgac<38>}\trput{5}
  }
 }
 \psset{linestyle=solid}\pstree{\CC{\hyperlink{spotgac<39>}{5}}\trput{5}}{
  \psset{linestyle=solid}\pstree{\CC{\hyperlink{spotgac<40>}{6}}\trput{2}}{
   \psset{linestyle=solid}\pstree{\CC{\hyperlink{spotgac<41>}{7}}\trput{4}}{
    \psset{linestyle=solid}\FAIL{spotgac<42>}\trput{6}
    \psset{linestyle=solid}\pstree{\CC{\hyperlink{spotgac<43>}{10}}\trput{9}}{
    }
   }
  }
 }
}
}
}
\only<44>{
\scalebox{ 0.40}{\psset{nodesep=0pt,levelsep=36pt}
\pstree{\CC{\hyperlink{spotgac<1>}{4}}}{
 \psset{linestyle=solid}\pstree{\CC{\hyperlink{spotgac<2>}{5}}\trput{2}}{
  \psset{linestyle=solid}\pstree{\CC{\hyperlink{spotgac<3>}{6}}\trput{4}}{
   \psset{linestyle=solid}\pstree{\CC{\hyperlink{spotgac<4>}{7}}\trput{5}}{
    \psset{linestyle=solid}\FAIL{spotgac<5>}\trput{6}
    \psset{linestyle=solid}\pstree{\CC{\hyperlink{spotgac<6>}{12}}\trput{9}}{
     \psset{linestyle=solid}\pstree{\CC{\hyperlink{spotgac<7>}{19}}\trput{6}}{
      \psset{linestyle=solid}\FAIL{spotgac<8>}\trput{6}
      \psset{linestyle=solid}\pstree{\CC{\hyperlink{spotgac<9>}{20}}\trput{9}}{
       \psset{linestyle=solid}\pstree{\CC{\hyperlink{spotgac<10>}{23}}\trput{6}}{
        \psset{linestyle=solid}\FAIL{spotgac<11>}\trput{2}
        \psset{linestyle=solid}\FAIL{spotgac<12>}\trput{8}
       }
       \psset{linestyle=solid}\FAIL{spotgac<13>}\trput{8}
      }
     }
     \psset{linestyle=solid}\FAIL{spotgac<14>}\trput{8}
    }
   }
   \psset{linestyle=solid}\FAIL{spotgac<15>}\trput{6}
  }
  \psset{linestyle=solid}\pstree{\CC{\hyperlink{spotgac<16>}{6}}\trput{6}}{
   \psset{linestyle=solid}\pstree{\CC{\hyperlink{spotgac<17>}{10}}\trput{4}}{
    \psset{linestyle=solid}\FAIL{spotgac<18>}\trput{4}
    \psset{linestyle=solid}\FAIL{spotgac<19>}\trput{6}
    \psset{linestyle=solid}\FAIL{spotgac<20>}\trput{8}
   }
   \psset{linestyle=solid}\pstree{\CC{\hyperlink{spotgac<21>}{7}}\trput{5}}{
    \psset{linestyle=solid}\pstree{\CC{\hyperlink{spotgac<22>}{10}}\trput{4}}{
     \psset{linestyle=solid}\FAIL{spotgac<23>}\trput{5}
     \psset{linestyle=solid}\FAIL{spotgac<24>}\trput{6}
     \psset{linestyle=solid}\FAIL{spotgac<25>}\trput{8}
    }
    \psset{linestyle=solid}\pstree{\CC{\hyperlink{spotgac<26>}{10}}\trput{9}}{
     \psset{linestyle=solid}\pstree{\CC{\hyperlink{spotgac<27>}{11}}\trput{5}}{
      \psset{linestyle=solid}\FAIL{spotgac<28>}\trput{6}
      \psset{linestyle=solid}\FAIL{spotgac<29>}\trput{8}
     }
     \psset{linestyle=solid}\FAIL{spotgac<30>}\trput{6}
     \psset{linestyle=solid}\FAIL{spotgac<31>}\trput{8}
    }
   }
  }
 }
 \psset{linestyle=solid}\pstree{\CC{\hyperlink{spotgac<32>}{5}}\trput{4}}{
  \psset{linestyle=solid}\pstree{\CC{\hyperlink{spotgac<33>}{6}}\trput{2}}{
   \psset{linestyle=solid}\FAIL{spotgac<34>}\trput{5}
   \psset{linestyle=solid}\FAIL{spotgac<35>}\trput{6}
  }
  \psset{linestyle=solid}\pstree{\CC{\hyperlink{spotgac<36>}{6}}\trput{6}}{
   \psset{linestyle=solid}\FAIL{spotgac<37>}\trput{2}
   \psset{linestyle=solid}\FAIL{spotgac<38>}\trput{5}
  }
 }
 \psset{linestyle=solid}\pstree{\CC{\hyperlink{spotgac<39>}{5}}\trput{5}}{
  \psset{linestyle=solid}\pstree{\CC{\hyperlink{spotgac<40>}{6}}\trput{2}}{
   \psset{linestyle=solid}\pstree{\CC{\hyperlink{spotgac<41>}{7}}\trput{4}}{
    \psset{linestyle=solid}\FAIL{spotgac<42>}\trput{6}
    \psset{linestyle=solid}\pstree{\CC{\hyperlink{spotgac<43>}{10}}\trput{9}}{
     \psset{linestyle=solid}\FAIL{spotgac<44>}\trput{6}
    }
   }
  }
 }
}
}
}
\only<45>{
\scalebox{ 0.40}{\psset{nodesep=0pt,levelsep=36pt}
\pstree{\CC{\hyperlink{spotgac<1>}{4}}}{
 \psset{linestyle=solid}\pstree{\CC{\hyperlink{spotgac<2>}{5}}\trput{2}}{
  \psset{linestyle=solid}\pstree{\CC{\hyperlink{spotgac<3>}{6}}\trput{4}}{
   \psset{linestyle=solid}\pstree{\CC{\hyperlink{spotgac<4>}{7}}\trput{5}}{
    \psset{linestyle=solid}\FAIL{spotgac<5>}\trput{6}
    \psset{linestyle=solid}\pstree{\CC{\hyperlink{spotgac<6>}{12}}\trput{9}}{
     \psset{linestyle=solid}\pstree{\CC{\hyperlink{spotgac<7>}{19}}\trput{6}}{
      \psset{linestyle=solid}\FAIL{spotgac<8>}\trput{6}
      \psset{linestyle=solid}\pstree{\CC{\hyperlink{spotgac<9>}{20}}\trput{9}}{
       \psset{linestyle=solid}\pstree{\CC{\hyperlink{spotgac<10>}{23}}\trput{6}}{
        \psset{linestyle=solid}\FAIL{spotgac<11>}\trput{2}
        \psset{linestyle=solid}\FAIL{spotgac<12>}\trput{8}
       }
       \psset{linestyle=solid}\FAIL{spotgac<13>}\trput{8}
      }
     }
     \psset{linestyle=solid}\FAIL{spotgac<14>}\trput{8}
    }
   }
   \psset{linestyle=solid}\FAIL{spotgac<15>}\trput{6}
  }
  \psset{linestyle=solid}\pstree{\CC{\hyperlink{spotgac<16>}{6}}\trput{6}}{
   \psset{linestyle=solid}\pstree{\CC{\hyperlink{spotgac<17>}{10}}\trput{4}}{
    \psset{linestyle=solid}\FAIL{spotgac<18>}\trput{4}
    \psset{linestyle=solid}\FAIL{spotgac<19>}\trput{6}
    \psset{linestyle=solid}\FAIL{spotgac<20>}\trput{8}
   }
   \psset{linestyle=solid}\pstree{\CC{\hyperlink{spotgac<21>}{7}}\trput{5}}{
    \psset{linestyle=solid}\pstree{\CC{\hyperlink{spotgac<22>}{10}}\trput{4}}{
     \psset{linestyle=solid}\FAIL{spotgac<23>}\trput{5}
     \psset{linestyle=solid}\FAIL{spotgac<24>}\trput{6}
     \psset{linestyle=solid}\FAIL{spotgac<25>}\trput{8}
    }
    \psset{linestyle=solid}\pstree{\CC{\hyperlink{spotgac<26>}{10}}\trput{9}}{
     \psset{linestyle=solid}\pstree{\CC{\hyperlink{spotgac<27>}{11}}\trput{5}}{
      \psset{linestyle=solid}\FAIL{spotgac<28>}\trput{6}
      \psset{linestyle=solid}\FAIL{spotgac<29>}\trput{8}
     }
     \psset{linestyle=solid}\FAIL{spotgac<30>}\trput{6}
     \psset{linestyle=solid}\FAIL{spotgac<31>}\trput{8}
    }
   }
  }
 }
 \psset{linestyle=solid}\pstree{\CC{\hyperlink{spotgac<32>}{5}}\trput{4}}{
  \psset{linestyle=solid}\pstree{\CC{\hyperlink{spotgac<33>}{6}}\trput{2}}{
   \psset{linestyle=solid}\FAIL{spotgac<34>}\trput{5}
   \psset{linestyle=solid}\FAIL{spotgac<35>}\trput{6}
  }
  \psset{linestyle=solid}\pstree{\CC{\hyperlink{spotgac<36>}{6}}\trput{6}}{
   \psset{linestyle=solid}\FAIL{spotgac<37>}\trput{2}
   \psset{linestyle=solid}\FAIL{spotgac<38>}\trput{5}
  }
 }
 \psset{linestyle=solid}\pstree{\CC{\hyperlink{spotgac<39>}{5}}\trput{5}}{
  \psset{linestyle=solid}\pstree{\CC{\hyperlink{spotgac<40>}{6}}\trput{2}}{
   \psset{linestyle=solid}\pstree{\CC{\hyperlink{spotgac<41>}{7}}\trput{4}}{
    \psset{linestyle=solid}\FAIL{spotgac<42>}\trput{6}
    \psset{linestyle=solid}\pstree{\CC{\hyperlink{spotgac<43>}{10}}\trput{9}}{
     \psset{linestyle=solid}\FAIL{spotgac<44>}\trput{6}
     \psset{linestyle=solid}\FAIL{spotgac<45>}\trput{8}
    }
   }
  }
 }
}
}
}
\only<46>{
\scalebox{ 0.40}{\psset{nodesep=0pt,levelsep=36pt}
\pstree{\CC{\hyperlink{spotgac<1>}{4}}}{
 \psset{linestyle=solid}\pstree{\CC{\hyperlink{spotgac<2>}{5}}\trput{2}}{
  \psset{linestyle=solid}\pstree{\CC{\hyperlink{spotgac<3>}{6}}\trput{4}}{
   \psset{linestyle=solid}\pstree{\CC{\hyperlink{spotgac<4>}{7}}\trput{5}}{
    \psset{linestyle=solid}\FAIL{spotgac<5>}\trput{6}
    \psset{linestyle=solid}\pstree{\CC{\hyperlink{spotgac<6>}{12}}\trput{9}}{
     \psset{linestyle=solid}\pstree{\CC{\hyperlink{spotgac<7>}{19}}\trput{6}}{
      \psset{linestyle=solid}\FAIL{spotgac<8>}\trput{6}
      \psset{linestyle=solid}\pstree{\CC{\hyperlink{spotgac<9>}{20}}\trput{9}}{
       \psset{linestyle=solid}\pstree{\CC{\hyperlink{spotgac<10>}{23}}\trput{6}}{
        \psset{linestyle=solid}\FAIL{spotgac<11>}\trput{2}
        \psset{linestyle=solid}\FAIL{spotgac<12>}\trput{8}
       }
       \psset{linestyle=solid}\FAIL{spotgac<13>}\trput{8}
      }
     }
     \psset{linestyle=solid}\FAIL{spotgac<14>}\trput{8}
    }
   }
   \psset{linestyle=solid}\FAIL{spotgac<15>}\trput{6}
  }
  \psset{linestyle=solid}\pstree{\CC{\hyperlink{spotgac<16>}{6}}\trput{6}}{
   \psset{linestyle=solid}\pstree{\CC{\hyperlink{spotgac<17>}{10}}\trput{4}}{
    \psset{linestyle=solid}\FAIL{spotgac<18>}\trput{4}
    \psset{linestyle=solid}\FAIL{spotgac<19>}\trput{6}
    \psset{linestyle=solid}\FAIL{spotgac<20>}\trput{8}
   }
   \psset{linestyle=solid}\pstree{\CC{\hyperlink{spotgac<21>}{7}}\trput{5}}{
    \psset{linestyle=solid}\pstree{\CC{\hyperlink{spotgac<22>}{10}}\trput{4}}{
     \psset{linestyle=solid}\FAIL{spotgac<23>}\trput{5}
     \psset{linestyle=solid}\FAIL{spotgac<24>}\trput{6}
     \psset{linestyle=solid}\FAIL{spotgac<25>}\trput{8}
    }
    \psset{linestyle=solid}\pstree{\CC{\hyperlink{spotgac<26>}{10}}\trput{9}}{
     \psset{linestyle=solid}\pstree{\CC{\hyperlink{spotgac<27>}{11}}\trput{5}}{
      \psset{linestyle=solid}\FAIL{spotgac<28>}\trput{6}
      \psset{linestyle=solid}\FAIL{spotgac<29>}\trput{8}
     }
     \psset{linestyle=solid}\FAIL{spotgac<30>}\trput{6}
     \psset{linestyle=solid}\FAIL{spotgac<31>}\trput{8}
    }
   }
  }
 }
 \psset{linestyle=solid}\pstree{\CC{\hyperlink{spotgac<32>}{5}}\trput{4}}{
  \psset{linestyle=solid}\pstree{\CC{\hyperlink{spotgac<33>}{6}}\trput{2}}{
   \psset{linestyle=solid}\FAIL{spotgac<34>}\trput{5}
   \psset{linestyle=solid}\FAIL{spotgac<35>}\trput{6}
  }
  \psset{linestyle=solid}\pstree{\CC{\hyperlink{spotgac<36>}{6}}\trput{6}}{
   \psset{linestyle=solid}\FAIL{spotgac<37>}\trput{2}
   \psset{linestyle=solid}\FAIL{spotgac<38>}\trput{5}
  }
 }
 \psset{linestyle=solid}\pstree{\CC{\hyperlink{spotgac<39>}{5}}\trput{5}}{
  \psset{linestyle=solid}\pstree{\CC{\hyperlink{spotgac<40>}{6}}\trput{2}}{
   \psset{linestyle=solid}\pstree{\CC{\hyperlink{spotgac<41>}{7}}\trput{4}}{
    \psset{linestyle=solid}\FAIL{spotgac<42>}\trput{6}
    \psset{linestyle=solid}\pstree{\CC{\hyperlink{spotgac<43>}{10}}\trput{9}}{
     \psset{linestyle=solid}\FAIL{spotgac<44>}\trput{6}
     \psset{linestyle=solid}\FAIL{spotgac<45>}\trput{8}
    }
   }
   \psset{linestyle=solid}\pstree{\CC{\hyperlink{spotgac<46>}{7}}\trput{6}}{
   }
  }
 }
}
}
}
\only<47>{
\scalebox{ 0.40}{\psset{nodesep=0pt,levelsep=36pt}
\pstree{\CC{\hyperlink{spotgac<1>}{4}}}{
 \psset{linestyle=solid}\pstree{\CC{\hyperlink{spotgac<2>}{5}}\trput{2}}{
  \psset{linestyle=solid}\pstree{\CC{\hyperlink{spotgac<3>}{6}}\trput{4}}{
   \psset{linestyle=solid}\pstree{\CC{\hyperlink{spotgac<4>}{7}}\trput{5}}{
    \psset{linestyle=solid}\FAIL{spotgac<5>}\trput{6}
    \psset{linestyle=solid}\pstree{\CC{\hyperlink{spotgac<6>}{12}}\trput{9}}{
     \psset{linestyle=solid}\pstree{\CC{\hyperlink{spotgac<7>}{19}}\trput{6}}{
      \psset{linestyle=solid}\FAIL{spotgac<8>}\trput{6}
      \psset{linestyle=solid}\pstree{\CC{\hyperlink{spotgac<9>}{20}}\trput{9}}{
       \psset{linestyle=solid}\pstree{\CC{\hyperlink{spotgac<10>}{23}}\trput{6}}{
        \psset{linestyle=solid}\FAIL{spotgac<11>}\trput{2}
        \psset{linestyle=solid}\FAIL{spotgac<12>}\trput{8}
       }
       \psset{linestyle=solid}\FAIL{spotgac<13>}\trput{8}
      }
     }
     \psset{linestyle=solid}\FAIL{spotgac<14>}\trput{8}
    }
   }
   \psset{linestyle=solid}\FAIL{spotgac<15>}\trput{6}
  }
  \psset{linestyle=solid}\pstree{\CC{\hyperlink{spotgac<16>}{6}}\trput{6}}{
   \psset{linestyle=solid}\pstree{\CC{\hyperlink{spotgac<17>}{10}}\trput{4}}{
    \psset{linestyle=solid}\FAIL{spotgac<18>}\trput{4}
    \psset{linestyle=solid}\FAIL{spotgac<19>}\trput{6}
    \psset{linestyle=solid}\FAIL{spotgac<20>}\trput{8}
   }
   \psset{linestyle=solid}\pstree{\CC{\hyperlink{spotgac<21>}{7}}\trput{5}}{
    \psset{linestyle=solid}\pstree{\CC{\hyperlink{spotgac<22>}{10}}\trput{4}}{
     \psset{linestyle=solid}\FAIL{spotgac<23>}\trput{5}
     \psset{linestyle=solid}\FAIL{spotgac<24>}\trput{6}
     \psset{linestyle=solid}\FAIL{spotgac<25>}\trput{8}
    }
    \psset{linestyle=solid}\pstree{\CC{\hyperlink{spotgac<26>}{10}}\trput{9}}{
     \psset{linestyle=solid}\pstree{\CC{\hyperlink{spotgac<27>}{11}}\trput{5}}{
      \psset{linestyle=solid}\FAIL{spotgac<28>}\trput{6}
      \psset{linestyle=solid}\FAIL{spotgac<29>}\trput{8}
     }
     \psset{linestyle=solid}\FAIL{spotgac<30>}\trput{6}
     \psset{linestyle=solid}\FAIL{spotgac<31>}\trput{8}
    }
   }
  }
 }
 \psset{linestyle=solid}\pstree{\CC{\hyperlink{spotgac<32>}{5}}\trput{4}}{
  \psset{linestyle=solid}\pstree{\CC{\hyperlink{spotgac<33>}{6}}\trput{2}}{
   \psset{linestyle=solid}\FAIL{spotgac<34>}\trput{5}
   \psset{linestyle=solid}\FAIL{spotgac<35>}\trput{6}
  }
  \psset{linestyle=solid}\pstree{\CC{\hyperlink{spotgac<36>}{6}}\trput{6}}{
   \psset{linestyle=solid}\FAIL{spotgac<37>}\trput{2}
   \psset{linestyle=solid}\FAIL{spotgac<38>}\trput{5}
  }
 }
 \psset{linestyle=solid}\pstree{\CC{\hyperlink{spotgac<39>}{5}}\trput{5}}{
  \psset{linestyle=solid}\pstree{\CC{\hyperlink{spotgac<40>}{6}}\trput{2}}{
   \psset{linestyle=solid}\pstree{\CC{\hyperlink{spotgac<41>}{7}}\trput{4}}{
    \psset{linestyle=solid}\FAIL{spotgac<42>}\trput{6}
    \psset{linestyle=solid}\pstree{\CC{\hyperlink{spotgac<43>}{10}}\trput{9}}{
     \psset{linestyle=solid}\FAIL{spotgac<44>}\trput{6}
     \psset{linestyle=solid}\FAIL{spotgac<45>}\trput{8}
    }
   }
   \psset{linestyle=solid}\pstree{\CC{\hyperlink{spotgac<46>}{7}}\trput{6}}{
    \psset{linestyle=solid}\FAIL{spotgac<47>}\trput{4}
   }
  }
 }
}
}
}
\only<48>{
\scalebox{ 0.40}{\psset{nodesep=0pt,levelsep=36pt}
\pstree{\CC{\hyperlink{spotgac<1>}{4}}}{
 \psset{linestyle=solid}\pstree{\CC{\hyperlink{spotgac<2>}{5}}\trput{2}}{
  \psset{linestyle=solid}\pstree{\CC{\hyperlink{spotgac<3>}{6}}\trput{4}}{
   \psset{linestyle=solid}\pstree{\CC{\hyperlink{spotgac<4>}{7}}\trput{5}}{
    \psset{linestyle=solid}\FAIL{spotgac<5>}\trput{6}
    \psset{linestyle=solid}\pstree{\CC{\hyperlink{spotgac<6>}{12}}\trput{9}}{
     \psset{linestyle=solid}\pstree{\CC{\hyperlink{spotgac<7>}{19}}\trput{6}}{
      \psset{linestyle=solid}\FAIL{spotgac<8>}\trput{6}
      \psset{linestyle=solid}\pstree{\CC{\hyperlink{spotgac<9>}{20}}\trput{9}}{
       \psset{linestyle=solid}\pstree{\CC{\hyperlink{spotgac<10>}{23}}\trput{6}}{
        \psset{linestyle=solid}\FAIL{spotgac<11>}\trput{2}
        \psset{linestyle=solid}\FAIL{spotgac<12>}\trput{8}
       }
       \psset{linestyle=solid}\FAIL{spotgac<13>}\trput{8}
      }
     }
     \psset{linestyle=solid}\FAIL{spotgac<14>}\trput{8}
    }
   }
   \psset{linestyle=solid}\FAIL{spotgac<15>}\trput{6}
  }
  \psset{linestyle=solid}\pstree{\CC{\hyperlink{spotgac<16>}{6}}\trput{6}}{
   \psset{linestyle=solid}\pstree{\CC{\hyperlink{spotgac<17>}{10}}\trput{4}}{
    \psset{linestyle=solid}\FAIL{spotgac<18>}\trput{4}
    \psset{linestyle=solid}\FAIL{spotgac<19>}\trput{6}
    \psset{linestyle=solid}\FAIL{spotgac<20>}\trput{8}
   }
   \psset{linestyle=solid}\pstree{\CC{\hyperlink{spotgac<21>}{7}}\trput{5}}{
    \psset{linestyle=solid}\pstree{\CC{\hyperlink{spotgac<22>}{10}}\trput{4}}{
     \psset{linestyle=solid}\FAIL{spotgac<23>}\trput{5}
     \psset{linestyle=solid}\FAIL{spotgac<24>}\trput{6}
     \psset{linestyle=solid}\FAIL{spotgac<25>}\trput{8}
    }
    \psset{linestyle=solid}\pstree{\CC{\hyperlink{spotgac<26>}{10}}\trput{9}}{
     \psset{linestyle=solid}\pstree{\CC{\hyperlink{spotgac<27>}{11}}\trput{5}}{
      \psset{linestyle=solid}\FAIL{spotgac<28>}\trput{6}
      \psset{linestyle=solid}\FAIL{spotgac<29>}\trput{8}
     }
     \psset{linestyle=solid}\FAIL{spotgac<30>}\trput{6}
     \psset{linestyle=solid}\FAIL{spotgac<31>}\trput{8}
    }
   }
  }
 }
 \psset{linestyle=solid}\pstree{\CC{\hyperlink{spotgac<32>}{5}}\trput{4}}{
  \psset{linestyle=solid}\pstree{\CC{\hyperlink{spotgac<33>}{6}}\trput{2}}{
   \psset{linestyle=solid}\FAIL{spotgac<34>}\trput{5}
   \psset{linestyle=solid}\FAIL{spotgac<35>}\trput{6}
  }
  \psset{linestyle=solid}\pstree{\CC{\hyperlink{spotgac<36>}{6}}\trput{6}}{
   \psset{linestyle=solid}\FAIL{spotgac<37>}\trput{2}
   \psset{linestyle=solid}\FAIL{spotgac<38>}\trput{5}
  }
 }
 \psset{linestyle=solid}\pstree{\CC{\hyperlink{spotgac<39>}{5}}\trput{5}}{
  \psset{linestyle=solid}\pstree{\CC{\hyperlink{spotgac<40>}{6}}\trput{2}}{
   \psset{linestyle=solid}\pstree{\CC{\hyperlink{spotgac<41>}{7}}\trput{4}}{
    \psset{linestyle=solid}\FAIL{spotgac<42>}\trput{6}
    \psset{linestyle=solid}\pstree{\CC{\hyperlink{spotgac<43>}{10}}\trput{9}}{
     \psset{linestyle=solid}\FAIL{spotgac<44>}\trput{6}
     \psset{linestyle=solid}\FAIL{spotgac<45>}\trput{8}
    }
   }
   \psset{linestyle=solid}\pstree{\CC{\hyperlink{spotgac<46>}{7}}\trput{6}}{
    \psset{linestyle=solid}\FAIL{spotgac<47>}\trput{4}
    \psset{linestyle=solid}\FAIL{spotgac<48>}\trput{9}
   }
  }
 }
}
}
}
\only<49>{
\scalebox{ 0.40}{\psset{nodesep=0pt,levelsep=36pt}
\pstree{\CC{\hyperlink{spotgac<1>}{4}}}{
 \psset{linestyle=solid}\pstree{\CC{\hyperlink{spotgac<2>}{5}}\trput{2}}{
  \psset{linestyle=solid}\pstree{\CC{\hyperlink{spotgac<3>}{6}}\trput{4}}{
   \psset{linestyle=solid}\pstree{\CC{\hyperlink{spotgac<4>}{7}}\trput{5}}{
    \psset{linestyle=solid}\FAIL{spotgac<5>}\trput{6}
    \psset{linestyle=solid}\pstree{\CC{\hyperlink{spotgac<6>}{12}}\trput{9}}{
     \psset{linestyle=solid}\pstree{\CC{\hyperlink{spotgac<7>}{19}}\trput{6}}{
      \psset{linestyle=solid}\FAIL{spotgac<8>}\trput{6}
      \psset{linestyle=solid}\pstree{\CC{\hyperlink{spotgac<9>}{20}}\trput{9}}{
       \psset{linestyle=solid}\pstree{\CC{\hyperlink{spotgac<10>}{23}}\trput{6}}{
        \psset{linestyle=solid}\FAIL{spotgac<11>}\trput{2}
        \psset{linestyle=solid}\FAIL{spotgac<12>}\trput{8}
       }
       \psset{linestyle=solid}\FAIL{spotgac<13>}\trput{8}
      }
     }
     \psset{linestyle=solid}\FAIL{spotgac<14>}\trput{8}
    }
   }
   \psset{linestyle=solid}\FAIL{spotgac<15>}\trput{6}
  }
  \psset{linestyle=solid}\pstree{\CC{\hyperlink{spotgac<16>}{6}}\trput{6}}{
   \psset{linestyle=solid}\pstree{\CC{\hyperlink{spotgac<17>}{10}}\trput{4}}{
    \psset{linestyle=solid}\FAIL{spotgac<18>}\trput{4}
    \psset{linestyle=solid}\FAIL{spotgac<19>}\trput{6}
    \psset{linestyle=solid}\FAIL{spotgac<20>}\trput{8}
   }
   \psset{linestyle=solid}\pstree{\CC{\hyperlink{spotgac<21>}{7}}\trput{5}}{
    \psset{linestyle=solid}\pstree{\CC{\hyperlink{spotgac<22>}{10}}\trput{4}}{
     \psset{linestyle=solid}\FAIL{spotgac<23>}\trput{5}
     \psset{linestyle=solid}\FAIL{spotgac<24>}\trput{6}
     \psset{linestyle=solid}\FAIL{spotgac<25>}\trput{8}
    }
    \psset{linestyle=solid}\pstree{\CC{\hyperlink{spotgac<26>}{10}}\trput{9}}{
     \psset{linestyle=solid}\pstree{\CC{\hyperlink{spotgac<27>}{11}}\trput{5}}{
      \psset{linestyle=solid}\FAIL{spotgac<28>}\trput{6}
      \psset{linestyle=solid}\FAIL{spotgac<29>}\trput{8}
     }
     \psset{linestyle=solid}\FAIL{spotgac<30>}\trput{6}
     \psset{linestyle=solid}\FAIL{spotgac<31>}\trput{8}
    }
   }
  }
 }
 \psset{linestyle=solid}\pstree{\CC{\hyperlink{spotgac<32>}{5}}\trput{4}}{
  \psset{linestyle=solid}\pstree{\CC{\hyperlink{spotgac<33>}{6}}\trput{2}}{
   \psset{linestyle=solid}\FAIL{spotgac<34>}\trput{5}
   \psset{linestyle=solid}\FAIL{spotgac<35>}\trput{6}
  }
  \psset{linestyle=solid}\pstree{\CC{\hyperlink{spotgac<36>}{6}}\trput{6}}{
   \psset{linestyle=solid}\FAIL{spotgac<37>}\trput{2}
   \psset{linestyle=solid}\FAIL{spotgac<38>}\trput{5}
  }
 }
 \psset{linestyle=solid}\pstree{\CC{\hyperlink{spotgac<39>}{5}}\trput{5}}{
  \psset{linestyle=solid}\pstree{\CC{\hyperlink{spotgac<40>}{6}}\trput{2}}{
   \psset{linestyle=solid}\pstree{\CC{\hyperlink{spotgac<41>}{7}}\trput{4}}{
    \psset{linestyle=solid}\FAIL{spotgac<42>}\trput{6}
    \psset{linestyle=solid}\pstree{\CC{\hyperlink{spotgac<43>}{10}}\trput{9}}{
     \psset{linestyle=solid}\FAIL{spotgac<44>}\trput{6}
     \psset{linestyle=solid}\FAIL{spotgac<45>}\trput{8}
    }
   }
   \psset{linestyle=solid}\pstree{\CC{\hyperlink{spotgac<46>}{7}}\trput{6}}{
    \psset{linestyle=solid}\FAIL{spotgac<47>}\trput{4}
    \psset{linestyle=solid}\FAIL{spotgac<48>}\trput{9}
   }
  }
  \psset{linestyle=solid}\FAIL{spotgac<49>}\trput{4}
 }
}
}
}
\only<50>{
\scalebox{ 0.40}{\psset{nodesep=0pt,levelsep=36pt}
\pstree{\CC{\hyperlink{spotgac<1>}{4}}}{
 \psset{linestyle=solid}\pstree{\CC{\hyperlink{spotgac<2>}{5}}\trput{2}}{
  \psset{linestyle=solid}\pstree{\CC{\hyperlink{spotgac<3>}{6}}\trput{4}}{
   \psset{linestyle=solid}\pstree{\CC{\hyperlink{spotgac<4>}{7}}\trput{5}}{
    \psset{linestyle=solid}\FAIL{spotgac<5>}\trput{6}
    \psset{linestyle=solid}\pstree{\CC{\hyperlink{spotgac<6>}{12}}\trput{9}}{
     \psset{linestyle=solid}\pstree{\CC{\hyperlink{spotgac<7>}{19}}\trput{6}}{
      \psset{linestyle=solid}\FAIL{spotgac<8>}\trput{6}
      \psset{linestyle=solid}\pstree{\CC{\hyperlink{spotgac<9>}{20}}\trput{9}}{
       \psset{linestyle=solid}\pstree{\CC{\hyperlink{spotgac<10>}{23}}\trput{6}}{
        \psset{linestyle=solid}\FAIL{spotgac<11>}\trput{2}
        \psset{linestyle=solid}\FAIL{spotgac<12>}\trput{8}
       }
       \psset{linestyle=solid}\FAIL{spotgac<13>}\trput{8}
      }
     }
     \psset{linestyle=solid}\FAIL{spotgac<14>}\trput{8}
    }
   }
   \psset{linestyle=solid}\FAIL{spotgac<15>}\trput{6}
  }
  \psset{linestyle=solid}\pstree{\CC{\hyperlink{spotgac<16>}{6}}\trput{6}}{
   \psset{linestyle=solid}\pstree{\CC{\hyperlink{spotgac<17>}{10}}\trput{4}}{
    \psset{linestyle=solid}\FAIL{spotgac<18>}\trput{4}
    \psset{linestyle=solid}\FAIL{spotgac<19>}\trput{6}
    \psset{linestyle=solid}\FAIL{spotgac<20>}\trput{8}
   }
   \psset{linestyle=solid}\pstree{\CC{\hyperlink{spotgac<21>}{7}}\trput{5}}{
    \psset{linestyle=solid}\pstree{\CC{\hyperlink{spotgac<22>}{10}}\trput{4}}{
     \psset{linestyle=solid}\FAIL{spotgac<23>}\trput{5}
     \psset{linestyle=solid}\FAIL{spotgac<24>}\trput{6}
     \psset{linestyle=solid}\FAIL{spotgac<25>}\trput{8}
    }
    \psset{linestyle=solid}\pstree{\CC{\hyperlink{spotgac<26>}{10}}\trput{9}}{
     \psset{linestyle=solid}\pstree{\CC{\hyperlink{spotgac<27>}{11}}\trput{5}}{
      \psset{linestyle=solid}\FAIL{spotgac<28>}\trput{6}
      \psset{linestyle=solid}\FAIL{spotgac<29>}\trput{8}
     }
     \psset{linestyle=solid}\FAIL{spotgac<30>}\trput{6}
     \psset{linestyle=solid}\FAIL{spotgac<31>}\trput{8}
    }
   }
  }
 }
 \psset{linestyle=solid}\pstree{\CC{\hyperlink{spotgac<32>}{5}}\trput{4}}{
  \psset{linestyle=solid}\pstree{\CC{\hyperlink{spotgac<33>}{6}}\trput{2}}{
   \psset{linestyle=solid}\FAIL{spotgac<34>}\trput{5}
   \psset{linestyle=solid}\FAIL{spotgac<35>}\trput{6}
  }
  \psset{linestyle=solid}\pstree{\CC{\hyperlink{spotgac<36>}{6}}\trput{6}}{
   \psset{linestyle=solid}\FAIL{spotgac<37>}\trput{2}
   \psset{linestyle=solid}\FAIL{spotgac<38>}\trput{5}
  }
 }
 \psset{linestyle=solid}\pstree{\CC{\hyperlink{spotgac<39>}{5}}\trput{5}}{
  \psset{linestyle=solid}\pstree{\CC{\hyperlink{spotgac<40>}{6}}\trput{2}}{
   \psset{linestyle=solid}\pstree{\CC{\hyperlink{spotgac<41>}{7}}\trput{4}}{
    \psset{linestyle=solid}\FAIL{spotgac<42>}\trput{6}
    \psset{linestyle=solid}\pstree{\CC{\hyperlink{spotgac<43>}{10}}\trput{9}}{
     \psset{linestyle=solid}\FAIL{spotgac<44>}\trput{6}
     \psset{linestyle=solid}\FAIL{spotgac<45>}\trput{8}
    }
   }
   \psset{linestyle=solid}\pstree{\CC{\hyperlink{spotgac<46>}{7}}\trput{6}}{
    \psset{linestyle=solid}\FAIL{spotgac<47>}\trput{4}
    \psset{linestyle=solid}\FAIL{spotgac<48>}\trput{9}
   }
  }
  \psset{linestyle=solid}\FAIL{spotgac<49>}\trput{4}
  \psset{linestyle=solid}\pstree{\CC{\hyperlink{spotgac<50>}{6}}\trput{6}}{
  }
 }
}
}
}
\only<51>{
\scalebox{ 0.40}{\psset{nodesep=0pt,levelsep=36pt}
\pstree{\CC{\hyperlink{spotgac<1>}{4}}}{
 \psset{linestyle=solid}\pstree{\CC{\hyperlink{spotgac<2>}{5}}\trput{2}}{
  \psset{linestyle=solid}\pstree{\CC{\hyperlink{spotgac<3>}{6}}\trput{4}}{
   \psset{linestyle=solid}\pstree{\CC{\hyperlink{spotgac<4>}{7}}\trput{5}}{
    \psset{linestyle=solid}\FAIL{spotgac<5>}\trput{6}
    \psset{linestyle=solid}\pstree{\CC{\hyperlink{spotgac<6>}{12}}\trput{9}}{
     \psset{linestyle=solid}\pstree{\CC{\hyperlink{spotgac<7>}{19}}\trput{6}}{
      \psset{linestyle=solid}\FAIL{spotgac<8>}\trput{6}
      \psset{linestyle=solid}\pstree{\CC{\hyperlink{spotgac<9>}{20}}\trput{9}}{
       \psset{linestyle=solid}\pstree{\CC{\hyperlink{spotgac<10>}{23}}\trput{6}}{
        \psset{linestyle=solid}\FAIL{spotgac<11>}\trput{2}
        \psset{linestyle=solid}\FAIL{spotgac<12>}\trput{8}
       }
       \psset{linestyle=solid}\FAIL{spotgac<13>}\trput{8}
      }
     }
     \psset{linestyle=solid}\FAIL{spotgac<14>}\trput{8}
    }
   }
   \psset{linestyle=solid}\FAIL{spotgac<15>}\trput{6}
  }
  \psset{linestyle=solid}\pstree{\CC{\hyperlink{spotgac<16>}{6}}\trput{6}}{
   \psset{linestyle=solid}\pstree{\CC{\hyperlink{spotgac<17>}{10}}\trput{4}}{
    \psset{linestyle=solid}\FAIL{spotgac<18>}\trput{4}
    \psset{linestyle=solid}\FAIL{spotgac<19>}\trput{6}
    \psset{linestyle=solid}\FAIL{spotgac<20>}\trput{8}
   }
   \psset{linestyle=solid}\pstree{\CC{\hyperlink{spotgac<21>}{7}}\trput{5}}{
    \psset{linestyle=solid}\pstree{\CC{\hyperlink{spotgac<22>}{10}}\trput{4}}{
     \psset{linestyle=solid}\FAIL{spotgac<23>}\trput{5}
     \psset{linestyle=solid}\FAIL{spotgac<24>}\trput{6}
     \psset{linestyle=solid}\FAIL{spotgac<25>}\trput{8}
    }
    \psset{linestyle=solid}\pstree{\CC{\hyperlink{spotgac<26>}{10}}\trput{9}}{
     \psset{linestyle=solid}\pstree{\CC{\hyperlink{spotgac<27>}{11}}\trput{5}}{
      \psset{linestyle=solid}\FAIL{spotgac<28>}\trput{6}
      \psset{linestyle=solid}\FAIL{spotgac<29>}\trput{8}
     }
     \psset{linestyle=solid}\FAIL{spotgac<30>}\trput{6}
     \psset{linestyle=solid}\FAIL{spotgac<31>}\trput{8}
    }
   }
  }
 }
 \psset{linestyle=solid}\pstree{\CC{\hyperlink{spotgac<32>}{5}}\trput{4}}{
  \psset{linestyle=solid}\pstree{\CC{\hyperlink{spotgac<33>}{6}}\trput{2}}{
   \psset{linestyle=solid}\FAIL{spotgac<34>}\trput{5}
   \psset{linestyle=solid}\FAIL{spotgac<35>}\trput{6}
  }
  \psset{linestyle=solid}\pstree{\CC{\hyperlink{spotgac<36>}{6}}\trput{6}}{
   \psset{linestyle=solid}\FAIL{spotgac<37>}\trput{2}
   \psset{linestyle=solid}\FAIL{spotgac<38>}\trput{5}
  }
 }
 \psset{linestyle=solid}\pstree{\CC{\hyperlink{spotgac<39>}{5}}\trput{5}}{
  \psset{linestyle=solid}\pstree{\CC{\hyperlink{spotgac<40>}{6}}\trput{2}}{
   \psset{linestyle=solid}\pstree{\CC{\hyperlink{spotgac<41>}{7}}\trput{4}}{
    \psset{linestyle=solid}\FAIL{spotgac<42>}\trput{6}
    \psset{linestyle=solid}\pstree{\CC{\hyperlink{spotgac<43>}{10}}\trput{9}}{
     \psset{linestyle=solid}\FAIL{spotgac<44>}\trput{6}
     \psset{linestyle=solid}\FAIL{spotgac<45>}\trput{8}
    }
   }
   \psset{linestyle=solid}\pstree{\CC{\hyperlink{spotgac<46>}{7}}\trput{6}}{
    \psset{linestyle=solid}\FAIL{spotgac<47>}\trput{4}
    \psset{linestyle=solid}\FAIL{spotgac<48>}\trput{9}
   }
  }
  \psset{linestyle=solid}\FAIL{spotgac<49>}\trput{4}
  \psset{linestyle=solid}\pstree{\CC{\hyperlink{spotgac<50>}{6}}\trput{6}}{
   \psset{linestyle=solid}\pstree{\CC{\hyperlink{spotgac<51>}{7}}\trput{2}}{
   }
  }
 }
}
}
}
\only<52>{
\scalebox{ 0.40}{\psset{nodesep=0pt,levelsep=36pt}
\pstree{\CC{\hyperlink{spotgac<1>}{4}}}{
 \psset{linestyle=solid}\pstree{\CC{\hyperlink{spotgac<2>}{5}}\trput{2}}{
  \psset{linestyle=solid}\pstree{\CC{\hyperlink{spotgac<3>}{6}}\trput{4}}{
   \psset{linestyle=solid}\pstree{\CC{\hyperlink{spotgac<4>}{7}}\trput{5}}{
    \psset{linestyle=solid}\FAIL{spotgac<5>}\trput{6}
    \psset{linestyle=solid}\pstree{\CC{\hyperlink{spotgac<6>}{12}}\trput{9}}{
     \psset{linestyle=solid}\pstree{\CC{\hyperlink{spotgac<7>}{19}}\trput{6}}{
      \psset{linestyle=solid}\FAIL{spotgac<8>}\trput{6}
      \psset{linestyle=solid}\pstree{\CC{\hyperlink{spotgac<9>}{20}}\trput{9}}{
       \psset{linestyle=solid}\pstree{\CC{\hyperlink{spotgac<10>}{23}}\trput{6}}{
        \psset{linestyle=solid}\FAIL{spotgac<11>}\trput{2}
        \psset{linestyle=solid}\FAIL{spotgac<12>}\trput{8}
       }
       \psset{linestyle=solid}\FAIL{spotgac<13>}\trput{8}
      }
     }
     \psset{linestyle=solid}\FAIL{spotgac<14>}\trput{8}
    }
   }
   \psset{linestyle=solid}\FAIL{spotgac<15>}\trput{6}
  }
  \psset{linestyle=solid}\pstree{\CC{\hyperlink{spotgac<16>}{6}}\trput{6}}{
   \psset{linestyle=solid}\pstree{\CC{\hyperlink{spotgac<17>}{10}}\trput{4}}{
    \psset{linestyle=solid}\FAIL{spotgac<18>}\trput{4}
    \psset{linestyle=solid}\FAIL{spotgac<19>}\trput{6}
    \psset{linestyle=solid}\FAIL{spotgac<20>}\trput{8}
   }
   \psset{linestyle=solid}\pstree{\CC{\hyperlink{spotgac<21>}{7}}\trput{5}}{
    \psset{linestyle=solid}\pstree{\CC{\hyperlink{spotgac<22>}{10}}\trput{4}}{
     \psset{linestyle=solid}\FAIL{spotgac<23>}\trput{5}
     \psset{linestyle=solid}\FAIL{spotgac<24>}\trput{6}
     \psset{linestyle=solid}\FAIL{spotgac<25>}\trput{8}
    }
    \psset{linestyle=solid}\pstree{\CC{\hyperlink{spotgac<26>}{10}}\trput{9}}{
     \psset{linestyle=solid}\pstree{\CC{\hyperlink{spotgac<27>}{11}}\trput{5}}{
      \psset{linestyle=solid}\FAIL{spotgac<28>}\trput{6}
      \psset{linestyle=solid}\FAIL{spotgac<29>}\trput{8}
     }
     \psset{linestyle=solid}\FAIL{spotgac<30>}\trput{6}
     \psset{linestyle=solid}\FAIL{spotgac<31>}\trput{8}
    }
   }
  }
 }
 \psset{linestyle=solid}\pstree{\CC{\hyperlink{spotgac<32>}{5}}\trput{4}}{
  \psset{linestyle=solid}\pstree{\CC{\hyperlink{spotgac<33>}{6}}\trput{2}}{
   \psset{linestyle=solid}\FAIL{spotgac<34>}\trput{5}
   \psset{linestyle=solid}\FAIL{spotgac<35>}\trput{6}
  }
  \psset{linestyle=solid}\pstree{\CC{\hyperlink{spotgac<36>}{6}}\trput{6}}{
   \psset{linestyle=solid}\FAIL{spotgac<37>}\trput{2}
   \psset{linestyle=solid}\FAIL{spotgac<38>}\trput{5}
  }
 }
 \psset{linestyle=solid}\pstree{\CC{\hyperlink{spotgac<39>}{5}}\trput{5}}{
  \psset{linestyle=solid}\pstree{\CC{\hyperlink{spotgac<40>}{6}}\trput{2}}{
   \psset{linestyle=solid}\pstree{\CC{\hyperlink{spotgac<41>}{7}}\trput{4}}{
    \psset{linestyle=solid}\FAIL{spotgac<42>}\trput{6}
    \psset{linestyle=solid}\pstree{\CC{\hyperlink{spotgac<43>}{10}}\trput{9}}{
     \psset{linestyle=solid}\FAIL{spotgac<44>}\trput{6}
     \psset{linestyle=solid}\FAIL{spotgac<45>}\trput{8}
    }
   }
   \psset{linestyle=solid}\pstree{\CC{\hyperlink{spotgac<46>}{7}}\trput{6}}{
    \psset{linestyle=solid}\FAIL{spotgac<47>}\trput{4}
    \psset{linestyle=solid}\FAIL{spotgac<48>}\trput{9}
   }
  }
  \psset{linestyle=solid}\FAIL{spotgac<49>}\trput{4}
  \psset{linestyle=solid}\pstree{\CC{\hyperlink{spotgac<50>}{6}}\trput{6}}{
   \psset{linestyle=solid}\pstree{\CC{\hyperlink{spotgac<51>}{7}}\trput{2}}{
    \psset{linestyle=solid}\FAIL{spotgac<52>}\trput{4}
   }
  }
 }
}
}
}
\only<53>{
\scalebox{ 0.40}{\psset{nodesep=0pt,levelsep=36pt}
\pstree{\CC{\hyperlink{spotgac<1>}{4}}}{
 \psset{linestyle=solid}\pstree{\CC{\hyperlink{spotgac<2>}{5}}\trput{2}}{
  \psset{linestyle=solid}\pstree{\CC{\hyperlink{spotgac<3>}{6}}\trput{4}}{
   \psset{linestyle=solid}\pstree{\CC{\hyperlink{spotgac<4>}{7}}\trput{5}}{
    \psset{linestyle=solid}\FAIL{spotgac<5>}\trput{6}
    \psset{linestyle=solid}\pstree{\CC{\hyperlink{spotgac<6>}{12}}\trput{9}}{
     \psset{linestyle=solid}\pstree{\CC{\hyperlink{spotgac<7>}{19}}\trput{6}}{
      \psset{linestyle=solid}\FAIL{spotgac<8>}\trput{6}
      \psset{linestyle=solid}\pstree{\CC{\hyperlink{spotgac<9>}{20}}\trput{9}}{
       \psset{linestyle=solid}\pstree{\CC{\hyperlink{spotgac<10>}{23}}\trput{6}}{
        \psset{linestyle=solid}\FAIL{spotgac<11>}\trput{2}
        \psset{linestyle=solid}\FAIL{spotgac<12>}\trput{8}
       }
       \psset{linestyle=solid}\FAIL{spotgac<13>}\trput{8}
      }
     }
     \psset{linestyle=solid}\FAIL{spotgac<14>}\trput{8}
    }
   }
   \psset{linestyle=solid}\FAIL{spotgac<15>}\trput{6}
  }
  \psset{linestyle=solid}\pstree{\CC{\hyperlink{spotgac<16>}{6}}\trput{6}}{
   \psset{linestyle=solid}\pstree{\CC{\hyperlink{spotgac<17>}{10}}\trput{4}}{
    \psset{linestyle=solid}\FAIL{spotgac<18>}\trput{4}
    \psset{linestyle=solid}\FAIL{spotgac<19>}\trput{6}
    \psset{linestyle=solid}\FAIL{spotgac<20>}\trput{8}
   }
   \psset{linestyle=solid}\pstree{\CC{\hyperlink{spotgac<21>}{7}}\trput{5}}{
    \psset{linestyle=solid}\pstree{\CC{\hyperlink{spotgac<22>}{10}}\trput{4}}{
     \psset{linestyle=solid}\FAIL{spotgac<23>}\trput{5}
     \psset{linestyle=solid}\FAIL{spotgac<24>}\trput{6}
     \psset{linestyle=solid}\FAIL{spotgac<25>}\trput{8}
    }
    \psset{linestyle=solid}\pstree{\CC{\hyperlink{spotgac<26>}{10}}\trput{9}}{
     \psset{linestyle=solid}\pstree{\CC{\hyperlink{spotgac<27>}{11}}\trput{5}}{
      \psset{linestyle=solid}\FAIL{spotgac<28>}\trput{6}
      \psset{linestyle=solid}\FAIL{spotgac<29>}\trput{8}
     }
     \psset{linestyle=solid}\FAIL{spotgac<30>}\trput{6}
     \psset{linestyle=solid}\FAIL{spotgac<31>}\trput{8}
    }
   }
  }
 }
 \psset{linestyle=solid}\pstree{\CC{\hyperlink{spotgac<32>}{5}}\trput{4}}{
  \psset{linestyle=solid}\pstree{\CC{\hyperlink{spotgac<33>}{6}}\trput{2}}{
   \psset{linestyle=solid}\FAIL{spotgac<34>}\trput{5}
   \psset{linestyle=solid}\FAIL{spotgac<35>}\trput{6}
  }
  \psset{linestyle=solid}\pstree{\CC{\hyperlink{spotgac<36>}{6}}\trput{6}}{
   \psset{linestyle=solid}\FAIL{spotgac<37>}\trput{2}
   \psset{linestyle=solid}\FAIL{spotgac<38>}\trput{5}
  }
 }
 \psset{linestyle=solid}\pstree{\CC{\hyperlink{spotgac<39>}{5}}\trput{5}}{
  \psset{linestyle=solid}\pstree{\CC{\hyperlink{spotgac<40>}{6}}\trput{2}}{
   \psset{linestyle=solid}\pstree{\CC{\hyperlink{spotgac<41>}{7}}\trput{4}}{
    \psset{linestyle=solid}\FAIL{spotgac<42>}\trput{6}
    \psset{linestyle=solid}\pstree{\CC{\hyperlink{spotgac<43>}{10}}\trput{9}}{
     \psset{linestyle=solid}\FAIL{spotgac<44>}\trput{6}
     \psset{linestyle=solid}\FAIL{spotgac<45>}\trput{8}
    }
   }
   \psset{linestyle=solid}\pstree{\CC{\hyperlink{spotgac<46>}{7}}\trput{6}}{
    \psset{linestyle=solid}\FAIL{spotgac<47>}\trput{4}
    \psset{linestyle=solid}\FAIL{spotgac<48>}\trput{9}
   }
  }
  \psset{linestyle=solid}\FAIL{spotgac<49>}\trput{4}
  \psset{linestyle=solid}\pstree{\CC{\hyperlink{spotgac<50>}{6}}\trput{6}}{
   \psset{linestyle=solid}\pstree{\CC{\hyperlink{spotgac<51>}{7}}\trput{2}}{
    \psset{linestyle=solid}\FAIL{spotgac<52>}\trput{4}
    \psset{linestyle=solid}\FAIL{spotgac<53>}\trput{9}
   }
  }
 }
}
}
}
\only<54>{
\scalebox{ 0.40}{\psset{nodesep=0pt,levelsep=36pt}
\pstree{\CC{\hyperlink{spotgac<1>}{4}}}{
 \psset{linestyle=solid}\pstree{\CC{\hyperlink{spotgac<2>}{5}}\trput{2}}{
  \psset{linestyle=solid}\pstree{\CC{\hyperlink{spotgac<3>}{6}}\trput{4}}{
   \psset{linestyle=solid}\pstree{\CC{\hyperlink{spotgac<4>}{7}}\trput{5}}{
    \psset{linestyle=solid}\FAIL{spotgac<5>}\trput{6}
    \psset{linestyle=solid}\pstree{\CC{\hyperlink{spotgac<6>}{12}}\trput{9}}{
     \psset{linestyle=solid}\pstree{\CC{\hyperlink{spotgac<7>}{19}}\trput{6}}{
      \psset{linestyle=solid}\FAIL{spotgac<8>}\trput{6}
      \psset{linestyle=solid}\pstree{\CC{\hyperlink{spotgac<9>}{20}}\trput{9}}{
       \psset{linestyle=solid}\pstree{\CC{\hyperlink{spotgac<10>}{23}}\trput{6}}{
        \psset{linestyle=solid}\FAIL{spotgac<11>}\trput{2}
        \psset{linestyle=solid}\FAIL{spotgac<12>}\trput{8}
       }
       \psset{linestyle=solid}\FAIL{spotgac<13>}\trput{8}
      }
     }
     \psset{linestyle=solid}\FAIL{spotgac<14>}\trput{8}
    }
   }
   \psset{linestyle=solid}\FAIL{spotgac<15>}\trput{6}
  }
  \psset{linestyle=solid}\pstree{\CC{\hyperlink{spotgac<16>}{6}}\trput{6}}{
   \psset{linestyle=solid}\pstree{\CC{\hyperlink{spotgac<17>}{10}}\trput{4}}{
    \psset{linestyle=solid}\FAIL{spotgac<18>}\trput{4}
    \psset{linestyle=solid}\FAIL{spotgac<19>}\trput{6}
    \psset{linestyle=solid}\FAIL{spotgac<20>}\trput{8}
   }
   \psset{linestyle=solid}\pstree{\CC{\hyperlink{spotgac<21>}{7}}\trput{5}}{
    \psset{linestyle=solid}\pstree{\CC{\hyperlink{spotgac<22>}{10}}\trput{4}}{
     \psset{linestyle=solid}\FAIL{spotgac<23>}\trput{5}
     \psset{linestyle=solid}\FAIL{spotgac<24>}\trput{6}
     \psset{linestyle=solid}\FAIL{spotgac<25>}\trput{8}
    }
    \psset{linestyle=solid}\pstree{\CC{\hyperlink{spotgac<26>}{10}}\trput{9}}{
     \psset{linestyle=solid}\pstree{\CC{\hyperlink{spotgac<27>}{11}}\trput{5}}{
      \psset{linestyle=solid}\FAIL{spotgac<28>}\trput{6}
      \psset{linestyle=solid}\FAIL{spotgac<29>}\trput{8}
     }
     \psset{linestyle=solid}\FAIL{spotgac<30>}\trput{6}
     \psset{linestyle=solid}\FAIL{spotgac<31>}\trput{8}
    }
   }
  }
 }
 \psset{linestyle=solid}\pstree{\CC{\hyperlink{spotgac<32>}{5}}\trput{4}}{
  \psset{linestyle=solid}\pstree{\CC{\hyperlink{spotgac<33>}{6}}\trput{2}}{
   \psset{linestyle=solid}\FAIL{spotgac<34>}\trput{5}
   \psset{linestyle=solid}\FAIL{spotgac<35>}\trput{6}
  }
  \psset{linestyle=solid}\pstree{\CC{\hyperlink{spotgac<36>}{6}}\trput{6}}{
   \psset{linestyle=solid}\FAIL{spotgac<37>}\trput{2}
   \psset{linestyle=solid}\FAIL{spotgac<38>}\trput{5}
  }
 }
 \psset{linestyle=solid}\pstree{\CC{\hyperlink{spotgac<39>}{5}}\trput{5}}{
  \psset{linestyle=solid}\pstree{\CC{\hyperlink{spotgac<40>}{6}}\trput{2}}{
   \psset{linestyle=solid}\pstree{\CC{\hyperlink{spotgac<41>}{7}}\trput{4}}{
    \psset{linestyle=solid}\FAIL{spotgac<42>}\trput{6}
    \psset{linestyle=solid}\pstree{\CC{\hyperlink{spotgac<43>}{10}}\trput{9}}{
     \psset{linestyle=solid}\FAIL{spotgac<44>}\trput{6}
     \psset{linestyle=solid}\FAIL{spotgac<45>}\trput{8}
    }
   }
   \psset{linestyle=solid}\pstree{\CC{\hyperlink{spotgac<46>}{7}}\trput{6}}{
    \psset{linestyle=solid}\FAIL{spotgac<47>}\trput{4}
    \psset{linestyle=solid}\FAIL{spotgac<48>}\trput{9}
   }
  }
  \psset{linestyle=solid}\FAIL{spotgac<49>}\trput{4}
  \psset{linestyle=solid}\pstree{\CC{\hyperlink{spotgac<50>}{6}}\trput{6}}{
   \psset{linestyle=solid}\pstree{\CC{\hyperlink{spotgac<51>}{7}}\trput{2}}{
    \psset{linestyle=solid}\FAIL{spotgac<52>}\trput{4}
    \psset{linestyle=solid}\FAIL{spotgac<53>}\trput{9}
   }
   \psset{linestyle=solid}\pstree{\CC{\hyperlink{spotgac<54>}{7}}\trput{4}}{
   }
  }
 }
}
}
}
\only<55>{
\scalebox{ 0.40}{\psset{nodesep=0pt,levelsep=36pt}
\pstree{\CC{\hyperlink{spotgac<1>}{4}}}{
 \psset{linestyle=solid}\pstree{\CC{\hyperlink{spotgac<2>}{5}}\trput{2}}{
  \psset{linestyle=solid}\pstree{\CC{\hyperlink{spotgac<3>}{6}}\trput{4}}{
   \psset{linestyle=solid}\pstree{\CC{\hyperlink{spotgac<4>}{7}}\trput{5}}{
    \psset{linestyle=solid}\FAIL{spotgac<5>}\trput{6}
    \psset{linestyle=solid}\pstree{\CC{\hyperlink{spotgac<6>}{12}}\trput{9}}{
     \psset{linestyle=solid}\pstree{\CC{\hyperlink{spotgac<7>}{19}}\trput{6}}{
      \psset{linestyle=solid}\FAIL{spotgac<8>}\trput{6}
      \psset{linestyle=solid}\pstree{\CC{\hyperlink{spotgac<9>}{20}}\trput{9}}{
       \psset{linestyle=solid}\pstree{\CC{\hyperlink{spotgac<10>}{23}}\trput{6}}{
        \psset{linestyle=solid}\FAIL{spotgac<11>}\trput{2}
        \psset{linestyle=solid}\FAIL{spotgac<12>}\trput{8}
       }
       \psset{linestyle=solid}\FAIL{spotgac<13>}\trput{8}
      }
     }
     \psset{linestyle=solid}\FAIL{spotgac<14>}\trput{8}
    }
   }
   \psset{linestyle=solid}\FAIL{spotgac<15>}\trput{6}
  }
  \psset{linestyle=solid}\pstree{\CC{\hyperlink{spotgac<16>}{6}}\trput{6}}{
   \psset{linestyle=solid}\pstree{\CC{\hyperlink{spotgac<17>}{10}}\trput{4}}{
    \psset{linestyle=solid}\FAIL{spotgac<18>}\trput{4}
    \psset{linestyle=solid}\FAIL{spotgac<19>}\trput{6}
    \psset{linestyle=solid}\FAIL{spotgac<20>}\trput{8}
   }
   \psset{linestyle=solid}\pstree{\CC{\hyperlink{spotgac<21>}{7}}\trput{5}}{
    \psset{linestyle=solid}\pstree{\CC{\hyperlink{spotgac<22>}{10}}\trput{4}}{
     \psset{linestyle=solid}\FAIL{spotgac<23>}\trput{5}
     \psset{linestyle=solid}\FAIL{spotgac<24>}\trput{6}
     \psset{linestyle=solid}\FAIL{spotgac<25>}\trput{8}
    }
    \psset{linestyle=solid}\pstree{\CC{\hyperlink{spotgac<26>}{10}}\trput{9}}{
     \psset{linestyle=solid}\pstree{\CC{\hyperlink{spotgac<27>}{11}}\trput{5}}{
      \psset{linestyle=solid}\FAIL{spotgac<28>}\trput{6}
      \psset{linestyle=solid}\FAIL{spotgac<29>}\trput{8}
     }
     \psset{linestyle=solid}\FAIL{spotgac<30>}\trput{6}
     \psset{linestyle=solid}\FAIL{spotgac<31>}\trput{8}
    }
   }
  }
 }
 \psset{linestyle=solid}\pstree{\CC{\hyperlink{spotgac<32>}{5}}\trput{4}}{
  \psset{linestyle=solid}\pstree{\CC{\hyperlink{spotgac<33>}{6}}\trput{2}}{
   \psset{linestyle=solid}\FAIL{spotgac<34>}\trput{5}
   \psset{linestyle=solid}\FAIL{spotgac<35>}\trput{6}
  }
  \psset{linestyle=solid}\pstree{\CC{\hyperlink{spotgac<36>}{6}}\trput{6}}{
   \psset{linestyle=solid}\FAIL{spotgac<37>}\trput{2}
   \psset{linestyle=solid}\FAIL{spotgac<38>}\trput{5}
  }
 }
 \psset{linestyle=solid}\pstree{\CC{\hyperlink{spotgac<39>}{5}}\trput{5}}{
  \psset{linestyle=solid}\pstree{\CC{\hyperlink{spotgac<40>}{6}}\trput{2}}{
   \psset{linestyle=solid}\pstree{\CC{\hyperlink{spotgac<41>}{7}}\trput{4}}{
    \psset{linestyle=solid}\FAIL{spotgac<42>}\trput{6}
    \psset{linestyle=solid}\pstree{\CC{\hyperlink{spotgac<43>}{10}}\trput{9}}{
     \psset{linestyle=solid}\FAIL{spotgac<44>}\trput{6}
     \psset{linestyle=solid}\FAIL{spotgac<45>}\trput{8}
    }
   }
   \psset{linestyle=solid}\pstree{\CC{\hyperlink{spotgac<46>}{7}}\trput{6}}{
    \psset{linestyle=solid}\FAIL{spotgac<47>}\trput{4}
    \psset{linestyle=solid}\FAIL{spotgac<48>}\trput{9}
   }
  }
  \psset{linestyle=solid}\FAIL{spotgac<49>}\trput{4}
  \psset{linestyle=solid}\pstree{\CC{\hyperlink{spotgac<50>}{6}}\trput{6}}{
   \psset{linestyle=solid}\pstree{\CC{\hyperlink{spotgac<51>}{7}}\trput{2}}{
    \psset{linestyle=solid}\FAIL{spotgac<52>}\trput{4}
    \psset{linestyle=solid}\FAIL{spotgac<53>}\trput{9}
   }
   \psset{linestyle=solid}\pstree{\CC{\hyperlink{spotgac<54>}{7}}\trput{4}}{
    \psset{linestyle=solid}\pstree{\CC{\hyperlink{spotgac<55>}{10}}\trput{2}}{
    }
   }
  }
 }
}
}
}
\only<56>{
\scalebox{ 0.40}{\psset{nodesep=0pt,levelsep=36pt}
\pstree{\CC{\hyperlink{spotgac<1>}{4}}}{
 \psset{linestyle=solid}\pstree{\CC{\hyperlink{spotgac<2>}{5}}\trput{2}}{
  \psset{linestyle=solid}\pstree{\CC{\hyperlink{spotgac<3>}{6}}\trput{4}}{
   \psset{linestyle=solid}\pstree{\CC{\hyperlink{spotgac<4>}{7}}\trput{5}}{
    \psset{linestyle=solid}\FAIL{spotgac<5>}\trput{6}
    \psset{linestyle=solid}\pstree{\CC{\hyperlink{spotgac<6>}{12}}\trput{9}}{
     \psset{linestyle=solid}\pstree{\CC{\hyperlink{spotgac<7>}{19}}\trput{6}}{
      \psset{linestyle=solid}\FAIL{spotgac<8>}\trput{6}
      \psset{linestyle=solid}\pstree{\CC{\hyperlink{spotgac<9>}{20}}\trput{9}}{
       \psset{linestyle=solid}\pstree{\CC{\hyperlink{spotgac<10>}{23}}\trput{6}}{
        \psset{linestyle=solid}\FAIL{spotgac<11>}\trput{2}
        \psset{linestyle=solid}\FAIL{spotgac<12>}\trput{8}
       }
       \psset{linestyle=solid}\FAIL{spotgac<13>}\trput{8}
      }
     }
     \psset{linestyle=solid}\FAIL{spotgac<14>}\trput{8}
    }
   }
   \psset{linestyle=solid}\FAIL{spotgac<15>}\trput{6}
  }
  \psset{linestyle=solid}\pstree{\CC{\hyperlink{spotgac<16>}{6}}\trput{6}}{
   \psset{linestyle=solid}\pstree{\CC{\hyperlink{spotgac<17>}{10}}\trput{4}}{
    \psset{linestyle=solid}\FAIL{spotgac<18>}\trput{4}
    \psset{linestyle=solid}\FAIL{spotgac<19>}\trput{6}
    \psset{linestyle=solid}\FAIL{spotgac<20>}\trput{8}
   }
   \psset{linestyle=solid}\pstree{\CC{\hyperlink{spotgac<21>}{7}}\trput{5}}{
    \psset{linestyle=solid}\pstree{\CC{\hyperlink{spotgac<22>}{10}}\trput{4}}{
     \psset{linestyle=solid}\FAIL{spotgac<23>}\trput{5}
     \psset{linestyle=solid}\FAIL{spotgac<24>}\trput{6}
     \psset{linestyle=solid}\FAIL{spotgac<25>}\trput{8}
    }
    \psset{linestyle=solid}\pstree{\CC{\hyperlink{spotgac<26>}{10}}\trput{9}}{
     \psset{linestyle=solid}\pstree{\CC{\hyperlink{spotgac<27>}{11}}\trput{5}}{
      \psset{linestyle=solid}\FAIL{spotgac<28>}\trput{6}
      \psset{linestyle=solid}\FAIL{spotgac<29>}\trput{8}
     }
     \psset{linestyle=solid}\FAIL{spotgac<30>}\trput{6}
     \psset{linestyle=solid}\FAIL{spotgac<31>}\trput{8}
    }
   }
  }
 }
 \psset{linestyle=solid}\pstree{\CC{\hyperlink{spotgac<32>}{5}}\trput{4}}{
  \psset{linestyle=solid}\pstree{\CC{\hyperlink{spotgac<33>}{6}}\trput{2}}{
   \psset{linestyle=solid}\FAIL{spotgac<34>}\trput{5}
   \psset{linestyle=solid}\FAIL{spotgac<35>}\trput{6}
  }
  \psset{linestyle=solid}\pstree{\CC{\hyperlink{spotgac<36>}{6}}\trput{6}}{
   \psset{linestyle=solid}\FAIL{spotgac<37>}\trput{2}
   \psset{linestyle=solid}\FAIL{spotgac<38>}\trput{5}
  }
 }
 \psset{linestyle=solid}\pstree{\CC{\hyperlink{spotgac<39>}{5}}\trput{5}}{
  \psset{linestyle=solid}\pstree{\CC{\hyperlink{spotgac<40>}{6}}\trput{2}}{
   \psset{linestyle=solid}\pstree{\CC{\hyperlink{spotgac<41>}{7}}\trput{4}}{
    \psset{linestyle=solid}\FAIL{spotgac<42>}\trput{6}
    \psset{linestyle=solid}\pstree{\CC{\hyperlink{spotgac<43>}{10}}\trput{9}}{
     \psset{linestyle=solid}\FAIL{spotgac<44>}\trput{6}
     \psset{linestyle=solid}\FAIL{spotgac<45>}\trput{8}
    }
   }
   \psset{linestyle=solid}\pstree{\CC{\hyperlink{spotgac<46>}{7}}\trput{6}}{
    \psset{linestyle=solid}\FAIL{spotgac<47>}\trput{4}
    \psset{linestyle=solid}\FAIL{spotgac<48>}\trput{9}
   }
  }
  \psset{linestyle=solid}\FAIL{spotgac<49>}\trput{4}
  \psset{linestyle=solid}\pstree{\CC{\hyperlink{spotgac<50>}{6}}\trput{6}}{
   \psset{linestyle=solid}\pstree{\CC{\hyperlink{spotgac<51>}{7}}\trput{2}}{
    \psset{linestyle=solid}\FAIL{spotgac<52>}\trput{4}
    \psset{linestyle=solid}\FAIL{spotgac<53>}\trput{9}
   }
   \psset{linestyle=solid}\pstree{\CC{\hyperlink{spotgac<54>}{7}}\trput{4}}{
    \psset{linestyle=solid}\pstree{\CC{\hyperlink{spotgac<55>}{10}}\trput{2}}{
     \psset{linestyle=solid}\pstree{\CC{\hyperlink{spotgac<56>}{12}}\trput{6}}{
     }
    }
   }
  }
 }
}
}
}
\only<57>{
\scalebox{ 0.40}{\psset{nodesep=0pt,levelsep=36pt}
\pstree{\CC{\hyperlink{spotgac<1>}{4}}}{
 \psset{linestyle=solid}\pstree{\CC{\hyperlink{spotgac<2>}{5}}\trput{2}}{
  \psset{linestyle=solid}\pstree{\CC{\hyperlink{spotgac<3>}{6}}\trput{4}}{
   \psset{linestyle=solid}\pstree{\CC{\hyperlink{spotgac<4>}{7}}\trput{5}}{
    \psset{linestyle=solid}\FAIL{spotgac<5>}\trput{6}
    \psset{linestyle=solid}\pstree{\CC{\hyperlink{spotgac<6>}{12}}\trput{9}}{
     \psset{linestyle=solid}\pstree{\CC{\hyperlink{spotgac<7>}{19}}\trput{6}}{
      \psset{linestyle=solid}\FAIL{spotgac<8>}\trput{6}
      \psset{linestyle=solid}\pstree{\CC{\hyperlink{spotgac<9>}{20}}\trput{9}}{
       \psset{linestyle=solid}\pstree{\CC{\hyperlink{spotgac<10>}{23}}\trput{6}}{
        \psset{linestyle=solid}\FAIL{spotgac<11>}\trput{2}
        \psset{linestyle=solid}\FAIL{spotgac<12>}\trput{8}
       }
       \psset{linestyle=solid}\FAIL{spotgac<13>}\trput{8}
      }
     }
     \psset{linestyle=solid}\FAIL{spotgac<14>}\trput{8}
    }
   }
   \psset{linestyle=solid}\FAIL{spotgac<15>}\trput{6}
  }
  \psset{linestyle=solid}\pstree{\CC{\hyperlink{spotgac<16>}{6}}\trput{6}}{
   \psset{linestyle=solid}\pstree{\CC{\hyperlink{spotgac<17>}{10}}\trput{4}}{
    \psset{linestyle=solid}\FAIL{spotgac<18>}\trput{4}
    \psset{linestyle=solid}\FAIL{spotgac<19>}\trput{6}
    \psset{linestyle=solid}\FAIL{spotgac<20>}\trput{8}
   }
   \psset{linestyle=solid}\pstree{\CC{\hyperlink{spotgac<21>}{7}}\trput{5}}{
    \psset{linestyle=solid}\pstree{\CC{\hyperlink{spotgac<22>}{10}}\trput{4}}{
     \psset{linestyle=solid}\FAIL{spotgac<23>}\trput{5}
     \psset{linestyle=solid}\FAIL{spotgac<24>}\trput{6}
     \psset{linestyle=solid}\FAIL{spotgac<25>}\trput{8}
    }
    \psset{linestyle=solid}\pstree{\CC{\hyperlink{spotgac<26>}{10}}\trput{9}}{
     \psset{linestyle=solid}\pstree{\CC{\hyperlink{spotgac<27>}{11}}\trput{5}}{
      \psset{linestyle=solid}\FAIL{spotgac<28>}\trput{6}
      \psset{linestyle=solid}\FAIL{spotgac<29>}\trput{8}
     }
     \psset{linestyle=solid}\FAIL{spotgac<30>}\trput{6}
     \psset{linestyle=solid}\FAIL{spotgac<31>}\trput{8}
    }
   }
  }
 }
 \psset{linestyle=solid}\pstree{\CC{\hyperlink{spotgac<32>}{5}}\trput{4}}{
  \psset{linestyle=solid}\pstree{\CC{\hyperlink{spotgac<33>}{6}}\trput{2}}{
   \psset{linestyle=solid}\FAIL{spotgac<34>}\trput{5}
   \psset{linestyle=solid}\FAIL{spotgac<35>}\trput{6}
  }
  \psset{linestyle=solid}\pstree{\CC{\hyperlink{spotgac<36>}{6}}\trput{6}}{
   \psset{linestyle=solid}\FAIL{spotgac<37>}\trput{2}
   \psset{linestyle=solid}\FAIL{spotgac<38>}\trput{5}
  }
 }
 \psset{linestyle=solid}\pstree{\CC{\hyperlink{spotgac<39>}{5}}\trput{5}}{
  \psset{linestyle=solid}\pstree{\CC{\hyperlink{spotgac<40>}{6}}\trput{2}}{
   \psset{linestyle=solid}\pstree{\CC{\hyperlink{spotgac<41>}{7}}\trput{4}}{
    \psset{linestyle=solid}\FAIL{spotgac<42>}\trput{6}
    \psset{linestyle=solid}\pstree{\CC{\hyperlink{spotgac<43>}{10}}\trput{9}}{
     \psset{linestyle=solid}\FAIL{spotgac<44>}\trput{6}
     \psset{linestyle=solid}\FAIL{spotgac<45>}\trput{8}
    }
   }
   \psset{linestyle=solid}\pstree{\CC{\hyperlink{spotgac<46>}{7}}\trput{6}}{
    \psset{linestyle=solid}\FAIL{spotgac<47>}\trput{4}
    \psset{linestyle=solid}\FAIL{spotgac<48>}\trput{9}
   }
  }
  \psset{linestyle=solid}\FAIL{spotgac<49>}\trput{4}
  \psset{linestyle=solid}\pstree{\CC{\hyperlink{spotgac<50>}{6}}\trput{6}}{
   \psset{linestyle=solid}\pstree{\CC{\hyperlink{spotgac<51>}{7}}\trput{2}}{
    \psset{linestyle=solid}\FAIL{spotgac<52>}\trput{4}
    \psset{linestyle=solid}\FAIL{spotgac<53>}\trput{9}
   }
   \psset{linestyle=solid}\pstree{\CC{\hyperlink{spotgac<54>}{7}}\trput{4}}{
    \psset{linestyle=solid}\pstree{\CC{\hyperlink{spotgac<55>}{10}}\trput{2}}{
     \psset{linestyle=solid}\pstree{\CC{\hyperlink{spotgac<56>}{12}}\trput{6}}{
      \psset{linestyle=solid}\FAIL{spotgac<57>}\trput{2}
     }
    }
   }
  }
 }
}
}
}
\only<58>{
\scalebox{ 0.40}{\psset{nodesep=0pt,levelsep=36pt}
\pstree{\CC{\hyperlink{spotgac<1>}{4}}}{
 \psset{linestyle=solid}\pstree{\CC{\hyperlink{spotgac<2>}{5}}\trput{2}}{
  \psset{linestyle=solid}\pstree{\CC{\hyperlink{spotgac<3>}{6}}\trput{4}}{
   \psset{linestyle=solid}\pstree{\CC{\hyperlink{spotgac<4>}{7}}\trput{5}}{
    \psset{linestyle=solid}\FAIL{spotgac<5>}\trput{6}
    \psset{linestyle=solid}\pstree{\CC{\hyperlink{spotgac<6>}{12}}\trput{9}}{
     \psset{linestyle=solid}\pstree{\CC{\hyperlink{spotgac<7>}{19}}\trput{6}}{
      \psset{linestyle=solid}\FAIL{spotgac<8>}\trput{6}
      \psset{linestyle=solid}\pstree{\CC{\hyperlink{spotgac<9>}{20}}\trput{9}}{
       \psset{linestyle=solid}\pstree{\CC{\hyperlink{spotgac<10>}{23}}\trput{6}}{
        \psset{linestyle=solid}\FAIL{spotgac<11>}\trput{2}
        \psset{linestyle=solid}\FAIL{spotgac<12>}\trput{8}
       }
       \psset{linestyle=solid}\FAIL{spotgac<13>}\trput{8}
      }
     }
     \psset{linestyle=solid}\FAIL{spotgac<14>}\trput{8}
    }
   }
   \psset{linestyle=solid}\FAIL{spotgac<15>}\trput{6}
  }
  \psset{linestyle=solid}\pstree{\CC{\hyperlink{spotgac<16>}{6}}\trput{6}}{
   \psset{linestyle=solid}\pstree{\CC{\hyperlink{spotgac<17>}{10}}\trput{4}}{
    \psset{linestyle=solid}\FAIL{spotgac<18>}\trput{4}
    \psset{linestyle=solid}\FAIL{spotgac<19>}\trput{6}
    \psset{linestyle=solid}\FAIL{spotgac<20>}\trput{8}
   }
   \psset{linestyle=solid}\pstree{\CC{\hyperlink{spotgac<21>}{7}}\trput{5}}{
    \psset{linestyle=solid}\pstree{\CC{\hyperlink{spotgac<22>}{10}}\trput{4}}{
     \psset{linestyle=solid}\FAIL{spotgac<23>}\trput{5}
     \psset{linestyle=solid}\FAIL{spotgac<24>}\trput{6}
     \psset{linestyle=solid}\FAIL{spotgac<25>}\trput{8}
    }
    \psset{linestyle=solid}\pstree{\CC{\hyperlink{spotgac<26>}{10}}\trput{9}}{
     \psset{linestyle=solid}\pstree{\CC{\hyperlink{spotgac<27>}{11}}\trput{5}}{
      \psset{linestyle=solid}\FAIL{spotgac<28>}\trput{6}
      \psset{linestyle=solid}\FAIL{spotgac<29>}\trput{8}
     }
     \psset{linestyle=solid}\FAIL{spotgac<30>}\trput{6}
     \psset{linestyle=solid}\FAIL{spotgac<31>}\trput{8}
    }
   }
  }
 }
 \psset{linestyle=solid}\pstree{\CC{\hyperlink{spotgac<32>}{5}}\trput{4}}{
  \psset{linestyle=solid}\pstree{\CC{\hyperlink{spotgac<33>}{6}}\trput{2}}{
   \psset{linestyle=solid}\FAIL{spotgac<34>}\trput{5}
   \psset{linestyle=solid}\FAIL{spotgac<35>}\trput{6}
  }
  \psset{linestyle=solid}\pstree{\CC{\hyperlink{spotgac<36>}{6}}\trput{6}}{
   \psset{linestyle=solid}\FAIL{spotgac<37>}\trput{2}
   \psset{linestyle=solid}\FAIL{spotgac<38>}\trput{5}
  }
 }
 \psset{linestyle=solid}\pstree{\CC{\hyperlink{spotgac<39>}{5}}\trput{5}}{
  \psset{linestyle=solid}\pstree{\CC{\hyperlink{spotgac<40>}{6}}\trput{2}}{
   \psset{linestyle=solid}\pstree{\CC{\hyperlink{spotgac<41>}{7}}\trput{4}}{
    \psset{linestyle=solid}\FAIL{spotgac<42>}\trput{6}
    \psset{linestyle=solid}\pstree{\CC{\hyperlink{spotgac<43>}{10}}\trput{9}}{
     \psset{linestyle=solid}\FAIL{spotgac<44>}\trput{6}
     \psset{linestyle=solid}\FAIL{spotgac<45>}\trput{8}
    }
   }
   \psset{linestyle=solid}\pstree{\CC{\hyperlink{spotgac<46>}{7}}\trput{6}}{
    \psset{linestyle=solid}\FAIL{spotgac<47>}\trput{4}
    \psset{linestyle=solid}\FAIL{spotgac<48>}\trput{9}
   }
  }
  \psset{linestyle=solid}\FAIL{spotgac<49>}\trput{4}
  \psset{linestyle=solid}\pstree{\CC{\hyperlink{spotgac<50>}{6}}\trput{6}}{
   \psset{linestyle=solid}\pstree{\CC{\hyperlink{spotgac<51>}{7}}\trput{2}}{
    \psset{linestyle=solid}\FAIL{spotgac<52>}\trput{4}
    \psset{linestyle=solid}\FAIL{spotgac<53>}\trput{9}
   }
   \psset{linestyle=solid}\pstree{\CC{\hyperlink{spotgac<54>}{7}}\trput{4}}{
    \psset{linestyle=solid}\pstree{\CC{\hyperlink{spotgac<55>}{10}}\trput{2}}{
     \psset{linestyle=solid}\pstree{\CC{\hyperlink{spotgac<56>}{12}}\trput{6}}{
      \psset{linestyle=solid}\FAIL{spotgac<57>}\trput{2}
      \psset{linestyle=solid}\FAIL{spotgac<58>}\trput{8}
     }
    }
   }
  }
 }
}
}
}
\only<59>{
\scalebox{ 0.40}{\psset{nodesep=0pt,levelsep=36pt}
\pstree{\CC{\hyperlink{spotgac<1>}{4}}}{
 \psset{linestyle=solid}\pstree{\CC{\hyperlink{spotgac<2>}{5}}\trput{2}}{
  \psset{linestyle=solid}\pstree{\CC{\hyperlink{spotgac<3>}{6}}\trput{4}}{
   \psset{linestyle=solid}\pstree{\CC{\hyperlink{spotgac<4>}{7}}\trput{5}}{
    \psset{linestyle=solid}\FAIL{spotgac<5>}\trput{6}
    \psset{linestyle=solid}\pstree{\CC{\hyperlink{spotgac<6>}{12}}\trput{9}}{
     \psset{linestyle=solid}\pstree{\CC{\hyperlink{spotgac<7>}{19}}\trput{6}}{
      \psset{linestyle=solid}\FAIL{spotgac<8>}\trput{6}
      \psset{linestyle=solid}\pstree{\CC{\hyperlink{spotgac<9>}{20}}\trput{9}}{
       \psset{linestyle=solid}\pstree{\CC{\hyperlink{spotgac<10>}{23}}\trput{6}}{
        \psset{linestyle=solid}\FAIL{spotgac<11>}\trput{2}
        \psset{linestyle=solid}\FAIL{spotgac<12>}\trput{8}
       }
       \psset{linestyle=solid}\FAIL{spotgac<13>}\trput{8}
      }
     }
     \psset{linestyle=solid}\FAIL{spotgac<14>}\trput{8}
    }
   }
   \psset{linestyle=solid}\FAIL{spotgac<15>}\trput{6}
  }
  \psset{linestyle=solid}\pstree{\CC{\hyperlink{spotgac<16>}{6}}\trput{6}}{
   \psset{linestyle=solid}\pstree{\CC{\hyperlink{spotgac<17>}{10}}\trput{4}}{
    \psset{linestyle=solid}\FAIL{spotgac<18>}\trput{4}
    \psset{linestyle=solid}\FAIL{spotgac<19>}\trput{6}
    \psset{linestyle=solid}\FAIL{spotgac<20>}\trput{8}
   }
   \psset{linestyle=solid}\pstree{\CC{\hyperlink{spotgac<21>}{7}}\trput{5}}{
    \psset{linestyle=solid}\pstree{\CC{\hyperlink{spotgac<22>}{10}}\trput{4}}{
     \psset{linestyle=solid}\FAIL{spotgac<23>}\trput{5}
     \psset{linestyle=solid}\FAIL{spotgac<24>}\trput{6}
     \psset{linestyle=solid}\FAIL{spotgac<25>}\trput{8}
    }
    \psset{linestyle=solid}\pstree{\CC{\hyperlink{spotgac<26>}{10}}\trput{9}}{
     \psset{linestyle=solid}\pstree{\CC{\hyperlink{spotgac<27>}{11}}\trput{5}}{
      \psset{linestyle=solid}\FAIL{spotgac<28>}\trput{6}
      \psset{linestyle=solid}\FAIL{spotgac<29>}\trput{8}
     }
     \psset{linestyle=solid}\FAIL{spotgac<30>}\trput{6}
     \psset{linestyle=solid}\FAIL{spotgac<31>}\trput{8}
    }
   }
  }
 }
 \psset{linestyle=solid}\pstree{\CC{\hyperlink{spotgac<32>}{5}}\trput{4}}{
  \psset{linestyle=solid}\pstree{\CC{\hyperlink{spotgac<33>}{6}}\trput{2}}{
   \psset{linestyle=solid}\FAIL{spotgac<34>}\trput{5}
   \psset{linestyle=solid}\FAIL{spotgac<35>}\trput{6}
  }
  \psset{linestyle=solid}\pstree{\CC{\hyperlink{spotgac<36>}{6}}\trput{6}}{
   \psset{linestyle=solid}\FAIL{spotgac<37>}\trput{2}
   \psset{linestyle=solid}\FAIL{spotgac<38>}\trput{5}
  }
 }
 \psset{linestyle=solid}\pstree{\CC{\hyperlink{spotgac<39>}{5}}\trput{5}}{
  \psset{linestyle=solid}\pstree{\CC{\hyperlink{spotgac<40>}{6}}\trput{2}}{
   \psset{linestyle=solid}\pstree{\CC{\hyperlink{spotgac<41>}{7}}\trput{4}}{
    \psset{linestyle=solid}\FAIL{spotgac<42>}\trput{6}
    \psset{linestyle=solid}\pstree{\CC{\hyperlink{spotgac<43>}{10}}\trput{9}}{
     \psset{linestyle=solid}\FAIL{spotgac<44>}\trput{6}
     \psset{linestyle=solid}\FAIL{spotgac<45>}\trput{8}
    }
   }
   \psset{linestyle=solid}\pstree{\CC{\hyperlink{spotgac<46>}{7}}\trput{6}}{
    \psset{linestyle=solid}\FAIL{spotgac<47>}\trput{4}
    \psset{linestyle=solid}\FAIL{spotgac<48>}\trput{9}
   }
  }
  \psset{linestyle=solid}\FAIL{spotgac<49>}\trput{4}
  \psset{linestyle=solid}\pstree{\CC{\hyperlink{spotgac<50>}{6}}\trput{6}}{
   \psset{linestyle=solid}\pstree{\CC{\hyperlink{spotgac<51>}{7}}\trput{2}}{
    \psset{linestyle=solid}\FAIL{spotgac<52>}\trput{4}
    \psset{linestyle=solid}\FAIL{spotgac<53>}\trput{9}
   }
   \psset{linestyle=solid}\pstree{\CC{\hyperlink{spotgac<54>}{7}}\trput{4}}{
    \psset{linestyle=solid}\pstree{\CC{\hyperlink{spotgac<55>}{10}}\trput{2}}{
     \psset{linestyle=solid}\pstree{\CC{\hyperlink{spotgac<56>}{12}}\trput{6}}{
      \psset{linestyle=solid}\FAIL{spotgac<57>}\trput{2}
      \psset{linestyle=solid}\FAIL{spotgac<58>}\trput{8}
     }
     \psset{linestyle=solid}\FAIL{spotgac<59>}\trput{8}
    }
   }
  }
 }
}
}
}
\only<60>{
\scalebox{ 0.40}{\psset{nodesep=0pt,levelsep=36pt}
\pstree{\CC{\hyperlink{spotgac<1>}{4}}}{
 \psset{linestyle=solid}\pstree{\CC{\hyperlink{spotgac<2>}{5}}\trput{2}}{
  \psset{linestyle=solid}\pstree{\CC{\hyperlink{spotgac<3>}{6}}\trput{4}}{
   \psset{linestyle=solid}\pstree{\CC{\hyperlink{spotgac<4>}{7}}\trput{5}}{
    \psset{linestyle=solid}\FAIL{spotgac<5>}\trput{6}
    \psset{linestyle=solid}\pstree{\CC{\hyperlink{spotgac<6>}{12}}\trput{9}}{
     \psset{linestyle=solid}\pstree{\CC{\hyperlink{spotgac<7>}{19}}\trput{6}}{
      \psset{linestyle=solid}\FAIL{spotgac<8>}\trput{6}
      \psset{linestyle=solid}\pstree{\CC{\hyperlink{spotgac<9>}{20}}\trput{9}}{
       \psset{linestyle=solid}\pstree{\CC{\hyperlink{spotgac<10>}{23}}\trput{6}}{
        \psset{linestyle=solid}\FAIL{spotgac<11>}\trput{2}
        \psset{linestyle=solid}\FAIL{spotgac<12>}\trput{8}
       }
       \psset{linestyle=solid}\FAIL{spotgac<13>}\trput{8}
      }
     }
     \psset{linestyle=solid}\FAIL{spotgac<14>}\trput{8}
    }
   }
   \psset{linestyle=solid}\FAIL{spotgac<15>}\trput{6}
  }
  \psset{linestyle=solid}\pstree{\CC{\hyperlink{spotgac<16>}{6}}\trput{6}}{
   \psset{linestyle=solid}\pstree{\CC{\hyperlink{spotgac<17>}{10}}\trput{4}}{
    \psset{linestyle=solid}\FAIL{spotgac<18>}\trput{4}
    \psset{linestyle=solid}\FAIL{spotgac<19>}\trput{6}
    \psset{linestyle=solid}\FAIL{spotgac<20>}\trput{8}
   }
   \psset{linestyle=solid}\pstree{\CC{\hyperlink{spotgac<21>}{7}}\trput{5}}{
    \psset{linestyle=solid}\pstree{\CC{\hyperlink{spotgac<22>}{10}}\trput{4}}{
     \psset{linestyle=solid}\FAIL{spotgac<23>}\trput{5}
     \psset{linestyle=solid}\FAIL{spotgac<24>}\trput{6}
     \psset{linestyle=solid}\FAIL{spotgac<25>}\trput{8}
    }
    \psset{linestyle=solid}\pstree{\CC{\hyperlink{spotgac<26>}{10}}\trput{9}}{
     \psset{linestyle=solid}\pstree{\CC{\hyperlink{spotgac<27>}{11}}\trput{5}}{
      \psset{linestyle=solid}\FAIL{spotgac<28>}\trput{6}
      \psset{linestyle=solid}\FAIL{spotgac<29>}\trput{8}
     }
     \psset{linestyle=solid}\FAIL{spotgac<30>}\trput{6}
     \psset{linestyle=solid}\FAIL{spotgac<31>}\trput{8}
    }
   }
  }
 }
 \psset{linestyle=solid}\pstree{\CC{\hyperlink{spotgac<32>}{5}}\trput{4}}{
  \psset{linestyle=solid}\pstree{\CC{\hyperlink{spotgac<33>}{6}}\trput{2}}{
   \psset{linestyle=solid}\FAIL{spotgac<34>}\trput{5}
   \psset{linestyle=solid}\FAIL{spotgac<35>}\trput{6}
  }
  \psset{linestyle=solid}\pstree{\CC{\hyperlink{spotgac<36>}{6}}\trput{6}}{
   \psset{linestyle=solid}\FAIL{spotgac<37>}\trput{2}
   \psset{linestyle=solid}\FAIL{spotgac<38>}\trput{5}
  }
 }
 \psset{linestyle=solid}\pstree{\CC{\hyperlink{spotgac<39>}{5}}\trput{5}}{
  \psset{linestyle=solid}\pstree{\CC{\hyperlink{spotgac<40>}{6}}\trput{2}}{
   \psset{linestyle=solid}\pstree{\CC{\hyperlink{spotgac<41>}{7}}\trput{4}}{
    \psset{linestyle=solid}\FAIL{spotgac<42>}\trput{6}
    \psset{linestyle=solid}\pstree{\CC{\hyperlink{spotgac<43>}{10}}\trput{9}}{
     \psset{linestyle=solid}\FAIL{spotgac<44>}\trput{6}
     \psset{linestyle=solid}\FAIL{spotgac<45>}\trput{8}
    }
   }
   \psset{linestyle=solid}\pstree{\CC{\hyperlink{spotgac<46>}{7}}\trput{6}}{
    \psset{linestyle=solid}\FAIL{spotgac<47>}\trput{4}
    \psset{linestyle=solid}\FAIL{spotgac<48>}\trput{9}
   }
  }
  \psset{linestyle=solid}\FAIL{spotgac<49>}\trput{4}
  \psset{linestyle=solid}\pstree{\CC{\hyperlink{spotgac<50>}{6}}\trput{6}}{
   \psset{linestyle=solid}\pstree{\CC{\hyperlink{spotgac<51>}{7}}\trput{2}}{
    \psset{linestyle=solid}\FAIL{spotgac<52>}\trput{4}
    \psset{linestyle=solid}\FAIL{spotgac<53>}\trput{9}
   }
   \psset{linestyle=solid}\pstree{\CC{\hyperlink{spotgac<54>}{7}}\trput{4}}{
    \psset{linestyle=solid}\pstree{\CC{\hyperlink{spotgac<55>}{10}}\trput{2}}{
     \psset{linestyle=solid}\pstree{\CC{\hyperlink{spotgac<56>}{12}}\trput{6}}{
      \psset{linestyle=solid}\FAIL{spotgac<57>}\trput{2}
      \psset{linestyle=solid}\FAIL{spotgac<58>}\trput{8}
     }
     \psset{linestyle=solid}\FAIL{spotgac<59>}\trput{8}
    }
    \psset{linestyle=solid}\pstree{\CC{\hyperlink{spotgac<60>}{10}}\trput{9}}{
    }
   }
  }
 }
}
}
}
\only<61>{
\scalebox{ 0.40}{\psset{nodesep=0pt,levelsep=36pt}
\pstree{\CC{\hyperlink{spotgac<1>}{4}}}{
 \psset{linestyle=solid}\pstree{\CC{\hyperlink{spotgac<2>}{5}}\trput{2}}{
  \psset{linestyle=solid}\pstree{\CC{\hyperlink{spotgac<3>}{6}}\trput{4}}{
   \psset{linestyle=solid}\pstree{\CC{\hyperlink{spotgac<4>}{7}}\trput{5}}{
    \psset{linestyle=solid}\FAIL{spotgac<5>}\trput{6}
    \psset{linestyle=solid}\pstree{\CC{\hyperlink{spotgac<6>}{12}}\trput{9}}{
     \psset{linestyle=solid}\pstree{\CC{\hyperlink{spotgac<7>}{19}}\trput{6}}{
      \psset{linestyle=solid}\FAIL{spotgac<8>}\trput{6}
      \psset{linestyle=solid}\pstree{\CC{\hyperlink{spotgac<9>}{20}}\trput{9}}{
       \psset{linestyle=solid}\pstree{\CC{\hyperlink{spotgac<10>}{23}}\trput{6}}{
        \psset{linestyle=solid}\FAIL{spotgac<11>}\trput{2}
        \psset{linestyle=solid}\FAIL{spotgac<12>}\trput{8}
       }
       \psset{linestyle=solid}\FAIL{spotgac<13>}\trput{8}
      }
     }
     \psset{linestyle=solid}\FAIL{spotgac<14>}\trput{8}
    }
   }
   \psset{linestyle=solid}\FAIL{spotgac<15>}\trput{6}
  }
  \psset{linestyle=solid}\pstree{\CC{\hyperlink{spotgac<16>}{6}}\trput{6}}{
   \psset{linestyle=solid}\pstree{\CC{\hyperlink{spotgac<17>}{10}}\trput{4}}{
    \psset{linestyle=solid}\FAIL{spotgac<18>}\trput{4}
    \psset{linestyle=solid}\FAIL{spotgac<19>}\trput{6}
    \psset{linestyle=solid}\FAIL{spotgac<20>}\trput{8}
   }
   \psset{linestyle=solid}\pstree{\CC{\hyperlink{spotgac<21>}{7}}\trput{5}}{
    \psset{linestyle=solid}\pstree{\CC{\hyperlink{spotgac<22>}{10}}\trput{4}}{
     \psset{linestyle=solid}\FAIL{spotgac<23>}\trput{5}
     \psset{linestyle=solid}\FAIL{spotgac<24>}\trput{6}
     \psset{linestyle=solid}\FAIL{spotgac<25>}\trput{8}
    }
    \psset{linestyle=solid}\pstree{\CC{\hyperlink{spotgac<26>}{10}}\trput{9}}{
     \psset{linestyle=solid}\pstree{\CC{\hyperlink{spotgac<27>}{11}}\trput{5}}{
      \psset{linestyle=solid}\FAIL{spotgac<28>}\trput{6}
      \psset{linestyle=solid}\FAIL{spotgac<29>}\trput{8}
     }
     \psset{linestyle=solid}\FAIL{spotgac<30>}\trput{6}
     \psset{linestyle=solid}\FAIL{spotgac<31>}\trput{8}
    }
   }
  }
 }
 \psset{linestyle=solid}\pstree{\CC{\hyperlink{spotgac<32>}{5}}\trput{4}}{
  \psset{linestyle=solid}\pstree{\CC{\hyperlink{spotgac<33>}{6}}\trput{2}}{
   \psset{linestyle=solid}\FAIL{spotgac<34>}\trput{5}
   \psset{linestyle=solid}\FAIL{spotgac<35>}\trput{6}
  }
  \psset{linestyle=solid}\pstree{\CC{\hyperlink{spotgac<36>}{6}}\trput{6}}{
   \psset{linestyle=solid}\FAIL{spotgac<37>}\trput{2}
   \psset{linestyle=solid}\FAIL{spotgac<38>}\trput{5}
  }
 }
 \psset{linestyle=solid}\pstree{\CC{\hyperlink{spotgac<39>}{5}}\trput{5}}{
  \psset{linestyle=solid}\pstree{\CC{\hyperlink{spotgac<40>}{6}}\trput{2}}{
   \psset{linestyle=solid}\pstree{\CC{\hyperlink{spotgac<41>}{7}}\trput{4}}{
    \psset{linestyle=solid}\FAIL{spotgac<42>}\trput{6}
    \psset{linestyle=solid}\pstree{\CC{\hyperlink{spotgac<43>}{10}}\trput{9}}{
     \psset{linestyle=solid}\FAIL{spotgac<44>}\trput{6}
     \psset{linestyle=solid}\FAIL{spotgac<45>}\trput{8}
    }
   }
   \psset{linestyle=solid}\pstree{\CC{\hyperlink{spotgac<46>}{7}}\trput{6}}{
    \psset{linestyle=solid}\FAIL{spotgac<47>}\trput{4}
    \psset{linestyle=solid}\FAIL{spotgac<48>}\trput{9}
   }
  }
  \psset{linestyle=solid}\FAIL{spotgac<49>}\trput{4}
  \psset{linestyle=solid}\pstree{\CC{\hyperlink{spotgac<50>}{6}}\trput{6}}{
   \psset{linestyle=solid}\pstree{\CC{\hyperlink{spotgac<51>}{7}}\trput{2}}{
    \psset{linestyle=solid}\FAIL{spotgac<52>}\trput{4}
    \psset{linestyle=solid}\FAIL{spotgac<53>}\trput{9}
   }
   \psset{linestyle=solid}\pstree{\CC{\hyperlink{spotgac<54>}{7}}\trput{4}}{
    \psset{linestyle=solid}\pstree{\CC{\hyperlink{spotgac<55>}{10}}\trput{2}}{
     \psset{linestyle=solid}\pstree{\CC{\hyperlink{spotgac<56>}{12}}\trput{6}}{
      \psset{linestyle=solid}\FAIL{spotgac<57>}\trput{2}
      \psset{linestyle=solid}\FAIL{spotgac<58>}\trput{8}
     }
     \psset{linestyle=solid}\FAIL{spotgac<59>}\trput{8}
    }
    \psset{linestyle=solid}\pstree{\CC{\hyperlink{spotgac<60>}{10}}\trput{9}}{
     \psset{linestyle=solid}\pstree{\CC{\hyperlink{spotgac<61>}{12}}\trput{6}}{
     }
    }
   }
  }
 }
}
}
}
\only<62>{
\scalebox{ 0.40}{\psset{nodesep=0pt,levelsep=36pt}
\pstree{\CC{\hyperlink{spotgac<1>}{4}}}{
 \psset{linestyle=solid}\pstree{\CC{\hyperlink{spotgac<2>}{5}}\trput{2}}{
  \psset{linestyle=solid}\pstree{\CC{\hyperlink{spotgac<3>}{6}}\trput{4}}{
   \psset{linestyle=solid}\pstree{\CC{\hyperlink{spotgac<4>}{7}}\trput{5}}{
    \psset{linestyle=solid}\FAIL{spotgac<5>}\trput{6}
    \psset{linestyle=solid}\pstree{\CC{\hyperlink{spotgac<6>}{12}}\trput{9}}{
     \psset{linestyle=solid}\pstree{\CC{\hyperlink{spotgac<7>}{19}}\trput{6}}{
      \psset{linestyle=solid}\FAIL{spotgac<8>}\trput{6}
      \psset{linestyle=solid}\pstree{\CC{\hyperlink{spotgac<9>}{20}}\trput{9}}{
       \psset{linestyle=solid}\pstree{\CC{\hyperlink{spotgac<10>}{23}}\trput{6}}{
        \psset{linestyle=solid}\FAIL{spotgac<11>}\trput{2}
        \psset{linestyle=solid}\FAIL{spotgac<12>}\trput{8}
       }
       \psset{linestyle=solid}\FAIL{spotgac<13>}\trput{8}
      }
     }
     \psset{linestyle=solid}\FAIL{spotgac<14>}\trput{8}
    }
   }
   \psset{linestyle=solid}\FAIL{spotgac<15>}\trput{6}
  }
  \psset{linestyle=solid}\pstree{\CC{\hyperlink{spotgac<16>}{6}}\trput{6}}{
   \psset{linestyle=solid}\pstree{\CC{\hyperlink{spotgac<17>}{10}}\trput{4}}{
    \psset{linestyle=solid}\FAIL{spotgac<18>}\trput{4}
    \psset{linestyle=solid}\FAIL{spotgac<19>}\trput{6}
    \psset{linestyle=solid}\FAIL{spotgac<20>}\trput{8}
   }
   \psset{linestyle=solid}\pstree{\CC{\hyperlink{spotgac<21>}{7}}\trput{5}}{
    \psset{linestyle=solid}\pstree{\CC{\hyperlink{spotgac<22>}{10}}\trput{4}}{
     \psset{linestyle=solid}\FAIL{spotgac<23>}\trput{5}
     \psset{linestyle=solid}\FAIL{spotgac<24>}\trput{6}
     \psset{linestyle=solid}\FAIL{spotgac<25>}\trput{8}
    }
    \psset{linestyle=solid}\pstree{\CC{\hyperlink{spotgac<26>}{10}}\trput{9}}{
     \psset{linestyle=solid}\pstree{\CC{\hyperlink{spotgac<27>}{11}}\trput{5}}{
      \psset{linestyle=solid}\FAIL{spotgac<28>}\trput{6}
      \psset{linestyle=solid}\FAIL{spotgac<29>}\trput{8}
     }
     \psset{linestyle=solid}\FAIL{spotgac<30>}\trput{6}
     \psset{linestyle=solid}\FAIL{spotgac<31>}\trput{8}
    }
   }
  }
 }
 \psset{linestyle=solid}\pstree{\CC{\hyperlink{spotgac<32>}{5}}\trput{4}}{
  \psset{linestyle=solid}\pstree{\CC{\hyperlink{spotgac<33>}{6}}\trput{2}}{
   \psset{linestyle=solid}\FAIL{spotgac<34>}\trput{5}
   \psset{linestyle=solid}\FAIL{spotgac<35>}\trput{6}
  }
  \psset{linestyle=solid}\pstree{\CC{\hyperlink{spotgac<36>}{6}}\trput{6}}{
   \psset{linestyle=solid}\FAIL{spotgac<37>}\trput{2}
   \psset{linestyle=solid}\FAIL{spotgac<38>}\trput{5}
  }
 }
 \psset{linestyle=solid}\pstree{\CC{\hyperlink{spotgac<39>}{5}}\trput{5}}{
  \psset{linestyle=solid}\pstree{\CC{\hyperlink{spotgac<40>}{6}}\trput{2}}{
   \psset{linestyle=solid}\pstree{\CC{\hyperlink{spotgac<41>}{7}}\trput{4}}{
    \psset{linestyle=solid}\FAIL{spotgac<42>}\trput{6}
    \psset{linestyle=solid}\pstree{\CC{\hyperlink{spotgac<43>}{10}}\trput{9}}{
     \psset{linestyle=solid}\FAIL{spotgac<44>}\trput{6}
     \psset{linestyle=solid}\FAIL{spotgac<45>}\trput{8}
    }
   }
   \psset{linestyle=solid}\pstree{\CC{\hyperlink{spotgac<46>}{7}}\trput{6}}{
    \psset{linestyle=solid}\FAIL{spotgac<47>}\trput{4}
    \psset{linestyle=solid}\FAIL{spotgac<48>}\trput{9}
   }
  }
  \psset{linestyle=solid}\FAIL{spotgac<49>}\trput{4}
  \psset{linestyle=solid}\pstree{\CC{\hyperlink{spotgac<50>}{6}}\trput{6}}{
   \psset{linestyle=solid}\pstree{\CC{\hyperlink{spotgac<51>}{7}}\trput{2}}{
    \psset{linestyle=solid}\FAIL{spotgac<52>}\trput{4}
    \psset{linestyle=solid}\FAIL{spotgac<53>}\trput{9}
   }
   \psset{linestyle=solid}\pstree{\CC{\hyperlink{spotgac<54>}{7}}\trput{4}}{
    \psset{linestyle=solid}\pstree{\CC{\hyperlink{spotgac<55>}{10}}\trput{2}}{
     \psset{linestyle=solid}\pstree{\CC{\hyperlink{spotgac<56>}{12}}\trput{6}}{
      \psset{linestyle=solid}\FAIL{spotgac<57>}\trput{2}
      \psset{linestyle=solid}\FAIL{spotgac<58>}\trput{8}
     }
     \psset{linestyle=solid}\FAIL{spotgac<59>}\trput{8}
    }
    \psset{linestyle=solid}\pstree{\CC{\hyperlink{spotgac<60>}{10}}\trput{9}}{
     \psset{linestyle=solid}\pstree{\CC{\hyperlink{spotgac<61>}{12}}\trput{6}}{
      \psset{linestyle=solid}\pstree{\CC{\hyperlink{spotgac<62>}{22}}\trput{2}}{
      }
     }
    }
   }
  }
 }
}
}
}
\only<63>{
\scalebox{ 0.40}{\psset{nodesep=0pt,levelsep=36pt}
\pstree{\CC{\hyperlink{spotgac<1>}{4}}}{
 \psset{linestyle=solid}\pstree{\CC{\hyperlink{spotgac<2>}{5}}\trput{2}}{
  \psset{linestyle=solid}\pstree{\CC{\hyperlink{spotgac<3>}{6}}\trput{4}}{
   \psset{linestyle=solid}\pstree{\CC{\hyperlink{spotgac<4>}{7}}\trput{5}}{
    \psset{linestyle=solid}\FAIL{spotgac<5>}\trput{6}
    \psset{linestyle=solid}\pstree{\CC{\hyperlink{spotgac<6>}{12}}\trput{9}}{
     \psset{linestyle=solid}\pstree{\CC{\hyperlink{spotgac<7>}{19}}\trput{6}}{
      \psset{linestyle=solid}\FAIL{spotgac<8>}\trput{6}
      \psset{linestyle=solid}\pstree{\CC{\hyperlink{spotgac<9>}{20}}\trput{9}}{
       \psset{linestyle=solid}\pstree{\CC{\hyperlink{spotgac<10>}{23}}\trput{6}}{
        \psset{linestyle=solid}\FAIL{spotgac<11>}\trput{2}
        \psset{linestyle=solid}\FAIL{spotgac<12>}\trput{8}
       }
       \psset{linestyle=solid}\FAIL{spotgac<13>}\trput{8}
      }
     }
     \psset{linestyle=solid}\FAIL{spotgac<14>}\trput{8}
    }
   }
   \psset{linestyle=solid}\FAIL{spotgac<15>}\trput{6}
  }
  \psset{linestyle=solid}\pstree{\CC{\hyperlink{spotgac<16>}{6}}\trput{6}}{
   \psset{linestyle=solid}\pstree{\CC{\hyperlink{spotgac<17>}{10}}\trput{4}}{
    \psset{linestyle=solid}\FAIL{spotgac<18>}\trput{4}
    \psset{linestyle=solid}\FAIL{spotgac<19>}\trput{6}
    \psset{linestyle=solid}\FAIL{spotgac<20>}\trput{8}
   }
   \psset{linestyle=solid}\pstree{\CC{\hyperlink{spotgac<21>}{7}}\trput{5}}{
    \psset{linestyle=solid}\pstree{\CC{\hyperlink{spotgac<22>}{10}}\trput{4}}{
     \psset{linestyle=solid}\FAIL{spotgac<23>}\trput{5}
     \psset{linestyle=solid}\FAIL{spotgac<24>}\trput{6}
     \psset{linestyle=solid}\FAIL{spotgac<25>}\trput{8}
    }
    \psset{linestyle=solid}\pstree{\CC{\hyperlink{spotgac<26>}{10}}\trput{9}}{
     \psset{linestyle=solid}\pstree{\CC{\hyperlink{spotgac<27>}{11}}\trput{5}}{
      \psset{linestyle=solid}\FAIL{spotgac<28>}\trput{6}
      \psset{linestyle=solid}\FAIL{spotgac<29>}\trput{8}
     }
     \psset{linestyle=solid}\FAIL{spotgac<30>}\trput{6}
     \psset{linestyle=solid}\FAIL{spotgac<31>}\trput{8}
    }
   }
  }
 }
 \psset{linestyle=solid}\pstree{\CC{\hyperlink{spotgac<32>}{5}}\trput{4}}{
  \psset{linestyle=solid}\pstree{\CC{\hyperlink{spotgac<33>}{6}}\trput{2}}{
   \psset{linestyle=solid}\FAIL{spotgac<34>}\trput{5}
   \psset{linestyle=solid}\FAIL{spotgac<35>}\trput{6}
  }
  \psset{linestyle=solid}\pstree{\CC{\hyperlink{spotgac<36>}{6}}\trput{6}}{
   \psset{linestyle=solid}\FAIL{spotgac<37>}\trput{2}
   \psset{linestyle=solid}\FAIL{spotgac<38>}\trput{5}
  }
 }
 \psset{linestyle=solid}\pstree{\CC{\hyperlink{spotgac<39>}{5}}\trput{5}}{
  \psset{linestyle=solid}\pstree{\CC{\hyperlink{spotgac<40>}{6}}\trput{2}}{
   \psset{linestyle=solid}\pstree{\CC{\hyperlink{spotgac<41>}{7}}\trput{4}}{
    \psset{linestyle=solid}\FAIL{spotgac<42>}\trput{6}
    \psset{linestyle=solid}\pstree{\CC{\hyperlink{spotgac<43>}{10}}\trput{9}}{
     \psset{linestyle=solid}\FAIL{spotgac<44>}\trput{6}
     \psset{linestyle=solid}\FAIL{spotgac<45>}\trput{8}
    }
   }
   \psset{linestyle=solid}\pstree{\CC{\hyperlink{spotgac<46>}{7}}\trput{6}}{
    \psset{linestyle=solid}\FAIL{spotgac<47>}\trput{4}
    \psset{linestyle=solid}\FAIL{spotgac<48>}\trput{9}
   }
  }
  \psset{linestyle=solid}\FAIL{spotgac<49>}\trput{4}
  \psset{linestyle=solid}\pstree{\CC{\hyperlink{spotgac<50>}{6}}\trput{6}}{
   \psset{linestyle=solid}\pstree{\CC{\hyperlink{spotgac<51>}{7}}\trput{2}}{
    \psset{linestyle=solid}\FAIL{spotgac<52>}\trput{4}
    \psset{linestyle=solid}\FAIL{spotgac<53>}\trput{9}
   }
   \psset{linestyle=solid}\pstree{\CC{\hyperlink{spotgac<54>}{7}}\trput{4}}{
    \psset{linestyle=solid}\pstree{\CC{\hyperlink{spotgac<55>}{10}}\trput{2}}{
     \psset{linestyle=solid}\pstree{\CC{\hyperlink{spotgac<56>}{12}}\trput{6}}{
      \psset{linestyle=solid}\FAIL{spotgac<57>}\trput{2}
      \psset{linestyle=solid}\FAIL{spotgac<58>}\trput{8}
     }
     \psset{linestyle=solid}\FAIL{spotgac<59>}\trput{8}
    }
    \psset{linestyle=solid}\pstree{\CC{\hyperlink{spotgac<60>}{10}}\trput{9}}{
     \psset{linestyle=solid}\pstree{\CC{\hyperlink{spotgac<61>}{12}}\trput{6}}{
      \psset{linestyle=solid}\pstree{\CC{\hyperlink{spotgac<62>}{22}}\trput{2}}{
       \psset{linestyle=solid}\pstree{\CC{\hyperlink{spotgac<63>}{24}}\trput{2}}{
       }
      }
     }
    }
   }
  }
 }
}
}
}
\only<64>{
\scalebox{ 0.40}{\psset{nodesep=0pt,levelsep=36pt}
\pstree{\CC{\hyperlink{spotgac<1>}{4}}}{
 \psset{linestyle=solid}\pstree{\CC{\hyperlink{spotgac<2>}{5}}\trput{2}}{
  \psset{linestyle=solid}\pstree{\CC{\hyperlink{spotgac<3>}{6}}\trput{4}}{
   \psset{linestyle=solid}\pstree{\CC{\hyperlink{spotgac<4>}{7}}\trput{5}}{
    \psset{linestyle=solid}\FAIL{spotgac<5>}\trput{6}
    \psset{linestyle=solid}\pstree{\CC{\hyperlink{spotgac<6>}{12}}\trput{9}}{
     \psset{linestyle=solid}\pstree{\CC{\hyperlink{spotgac<7>}{19}}\trput{6}}{
      \psset{linestyle=solid}\FAIL{spotgac<8>}\trput{6}
      \psset{linestyle=solid}\pstree{\CC{\hyperlink{spotgac<9>}{20}}\trput{9}}{
       \psset{linestyle=solid}\pstree{\CC{\hyperlink{spotgac<10>}{23}}\trput{6}}{
        \psset{linestyle=solid}\FAIL{spotgac<11>}\trput{2}
        \psset{linestyle=solid}\FAIL{spotgac<12>}\trput{8}
       }
       \psset{linestyle=solid}\FAIL{spotgac<13>}\trput{8}
      }
     }
     \psset{linestyle=solid}\FAIL{spotgac<14>}\trput{8}
    }
   }
   \psset{linestyle=solid}\FAIL{spotgac<15>}\trput{6}
  }
  \psset{linestyle=solid}\pstree{\CC{\hyperlink{spotgac<16>}{6}}\trput{6}}{
   \psset{linestyle=solid}\pstree{\CC{\hyperlink{spotgac<17>}{10}}\trput{4}}{
    \psset{linestyle=solid}\FAIL{spotgac<18>}\trput{4}
    \psset{linestyle=solid}\FAIL{spotgac<19>}\trput{6}
    \psset{linestyle=solid}\FAIL{spotgac<20>}\trput{8}
   }
   \psset{linestyle=solid}\pstree{\CC{\hyperlink{spotgac<21>}{7}}\trput{5}}{
    \psset{linestyle=solid}\pstree{\CC{\hyperlink{spotgac<22>}{10}}\trput{4}}{
     \psset{linestyle=solid}\FAIL{spotgac<23>}\trput{5}
     \psset{linestyle=solid}\FAIL{spotgac<24>}\trput{6}
     \psset{linestyle=solid}\FAIL{spotgac<25>}\trput{8}
    }
    \psset{linestyle=solid}\pstree{\CC{\hyperlink{spotgac<26>}{10}}\trput{9}}{
     \psset{linestyle=solid}\pstree{\CC{\hyperlink{spotgac<27>}{11}}\trput{5}}{
      \psset{linestyle=solid}\FAIL{spotgac<28>}\trput{6}
      \psset{linestyle=solid}\FAIL{spotgac<29>}\trput{8}
     }
     \psset{linestyle=solid}\FAIL{spotgac<30>}\trput{6}
     \psset{linestyle=solid}\FAIL{spotgac<31>}\trput{8}
    }
   }
  }
 }
 \psset{linestyle=solid}\pstree{\CC{\hyperlink{spotgac<32>}{5}}\trput{4}}{
  \psset{linestyle=solid}\pstree{\CC{\hyperlink{spotgac<33>}{6}}\trput{2}}{
   \psset{linestyle=solid}\FAIL{spotgac<34>}\trput{5}
   \psset{linestyle=solid}\FAIL{spotgac<35>}\trput{6}
  }
  \psset{linestyle=solid}\pstree{\CC{\hyperlink{spotgac<36>}{6}}\trput{6}}{
   \psset{linestyle=solid}\FAIL{spotgac<37>}\trput{2}
   \psset{linestyle=solid}\FAIL{spotgac<38>}\trput{5}
  }
 }
 \psset{linestyle=solid}\pstree{\CC{\hyperlink{spotgac<39>}{5}}\trput{5}}{
  \psset{linestyle=solid}\pstree{\CC{\hyperlink{spotgac<40>}{6}}\trput{2}}{
   \psset{linestyle=solid}\pstree{\CC{\hyperlink{spotgac<41>}{7}}\trput{4}}{
    \psset{linestyle=solid}\FAIL{spotgac<42>}\trput{6}
    \psset{linestyle=solid}\pstree{\CC{\hyperlink{spotgac<43>}{10}}\trput{9}}{
     \psset{linestyle=solid}\FAIL{spotgac<44>}\trput{6}
     \psset{linestyle=solid}\FAIL{spotgac<45>}\trput{8}
    }
   }
   \psset{linestyle=solid}\pstree{\CC{\hyperlink{spotgac<46>}{7}}\trput{6}}{
    \psset{linestyle=solid}\FAIL{spotgac<47>}\trput{4}
    \psset{linestyle=solid}\FAIL{spotgac<48>}\trput{9}
   }
  }
  \psset{linestyle=solid}\FAIL{spotgac<49>}\trput{4}
  \psset{linestyle=solid}\pstree{\CC{\hyperlink{spotgac<50>}{6}}\trput{6}}{
   \psset{linestyle=solid}\pstree{\CC{\hyperlink{spotgac<51>}{7}}\trput{2}}{
    \psset{linestyle=solid}\FAIL{spotgac<52>}\trput{4}
    \psset{linestyle=solid}\FAIL{spotgac<53>}\trput{9}
   }
   \psset{linestyle=solid}\pstree{\CC{\hyperlink{spotgac<54>}{7}}\trput{4}}{
    \psset{linestyle=solid}\pstree{\CC{\hyperlink{spotgac<55>}{10}}\trput{2}}{
     \psset{linestyle=solid}\pstree{\CC{\hyperlink{spotgac<56>}{12}}\trput{6}}{
      \psset{linestyle=solid}\FAIL{spotgac<57>}\trput{2}
      \psset{linestyle=solid}\FAIL{spotgac<58>}\trput{8}
     }
     \psset{linestyle=solid}\FAIL{spotgac<59>}\trput{8}
    }
    \psset{linestyle=solid}\pstree{\CC{\hyperlink{spotgac<60>}{10}}\trput{9}}{
     \psset{linestyle=solid}\pstree{\CC{\hyperlink{spotgac<61>}{12}}\trput{6}}{
      \psset{linestyle=solid}\pstree{\CC{\hyperlink{spotgac<62>}{22}}\trput{2}}{
       \psset{linestyle=solid}\pstree{\CC{\hyperlink{spotgac<63>}{24}}\trput{2}}{
        \psset{linestyle=solid}\pstree{\CC{\hyperlink{spotgac<64>}{31}}\trput{1}}{
        }
       }
      }
     }
    }
   }
  }
 }
}
}
}
\only<65>{
\scalebox{ 0.40}{\psset{nodesep=0pt,levelsep=36pt}
\pstree{\CC{\hyperlink{spotgac<1>}{4}}}{
 \psset{linestyle=solid}\pstree{\CC{\hyperlink{spotgac<2>}{5}}\trput{2}}{
  \psset{linestyle=solid}\pstree{\CC{\hyperlink{spotgac<3>}{6}}\trput{4}}{
   \psset{linestyle=solid}\pstree{\CC{\hyperlink{spotgac<4>}{7}}\trput{5}}{
    \psset{linestyle=solid}\FAIL{spotgac<5>}\trput{6}
    \psset{linestyle=solid}\pstree{\CC{\hyperlink{spotgac<6>}{12}}\trput{9}}{
     \psset{linestyle=solid}\pstree{\CC{\hyperlink{spotgac<7>}{19}}\trput{6}}{
      \psset{linestyle=solid}\FAIL{spotgac<8>}\trput{6}
      \psset{linestyle=solid}\pstree{\CC{\hyperlink{spotgac<9>}{20}}\trput{9}}{
       \psset{linestyle=solid}\pstree{\CC{\hyperlink{spotgac<10>}{23}}\trput{6}}{
        \psset{linestyle=solid}\FAIL{spotgac<11>}\trput{2}
        \psset{linestyle=solid}\FAIL{spotgac<12>}\trput{8}
       }
       \psset{linestyle=solid}\FAIL{spotgac<13>}\trput{8}
      }
     }
     \psset{linestyle=solid}\FAIL{spotgac<14>}\trput{8}
    }
   }
   \psset{linestyle=solid}\FAIL{spotgac<15>}\trput{6}
  }
  \psset{linestyle=solid}\pstree{\CC{\hyperlink{spotgac<16>}{6}}\trput{6}}{
   \psset{linestyle=solid}\pstree{\CC{\hyperlink{spotgac<17>}{10}}\trput{4}}{
    \psset{linestyle=solid}\FAIL{spotgac<18>}\trput{4}
    \psset{linestyle=solid}\FAIL{spotgac<19>}\trput{6}
    \psset{linestyle=solid}\FAIL{spotgac<20>}\trput{8}
   }
   \psset{linestyle=solid}\pstree{\CC{\hyperlink{spotgac<21>}{7}}\trput{5}}{
    \psset{linestyle=solid}\pstree{\CC{\hyperlink{spotgac<22>}{10}}\trput{4}}{
     \psset{linestyle=solid}\FAIL{spotgac<23>}\trput{5}
     \psset{linestyle=solid}\FAIL{spotgac<24>}\trput{6}
     \psset{linestyle=solid}\FAIL{spotgac<25>}\trput{8}
    }
    \psset{linestyle=solid}\pstree{\CC{\hyperlink{spotgac<26>}{10}}\trput{9}}{
     \psset{linestyle=solid}\pstree{\CC{\hyperlink{spotgac<27>}{11}}\trput{5}}{
      \psset{linestyle=solid}\FAIL{spotgac<28>}\trput{6}
      \psset{linestyle=solid}\FAIL{spotgac<29>}\trput{8}
     }
     \psset{linestyle=solid}\FAIL{spotgac<30>}\trput{6}
     \psset{linestyle=solid}\FAIL{spotgac<31>}\trput{8}
    }
   }
  }
 }
 \psset{linestyle=solid}\pstree{\CC{\hyperlink{spotgac<32>}{5}}\trput{4}}{
  \psset{linestyle=solid}\pstree{\CC{\hyperlink{spotgac<33>}{6}}\trput{2}}{
   \psset{linestyle=solid}\FAIL{spotgac<34>}\trput{5}
   \psset{linestyle=solid}\FAIL{spotgac<35>}\trput{6}
  }
  \psset{linestyle=solid}\pstree{\CC{\hyperlink{spotgac<36>}{6}}\trput{6}}{
   \psset{linestyle=solid}\FAIL{spotgac<37>}\trput{2}
   \psset{linestyle=solid}\FAIL{spotgac<38>}\trput{5}
  }
 }
 \psset{linestyle=solid}\pstree{\CC{\hyperlink{spotgac<39>}{5}}\trput{5}}{
  \psset{linestyle=solid}\pstree{\CC{\hyperlink{spotgac<40>}{6}}\trput{2}}{
   \psset{linestyle=solid}\pstree{\CC{\hyperlink{spotgac<41>}{7}}\trput{4}}{
    \psset{linestyle=solid}\FAIL{spotgac<42>}\trput{6}
    \psset{linestyle=solid}\pstree{\CC{\hyperlink{spotgac<43>}{10}}\trput{9}}{
     \psset{linestyle=solid}\FAIL{spotgac<44>}\trput{6}
     \psset{linestyle=solid}\FAIL{spotgac<45>}\trput{8}
    }
   }
   \psset{linestyle=solid}\pstree{\CC{\hyperlink{spotgac<46>}{7}}\trput{6}}{
    \psset{linestyle=solid}\FAIL{spotgac<47>}\trput{4}
    \psset{linestyle=solid}\FAIL{spotgac<48>}\trput{9}
   }
  }
  \psset{linestyle=solid}\FAIL{spotgac<49>}\trput{4}
  \psset{linestyle=solid}\pstree{\CC{\hyperlink{spotgac<50>}{6}}\trput{6}}{
   \psset{linestyle=solid}\pstree{\CC{\hyperlink{spotgac<51>}{7}}\trput{2}}{
    \psset{linestyle=solid}\FAIL{spotgac<52>}\trput{4}
    \psset{linestyle=solid}\FAIL{spotgac<53>}\trput{9}
   }
   \psset{linestyle=solid}\pstree{\CC{\hyperlink{spotgac<54>}{7}}\trput{4}}{
    \psset{linestyle=solid}\pstree{\CC{\hyperlink{spotgac<55>}{10}}\trput{2}}{
     \psset{linestyle=solid}\pstree{\CC{\hyperlink{spotgac<56>}{12}}\trput{6}}{
      \psset{linestyle=solid}\FAIL{spotgac<57>}\trput{2}
      \psset{linestyle=solid}\FAIL{spotgac<58>}\trput{8}
     }
     \psset{linestyle=solid}\FAIL{spotgac<59>}\trput{8}
    }
    \psset{linestyle=solid}\pstree{\CC{\hyperlink{spotgac<60>}{10}}\trput{9}}{
     \psset{linestyle=solid}\pstree{\CC{\hyperlink{spotgac<61>}{12}}\trput{6}}{
      \psset{linestyle=solid}\pstree{\CC{\hyperlink{spotgac<62>}{22}}\trput{2}}{
       \psset{linestyle=solid}\pstree{\CC{\hyperlink{spotgac<63>}{24}}\trput{2}}{
        \psset{linestyle=solid}\pstree{\CC{\hyperlink{spotgac<64>}{31}}\trput{1}}{
         \psset{linestyle=solid}\FAIL{spotgac<65>}\trput{4}
        }
       }
      }
     }
    }
   }
  }
 }
}
}
}
\only<66>{
\scalebox{ 0.40}{\psset{nodesep=0pt,levelsep=36pt}
\pstree{\CC{\hyperlink{spotgac<1>}{4}}}{
 \psset{linestyle=solid}\pstree{\CC{\hyperlink{spotgac<2>}{5}}\trput{2}}{
  \psset{linestyle=solid}\pstree{\CC{\hyperlink{spotgac<3>}{6}}\trput{4}}{
   \psset{linestyle=solid}\pstree{\CC{\hyperlink{spotgac<4>}{7}}\trput{5}}{
    \psset{linestyle=solid}\FAIL{spotgac<5>}\trput{6}
    \psset{linestyle=solid}\pstree{\CC{\hyperlink{spotgac<6>}{12}}\trput{9}}{
     \psset{linestyle=solid}\pstree{\CC{\hyperlink{spotgac<7>}{19}}\trput{6}}{
      \psset{linestyle=solid}\FAIL{spotgac<8>}\trput{6}
      \psset{linestyle=solid}\pstree{\CC{\hyperlink{spotgac<9>}{20}}\trput{9}}{
       \psset{linestyle=solid}\pstree{\CC{\hyperlink{spotgac<10>}{23}}\trput{6}}{
        \psset{linestyle=solid}\FAIL{spotgac<11>}\trput{2}
        \psset{linestyle=solid}\FAIL{spotgac<12>}\trput{8}
       }
       \psset{linestyle=solid}\FAIL{spotgac<13>}\trput{8}
      }
     }
     \psset{linestyle=solid}\FAIL{spotgac<14>}\trput{8}
    }
   }
   \psset{linestyle=solid}\FAIL{spotgac<15>}\trput{6}
  }
  \psset{linestyle=solid}\pstree{\CC{\hyperlink{spotgac<16>}{6}}\trput{6}}{
   \psset{linestyle=solid}\pstree{\CC{\hyperlink{spotgac<17>}{10}}\trput{4}}{
    \psset{linestyle=solid}\FAIL{spotgac<18>}\trput{4}
    \psset{linestyle=solid}\FAIL{spotgac<19>}\trput{6}
    \psset{linestyle=solid}\FAIL{spotgac<20>}\trput{8}
   }
   \psset{linestyle=solid}\pstree{\CC{\hyperlink{spotgac<21>}{7}}\trput{5}}{
    \psset{linestyle=solid}\pstree{\CC{\hyperlink{spotgac<22>}{10}}\trput{4}}{
     \psset{linestyle=solid}\FAIL{spotgac<23>}\trput{5}
     \psset{linestyle=solid}\FAIL{spotgac<24>}\trput{6}
     \psset{linestyle=solid}\FAIL{spotgac<25>}\trput{8}
    }
    \psset{linestyle=solid}\pstree{\CC{\hyperlink{spotgac<26>}{10}}\trput{9}}{
     \psset{linestyle=solid}\pstree{\CC{\hyperlink{spotgac<27>}{11}}\trput{5}}{
      \psset{linestyle=solid}\FAIL{spotgac<28>}\trput{6}
      \psset{linestyle=solid}\FAIL{spotgac<29>}\trput{8}
     }
     \psset{linestyle=solid}\FAIL{spotgac<30>}\trput{6}
     \psset{linestyle=solid}\FAIL{spotgac<31>}\trput{8}
    }
   }
  }
 }
 \psset{linestyle=solid}\pstree{\CC{\hyperlink{spotgac<32>}{5}}\trput{4}}{
  \psset{linestyle=solid}\pstree{\CC{\hyperlink{spotgac<33>}{6}}\trput{2}}{
   \psset{linestyle=solid}\FAIL{spotgac<34>}\trput{5}
   \psset{linestyle=solid}\FAIL{spotgac<35>}\trput{6}
  }
  \psset{linestyle=solid}\pstree{\CC{\hyperlink{spotgac<36>}{6}}\trput{6}}{
   \psset{linestyle=solid}\FAIL{spotgac<37>}\trput{2}
   \psset{linestyle=solid}\FAIL{spotgac<38>}\trput{5}
  }
 }
 \psset{linestyle=solid}\pstree{\CC{\hyperlink{spotgac<39>}{5}}\trput{5}}{
  \psset{linestyle=solid}\pstree{\CC{\hyperlink{spotgac<40>}{6}}\trput{2}}{
   \psset{linestyle=solid}\pstree{\CC{\hyperlink{spotgac<41>}{7}}\trput{4}}{
    \psset{linestyle=solid}\FAIL{spotgac<42>}\trput{6}
    \psset{linestyle=solid}\pstree{\CC{\hyperlink{spotgac<43>}{10}}\trput{9}}{
     \psset{linestyle=solid}\FAIL{spotgac<44>}\trput{6}
     \psset{linestyle=solid}\FAIL{spotgac<45>}\trput{8}
    }
   }
   \psset{linestyle=solid}\pstree{\CC{\hyperlink{spotgac<46>}{7}}\trput{6}}{
    \psset{linestyle=solid}\FAIL{spotgac<47>}\trput{4}
    \psset{linestyle=solid}\FAIL{spotgac<48>}\trput{9}
   }
  }
  \psset{linestyle=solid}\FAIL{spotgac<49>}\trput{4}
  \psset{linestyle=solid}\pstree{\CC{\hyperlink{spotgac<50>}{6}}\trput{6}}{
   \psset{linestyle=solid}\pstree{\CC{\hyperlink{spotgac<51>}{7}}\trput{2}}{
    \psset{linestyle=solid}\FAIL{spotgac<52>}\trput{4}
    \psset{linestyle=solid}\FAIL{spotgac<53>}\trput{9}
   }
   \psset{linestyle=solid}\pstree{\CC{\hyperlink{spotgac<54>}{7}}\trput{4}}{
    \psset{linestyle=solid}\pstree{\CC{\hyperlink{spotgac<55>}{10}}\trput{2}}{
     \psset{linestyle=solid}\pstree{\CC{\hyperlink{spotgac<56>}{12}}\trput{6}}{
      \psset{linestyle=solid}\FAIL{spotgac<57>}\trput{2}
      \psset{linestyle=solid}\FAIL{spotgac<58>}\trput{8}
     }
     \psset{linestyle=solid}\FAIL{spotgac<59>}\trput{8}
    }
    \psset{linestyle=solid}\pstree{\CC{\hyperlink{spotgac<60>}{10}}\trput{9}}{
     \psset{linestyle=solid}\pstree{\CC{\hyperlink{spotgac<61>}{12}}\trput{6}}{
      \psset{linestyle=solid}\pstree{\CC{\hyperlink{spotgac<62>}{22}}\trput{2}}{
       \psset{linestyle=solid}\pstree{\CC{\hyperlink{spotgac<63>}{24}}\trput{2}}{
        \psset{linestyle=solid}\pstree{\CC{\hyperlink{spotgac<64>}{31}}\trput{1}}{
         \psset{linestyle=solid}\FAIL{spotgac<65>}\trput{4}
         \psset{linestyle=solid}\SUCC{spotgac<66>}\trput{8}
        }
       }
      }
     }
    }
   }
  }
 }
}
}
}
\only<67>{
\scalebox{ 0.40}{\psset{nodesep=0pt,levelsep=36pt}
\pstree{\CC{\hyperlink{spotgac<1>}{4}}}{
 \psset{linestyle=solid}\pstree{\CC{\hyperlink{spotgac<2>}{5}}\trput{2}}{
  \psset{linestyle=solid}\pstree{\CC{\hyperlink{spotgac<3>}{6}}\trput{4}}{
   \psset{linestyle=solid}\pstree{\CC{\hyperlink{spotgac<4>}{7}}\trput{5}}{
    \psset{linestyle=solid}\FAIL{spotgac<5>}\trput{6}
    \psset{linestyle=solid}\pstree{\CC{\hyperlink{spotgac<6>}{12}}\trput{9}}{
     \psset{linestyle=solid}\pstree{\CC{\hyperlink{spotgac<7>}{19}}\trput{6}}{
      \psset{linestyle=solid}\FAIL{spotgac<8>}\trput{6}
      \psset{linestyle=solid}\pstree{\CC{\hyperlink{spotgac<9>}{20}}\trput{9}}{
       \psset{linestyle=solid}\pstree{\CC{\hyperlink{spotgac<10>}{23}}\trput{6}}{
        \psset{linestyle=solid}\FAIL{spotgac<11>}\trput{2}
        \psset{linestyle=solid}\FAIL{spotgac<12>}\trput{8}
       }
       \psset{linestyle=solid}\FAIL{spotgac<13>}\trput{8}
      }
     }
     \psset{linestyle=solid}\FAIL{spotgac<14>}\trput{8}
    }
   }
   \psset{linestyle=solid}\FAIL{spotgac<15>}\trput{6}
  }
  \psset{linestyle=solid}\pstree{\CC{\hyperlink{spotgac<16>}{6}}\trput{6}}{
   \psset{linestyle=solid}\pstree{\CC{\hyperlink{spotgac<17>}{10}}\trput{4}}{
    \psset{linestyle=solid}\FAIL{spotgac<18>}\trput{4}
    \psset{linestyle=solid}\FAIL{spotgac<19>}\trput{6}
    \psset{linestyle=solid}\FAIL{spotgac<20>}\trput{8}
   }
   \psset{linestyle=solid}\pstree{\CC{\hyperlink{spotgac<21>}{7}}\trput{5}}{
    \psset{linestyle=solid}\pstree{\CC{\hyperlink{spotgac<22>}{10}}\trput{4}}{
     \psset{linestyle=solid}\FAIL{spotgac<23>}\trput{5}
     \psset{linestyle=solid}\FAIL{spotgac<24>}\trput{6}
     \psset{linestyle=solid}\FAIL{spotgac<25>}\trput{8}
    }
    \psset{linestyle=solid}\pstree{\CC{\hyperlink{spotgac<26>}{10}}\trput{9}}{
     \psset{linestyle=solid}\pstree{\CC{\hyperlink{spotgac<27>}{11}}\trput{5}}{
      \psset{linestyle=solid}\FAIL{spotgac<28>}\trput{6}
      \psset{linestyle=solid}\FAIL{spotgac<29>}\trput{8}
     }
     \psset{linestyle=solid}\FAIL{spotgac<30>}\trput{6}
     \psset{linestyle=solid}\FAIL{spotgac<31>}\trput{8}
    }
   }
  }
 }
 \psset{linestyle=solid}\pstree{\CC{\hyperlink{spotgac<32>}{5}}\trput{4}}{
  \psset{linestyle=solid}\pstree{\CC{\hyperlink{spotgac<33>}{6}}\trput{2}}{
   \psset{linestyle=solid}\FAIL{spotgac<34>}\trput{5}
   \psset{linestyle=solid}\FAIL{spotgac<35>}\trput{6}
  }
  \psset{linestyle=solid}\pstree{\CC{\hyperlink{spotgac<36>}{6}}\trput{6}}{
   \psset{linestyle=solid}\FAIL{spotgac<37>}\trput{2}
   \psset{linestyle=solid}\FAIL{spotgac<38>}\trput{5}
  }
 }
 \psset{linestyle=solid}\pstree{\CC{\hyperlink{spotgac<39>}{5}}\trput{5}}{
  \psset{linestyle=solid}\pstree{\CC{\hyperlink{spotgac<40>}{6}}\trput{2}}{
   \psset{linestyle=solid}\pstree{\CC{\hyperlink{spotgac<41>}{7}}\trput{4}}{
    \psset{linestyle=solid}\FAIL{spotgac<42>}\trput{6}
    \psset{linestyle=solid}\pstree{\CC{\hyperlink{spotgac<43>}{10}}\trput{9}}{
     \psset{linestyle=solid}\FAIL{spotgac<44>}\trput{6}
     \psset{linestyle=solid}\FAIL{spotgac<45>}\trput{8}
    }
   }
   \psset{linestyle=solid}\pstree{\CC{\hyperlink{spotgac<46>}{7}}\trput{6}}{
    \psset{linestyle=solid}\FAIL{spotgac<47>}\trput{4}
    \psset{linestyle=solid}\FAIL{spotgac<48>}\trput{9}
   }
  }
  \psset{linestyle=solid}\FAIL{spotgac<49>}\trput{4}
  \psset{linestyle=solid}\pstree{\CC{\hyperlink{spotgac<50>}{6}}\trput{6}}{
   \psset{linestyle=solid}\pstree{\CC{\hyperlink{spotgac<51>}{7}}\trput{2}}{
    \psset{linestyle=solid}\FAIL{spotgac<52>}\trput{4}
    \psset{linestyle=solid}\FAIL{spotgac<53>}\trput{9}
   }
   \psset{linestyle=solid}\pstree{\CC{\hyperlink{spotgac<54>}{7}}\trput{4}}{
    \psset{linestyle=solid}\pstree{\CC{\hyperlink{spotgac<55>}{10}}\trput{2}}{
     \psset{linestyle=solid}\pstree{\CC{\hyperlink{spotgac<56>}{12}}\trput{6}}{
      \psset{linestyle=solid}\FAIL{spotgac<57>}\trput{2}
      \psset{linestyle=solid}\FAIL{spotgac<58>}\trput{8}
     }
     \psset{linestyle=solid}\FAIL{spotgac<59>}\trput{8}
    }
    \psset{linestyle=solid}\pstree{\CC{\hyperlink{spotgac<60>}{10}}\trput{9}}{
     \psset{linestyle=solid}\pstree{\CC{\hyperlink{spotgac<61>}{12}}\trput{6}}{
      \psset{linestyle=solid}\pstree{\CC{\hyperlink{spotgac<62>}{22}}\trput{2}}{
       \psset{linestyle=solid}\pstree{\CC{\hyperlink{spotgac<63>}{24}}\trput{2}}{
        \psset{linestyle=solid}\pstree{\CC{\hyperlink{spotgac<64>}{31}}\trput{1}}{
         \psset{linestyle=solid}\FAIL{spotgac<65>}\trput{4}
         \psset{linestyle=solid}\SUCC{spotgac<66>}\trput{8}
        }
        \psset{linestyle=solid}\FAIL{spotgac<67>}\trput{7}
       }
      }
     }
    }
   }
  }
 }
}
}
}
\only<68>{
\scalebox{ 0.40}{\psset{nodesep=0pt,levelsep=36pt}
\pstree{\CC{\hyperlink{spotgac<1>}{4}}}{
 \psset{linestyle=solid}\pstree{\CC{\hyperlink{spotgac<2>}{5}}\trput{2}}{
  \psset{linestyle=solid}\pstree{\CC{\hyperlink{spotgac<3>}{6}}\trput{4}}{
   \psset{linestyle=solid}\pstree{\CC{\hyperlink{spotgac<4>}{7}}\trput{5}}{
    \psset{linestyle=solid}\FAIL{spotgac<5>}\trput{6}
    \psset{linestyle=solid}\pstree{\CC{\hyperlink{spotgac<6>}{12}}\trput{9}}{
     \psset{linestyle=solid}\pstree{\CC{\hyperlink{spotgac<7>}{19}}\trput{6}}{
      \psset{linestyle=solid}\FAIL{spotgac<8>}\trput{6}
      \psset{linestyle=solid}\pstree{\CC{\hyperlink{spotgac<9>}{20}}\trput{9}}{
       \psset{linestyle=solid}\pstree{\CC{\hyperlink{spotgac<10>}{23}}\trput{6}}{
        \psset{linestyle=solid}\FAIL{spotgac<11>}\trput{2}
        \psset{linestyle=solid}\FAIL{spotgac<12>}\trput{8}
       }
       \psset{linestyle=solid}\FAIL{spotgac<13>}\trput{8}
      }
     }
     \psset{linestyle=solid}\FAIL{spotgac<14>}\trput{8}
    }
   }
   \psset{linestyle=solid}\FAIL{spotgac<15>}\trput{6}
  }
  \psset{linestyle=solid}\pstree{\CC{\hyperlink{spotgac<16>}{6}}\trput{6}}{
   \psset{linestyle=solid}\pstree{\CC{\hyperlink{spotgac<17>}{10}}\trput{4}}{
    \psset{linestyle=solid}\FAIL{spotgac<18>}\trput{4}
    \psset{linestyle=solid}\FAIL{spotgac<19>}\trput{6}
    \psset{linestyle=solid}\FAIL{spotgac<20>}\trput{8}
   }
   \psset{linestyle=solid}\pstree{\CC{\hyperlink{spotgac<21>}{7}}\trput{5}}{
    \psset{linestyle=solid}\pstree{\CC{\hyperlink{spotgac<22>}{10}}\trput{4}}{
     \psset{linestyle=solid}\FAIL{spotgac<23>}\trput{5}
     \psset{linestyle=solid}\FAIL{spotgac<24>}\trput{6}
     \psset{linestyle=solid}\FAIL{spotgac<25>}\trput{8}
    }
    \psset{linestyle=solid}\pstree{\CC{\hyperlink{spotgac<26>}{10}}\trput{9}}{
     \psset{linestyle=solid}\pstree{\CC{\hyperlink{spotgac<27>}{11}}\trput{5}}{
      \psset{linestyle=solid}\FAIL{spotgac<28>}\trput{6}
      \psset{linestyle=solid}\FAIL{spotgac<29>}\trput{8}
     }
     \psset{linestyle=solid}\FAIL{spotgac<30>}\trput{6}
     \psset{linestyle=solid}\FAIL{spotgac<31>}\trput{8}
    }
   }
  }
 }
 \psset{linestyle=solid}\pstree{\CC{\hyperlink{spotgac<32>}{5}}\trput{4}}{
  \psset{linestyle=solid}\pstree{\CC{\hyperlink{spotgac<33>}{6}}\trput{2}}{
   \psset{linestyle=solid}\FAIL{spotgac<34>}\trput{5}
   \psset{linestyle=solid}\FAIL{spotgac<35>}\trput{6}
  }
  \psset{linestyle=solid}\pstree{\CC{\hyperlink{spotgac<36>}{6}}\trput{6}}{
   \psset{linestyle=solid}\FAIL{spotgac<37>}\trput{2}
   \psset{linestyle=solid}\FAIL{spotgac<38>}\trput{5}
  }
 }
 \psset{linestyle=solid}\pstree{\CC{\hyperlink{spotgac<39>}{5}}\trput{5}}{
  \psset{linestyle=solid}\pstree{\CC{\hyperlink{spotgac<40>}{6}}\trput{2}}{
   \psset{linestyle=solid}\pstree{\CC{\hyperlink{spotgac<41>}{7}}\trput{4}}{
    \psset{linestyle=solid}\FAIL{spotgac<42>}\trput{6}
    \psset{linestyle=solid}\pstree{\CC{\hyperlink{spotgac<43>}{10}}\trput{9}}{
     \psset{linestyle=solid}\FAIL{spotgac<44>}\trput{6}
     \psset{linestyle=solid}\FAIL{spotgac<45>}\trput{8}
    }
   }
   \psset{linestyle=solid}\pstree{\CC{\hyperlink{spotgac<46>}{7}}\trput{6}}{
    \psset{linestyle=solid}\FAIL{spotgac<47>}\trput{4}
    \psset{linestyle=solid}\FAIL{spotgac<48>}\trput{9}
   }
  }
  \psset{linestyle=solid}\FAIL{spotgac<49>}\trput{4}
  \psset{linestyle=solid}\pstree{\CC{\hyperlink{spotgac<50>}{6}}\trput{6}}{
   \psset{linestyle=solid}\pstree{\CC{\hyperlink{spotgac<51>}{7}}\trput{2}}{
    \psset{linestyle=solid}\FAIL{spotgac<52>}\trput{4}
    \psset{linestyle=solid}\FAIL{spotgac<53>}\trput{9}
   }
   \psset{linestyle=solid}\pstree{\CC{\hyperlink{spotgac<54>}{7}}\trput{4}}{
    \psset{linestyle=solid}\pstree{\CC{\hyperlink{spotgac<55>}{10}}\trput{2}}{
     \psset{linestyle=solid}\pstree{\CC{\hyperlink{spotgac<56>}{12}}\trput{6}}{
      \psset{linestyle=solid}\FAIL{spotgac<57>}\trput{2}
      \psset{linestyle=solid}\FAIL{spotgac<58>}\trput{8}
     }
     \psset{linestyle=solid}\FAIL{spotgac<59>}\trput{8}
    }
    \psset{linestyle=solid}\pstree{\CC{\hyperlink{spotgac<60>}{10}}\trput{9}}{
     \psset{linestyle=solid}\pstree{\CC{\hyperlink{spotgac<61>}{12}}\trput{6}}{
      \psset{linestyle=solid}\pstree{\CC{\hyperlink{spotgac<62>}{22}}\trput{2}}{
       \psset{linestyle=solid}\pstree{\CC{\hyperlink{spotgac<63>}{24}}\trput{2}}{
        \psset{linestyle=solid}\pstree{\CC{\hyperlink{spotgac<64>}{31}}\trput{1}}{
         \psset{linestyle=solid}\FAIL{spotgac<65>}\trput{4}
         \psset{linestyle=solid}\SUCC{spotgac<66>}\trput{8}
        }
        \psset{linestyle=solid}\FAIL{spotgac<67>}\trput{7}
       }
       \psset{linestyle=solid}\FAIL{spotgac<68>}\trput{4}
      }
     }
    }
   }
  }
 }
}
}
}
\only<69>{
\scalebox{ 0.40}{\psset{nodesep=0pt,levelsep=36pt}
\pstree{\CC{\hyperlink{spotgac<1>}{4}}}{
 \psset{linestyle=solid}\pstree{\CC{\hyperlink{spotgac<2>}{5}}\trput{2}}{
  \psset{linestyle=solid}\pstree{\CC{\hyperlink{spotgac<3>}{6}}\trput{4}}{
   \psset{linestyle=solid}\pstree{\CC{\hyperlink{spotgac<4>}{7}}\trput{5}}{
    \psset{linestyle=solid}\FAIL{spotgac<5>}\trput{6}
    \psset{linestyle=solid}\pstree{\CC{\hyperlink{spotgac<6>}{12}}\trput{9}}{
     \psset{linestyle=solid}\pstree{\CC{\hyperlink{spotgac<7>}{19}}\trput{6}}{
      \psset{linestyle=solid}\FAIL{spotgac<8>}\trput{6}
      \psset{linestyle=solid}\pstree{\CC{\hyperlink{spotgac<9>}{20}}\trput{9}}{
       \psset{linestyle=solid}\pstree{\CC{\hyperlink{spotgac<10>}{23}}\trput{6}}{
        \psset{linestyle=solid}\FAIL{spotgac<11>}\trput{2}
        \psset{linestyle=solid}\FAIL{spotgac<12>}\trput{8}
       }
       \psset{linestyle=solid}\FAIL{spotgac<13>}\trput{8}
      }
     }
     \psset{linestyle=solid}\FAIL{spotgac<14>}\trput{8}
    }
   }
   \psset{linestyle=solid}\FAIL{spotgac<15>}\trput{6}
  }
  \psset{linestyle=solid}\pstree{\CC{\hyperlink{spotgac<16>}{6}}\trput{6}}{
   \psset{linestyle=solid}\pstree{\CC{\hyperlink{spotgac<17>}{10}}\trput{4}}{
    \psset{linestyle=solid}\FAIL{spotgac<18>}\trput{4}
    \psset{linestyle=solid}\FAIL{spotgac<19>}\trput{6}
    \psset{linestyle=solid}\FAIL{spotgac<20>}\trput{8}
   }
   \psset{linestyle=solid}\pstree{\CC{\hyperlink{spotgac<21>}{7}}\trput{5}}{
    \psset{linestyle=solid}\pstree{\CC{\hyperlink{spotgac<22>}{10}}\trput{4}}{
     \psset{linestyle=solid}\FAIL{spotgac<23>}\trput{5}
     \psset{linestyle=solid}\FAIL{spotgac<24>}\trput{6}
     \psset{linestyle=solid}\FAIL{spotgac<25>}\trput{8}
    }
    \psset{linestyle=solid}\pstree{\CC{\hyperlink{spotgac<26>}{10}}\trput{9}}{
     \psset{linestyle=solid}\pstree{\CC{\hyperlink{spotgac<27>}{11}}\trput{5}}{
      \psset{linestyle=solid}\FAIL{spotgac<28>}\trput{6}
      \psset{linestyle=solid}\FAIL{spotgac<29>}\trput{8}
     }
     \psset{linestyle=solid}\FAIL{spotgac<30>}\trput{6}
     \psset{linestyle=solid}\FAIL{spotgac<31>}\trput{8}
    }
   }
  }
 }
 \psset{linestyle=solid}\pstree{\CC{\hyperlink{spotgac<32>}{5}}\trput{4}}{
  \psset{linestyle=solid}\pstree{\CC{\hyperlink{spotgac<33>}{6}}\trput{2}}{
   \psset{linestyle=solid}\FAIL{spotgac<34>}\trput{5}
   \psset{linestyle=solid}\FAIL{spotgac<35>}\trput{6}
  }
  \psset{linestyle=solid}\pstree{\CC{\hyperlink{spotgac<36>}{6}}\trput{6}}{
   \psset{linestyle=solid}\FAIL{spotgac<37>}\trput{2}
   \psset{linestyle=solid}\FAIL{spotgac<38>}\trput{5}
  }
 }
 \psset{linestyle=solid}\pstree{\CC{\hyperlink{spotgac<39>}{5}}\trput{5}}{
  \psset{linestyle=solid}\pstree{\CC{\hyperlink{spotgac<40>}{6}}\trput{2}}{
   \psset{linestyle=solid}\pstree{\CC{\hyperlink{spotgac<41>}{7}}\trput{4}}{
    \psset{linestyle=solid}\FAIL{spotgac<42>}\trput{6}
    \psset{linestyle=solid}\pstree{\CC{\hyperlink{spotgac<43>}{10}}\trput{9}}{
     \psset{linestyle=solid}\FAIL{spotgac<44>}\trput{6}
     \psset{linestyle=solid}\FAIL{spotgac<45>}\trput{8}
    }
   }
   \psset{linestyle=solid}\pstree{\CC{\hyperlink{spotgac<46>}{7}}\trput{6}}{
    \psset{linestyle=solid}\FAIL{spotgac<47>}\trput{4}
    \psset{linestyle=solid}\FAIL{spotgac<48>}\trput{9}
   }
  }
  \psset{linestyle=solid}\FAIL{spotgac<49>}\trput{4}
  \psset{linestyle=solid}\pstree{\CC{\hyperlink{spotgac<50>}{6}}\trput{6}}{
   \psset{linestyle=solid}\pstree{\CC{\hyperlink{spotgac<51>}{7}}\trput{2}}{
    \psset{linestyle=solid}\FAIL{spotgac<52>}\trput{4}
    \psset{linestyle=solid}\FAIL{spotgac<53>}\trput{9}
   }
   \psset{linestyle=solid}\pstree{\CC{\hyperlink{spotgac<54>}{7}}\trput{4}}{
    \psset{linestyle=solid}\pstree{\CC{\hyperlink{spotgac<55>}{10}}\trput{2}}{
     \psset{linestyle=solid}\pstree{\CC{\hyperlink{spotgac<56>}{12}}\trput{6}}{
      \psset{linestyle=solid}\FAIL{spotgac<57>}\trput{2}
      \psset{linestyle=solid}\FAIL{spotgac<58>}\trput{8}
     }
     \psset{linestyle=solid}\FAIL{spotgac<59>}\trput{8}
    }
    \psset{linestyle=solid}\pstree{\CC{\hyperlink{spotgac<60>}{10}}\trput{9}}{
     \psset{linestyle=solid}\pstree{\CC{\hyperlink{spotgac<61>}{12}}\trput{6}}{
      \psset{linestyle=solid}\pstree{\CC{\hyperlink{spotgac<62>}{22}}\trput{2}}{
       \psset{linestyle=solid}\pstree{\CC{\hyperlink{spotgac<63>}{24}}\trput{2}}{
        \psset{linestyle=solid}\pstree{\CC{\hyperlink{spotgac<64>}{31}}\trput{1}}{
         \psset{linestyle=solid}\FAIL{spotgac<65>}\trput{4}
         \psset{linestyle=solid}\SUCC{spotgac<66>}\trput{8}
        }
        \psset{linestyle=solid}\FAIL{spotgac<67>}\trput{7}
       }
       \psset{linestyle=solid}\FAIL{spotgac<68>}\trput{4}
      }
      \psset{linestyle=solid}\pstree{\CC{\hyperlink{spotgac<69>}{22}}\trput{8}}{
      }
     }
    }
   }
  }
 }
}
}
}
\only<70>{
\scalebox{ 0.40}{\psset{nodesep=0pt,levelsep=36pt}
\pstree{\CC{\hyperlink{spotgac<1>}{4}}}{
 \psset{linestyle=solid}\pstree{\CC{\hyperlink{spotgac<2>}{5}}\trput{2}}{
  \psset{linestyle=solid}\pstree{\CC{\hyperlink{spotgac<3>}{6}}\trput{4}}{
   \psset{linestyle=solid}\pstree{\CC{\hyperlink{spotgac<4>}{7}}\trput{5}}{
    \psset{linestyle=solid}\FAIL{spotgac<5>}\trput{6}
    \psset{linestyle=solid}\pstree{\CC{\hyperlink{spotgac<6>}{12}}\trput{9}}{
     \psset{linestyle=solid}\pstree{\CC{\hyperlink{spotgac<7>}{19}}\trput{6}}{
      \psset{linestyle=solid}\FAIL{spotgac<8>}\trput{6}
      \psset{linestyle=solid}\pstree{\CC{\hyperlink{spotgac<9>}{20}}\trput{9}}{
       \psset{linestyle=solid}\pstree{\CC{\hyperlink{spotgac<10>}{23}}\trput{6}}{
        \psset{linestyle=solid}\FAIL{spotgac<11>}\trput{2}
        \psset{linestyle=solid}\FAIL{spotgac<12>}\trput{8}
       }
       \psset{linestyle=solid}\FAIL{spotgac<13>}\trput{8}
      }
     }
     \psset{linestyle=solid}\FAIL{spotgac<14>}\trput{8}
    }
   }
   \psset{linestyle=solid}\FAIL{spotgac<15>}\trput{6}
  }
  \psset{linestyle=solid}\pstree{\CC{\hyperlink{spotgac<16>}{6}}\trput{6}}{
   \psset{linestyle=solid}\pstree{\CC{\hyperlink{spotgac<17>}{10}}\trput{4}}{
    \psset{linestyle=solid}\FAIL{spotgac<18>}\trput{4}
    \psset{linestyle=solid}\FAIL{spotgac<19>}\trput{6}
    \psset{linestyle=solid}\FAIL{spotgac<20>}\trput{8}
   }
   \psset{linestyle=solid}\pstree{\CC{\hyperlink{spotgac<21>}{7}}\trput{5}}{
    \psset{linestyle=solid}\pstree{\CC{\hyperlink{spotgac<22>}{10}}\trput{4}}{
     \psset{linestyle=solid}\FAIL{spotgac<23>}\trput{5}
     \psset{linestyle=solid}\FAIL{spotgac<24>}\trput{6}
     \psset{linestyle=solid}\FAIL{spotgac<25>}\trput{8}
    }
    \psset{linestyle=solid}\pstree{\CC{\hyperlink{spotgac<26>}{10}}\trput{9}}{
     \psset{linestyle=solid}\pstree{\CC{\hyperlink{spotgac<27>}{11}}\trput{5}}{
      \psset{linestyle=solid}\FAIL{spotgac<28>}\trput{6}
      \psset{linestyle=solid}\FAIL{spotgac<29>}\trput{8}
     }
     \psset{linestyle=solid}\FAIL{spotgac<30>}\trput{6}
     \psset{linestyle=solid}\FAIL{spotgac<31>}\trput{8}
    }
   }
  }
 }
 \psset{linestyle=solid}\pstree{\CC{\hyperlink{spotgac<32>}{5}}\trput{4}}{
  \psset{linestyle=solid}\pstree{\CC{\hyperlink{spotgac<33>}{6}}\trput{2}}{
   \psset{linestyle=solid}\FAIL{spotgac<34>}\trput{5}
   \psset{linestyle=solid}\FAIL{spotgac<35>}\trput{6}
  }
  \psset{linestyle=solid}\pstree{\CC{\hyperlink{spotgac<36>}{6}}\trput{6}}{
   \psset{linestyle=solid}\FAIL{spotgac<37>}\trput{2}
   \psset{linestyle=solid}\FAIL{spotgac<38>}\trput{5}
  }
 }
 \psset{linestyle=solid}\pstree{\CC{\hyperlink{spotgac<39>}{5}}\trput{5}}{
  \psset{linestyle=solid}\pstree{\CC{\hyperlink{spotgac<40>}{6}}\trput{2}}{
   \psset{linestyle=solid}\pstree{\CC{\hyperlink{spotgac<41>}{7}}\trput{4}}{
    \psset{linestyle=solid}\FAIL{spotgac<42>}\trput{6}
    \psset{linestyle=solid}\pstree{\CC{\hyperlink{spotgac<43>}{10}}\trput{9}}{
     \psset{linestyle=solid}\FAIL{spotgac<44>}\trput{6}
     \psset{linestyle=solid}\FAIL{spotgac<45>}\trput{8}
    }
   }
   \psset{linestyle=solid}\pstree{\CC{\hyperlink{spotgac<46>}{7}}\trput{6}}{
    \psset{linestyle=solid}\FAIL{spotgac<47>}\trput{4}
    \psset{linestyle=solid}\FAIL{spotgac<48>}\trput{9}
   }
  }
  \psset{linestyle=solid}\FAIL{spotgac<49>}\trput{4}
  \psset{linestyle=solid}\pstree{\CC{\hyperlink{spotgac<50>}{6}}\trput{6}}{
   \psset{linestyle=solid}\pstree{\CC{\hyperlink{spotgac<51>}{7}}\trput{2}}{
    \psset{linestyle=solid}\FAIL{spotgac<52>}\trput{4}
    \psset{linestyle=solid}\FAIL{spotgac<53>}\trput{9}
   }
   \psset{linestyle=solid}\pstree{\CC{\hyperlink{spotgac<54>}{7}}\trput{4}}{
    \psset{linestyle=solid}\pstree{\CC{\hyperlink{spotgac<55>}{10}}\trput{2}}{
     \psset{linestyle=solid}\pstree{\CC{\hyperlink{spotgac<56>}{12}}\trput{6}}{
      \psset{linestyle=solid}\FAIL{spotgac<57>}\trput{2}
      \psset{linestyle=solid}\FAIL{spotgac<58>}\trput{8}
     }
     \psset{linestyle=solid}\FAIL{spotgac<59>}\trput{8}
    }
    \psset{linestyle=solid}\pstree{\CC{\hyperlink{spotgac<60>}{10}}\trput{9}}{
     \psset{linestyle=solid}\pstree{\CC{\hyperlink{spotgac<61>}{12}}\trput{6}}{
      \psset{linestyle=solid}\pstree{\CC{\hyperlink{spotgac<62>}{22}}\trput{2}}{
       \psset{linestyle=solid}\pstree{\CC{\hyperlink{spotgac<63>}{24}}\trput{2}}{
        \psset{linestyle=solid}\pstree{\CC{\hyperlink{spotgac<64>}{31}}\trput{1}}{
         \psset{linestyle=solid}\FAIL{spotgac<65>}\trput{4}
         \psset{linestyle=solid}\SUCC{spotgac<66>}\trput{8}
        }
        \psset{linestyle=solid}\FAIL{spotgac<67>}\trput{7}
       }
       \psset{linestyle=solid}\FAIL{spotgac<68>}\trput{4}
      }
      \psset{linestyle=solid}\pstree{\CC{\hyperlink{spotgac<69>}{22}}\trput{8}}{
       \psset{linestyle=solid}\pstree{\CC{\hyperlink{spotgac<70>}{24}}\trput{2}}{
       }
      }
     }
    }
   }
  }
 }
}
}
}
\only<71>{
\scalebox{ 0.30}{\psset{nodesep=0pt,levelsep=36pt}
\pstree{\CC{\hyperlink{spotgac<1>}{4}}}{
 \psset{linestyle=solid}\pstree{\CC{\hyperlink{spotgac<2>}{5}}\trput{2}}{
  \psset{linestyle=solid}\pstree{\CC{\hyperlink{spotgac<3>}{6}}\trput{4}}{
   \psset{linestyle=solid}\pstree{\CC{\hyperlink{spotgac<4>}{7}}\trput{5}}{
    \psset{linestyle=solid}\FAIL{spotgac<5>}\trput{6}
    \psset{linestyle=solid}\pstree{\CC{\hyperlink{spotgac<6>}{12}}\trput{9}}{
     \psset{linestyle=solid}\pstree{\CC{\hyperlink{spotgac<7>}{19}}\trput{6}}{
      \psset{linestyle=solid}\FAIL{spotgac<8>}\trput{6}
      \psset{linestyle=solid}\pstree{\CC{\hyperlink{spotgac<9>}{20}}\trput{9}}{
       \psset{linestyle=solid}\pstree{\CC{\hyperlink{spotgac<10>}{23}}\trput{6}}{
        \psset{linestyle=solid}\FAIL{spotgac<11>}\trput{2}
        \psset{linestyle=solid}\FAIL{spotgac<12>}\trput{8}
       }
       \psset{linestyle=solid}\FAIL{spotgac<13>}\trput{8}
      }
     }
     \psset{linestyle=solid}\FAIL{spotgac<14>}\trput{8}
    }
   }
   \psset{linestyle=solid}\FAIL{spotgac<15>}\trput{6}
  }
  \psset{linestyle=solid}\pstree{\CC{\hyperlink{spotgac<16>}{6}}\trput{6}}{
   \psset{linestyle=solid}\pstree{\CC{\hyperlink{spotgac<17>}{10}}\trput{4}}{
    \psset{linestyle=solid}\FAIL{spotgac<18>}\trput{4}
    \psset{linestyle=solid}\FAIL{spotgac<19>}\trput{6}
    \psset{linestyle=solid}\FAIL{spotgac<20>}\trput{8}
   }
   \psset{linestyle=solid}\pstree{\CC{\hyperlink{spotgac<21>}{7}}\trput{5}}{
    \psset{linestyle=solid}\pstree{\CC{\hyperlink{spotgac<22>}{10}}\trput{4}}{
     \psset{linestyle=solid}\FAIL{spotgac<23>}\trput{5}
     \psset{linestyle=solid}\FAIL{spotgac<24>}\trput{6}
     \psset{linestyle=solid}\FAIL{spotgac<25>}\trput{8}
    }
    \psset{linestyle=solid}\pstree{\CC{\hyperlink{spotgac<26>}{10}}\trput{9}}{
     \psset{linestyle=solid}\pstree{\CC{\hyperlink{spotgac<27>}{11}}\trput{5}}{
      \psset{linestyle=solid}\FAIL{spotgac<28>}\trput{6}
      \psset{linestyle=solid}\FAIL{spotgac<29>}\trput{8}
     }
     \psset{linestyle=solid}\FAIL{spotgac<30>}\trput{6}
     \psset{linestyle=solid}\FAIL{spotgac<31>}\trput{8}
    }
   }
  }
 }
 \psset{linestyle=solid}\pstree{\CC{\hyperlink{spotgac<32>}{5}}\trput{4}}{
  \psset{linestyle=solid}\pstree{\CC{\hyperlink{spotgac<33>}{6}}\trput{2}}{
   \psset{linestyle=solid}\FAIL{spotgac<34>}\trput{5}
   \psset{linestyle=solid}\FAIL{spotgac<35>}\trput{6}
  }
  \psset{linestyle=solid}\pstree{\CC{\hyperlink{spotgac<36>}{6}}\trput{6}}{
   \psset{linestyle=solid}\FAIL{spotgac<37>}\trput{2}
   \psset{linestyle=solid}\FAIL{spotgac<38>}\trput{5}
  }
 }
 \psset{linestyle=solid}\pstree{\CC{\hyperlink{spotgac<39>}{5}}\trput{5}}{
  \psset{linestyle=solid}\pstree{\CC{\hyperlink{spotgac<40>}{6}}\trput{2}}{
   \psset{linestyle=solid}\pstree{\CC{\hyperlink{spotgac<41>}{7}}\trput{4}}{
    \psset{linestyle=solid}\FAIL{spotgac<42>}\trput{6}
    \psset{linestyle=solid}\pstree{\CC{\hyperlink{spotgac<43>}{10}}\trput{9}}{
     \psset{linestyle=solid}\FAIL{spotgac<44>}\trput{6}
     \psset{linestyle=solid}\FAIL{spotgac<45>}\trput{8}
    }
   }
   \psset{linestyle=solid}\pstree{\CC{\hyperlink{spotgac<46>}{7}}\trput{6}}{
    \psset{linestyle=solid}\FAIL{spotgac<47>}\trput{4}
    \psset{linestyle=solid}\FAIL{spotgac<48>}\trput{9}
   }
  }
  \psset{linestyle=solid}\FAIL{spotgac<49>}\trput{4}
  \psset{linestyle=solid}\pstree{\CC{\hyperlink{spotgac<50>}{6}}\trput{6}}{
   \psset{linestyle=solid}\pstree{\CC{\hyperlink{spotgac<51>}{7}}\trput{2}}{
    \psset{linestyle=solid}\FAIL{spotgac<52>}\trput{4}
    \psset{linestyle=solid}\FAIL{spotgac<53>}\trput{9}
   }
   \psset{linestyle=solid}\pstree{\CC{\hyperlink{spotgac<54>}{7}}\trput{4}}{
    \psset{linestyle=solid}\pstree{\CC{\hyperlink{spotgac<55>}{10}}\trput{2}}{
     \psset{linestyle=solid}\pstree{\CC{\hyperlink{spotgac<56>}{12}}\trput{6}}{
      \psset{linestyle=solid}\FAIL{spotgac<57>}\trput{2}
      \psset{linestyle=solid}\FAIL{spotgac<58>}\trput{8}
     }
     \psset{linestyle=solid}\FAIL{spotgac<59>}\trput{8}
    }
    \psset{linestyle=solid}\pstree{\CC{\hyperlink{spotgac<60>}{10}}\trput{9}}{
     \psset{linestyle=solid}\pstree{\CC{\hyperlink{spotgac<61>}{12}}\trput{6}}{
      \psset{linestyle=solid}\pstree{\CC{\hyperlink{spotgac<62>}{22}}\trput{2}}{
       \psset{linestyle=solid}\pstree{\CC{\hyperlink{spotgac<63>}{24}}\trput{2}}{
        \psset{linestyle=solid}\pstree{\CC{\hyperlink{spotgac<64>}{31}}\trput{1}}{
         \psset{linestyle=solid}\FAIL{spotgac<65>}\trput{4}
         \psset{linestyle=solid}\SUCC{spotgac<66>}\trput{8}
        }
        \psset{linestyle=solid}\FAIL{spotgac<67>}\trput{7}
       }
       \psset{linestyle=solid}\FAIL{spotgac<68>}\trput{4}
      }
      \psset{linestyle=solid}\pstree{\CC{\hyperlink{spotgac<69>}{22}}\trput{8}}{
       \psset{linestyle=solid}\pstree{\CC{\hyperlink{spotgac<70>}{24}}\trput{2}}{
        \psset{linestyle=solid}\FAIL{spotgac<71>}\trput{1}
       }
      }
     }
    }
   }
  }
 }
}
}
}
\only<72>{
\scalebox{ 0.30}{\psset{nodesep=0pt,levelsep=36pt}
\pstree{\CC{\hyperlink{spotgac<1>}{4}}}{
 \psset{linestyle=solid}\pstree{\CC{\hyperlink{spotgac<2>}{5}}\trput{2}}{
  \psset{linestyle=solid}\pstree{\CC{\hyperlink{spotgac<3>}{6}}\trput{4}}{
   \psset{linestyle=solid}\pstree{\CC{\hyperlink{spotgac<4>}{7}}\trput{5}}{
    \psset{linestyle=solid}\FAIL{spotgac<5>}\trput{6}
    \psset{linestyle=solid}\pstree{\CC{\hyperlink{spotgac<6>}{12}}\trput{9}}{
     \psset{linestyle=solid}\pstree{\CC{\hyperlink{spotgac<7>}{19}}\trput{6}}{
      \psset{linestyle=solid}\FAIL{spotgac<8>}\trput{6}
      \psset{linestyle=solid}\pstree{\CC{\hyperlink{spotgac<9>}{20}}\trput{9}}{
       \psset{linestyle=solid}\pstree{\CC{\hyperlink{spotgac<10>}{23}}\trput{6}}{
        \psset{linestyle=solid}\FAIL{spotgac<11>}\trput{2}
        \psset{linestyle=solid}\FAIL{spotgac<12>}\trput{8}
       }
       \psset{linestyle=solid}\FAIL{spotgac<13>}\trput{8}
      }
     }
     \psset{linestyle=solid}\FAIL{spotgac<14>}\trput{8}
    }
   }
   \psset{linestyle=solid}\FAIL{spotgac<15>}\trput{6}
  }
  \psset{linestyle=solid}\pstree{\CC{\hyperlink{spotgac<16>}{6}}\trput{6}}{
   \psset{linestyle=solid}\pstree{\CC{\hyperlink{spotgac<17>}{10}}\trput{4}}{
    \psset{linestyle=solid}\FAIL{spotgac<18>}\trput{4}
    \psset{linestyle=solid}\FAIL{spotgac<19>}\trput{6}
    \psset{linestyle=solid}\FAIL{spotgac<20>}\trput{8}
   }
   \psset{linestyle=solid}\pstree{\CC{\hyperlink{spotgac<21>}{7}}\trput{5}}{
    \psset{linestyle=solid}\pstree{\CC{\hyperlink{spotgac<22>}{10}}\trput{4}}{
     \psset{linestyle=solid}\FAIL{spotgac<23>}\trput{5}
     \psset{linestyle=solid}\FAIL{spotgac<24>}\trput{6}
     \psset{linestyle=solid}\FAIL{spotgac<25>}\trput{8}
    }
    \psset{linestyle=solid}\pstree{\CC{\hyperlink{spotgac<26>}{10}}\trput{9}}{
     \psset{linestyle=solid}\pstree{\CC{\hyperlink{spotgac<27>}{11}}\trput{5}}{
      \psset{linestyle=solid}\FAIL{spotgac<28>}\trput{6}
      \psset{linestyle=solid}\FAIL{spotgac<29>}\trput{8}
     }
     \psset{linestyle=solid}\FAIL{spotgac<30>}\trput{6}
     \psset{linestyle=solid}\FAIL{spotgac<31>}\trput{8}
    }
   }
  }
 }
 \psset{linestyle=solid}\pstree{\CC{\hyperlink{spotgac<32>}{5}}\trput{4}}{
  \psset{linestyle=solid}\pstree{\CC{\hyperlink{spotgac<33>}{6}}\trput{2}}{
   \psset{linestyle=solid}\FAIL{spotgac<34>}\trput{5}
   \psset{linestyle=solid}\FAIL{spotgac<35>}\trput{6}
  }
  \psset{linestyle=solid}\pstree{\CC{\hyperlink{spotgac<36>}{6}}\trput{6}}{
   \psset{linestyle=solid}\FAIL{spotgac<37>}\trput{2}
   \psset{linestyle=solid}\FAIL{spotgac<38>}\trput{5}
  }
 }
 \psset{linestyle=solid}\pstree{\CC{\hyperlink{spotgac<39>}{5}}\trput{5}}{
  \psset{linestyle=solid}\pstree{\CC{\hyperlink{spotgac<40>}{6}}\trput{2}}{
   \psset{linestyle=solid}\pstree{\CC{\hyperlink{spotgac<41>}{7}}\trput{4}}{
    \psset{linestyle=solid}\FAIL{spotgac<42>}\trput{6}
    \psset{linestyle=solid}\pstree{\CC{\hyperlink{spotgac<43>}{10}}\trput{9}}{
     \psset{linestyle=solid}\FAIL{spotgac<44>}\trput{6}
     \psset{linestyle=solid}\FAIL{spotgac<45>}\trput{8}
    }
   }
   \psset{linestyle=solid}\pstree{\CC{\hyperlink{spotgac<46>}{7}}\trput{6}}{
    \psset{linestyle=solid}\FAIL{spotgac<47>}\trput{4}
    \psset{linestyle=solid}\FAIL{spotgac<48>}\trput{9}
   }
  }
  \psset{linestyle=solid}\FAIL{spotgac<49>}\trput{4}
  \psset{linestyle=solid}\pstree{\CC{\hyperlink{spotgac<50>}{6}}\trput{6}}{
   \psset{linestyle=solid}\pstree{\CC{\hyperlink{spotgac<51>}{7}}\trput{2}}{
    \psset{linestyle=solid}\FAIL{spotgac<52>}\trput{4}
    \psset{linestyle=solid}\FAIL{spotgac<53>}\trput{9}
   }
   \psset{linestyle=solid}\pstree{\CC{\hyperlink{spotgac<54>}{7}}\trput{4}}{
    \psset{linestyle=solid}\pstree{\CC{\hyperlink{spotgac<55>}{10}}\trput{2}}{
     \psset{linestyle=solid}\pstree{\CC{\hyperlink{spotgac<56>}{12}}\trput{6}}{
      \psset{linestyle=solid}\FAIL{spotgac<57>}\trput{2}
      \psset{linestyle=solid}\FAIL{spotgac<58>}\trput{8}
     }
     \psset{linestyle=solid}\FAIL{spotgac<59>}\trput{8}
    }
    \psset{linestyle=solid}\pstree{\CC{\hyperlink{spotgac<60>}{10}}\trput{9}}{
     \psset{linestyle=solid}\pstree{\CC{\hyperlink{spotgac<61>}{12}}\trput{6}}{
      \psset{linestyle=solid}\pstree{\CC{\hyperlink{spotgac<62>}{22}}\trput{2}}{
       \psset{linestyle=solid}\pstree{\CC{\hyperlink{spotgac<63>}{24}}\trput{2}}{
        \psset{linestyle=solid}\pstree{\CC{\hyperlink{spotgac<64>}{31}}\trput{1}}{
         \psset{linestyle=solid}\FAIL{spotgac<65>}\trput{4}
         \psset{linestyle=solid}\SUCC{spotgac<66>}\trput{8}
        }
        \psset{linestyle=solid}\FAIL{spotgac<67>}\trput{7}
       }
       \psset{linestyle=solid}\FAIL{spotgac<68>}\trput{4}
      }
      \psset{linestyle=solid}\pstree{\CC{\hyperlink{spotgac<69>}{22}}\trput{8}}{
       \psset{linestyle=solid}\pstree{\CC{\hyperlink{spotgac<70>}{24}}\trput{2}}{
        \psset{linestyle=solid}\FAIL{spotgac<71>}\trput{1}
        \psset{linestyle=solid}\FAIL{spotgac<72>}\trput{7}
       }
      }
     }
    }
   }
  }
 }
}
}
}
\only<73>{
\scalebox{ 0.30}{\psset{nodesep=0pt,levelsep=36pt}
\pstree{\CC{\hyperlink{spotgac<1>}{4}}}{
 \psset{linestyle=solid}\pstree{\CC{\hyperlink{spotgac<2>}{5}}\trput{2}}{
  \psset{linestyle=solid}\pstree{\CC{\hyperlink{spotgac<3>}{6}}\trput{4}}{
   \psset{linestyle=solid}\pstree{\CC{\hyperlink{spotgac<4>}{7}}\trput{5}}{
    \psset{linestyle=solid}\FAIL{spotgac<5>}\trput{6}
    \psset{linestyle=solid}\pstree{\CC{\hyperlink{spotgac<6>}{12}}\trput{9}}{
     \psset{linestyle=solid}\pstree{\CC{\hyperlink{spotgac<7>}{19}}\trput{6}}{
      \psset{linestyle=solid}\FAIL{spotgac<8>}\trput{6}
      \psset{linestyle=solid}\pstree{\CC{\hyperlink{spotgac<9>}{20}}\trput{9}}{
       \psset{linestyle=solid}\pstree{\CC{\hyperlink{spotgac<10>}{23}}\trput{6}}{
        \psset{linestyle=solid}\FAIL{spotgac<11>}\trput{2}
        \psset{linestyle=solid}\FAIL{spotgac<12>}\trput{8}
       }
       \psset{linestyle=solid}\FAIL{spotgac<13>}\trput{8}
      }
     }
     \psset{linestyle=solid}\FAIL{spotgac<14>}\trput{8}
    }
   }
   \psset{linestyle=solid}\FAIL{spotgac<15>}\trput{6}
  }
  \psset{linestyle=solid}\pstree{\CC{\hyperlink{spotgac<16>}{6}}\trput{6}}{
   \psset{linestyle=solid}\pstree{\CC{\hyperlink{spotgac<17>}{10}}\trput{4}}{
    \psset{linestyle=solid}\FAIL{spotgac<18>}\trput{4}
    \psset{linestyle=solid}\FAIL{spotgac<19>}\trput{6}
    \psset{linestyle=solid}\FAIL{spotgac<20>}\trput{8}
   }
   \psset{linestyle=solid}\pstree{\CC{\hyperlink{spotgac<21>}{7}}\trput{5}}{
    \psset{linestyle=solid}\pstree{\CC{\hyperlink{spotgac<22>}{10}}\trput{4}}{
     \psset{linestyle=solid}\FAIL{spotgac<23>}\trput{5}
     \psset{linestyle=solid}\FAIL{spotgac<24>}\trput{6}
     \psset{linestyle=solid}\FAIL{spotgac<25>}\trput{8}
    }
    \psset{linestyle=solid}\pstree{\CC{\hyperlink{spotgac<26>}{10}}\trput{9}}{
     \psset{linestyle=solid}\pstree{\CC{\hyperlink{spotgac<27>}{11}}\trput{5}}{
      \psset{linestyle=solid}\FAIL{spotgac<28>}\trput{6}
      \psset{linestyle=solid}\FAIL{spotgac<29>}\trput{8}
     }
     \psset{linestyle=solid}\FAIL{spotgac<30>}\trput{6}
     \psset{linestyle=solid}\FAIL{spotgac<31>}\trput{8}
    }
   }
  }
 }
 \psset{linestyle=solid}\pstree{\CC{\hyperlink{spotgac<32>}{5}}\trput{4}}{
  \psset{linestyle=solid}\pstree{\CC{\hyperlink{spotgac<33>}{6}}\trput{2}}{
   \psset{linestyle=solid}\FAIL{spotgac<34>}\trput{5}
   \psset{linestyle=solid}\FAIL{spotgac<35>}\trput{6}
  }
  \psset{linestyle=solid}\pstree{\CC{\hyperlink{spotgac<36>}{6}}\trput{6}}{
   \psset{linestyle=solid}\FAIL{spotgac<37>}\trput{2}
   \psset{linestyle=solid}\FAIL{spotgac<38>}\trput{5}
  }
 }
 \psset{linestyle=solid}\pstree{\CC{\hyperlink{spotgac<39>}{5}}\trput{5}}{
  \psset{linestyle=solid}\pstree{\CC{\hyperlink{spotgac<40>}{6}}\trput{2}}{
   \psset{linestyle=solid}\pstree{\CC{\hyperlink{spotgac<41>}{7}}\trput{4}}{
    \psset{linestyle=solid}\FAIL{spotgac<42>}\trput{6}
    \psset{linestyle=solid}\pstree{\CC{\hyperlink{spotgac<43>}{10}}\trput{9}}{
     \psset{linestyle=solid}\FAIL{spotgac<44>}\trput{6}
     \psset{linestyle=solid}\FAIL{spotgac<45>}\trput{8}
    }
   }
   \psset{linestyle=solid}\pstree{\CC{\hyperlink{spotgac<46>}{7}}\trput{6}}{
    \psset{linestyle=solid}\FAIL{spotgac<47>}\trput{4}
    \psset{linestyle=solid}\FAIL{spotgac<48>}\trput{9}
   }
  }
  \psset{linestyle=solid}\FAIL{spotgac<49>}\trput{4}
  \psset{linestyle=solid}\pstree{\CC{\hyperlink{spotgac<50>}{6}}\trput{6}}{
   \psset{linestyle=solid}\pstree{\CC{\hyperlink{spotgac<51>}{7}}\trput{2}}{
    \psset{linestyle=solid}\FAIL{spotgac<52>}\trput{4}
    \psset{linestyle=solid}\FAIL{spotgac<53>}\trput{9}
   }
   \psset{linestyle=solid}\pstree{\CC{\hyperlink{spotgac<54>}{7}}\trput{4}}{
    \psset{linestyle=solid}\pstree{\CC{\hyperlink{spotgac<55>}{10}}\trput{2}}{
     \psset{linestyle=solid}\pstree{\CC{\hyperlink{spotgac<56>}{12}}\trput{6}}{
      \psset{linestyle=solid}\FAIL{spotgac<57>}\trput{2}
      \psset{linestyle=solid}\FAIL{spotgac<58>}\trput{8}
     }
     \psset{linestyle=solid}\FAIL{spotgac<59>}\trput{8}
    }
    \psset{linestyle=solid}\pstree{\CC{\hyperlink{spotgac<60>}{10}}\trput{9}}{
     \psset{linestyle=solid}\pstree{\CC{\hyperlink{spotgac<61>}{12}}\trput{6}}{
      \psset{linestyle=solid}\pstree{\CC{\hyperlink{spotgac<62>}{22}}\trput{2}}{
       \psset{linestyle=solid}\pstree{\CC{\hyperlink{spotgac<63>}{24}}\trput{2}}{
        \psset{linestyle=solid}\pstree{\CC{\hyperlink{spotgac<64>}{31}}\trput{1}}{
         \psset{linestyle=solid}\FAIL{spotgac<65>}\trput{4}
         \psset{linestyle=solid}\SUCC{spotgac<66>}\trput{8}
        }
        \psset{linestyle=solid}\FAIL{spotgac<67>}\trput{7}
       }
       \psset{linestyle=solid}\FAIL{spotgac<68>}\trput{4}
      }
      \psset{linestyle=solid}\pstree{\CC{\hyperlink{spotgac<69>}{22}}\trput{8}}{
       \psset{linestyle=solid}\pstree{\CC{\hyperlink{spotgac<70>}{24}}\trput{2}}{
        \psset{linestyle=solid}\FAIL{spotgac<71>}\trput{1}
        \psset{linestyle=solid}\FAIL{spotgac<72>}\trput{7}
       }
       \psset{linestyle=solid}\FAIL{spotgac<73>}\trput{4}
      }
     }
    }
   }
  }
 }
}
}
}
\only<74>{
\scalebox{ 0.30}{\psset{nodesep=0pt,levelsep=36pt}
\pstree{\CC{\hyperlink{spotgac<1>}{4}}}{
 \psset{linestyle=solid}\pstree{\CC{\hyperlink{spotgac<2>}{5}}\trput{2}}{
  \psset{linestyle=solid}\pstree{\CC{\hyperlink{spotgac<3>}{6}}\trput{4}}{
   \psset{linestyle=solid}\pstree{\CC{\hyperlink{spotgac<4>}{7}}\trput{5}}{
    \psset{linestyle=solid}\FAIL{spotgac<5>}\trput{6}
    \psset{linestyle=solid}\pstree{\CC{\hyperlink{spotgac<6>}{12}}\trput{9}}{
     \psset{linestyle=solid}\pstree{\CC{\hyperlink{spotgac<7>}{19}}\trput{6}}{
      \psset{linestyle=solid}\FAIL{spotgac<8>}\trput{6}
      \psset{linestyle=solid}\pstree{\CC{\hyperlink{spotgac<9>}{20}}\trput{9}}{
       \psset{linestyle=solid}\pstree{\CC{\hyperlink{spotgac<10>}{23}}\trput{6}}{
        \psset{linestyle=solid}\FAIL{spotgac<11>}\trput{2}
        \psset{linestyle=solid}\FAIL{spotgac<12>}\trput{8}
       }
       \psset{linestyle=solid}\FAIL{spotgac<13>}\trput{8}
      }
     }
     \psset{linestyle=solid}\FAIL{spotgac<14>}\trput{8}
    }
   }
   \psset{linestyle=solid}\FAIL{spotgac<15>}\trput{6}
  }
  \psset{linestyle=solid}\pstree{\CC{\hyperlink{spotgac<16>}{6}}\trput{6}}{
   \psset{linestyle=solid}\pstree{\CC{\hyperlink{spotgac<17>}{10}}\trput{4}}{
    \psset{linestyle=solid}\FAIL{spotgac<18>}\trput{4}
    \psset{linestyle=solid}\FAIL{spotgac<19>}\trput{6}
    \psset{linestyle=solid}\FAIL{spotgac<20>}\trput{8}
   }
   \psset{linestyle=solid}\pstree{\CC{\hyperlink{spotgac<21>}{7}}\trput{5}}{
    \psset{linestyle=solid}\pstree{\CC{\hyperlink{spotgac<22>}{10}}\trput{4}}{
     \psset{linestyle=solid}\FAIL{spotgac<23>}\trput{5}
     \psset{linestyle=solid}\FAIL{spotgac<24>}\trput{6}
     \psset{linestyle=solid}\FAIL{spotgac<25>}\trput{8}
    }
    \psset{linestyle=solid}\pstree{\CC{\hyperlink{spotgac<26>}{10}}\trput{9}}{
     \psset{linestyle=solid}\pstree{\CC{\hyperlink{spotgac<27>}{11}}\trput{5}}{
      \psset{linestyle=solid}\FAIL{spotgac<28>}\trput{6}
      \psset{linestyle=solid}\FAIL{spotgac<29>}\trput{8}
     }
     \psset{linestyle=solid}\FAIL{spotgac<30>}\trput{6}
     \psset{linestyle=solid}\FAIL{spotgac<31>}\trput{8}
    }
   }
  }
 }
 \psset{linestyle=solid}\pstree{\CC{\hyperlink{spotgac<32>}{5}}\trput{4}}{
  \psset{linestyle=solid}\pstree{\CC{\hyperlink{spotgac<33>}{6}}\trput{2}}{
   \psset{linestyle=solid}\FAIL{spotgac<34>}\trput{5}
   \psset{linestyle=solid}\FAIL{spotgac<35>}\trput{6}
  }
  \psset{linestyle=solid}\pstree{\CC{\hyperlink{spotgac<36>}{6}}\trput{6}}{
   \psset{linestyle=solid}\FAIL{spotgac<37>}\trput{2}
   \psset{linestyle=solid}\FAIL{spotgac<38>}\trput{5}
  }
 }
 \psset{linestyle=solid}\pstree{\CC{\hyperlink{spotgac<39>}{5}}\trput{5}}{
  \psset{linestyle=solid}\pstree{\CC{\hyperlink{spotgac<40>}{6}}\trput{2}}{
   \psset{linestyle=solid}\pstree{\CC{\hyperlink{spotgac<41>}{7}}\trput{4}}{
    \psset{linestyle=solid}\FAIL{spotgac<42>}\trput{6}
    \psset{linestyle=solid}\pstree{\CC{\hyperlink{spotgac<43>}{10}}\trput{9}}{
     \psset{linestyle=solid}\FAIL{spotgac<44>}\trput{6}
     \psset{linestyle=solid}\FAIL{spotgac<45>}\trput{8}
    }
   }
   \psset{linestyle=solid}\pstree{\CC{\hyperlink{spotgac<46>}{7}}\trput{6}}{
    \psset{linestyle=solid}\FAIL{spotgac<47>}\trput{4}
    \psset{linestyle=solid}\FAIL{spotgac<48>}\trput{9}
   }
  }
  \psset{linestyle=solid}\FAIL{spotgac<49>}\trput{4}
  \psset{linestyle=solid}\pstree{\CC{\hyperlink{spotgac<50>}{6}}\trput{6}}{
   \psset{linestyle=solid}\pstree{\CC{\hyperlink{spotgac<51>}{7}}\trput{2}}{
    \psset{linestyle=solid}\FAIL{spotgac<52>}\trput{4}
    \psset{linestyle=solid}\FAIL{spotgac<53>}\trput{9}
   }
   \psset{linestyle=solid}\pstree{\CC{\hyperlink{spotgac<54>}{7}}\trput{4}}{
    \psset{linestyle=solid}\pstree{\CC{\hyperlink{spotgac<55>}{10}}\trput{2}}{
     \psset{linestyle=solid}\pstree{\CC{\hyperlink{spotgac<56>}{12}}\trput{6}}{
      \psset{linestyle=solid}\FAIL{spotgac<57>}\trput{2}
      \psset{linestyle=solid}\FAIL{spotgac<58>}\trput{8}
     }
     \psset{linestyle=solid}\FAIL{spotgac<59>}\trput{8}
    }
    \psset{linestyle=solid}\pstree{\CC{\hyperlink{spotgac<60>}{10}}\trput{9}}{
     \psset{linestyle=solid}\pstree{\CC{\hyperlink{spotgac<61>}{12}}\trput{6}}{
      \psset{linestyle=solid}\pstree{\CC{\hyperlink{spotgac<62>}{22}}\trput{2}}{
       \psset{linestyle=solid}\pstree{\CC{\hyperlink{spotgac<63>}{24}}\trput{2}}{
        \psset{linestyle=solid}\pstree{\CC{\hyperlink{spotgac<64>}{31}}\trput{1}}{
         \psset{linestyle=solid}\FAIL{spotgac<65>}\trput{4}
         \psset{linestyle=solid}\SUCC{spotgac<66>}\trput{8}
        }
        \psset{linestyle=solid}\FAIL{spotgac<67>}\trput{7}
       }
       \psset{linestyle=solid}\FAIL{spotgac<68>}\trput{4}
      }
      \psset{linestyle=solid}\pstree{\CC{\hyperlink{spotgac<69>}{22}}\trput{8}}{
       \psset{linestyle=solid}\pstree{\CC{\hyperlink{spotgac<70>}{24}}\trput{2}}{
        \psset{linestyle=solid}\FAIL{spotgac<71>}\trput{1}
        \psset{linestyle=solid}\FAIL{spotgac<72>}\trput{7}
       }
       \psset{linestyle=solid}\FAIL{spotgac<73>}\trput{4}
      }
     }
     \psset{linestyle=solid}\FAIL{spotgac<74>}\trput{8}
    }
   }
  }
 }
}
}
}
\only<75>{
\scalebox{ 0.30}{\psset{nodesep=0pt,levelsep=36pt}
\pstree{\CC{\hyperlink{spotgac<1>}{4}}}{
 \psset{linestyle=solid}\pstree{\CC{\hyperlink{spotgac<2>}{5}}\trput{2}}{
  \psset{linestyle=solid}\pstree{\CC{\hyperlink{spotgac<3>}{6}}\trput{4}}{
   \psset{linestyle=solid}\pstree{\CC{\hyperlink{spotgac<4>}{7}}\trput{5}}{
    \psset{linestyle=solid}\FAIL{spotgac<5>}\trput{6}
    \psset{linestyle=solid}\pstree{\CC{\hyperlink{spotgac<6>}{12}}\trput{9}}{
     \psset{linestyle=solid}\pstree{\CC{\hyperlink{spotgac<7>}{19}}\trput{6}}{
      \psset{linestyle=solid}\FAIL{spotgac<8>}\trput{6}
      \psset{linestyle=solid}\pstree{\CC{\hyperlink{spotgac<9>}{20}}\trput{9}}{
       \psset{linestyle=solid}\pstree{\CC{\hyperlink{spotgac<10>}{23}}\trput{6}}{
        \psset{linestyle=solid}\FAIL{spotgac<11>}\trput{2}
        \psset{linestyle=solid}\FAIL{spotgac<12>}\trput{8}
       }
       \psset{linestyle=solid}\FAIL{spotgac<13>}\trput{8}
      }
     }
     \psset{linestyle=solid}\FAIL{spotgac<14>}\trput{8}
    }
   }
   \psset{linestyle=solid}\FAIL{spotgac<15>}\trput{6}
  }
  \psset{linestyle=solid}\pstree{\CC{\hyperlink{spotgac<16>}{6}}\trput{6}}{
   \psset{linestyle=solid}\pstree{\CC{\hyperlink{spotgac<17>}{10}}\trput{4}}{
    \psset{linestyle=solid}\FAIL{spotgac<18>}\trput{4}
    \psset{linestyle=solid}\FAIL{spotgac<19>}\trput{6}
    \psset{linestyle=solid}\FAIL{spotgac<20>}\trput{8}
   }
   \psset{linestyle=solid}\pstree{\CC{\hyperlink{spotgac<21>}{7}}\trput{5}}{
    \psset{linestyle=solid}\pstree{\CC{\hyperlink{spotgac<22>}{10}}\trput{4}}{
     \psset{linestyle=solid}\FAIL{spotgac<23>}\trput{5}
     \psset{linestyle=solid}\FAIL{spotgac<24>}\trput{6}
     \psset{linestyle=solid}\FAIL{spotgac<25>}\trput{8}
    }
    \psset{linestyle=solid}\pstree{\CC{\hyperlink{spotgac<26>}{10}}\trput{9}}{
     \psset{linestyle=solid}\pstree{\CC{\hyperlink{spotgac<27>}{11}}\trput{5}}{
      \psset{linestyle=solid}\FAIL{spotgac<28>}\trput{6}
      \psset{linestyle=solid}\FAIL{spotgac<29>}\trput{8}
     }
     \psset{linestyle=solid}\FAIL{spotgac<30>}\trput{6}
     \psset{linestyle=solid}\FAIL{spotgac<31>}\trput{8}
    }
   }
  }
 }
 \psset{linestyle=solid}\pstree{\CC{\hyperlink{spotgac<32>}{5}}\trput{4}}{
  \psset{linestyle=solid}\pstree{\CC{\hyperlink{spotgac<33>}{6}}\trput{2}}{
   \psset{linestyle=solid}\FAIL{spotgac<34>}\trput{5}
   \psset{linestyle=solid}\FAIL{spotgac<35>}\trput{6}
  }
  \psset{linestyle=solid}\pstree{\CC{\hyperlink{spotgac<36>}{6}}\trput{6}}{
   \psset{linestyle=solid}\FAIL{spotgac<37>}\trput{2}
   \psset{linestyle=solid}\FAIL{spotgac<38>}\trput{5}
  }
 }
 \psset{linestyle=solid}\pstree{\CC{\hyperlink{spotgac<39>}{5}}\trput{5}}{
  \psset{linestyle=solid}\pstree{\CC{\hyperlink{spotgac<40>}{6}}\trput{2}}{
   \psset{linestyle=solid}\pstree{\CC{\hyperlink{spotgac<41>}{7}}\trput{4}}{
    \psset{linestyle=solid}\FAIL{spotgac<42>}\trput{6}
    \psset{linestyle=solid}\pstree{\CC{\hyperlink{spotgac<43>}{10}}\trput{9}}{
     \psset{linestyle=solid}\FAIL{spotgac<44>}\trput{6}
     \psset{linestyle=solid}\FAIL{spotgac<45>}\trput{8}
    }
   }
   \psset{linestyle=solid}\pstree{\CC{\hyperlink{spotgac<46>}{7}}\trput{6}}{
    \psset{linestyle=solid}\FAIL{spotgac<47>}\trput{4}
    \psset{linestyle=solid}\FAIL{spotgac<48>}\trput{9}
   }
  }
  \psset{linestyle=solid}\FAIL{spotgac<49>}\trput{4}
  \psset{linestyle=solid}\pstree{\CC{\hyperlink{spotgac<50>}{6}}\trput{6}}{
   \psset{linestyle=solid}\pstree{\CC{\hyperlink{spotgac<51>}{7}}\trput{2}}{
    \psset{linestyle=solid}\FAIL{spotgac<52>}\trput{4}
    \psset{linestyle=solid}\FAIL{spotgac<53>}\trput{9}
   }
   \psset{linestyle=solid}\pstree{\CC{\hyperlink{spotgac<54>}{7}}\trput{4}}{
    \psset{linestyle=solid}\pstree{\CC{\hyperlink{spotgac<55>}{10}}\trput{2}}{
     \psset{linestyle=solid}\pstree{\CC{\hyperlink{spotgac<56>}{12}}\trput{6}}{
      \psset{linestyle=solid}\FAIL{spotgac<57>}\trput{2}
      \psset{linestyle=solid}\FAIL{spotgac<58>}\trput{8}
     }
     \psset{linestyle=solid}\FAIL{spotgac<59>}\trput{8}
    }
    \psset{linestyle=solid}\pstree{\CC{\hyperlink{spotgac<60>}{10}}\trput{9}}{
     \psset{linestyle=solid}\pstree{\CC{\hyperlink{spotgac<61>}{12}}\trput{6}}{
      \psset{linestyle=solid}\pstree{\CC{\hyperlink{spotgac<62>}{22}}\trput{2}}{
       \psset{linestyle=solid}\pstree{\CC{\hyperlink{spotgac<63>}{24}}\trput{2}}{
        \psset{linestyle=solid}\pstree{\CC{\hyperlink{spotgac<64>}{31}}\trput{1}}{
         \psset{linestyle=solid}\FAIL{spotgac<65>}\trput{4}
         \psset{linestyle=solid}\SUCC{spotgac<66>}\trput{8}
        }
        \psset{linestyle=solid}\FAIL{spotgac<67>}\trput{7}
       }
       \psset{linestyle=solid}\FAIL{spotgac<68>}\trput{4}
      }
      \psset{linestyle=solid}\pstree{\CC{\hyperlink{spotgac<69>}{22}}\trput{8}}{
       \psset{linestyle=solid}\pstree{\CC{\hyperlink{spotgac<70>}{24}}\trput{2}}{
        \psset{linestyle=solid}\FAIL{spotgac<71>}\trput{1}
        \psset{linestyle=solid}\FAIL{spotgac<72>}\trput{7}
       }
       \psset{linestyle=solid}\FAIL{spotgac<73>}\trput{4}
      }
     }
     \psset{linestyle=solid}\FAIL{spotgac<74>}\trput{8}
    }
   }
  }
 }
 \psset{linestyle=solid}\pstree{\CC{\hyperlink{spotgac<75>}{5}}\trput{6}}{
 }
}
}
}
\only<76>{
\scalebox{ 0.30}{\psset{nodesep=0pt,levelsep=36pt}
\pstree{\CC{\hyperlink{spotgac<1>}{4}}}{
 \psset{linestyle=solid}\pstree{\CC{\hyperlink{spotgac<2>}{5}}\trput{2}}{
  \psset{linestyle=solid}\pstree{\CC{\hyperlink{spotgac<3>}{6}}\trput{4}}{
   \psset{linestyle=solid}\pstree{\CC{\hyperlink{spotgac<4>}{7}}\trput{5}}{
    \psset{linestyle=solid}\FAIL{spotgac<5>}\trput{6}
    \psset{linestyle=solid}\pstree{\CC{\hyperlink{spotgac<6>}{12}}\trput{9}}{
     \psset{linestyle=solid}\pstree{\CC{\hyperlink{spotgac<7>}{19}}\trput{6}}{
      \psset{linestyle=solid}\FAIL{spotgac<8>}\trput{6}
      \psset{linestyle=solid}\pstree{\CC{\hyperlink{spotgac<9>}{20}}\trput{9}}{
       \psset{linestyle=solid}\pstree{\CC{\hyperlink{spotgac<10>}{23}}\trput{6}}{
        \psset{linestyle=solid}\FAIL{spotgac<11>}\trput{2}
        \psset{linestyle=solid}\FAIL{spotgac<12>}\trput{8}
       }
       \psset{linestyle=solid}\FAIL{spotgac<13>}\trput{8}
      }
     }
     \psset{linestyle=solid}\FAIL{spotgac<14>}\trput{8}
    }
   }
   \psset{linestyle=solid}\FAIL{spotgac<15>}\trput{6}
  }
  \psset{linestyle=solid}\pstree{\CC{\hyperlink{spotgac<16>}{6}}\trput{6}}{
   \psset{linestyle=solid}\pstree{\CC{\hyperlink{spotgac<17>}{10}}\trput{4}}{
    \psset{linestyle=solid}\FAIL{spotgac<18>}\trput{4}
    \psset{linestyle=solid}\FAIL{spotgac<19>}\trput{6}
    \psset{linestyle=solid}\FAIL{spotgac<20>}\trput{8}
   }
   \psset{linestyle=solid}\pstree{\CC{\hyperlink{spotgac<21>}{7}}\trput{5}}{
    \psset{linestyle=solid}\pstree{\CC{\hyperlink{spotgac<22>}{10}}\trput{4}}{
     \psset{linestyle=solid}\FAIL{spotgac<23>}\trput{5}
     \psset{linestyle=solid}\FAIL{spotgac<24>}\trput{6}
     \psset{linestyle=solid}\FAIL{spotgac<25>}\trput{8}
    }
    \psset{linestyle=solid}\pstree{\CC{\hyperlink{spotgac<26>}{10}}\trput{9}}{
     \psset{linestyle=solid}\pstree{\CC{\hyperlink{spotgac<27>}{11}}\trput{5}}{
      \psset{linestyle=solid}\FAIL{spotgac<28>}\trput{6}
      \psset{linestyle=solid}\FAIL{spotgac<29>}\trput{8}
     }
     \psset{linestyle=solid}\FAIL{spotgac<30>}\trput{6}
     \psset{linestyle=solid}\FAIL{spotgac<31>}\trput{8}
    }
   }
  }
 }
 \psset{linestyle=solid}\pstree{\CC{\hyperlink{spotgac<32>}{5}}\trput{4}}{
  \psset{linestyle=solid}\pstree{\CC{\hyperlink{spotgac<33>}{6}}\trput{2}}{
   \psset{linestyle=solid}\FAIL{spotgac<34>}\trput{5}
   \psset{linestyle=solid}\FAIL{spotgac<35>}\trput{6}
  }
  \psset{linestyle=solid}\pstree{\CC{\hyperlink{spotgac<36>}{6}}\trput{6}}{
   \psset{linestyle=solid}\FAIL{spotgac<37>}\trput{2}
   \psset{linestyle=solid}\FAIL{spotgac<38>}\trput{5}
  }
 }
 \psset{linestyle=solid}\pstree{\CC{\hyperlink{spotgac<39>}{5}}\trput{5}}{
  \psset{linestyle=solid}\pstree{\CC{\hyperlink{spotgac<40>}{6}}\trput{2}}{
   \psset{linestyle=solid}\pstree{\CC{\hyperlink{spotgac<41>}{7}}\trput{4}}{
    \psset{linestyle=solid}\FAIL{spotgac<42>}\trput{6}
    \psset{linestyle=solid}\pstree{\CC{\hyperlink{spotgac<43>}{10}}\trput{9}}{
     \psset{linestyle=solid}\FAIL{spotgac<44>}\trput{6}
     \psset{linestyle=solid}\FAIL{spotgac<45>}\trput{8}
    }
   }
   \psset{linestyle=solid}\pstree{\CC{\hyperlink{spotgac<46>}{7}}\trput{6}}{
    \psset{linestyle=solid}\FAIL{spotgac<47>}\trput{4}
    \psset{linestyle=solid}\FAIL{spotgac<48>}\trput{9}
   }
  }
  \psset{linestyle=solid}\FAIL{spotgac<49>}\trput{4}
  \psset{linestyle=solid}\pstree{\CC{\hyperlink{spotgac<50>}{6}}\trput{6}}{
   \psset{linestyle=solid}\pstree{\CC{\hyperlink{spotgac<51>}{7}}\trput{2}}{
    \psset{linestyle=solid}\FAIL{spotgac<52>}\trput{4}
    \psset{linestyle=solid}\FAIL{spotgac<53>}\trput{9}
   }
   \psset{linestyle=solid}\pstree{\CC{\hyperlink{spotgac<54>}{7}}\trput{4}}{
    \psset{linestyle=solid}\pstree{\CC{\hyperlink{spotgac<55>}{10}}\trput{2}}{
     \psset{linestyle=solid}\pstree{\CC{\hyperlink{spotgac<56>}{12}}\trput{6}}{
      \psset{linestyle=solid}\FAIL{spotgac<57>}\trput{2}
      \psset{linestyle=solid}\FAIL{spotgac<58>}\trput{8}
     }
     \psset{linestyle=solid}\FAIL{spotgac<59>}\trput{8}
    }
    \psset{linestyle=solid}\pstree{\CC{\hyperlink{spotgac<60>}{10}}\trput{9}}{
     \psset{linestyle=solid}\pstree{\CC{\hyperlink{spotgac<61>}{12}}\trput{6}}{
      \psset{linestyle=solid}\pstree{\CC{\hyperlink{spotgac<62>}{22}}\trput{2}}{
       \psset{linestyle=solid}\pstree{\CC{\hyperlink{spotgac<63>}{24}}\trput{2}}{
        \psset{linestyle=solid}\pstree{\CC{\hyperlink{spotgac<64>}{31}}\trput{1}}{
         \psset{linestyle=solid}\FAIL{spotgac<65>}\trput{4}
         \psset{linestyle=solid}\SUCC{spotgac<66>}\trput{8}
        }
        \psset{linestyle=solid}\FAIL{spotgac<67>}\trput{7}
       }
       \psset{linestyle=solid}\FAIL{spotgac<68>}\trput{4}
      }
      \psset{linestyle=solid}\pstree{\CC{\hyperlink{spotgac<69>}{22}}\trput{8}}{
       \psset{linestyle=solid}\pstree{\CC{\hyperlink{spotgac<70>}{24}}\trput{2}}{
        \psset{linestyle=solid}\FAIL{spotgac<71>}\trput{1}
        \psset{linestyle=solid}\FAIL{spotgac<72>}\trput{7}
       }
       \psset{linestyle=solid}\FAIL{spotgac<73>}\trput{4}
      }
     }
     \psset{linestyle=solid}\FAIL{spotgac<74>}\trput{8}
    }
   }
  }
 }
 \psset{linestyle=solid}\pstree{\CC{\hyperlink{spotgac<75>}{5}}\trput{6}}{
  \psset{linestyle=solid}\pstree{\CC{\hyperlink{spotgac<76>}{6}}\trput{2}}{
  }
 }
}
}
}
\only<77>{
\scalebox{ 0.30}{\psset{nodesep=0pt,levelsep=36pt}
\pstree{\CC{\hyperlink{spotgac<1>}{4}}}{
 \psset{linestyle=solid}\pstree{\CC{\hyperlink{spotgac<2>}{5}}\trput{2}}{
  \psset{linestyle=solid}\pstree{\CC{\hyperlink{spotgac<3>}{6}}\trput{4}}{
   \psset{linestyle=solid}\pstree{\CC{\hyperlink{spotgac<4>}{7}}\trput{5}}{
    \psset{linestyle=solid}\FAIL{spotgac<5>}\trput{6}
    \psset{linestyle=solid}\pstree{\CC{\hyperlink{spotgac<6>}{12}}\trput{9}}{
     \psset{linestyle=solid}\pstree{\CC{\hyperlink{spotgac<7>}{19}}\trput{6}}{
      \psset{linestyle=solid}\FAIL{spotgac<8>}\trput{6}
      \psset{linestyle=solid}\pstree{\CC{\hyperlink{spotgac<9>}{20}}\trput{9}}{
       \psset{linestyle=solid}\pstree{\CC{\hyperlink{spotgac<10>}{23}}\trput{6}}{
        \psset{linestyle=solid}\FAIL{spotgac<11>}\trput{2}
        \psset{linestyle=solid}\FAIL{spotgac<12>}\trput{8}
       }
       \psset{linestyle=solid}\FAIL{spotgac<13>}\trput{8}
      }
     }
     \psset{linestyle=solid}\FAIL{spotgac<14>}\trput{8}
    }
   }
   \psset{linestyle=solid}\FAIL{spotgac<15>}\trput{6}
  }
  \psset{linestyle=solid}\pstree{\CC{\hyperlink{spotgac<16>}{6}}\trput{6}}{
   \psset{linestyle=solid}\pstree{\CC{\hyperlink{spotgac<17>}{10}}\trput{4}}{
    \psset{linestyle=solid}\FAIL{spotgac<18>}\trput{4}
    \psset{linestyle=solid}\FAIL{spotgac<19>}\trput{6}
    \psset{linestyle=solid}\FAIL{spotgac<20>}\trput{8}
   }
   \psset{linestyle=solid}\pstree{\CC{\hyperlink{spotgac<21>}{7}}\trput{5}}{
    \psset{linestyle=solid}\pstree{\CC{\hyperlink{spotgac<22>}{10}}\trput{4}}{
     \psset{linestyle=solid}\FAIL{spotgac<23>}\trput{5}
     \psset{linestyle=solid}\FAIL{spotgac<24>}\trput{6}
     \psset{linestyle=solid}\FAIL{spotgac<25>}\trput{8}
    }
    \psset{linestyle=solid}\pstree{\CC{\hyperlink{spotgac<26>}{10}}\trput{9}}{
     \psset{linestyle=solid}\pstree{\CC{\hyperlink{spotgac<27>}{11}}\trput{5}}{
      \psset{linestyle=solid}\FAIL{spotgac<28>}\trput{6}
      \psset{linestyle=solid}\FAIL{spotgac<29>}\trput{8}
     }
     \psset{linestyle=solid}\FAIL{spotgac<30>}\trput{6}
     \psset{linestyle=solid}\FAIL{spotgac<31>}\trput{8}
    }
   }
  }
 }
 \psset{linestyle=solid}\pstree{\CC{\hyperlink{spotgac<32>}{5}}\trput{4}}{
  \psset{linestyle=solid}\pstree{\CC{\hyperlink{spotgac<33>}{6}}\trput{2}}{
   \psset{linestyle=solid}\FAIL{spotgac<34>}\trput{5}
   \psset{linestyle=solid}\FAIL{spotgac<35>}\trput{6}
  }
  \psset{linestyle=solid}\pstree{\CC{\hyperlink{spotgac<36>}{6}}\trput{6}}{
   \psset{linestyle=solid}\FAIL{spotgac<37>}\trput{2}
   \psset{linestyle=solid}\FAIL{spotgac<38>}\trput{5}
  }
 }
 \psset{linestyle=solid}\pstree{\CC{\hyperlink{spotgac<39>}{5}}\trput{5}}{
  \psset{linestyle=solid}\pstree{\CC{\hyperlink{spotgac<40>}{6}}\trput{2}}{
   \psset{linestyle=solid}\pstree{\CC{\hyperlink{spotgac<41>}{7}}\trput{4}}{
    \psset{linestyle=solid}\FAIL{spotgac<42>}\trput{6}
    \psset{linestyle=solid}\pstree{\CC{\hyperlink{spotgac<43>}{10}}\trput{9}}{
     \psset{linestyle=solid}\FAIL{spotgac<44>}\trput{6}
     \psset{linestyle=solid}\FAIL{spotgac<45>}\trput{8}
    }
   }
   \psset{linestyle=solid}\pstree{\CC{\hyperlink{spotgac<46>}{7}}\trput{6}}{
    \psset{linestyle=solid}\FAIL{spotgac<47>}\trput{4}
    \psset{linestyle=solid}\FAIL{spotgac<48>}\trput{9}
   }
  }
  \psset{linestyle=solid}\FAIL{spotgac<49>}\trput{4}
  \psset{linestyle=solid}\pstree{\CC{\hyperlink{spotgac<50>}{6}}\trput{6}}{
   \psset{linestyle=solid}\pstree{\CC{\hyperlink{spotgac<51>}{7}}\trput{2}}{
    \psset{linestyle=solid}\FAIL{spotgac<52>}\trput{4}
    \psset{linestyle=solid}\FAIL{spotgac<53>}\trput{9}
   }
   \psset{linestyle=solid}\pstree{\CC{\hyperlink{spotgac<54>}{7}}\trput{4}}{
    \psset{linestyle=solid}\pstree{\CC{\hyperlink{spotgac<55>}{10}}\trput{2}}{
     \psset{linestyle=solid}\pstree{\CC{\hyperlink{spotgac<56>}{12}}\trput{6}}{
      \psset{linestyle=solid}\FAIL{spotgac<57>}\trput{2}
      \psset{linestyle=solid}\FAIL{spotgac<58>}\trput{8}
     }
     \psset{linestyle=solid}\FAIL{spotgac<59>}\trput{8}
    }
    \psset{linestyle=solid}\pstree{\CC{\hyperlink{spotgac<60>}{10}}\trput{9}}{
     \psset{linestyle=solid}\pstree{\CC{\hyperlink{spotgac<61>}{12}}\trput{6}}{
      \psset{linestyle=solid}\pstree{\CC{\hyperlink{spotgac<62>}{22}}\trput{2}}{
       \psset{linestyle=solid}\pstree{\CC{\hyperlink{spotgac<63>}{24}}\trput{2}}{
        \psset{linestyle=solid}\pstree{\CC{\hyperlink{spotgac<64>}{31}}\trput{1}}{
         \psset{linestyle=solid}\FAIL{spotgac<65>}\trput{4}
         \psset{linestyle=solid}\SUCC{spotgac<66>}\trput{8}
        }
        \psset{linestyle=solid}\FAIL{spotgac<67>}\trput{7}
       }
       \psset{linestyle=solid}\FAIL{spotgac<68>}\trput{4}
      }
      \psset{linestyle=solid}\pstree{\CC{\hyperlink{spotgac<69>}{22}}\trput{8}}{
       \psset{linestyle=solid}\pstree{\CC{\hyperlink{spotgac<70>}{24}}\trput{2}}{
        \psset{linestyle=solid}\FAIL{spotgac<71>}\trput{1}
        \psset{linestyle=solid}\FAIL{spotgac<72>}\trput{7}
       }
       \psset{linestyle=solid}\FAIL{spotgac<73>}\trput{4}
      }
     }
     \psset{linestyle=solid}\FAIL{spotgac<74>}\trput{8}
    }
   }
  }
 }
 \psset{linestyle=solid}\pstree{\CC{\hyperlink{spotgac<75>}{5}}\trput{6}}{
  \psset{linestyle=solid}\pstree{\CC{\hyperlink{spotgac<76>}{6}}\trput{2}}{
   \psset{linestyle=solid}\FAIL{spotgac<77>}\trput{4}
  }
 }
}
}
}
\only<78>{
\scalebox{ 0.30}{\psset{nodesep=0pt,levelsep=36pt}
\pstree{\CC{\hyperlink{spotgac<1>}{4}}}{
 \psset{linestyle=solid}\pstree{\CC{\hyperlink{spotgac<2>}{5}}\trput{2}}{
  \psset{linestyle=solid}\pstree{\CC{\hyperlink{spotgac<3>}{6}}\trput{4}}{
   \psset{linestyle=solid}\pstree{\CC{\hyperlink{spotgac<4>}{7}}\trput{5}}{
    \psset{linestyle=solid}\FAIL{spotgac<5>}\trput{6}
    \psset{linestyle=solid}\pstree{\CC{\hyperlink{spotgac<6>}{12}}\trput{9}}{
     \psset{linestyle=solid}\pstree{\CC{\hyperlink{spotgac<7>}{19}}\trput{6}}{
      \psset{linestyle=solid}\FAIL{spotgac<8>}\trput{6}
      \psset{linestyle=solid}\pstree{\CC{\hyperlink{spotgac<9>}{20}}\trput{9}}{
       \psset{linestyle=solid}\pstree{\CC{\hyperlink{spotgac<10>}{23}}\trput{6}}{
        \psset{linestyle=solid}\FAIL{spotgac<11>}\trput{2}
        \psset{linestyle=solid}\FAIL{spotgac<12>}\trput{8}
       }
       \psset{linestyle=solid}\FAIL{spotgac<13>}\trput{8}
      }
     }
     \psset{linestyle=solid}\FAIL{spotgac<14>}\trput{8}
    }
   }
   \psset{linestyle=solid}\FAIL{spotgac<15>}\trput{6}
  }
  \psset{linestyle=solid}\pstree{\CC{\hyperlink{spotgac<16>}{6}}\trput{6}}{
   \psset{linestyle=solid}\pstree{\CC{\hyperlink{spotgac<17>}{10}}\trput{4}}{
    \psset{linestyle=solid}\FAIL{spotgac<18>}\trput{4}
    \psset{linestyle=solid}\FAIL{spotgac<19>}\trput{6}
    \psset{linestyle=solid}\FAIL{spotgac<20>}\trput{8}
   }
   \psset{linestyle=solid}\pstree{\CC{\hyperlink{spotgac<21>}{7}}\trput{5}}{
    \psset{linestyle=solid}\pstree{\CC{\hyperlink{spotgac<22>}{10}}\trput{4}}{
     \psset{linestyle=solid}\FAIL{spotgac<23>}\trput{5}
     \psset{linestyle=solid}\FAIL{spotgac<24>}\trput{6}
     \psset{linestyle=solid}\FAIL{spotgac<25>}\trput{8}
    }
    \psset{linestyle=solid}\pstree{\CC{\hyperlink{spotgac<26>}{10}}\trput{9}}{
     \psset{linestyle=solid}\pstree{\CC{\hyperlink{spotgac<27>}{11}}\trput{5}}{
      \psset{linestyle=solid}\FAIL{spotgac<28>}\trput{6}
      \psset{linestyle=solid}\FAIL{spotgac<29>}\trput{8}
     }
     \psset{linestyle=solid}\FAIL{spotgac<30>}\trput{6}
     \psset{linestyle=solid}\FAIL{spotgac<31>}\trput{8}
    }
   }
  }
 }
 \psset{linestyle=solid}\pstree{\CC{\hyperlink{spotgac<32>}{5}}\trput{4}}{
  \psset{linestyle=solid}\pstree{\CC{\hyperlink{spotgac<33>}{6}}\trput{2}}{
   \psset{linestyle=solid}\FAIL{spotgac<34>}\trput{5}
   \psset{linestyle=solid}\FAIL{spotgac<35>}\trput{6}
  }
  \psset{linestyle=solid}\pstree{\CC{\hyperlink{spotgac<36>}{6}}\trput{6}}{
   \psset{linestyle=solid}\FAIL{spotgac<37>}\trput{2}
   \psset{linestyle=solid}\FAIL{spotgac<38>}\trput{5}
  }
 }
 \psset{linestyle=solid}\pstree{\CC{\hyperlink{spotgac<39>}{5}}\trput{5}}{
  \psset{linestyle=solid}\pstree{\CC{\hyperlink{spotgac<40>}{6}}\trput{2}}{
   \psset{linestyle=solid}\pstree{\CC{\hyperlink{spotgac<41>}{7}}\trput{4}}{
    \psset{linestyle=solid}\FAIL{spotgac<42>}\trput{6}
    \psset{linestyle=solid}\pstree{\CC{\hyperlink{spotgac<43>}{10}}\trput{9}}{
     \psset{linestyle=solid}\FAIL{spotgac<44>}\trput{6}
     \psset{linestyle=solid}\FAIL{spotgac<45>}\trput{8}
    }
   }
   \psset{linestyle=solid}\pstree{\CC{\hyperlink{spotgac<46>}{7}}\trput{6}}{
    \psset{linestyle=solid}\FAIL{spotgac<47>}\trput{4}
    \psset{linestyle=solid}\FAIL{spotgac<48>}\trput{9}
   }
  }
  \psset{linestyle=solid}\FAIL{spotgac<49>}\trput{4}
  \psset{linestyle=solid}\pstree{\CC{\hyperlink{spotgac<50>}{6}}\trput{6}}{
   \psset{linestyle=solid}\pstree{\CC{\hyperlink{spotgac<51>}{7}}\trput{2}}{
    \psset{linestyle=solid}\FAIL{spotgac<52>}\trput{4}
    \psset{linestyle=solid}\FAIL{spotgac<53>}\trput{9}
   }
   \psset{linestyle=solid}\pstree{\CC{\hyperlink{spotgac<54>}{7}}\trput{4}}{
    \psset{linestyle=solid}\pstree{\CC{\hyperlink{spotgac<55>}{10}}\trput{2}}{
     \psset{linestyle=solid}\pstree{\CC{\hyperlink{spotgac<56>}{12}}\trput{6}}{
      \psset{linestyle=solid}\FAIL{spotgac<57>}\trput{2}
      \psset{linestyle=solid}\FAIL{spotgac<58>}\trput{8}
     }
     \psset{linestyle=solid}\FAIL{spotgac<59>}\trput{8}
    }
    \psset{linestyle=solid}\pstree{\CC{\hyperlink{spotgac<60>}{10}}\trput{9}}{
     \psset{linestyle=solid}\pstree{\CC{\hyperlink{spotgac<61>}{12}}\trput{6}}{
      \psset{linestyle=solid}\pstree{\CC{\hyperlink{spotgac<62>}{22}}\trput{2}}{
       \psset{linestyle=solid}\pstree{\CC{\hyperlink{spotgac<63>}{24}}\trput{2}}{
        \psset{linestyle=solid}\pstree{\CC{\hyperlink{spotgac<64>}{31}}\trput{1}}{
         \psset{linestyle=solid}\FAIL{spotgac<65>}\trput{4}
         \psset{linestyle=solid}\SUCC{spotgac<66>}\trput{8}
        }
        \psset{linestyle=solid}\FAIL{spotgac<67>}\trput{7}
       }
       \psset{linestyle=solid}\FAIL{spotgac<68>}\trput{4}
      }
      \psset{linestyle=solid}\pstree{\CC{\hyperlink{spotgac<69>}{22}}\trput{8}}{
       \psset{linestyle=solid}\pstree{\CC{\hyperlink{spotgac<70>}{24}}\trput{2}}{
        \psset{linestyle=solid}\FAIL{spotgac<71>}\trput{1}
        \psset{linestyle=solid}\FAIL{spotgac<72>}\trput{7}
       }
       \psset{linestyle=solid}\FAIL{spotgac<73>}\trput{4}
      }
     }
     \psset{linestyle=solid}\FAIL{spotgac<74>}\trput{8}
    }
   }
  }
 }
 \psset{linestyle=solid}\pstree{\CC{\hyperlink{spotgac<75>}{5}}\trput{6}}{
  \psset{linestyle=solid}\pstree{\CC{\hyperlink{spotgac<76>}{6}}\trput{2}}{
   \psset{linestyle=solid}\FAIL{spotgac<77>}\trput{4}
   \psset{linestyle=solid}\pstree{\CC{\hyperlink{spotgac<78>}{7}}\trput{5}}{
   }
  }
 }
}
}
}
\only<79>{
\scalebox{ 0.30}{\psset{nodesep=0pt,levelsep=36pt}
\pstree{\CC{\hyperlink{spotgac<1>}{4}}}{
 \psset{linestyle=solid}\pstree{\CC{\hyperlink{spotgac<2>}{5}}\trput{2}}{
  \psset{linestyle=solid}\pstree{\CC{\hyperlink{spotgac<3>}{6}}\trput{4}}{
   \psset{linestyle=solid}\pstree{\CC{\hyperlink{spotgac<4>}{7}}\trput{5}}{
    \psset{linestyle=solid}\FAIL{spotgac<5>}\trput{6}
    \psset{linestyle=solid}\pstree{\CC{\hyperlink{spotgac<6>}{12}}\trput{9}}{
     \psset{linestyle=solid}\pstree{\CC{\hyperlink{spotgac<7>}{19}}\trput{6}}{
      \psset{linestyle=solid}\FAIL{spotgac<8>}\trput{6}
      \psset{linestyle=solid}\pstree{\CC{\hyperlink{spotgac<9>}{20}}\trput{9}}{
       \psset{linestyle=solid}\pstree{\CC{\hyperlink{spotgac<10>}{23}}\trput{6}}{
        \psset{linestyle=solid}\FAIL{spotgac<11>}\trput{2}
        \psset{linestyle=solid}\FAIL{spotgac<12>}\trput{8}
       }
       \psset{linestyle=solid}\FAIL{spotgac<13>}\trput{8}
      }
     }
     \psset{linestyle=solid}\FAIL{spotgac<14>}\trput{8}
    }
   }
   \psset{linestyle=solid}\FAIL{spotgac<15>}\trput{6}
  }
  \psset{linestyle=solid}\pstree{\CC{\hyperlink{spotgac<16>}{6}}\trput{6}}{
   \psset{linestyle=solid}\pstree{\CC{\hyperlink{spotgac<17>}{10}}\trput{4}}{
    \psset{linestyle=solid}\FAIL{spotgac<18>}\trput{4}
    \psset{linestyle=solid}\FAIL{spotgac<19>}\trput{6}
    \psset{linestyle=solid}\FAIL{spotgac<20>}\trput{8}
   }
   \psset{linestyle=solid}\pstree{\CC{\hyperlink{spotgac<21>}{7}}\trput{5}}{
    \psset{linestyle=solid}\pstree{\CC{\hyperlink{spotgac<22>}{10}}\trput{4}}{
     \psset{linestyle=solid}\FAIL{spotgac<23>}\trput{5}
     \psset{linestyle=solid}\FAIL{spotgac<24>}\trput{6}
     \psset{linestyle=solid}\FAIL{spotgac<25>}\trput{8}
    }
    \psset{linestyle=solid}\pstree{\CC{\hyperlink{spotgac<26>}{10}}\trput{9}}{
     \psset{linestyle=solid}\pstree{\CC{\hyperlink{spotgac<27>}{11}}\trput{5}}{
      \psset{linestyle=solid}\FAIL{spotgac<28>}\trput{6}
      \psset{linestyle=solid}\FAIL{spotgac<29>}\trput{8}
     }
     \psset{linestyle=solid}\FAIL{spotgac<30>}\trput{6}
     \psset{linestyle=solid}\FAIL{spotgac<31>}\trput{8}
    }
   }
  }
 }
 \psset{linestyle=solid}\pstree{\CC{\hyperlink{spotgac<32>}{5}}\trput{4}}{
  \psset{linestyle=solid}\pstree{\CC{\hyperlink{spotgac<33>}{6}}\trput{2}}{
   \psset{linestyle=solid}\FAIL{spotgac<34>}\trput{5}
   \psset{linestyle=solid}\FAIL{spotgac<35>}\trput{6}
  }
  \psset{linestyle=solid}\pstree{\CC{\hyperlink{spotgac<36>}{6}}\trput{6}}{
   \psset{linestyle=solid}\FAIL{spotgac<37>}\trput{2}
   \psset{linestyle=solid}\FAIL{spotgac<38>}\trput{5}
  }
 }
 \psset{linestyle=solid}\pstree{\CC{\hyperlink{spotgac<39>}{5}}\trput{5}}{
  \psset{linestyle=solid}\pstree{\CC{\hyperlink{spotgac<40>}{6}}\trput{2}}{
   \psset{linestyle=solid}\pstree{\CC{\hyperlink{spotgac<41>}{7}}\trput{4}}{
    \psset{linestyle=solid}\FAIL{spotgac<42>}\trput{6}
    \psset{linestyle=solid}\pstree{\CC{\hyperlink{spotgac<43>}{10}}\trput{9}}{
     \psset{linestyle=solid}\FAIL{spotgac<44>}\trput{6}
     \psset{linestyle=solid}\FAIL{spotgac<45>}\trput{8}
    }
   }
   \psset{linestyle=solid}\pstree{\CC{\hyperlink{spotgac<46>}{7}}\trput{6}}{
    \psset{linestyle=solid}\FAIL{spotgac<47>}\trput{4}
    \psset{linestyle=solid}\FAIL{spotgac<48>}\trput{9}
   }
  }
  \psset{linestyle=solid}\FAIL{spotgac<49>}\trput{4}
  \psset{linestyle=solid}\pstree{\CC{\hyperlink{spotgac<50>}{6}}\trput{6}}{
   \psset{linestyle=solid}\pstree{\CC{\hyperlink{spotgac<51>}{7}}\trput{2}}{
    \psset{linestyle=solid}\FAIL{spotgac<52>}\trput{4}
    \psset{linestyle=solid}\FAIL{spotgac<53>}\trput{9}
   }
   \psset{linestyle=solid}\pstree{\CC{\hyperlink{spotgac<54>}{7}}\trput{4}}{
    \psset{linestyle=solid}\pstree{\CC{\hyperlink{spotgac<55>}{10}}\trput{2}}{
     \psset{linestyle=solid}\pstree{\CC{\hyperlink{spotgac<56>}{12}}\trput{6}}{
      \psset{linestyle=solid}\FAIL{spotgac<57>}\trput{2}
      \psset{linestyle=solid}\FAIL{spotgac<58>}\trput{8}
     }
     \psset{linestyle=solid}\FAIL{spotgac<59>}\trput{8}
    }
    \psset{linestyle=solid}\pstree{\CC{\hyperlink{spotgac<60>}{10}}\trput{9}}{
     \psset{linestyle=solid}\pstree{\CC{\hyperlink{spotgac<61>}{12}}\trput{6}}{
      \psset{linestyle=solid}\pstree{\CC{\hyperlink{spotgac<62>}{22}}\trput{2}}{
       \psset{linestyle=solid}\pstree{\CC{\hyperlink{spotgac<63>}{24}}\trput{2}}{
        \psset{linestyle=solid}\pstree{\CC{\hyperlink{spotgac<64>}{31}}\trput{1}}{
         \psset{linestyle=solid}\FAIL{spotgac<65>}\trput{4}
         \psset{linestyle=solid}\SUCC{spotgac<66>}\trput{8}
        }
        \psset{linestyle=solid}\FAIL{spotgac<67>}\trput{7}
       }
       \psset{linestyle=solid}\FAIL{spotgac<68>}\trput{4}
      }
      \psset{linestyle=solid}\pstree{\CC{\hyperlink{spotgac<69>}{22}}\trput{8}}{
       \psset{linestyle=solid}\pstree{\CC{\hyperlink{spotgac<70>}{24}}\trput{2}}{
        \psset{linestyle=solid}\FAIL{spotgac<71>}\trput{1}
        \psset{linestyle=solid}\FAIL{spotgac<72>}\trput{7}
       }
       \psset{linestyle=solid}\FAIL{spotgac<73>}\trput{4}
      }
     }
     \psset{linestyle=solid}\FAIL{spotgac<74>}\trput{8}
    }
   }
  }
 }
 \psset{linestyle=solid}\pstree{\CC{\hyperlink{spotgac<75>}{5}}\trput{6}}{
  \psset{linestyle=solid}\pstree{\CC{\hyperlink{spotgac<76>}{6}}\trput{2}}{
   \psset{linestyle=solid}\FAIL{spotgac<77>}\trput{4}
   \psset{linestyle=solid}\pstree{\CC{\hyperlink{spotgac<78>}{7}}\trput{5}}{
    \psset{linestyle=solid}\FAIL{spotgac<79>}\trput{4}
   }
  }
 }
}
}
}
\only<80>{
\scalebox{ 0.30}{\psset{nodesep=0pt,levelsep=36pt}
\pstree{\CC{\hyperlink{spotgac<1>}{4}}}{
 \psset{linestyle=solid}\pstree{\CC{\hyperlink{spotgac<2>}{5}}\trput{2}}{
  \psset{linestyle=solid}\pstree{\CC{\hyperlink{spotgac<3>}{6}}\trput{4}}{
   \psset{linestyle=solid}\pstree{\CC{\hyperlink{spotgac<4>}{7}}\trput{5}}{
    \psset{linestyle=solid}\FAIL{spotgac<5>}\trput{6}
    \psset{linestyle=solid}\pstree{\CC{\hyperlink{spotgac<6>}{12}}\trput{9}}{
     \psset{linestyle=solid}\pstree{\CC{\hyperlink{spotgac<7>}{19}}\trput{6}}{
      \psset{linestyle=solid}\FAIL{spotgac<8>}\trput{6}
      \psset{linestyle=solid}\pstree{\CC{\hyperlink{spotgac<9>}{20}}\trput{9}}{
       \psset{linestyle=solid}\pstree{\CC{\hyperlink{spotgac<10>}{23}}\trput{6}}{
        \psset{linestyle=solid}\FAIL{spotgac<11>}\trput{2}
        \psset{linestyle=solid}\FAIL{spotgac<12>}\trput{8}
       }
       \psset{linestyle=solid}\FAIL{spotgac<13>}\trput{8}
      }
     }
     \psset{linestyle=solid}\FAIL{spotgac<14>}\trput{8}
    }
   }
   \psset{linestyle=solid}\FAIL{spotgac<15>}\trput{6}
  }
  \psset{linestyle=solid}\pstree{\CC{\hyperlink{spotgac<16>}{6}}\trput{6}}{
   \psset{linestyle=solid}\pstree{\CC{\hyperlink{spotgac<17>}{10}}\trput{4}}{
    \psset{linestyle=solid}\FAIL{spotgac<18>}\trput{4}
    \psset{linestyle=solid}\FAIL{spotgac<19>}\trput{6}
    \psset{linestyle=solid}\FAIL{spotgac<20>}\trput{8}
   }
   \psset{linestyle=solid}\pstree{\CC{\hyperlink{spotgac<21>}{7}}\trput{5}}{
    \psset{linestyle=solid}\pstree{\CC{\hyperlink{spotgac<22>}{10}}\trput{4}}{
     \psset{linestyle=solid}\FAIL{spotgac<23>}\trput{5}
     \psset{linestyle=solid}\FAIL{spotgac<24>}\trput{6}
     \psset{linestyle=solid}\FAIL{spotgac<25>}\trput{8}
    }
    \psset{linestyle=solid}\pstree{\CC{\hyperlink{spotgac<26>}{10}}\trput{9}}{
     \psset{linestyle=solid}\pstree{\CC{\hyperlink{spotgac<27>}{11}}\trput{5}}{
      \psset{linestyle=solid}\FAIL{spotgac<28>}\trput{6}
      \psset{linestyle=solid}\FAIL{spotgac<29>}\trput{8}
     }
     \psset{linestyle=solid}\FAIL{spotgac<30>}\trput{6}
     \psset{linestyle=solid}\FAIL{spotgac<31>}\trput{8}
    }
   }
  }
 }
 \psset{linestyle=solid}\pstree{\CC{\hyperlink{spotgac<32>}{5}}\trput{4}}{
  \psset{linestyle=solid}\pstree{\CC{\hyperlink{spotgac<33>}{6}}\trput{2}}{
   \psset{linestyle=solid}\FAIL{spotgac<34>}\trput{5}
   \psset{linestyle=solid}\FAIL{spotgac<35>}\trput{6}
  }
  \psset{linestyle=solid}\pstree{\CC{\hyperlink{spotgac<36>}{6}}\trput{6}}{
   \psset{linestyle=solid}\FAIL{spotgac<37>}\trput{2}
   \psset{linestyle=solid}\FAIL{spotgac<38>}\trput{5}
  }
 }
 \psset{linestyle=solid}\pstree{\CC{\hyperlink{spotgac<39>}{5}}\trput{5}}{
  \psset{linestyle=solid}\pstree{\CC{\hyperlink{spotgac<40>}{6}}\trput{2}}{
   \psset{linestyle=solid}\pstree{\CC{\hyperlink{spotgac<41>}{7}}\trput{4}}{
    \psset{linestyle=solid}\FAIL{spotgac<42>}\trput{6}
    \psset{linestyle=solid}\pstree{\CC{\hyperlink{spotgac<43>}{10}}\trput{9}}{
     \psset{linestyle=solid}\FAIL{spotgac<44>}\trput{6}
     \psset{linestyle=solid}\FAIL{spotgac<45>}\trput{8}
    }
   }
   \psset{linestyle=solid}\pstree{\CC{\hyperlink{spotgac<46>}{7}}\trput{6}}{
    \psset{linestyle=solid}\FAIL{spotgac<47>}\trput{4}
    \psset{linestyle=solid}\FAIL{spotgac<48>}\trput{9}
   }
  }
  \psset{linestyle=solid}\FAIL{spotgac<49>}\trput{4}
  \psset{linestyle=solid}\pstree{\CC{\hyperlink{spotgac<50>}{6}}\trput{6}}{
   \psset{linestyle=solid}\pstree{\CC{\hyperlink{spotgac<51>}{7}}\trput{2}}{
    \psset{linestyle=solid}\FAIL{spotgac<52>}\trput{4}
    \psset{linestyle=solid}\FAIL{spotgac<53>}\trput{9}
   }
   \psset{linestyle=solid}\pstree{\CC{\hyperlink{spotgac<54>}{7}}\trput{4}}{
    \psset{linestyle=solid}\pstree{\CC{\hyperlink{spotgac<55>}{10}}\trput{2}}{
     \psset{linestyle=solid}\pstree{\CC{\hyperlink{spotgac<56>}{12}}\trput{6}}{
      \psset{linestyle=solid}\FAIL{spotgac<57>}\trput{2}
      \psset{linestyle=solid}\FAIL{spotgac<58>}\trput{8}
     }
     \psset{linestyle=solid}\FAIL{spotgac<59>}\trput{8}
    }
    \psset{linestyle=solid}\pstree{\CC{\hyperlink{spotgac<60>}{10}}\trput{9}}{
     \psset{linestyle=solid}\pstree{\CC{\hyperlink{spotgac<61>}{12}}\trput{6}}{
      \psset{linestyle=solid}\pstree{\CC{\hyperlink{spotgac<62>}{22}}\trput{2}}{
       \psset{linestyle=solid}\pstree{\CC{\hyperlink{spotgac<63>}{24}}\trput{2}}{
        \psset{linestyle=solid}\pstree{\CC{\hyperlink{spotgac<64>}{31}}\trput{1}}{
         \psset{linestyle=solid}\FAIL{spotgac<65>}\trput{4}
         \psset{linestyle=solid}\SUCC{spotgac<66>}\trput{8}
        }
        \psset{linestyle=solid}\FAIL{spotgac<67>}\trput{7}
       }
       \psset{linestyle=solid}\FAIL{spotgac<68>}\trput{4}
      }
      \psset{linestyle=solid}\pstree{\CC{\hyperlink{spotgac<69>}{22}}\trput{8}}{
       \psset{linestyle=solid}\pstree{\CC{\hyperlink{spotgac<70>}{24}}\trput{2}}{
        \psset{linestyle=solid}\FAIL{spotgac<71>}\trput{1}
        \psset{linestyle=solid}\FAIL{spotgac<72>}\trput{7}
       }
       \psset{linestyle=solid}\FAIL{spotgac<73>}\trput{4}
      }
     }
     \psset{linestyle=solid}\FAIL{spotgac<74>}\trput{8}
    }
   }
  }
 }
 \psset{linestyle=solid}\pstree{\CC{\hyperlink{spotgac<75>}{5}}\trput{6}}{
  \psset{linestyle=solid}\pstree{\CC{\hyperlink{spotgac<76>}{6}}\trput{2}}{
   \psset{linestyle=solid}\FAIL{spotgac<77>}\trput{4}
   \psset{linestyle=solid}\pstree{\CC{\hyperlink{spotgac<78>}{7}}\trput{5}}{
    \psset{linestyle=solid}\FAIL{spotgac<79>}\trput{4}
    \psset{linestyle=solid}\FAIL{spotgac<80>}\trput{9}
   }
  }
 }
}
}
}
\only<81>{
\scalebox{ 0.30}{\psset{nodesep=0pt,levelsep=36pt}
\pstree{\CC{\hyperlink{spotgac<1>}{4}}}{
 \psset{linestyle=solid}\pstree{\CC{\hyperlink{spotgac<2>}{5}}\trput{2}}{
  \psset{linestyle=solid}\pstree{\CC{\hyperlink{spotgac<3>}{6}}\trput{4}}{
   \psset{linestyle=solid}\pstree{\CC{\hyperlink{spotgac<4>}{7}}\trput{5}}{
    \psset{linestyle=solid}\FAIL{spotgac<5>}\trput{6}
    \psset{linestyle=solid}\pstree{\CC{\hyperlink{spotgac<6>}{12}}\trput{9}}{
     \psset{linestyle=solid}\pstree{\CC{\hyperlink{spotgac<7>}{19}}\trput{6}}{
      \psset{linestyle=solid}\FAIL{spotgac<8>}\trput{6}
      \psset{linestyle=solid}\pstree{\CC{\hyperlink{spotgac<9>}{20}}\trput{9}}{
       \psset{linestyle=solid}\pstree{\CC{\hyperlink{spotgac<10>}{23}}\trput{6}}{
        \psset{linestyle=solid}\FAIL{spotgac<11>}\trput{2}
        \psset{linestyle=solid}\FAIL{spotgac<12>}\trput{8}
       }
       \psset{linestyle=solid}\FAIL{spotgac<13>}\trput{8}
      }
     }
     \psset{linestyle=solid}\FAIL{spotgac<14>}\trput{8}
    }
   }
   \psset{linestyle=solid}\FAIL{spotgac<15>}\trput{6}
  }
  \psset{linestyle=solid}\pstree{\CC{\hyperlink{spotgac<16>}{6}}\trput{6}}{
   \psset{linestyle=solid}\pstree{\CC{\hyperlink{spotgac<17>}{10}}\trput{4}}{
    \psset{linestyle=solid}\FAIL{spotgac<18>}\trput{4}
    \psset{linestyle=solid}\FAIL{spotgac<19>}\trput{6}
    \psset{linestyle=solid}\FAIL{spotgac<20>}\trput{8}
   }
   \psset{linestyle=solid}\pstree{\CC{\hyperlink{spotgac<21>}{7}}\trput{5}}{
    \psset{linestyle=solid}\pstree{\CC{\hyperlink{spotgac<22>}{10}}\trput{4}}{
     \psset{linestyle=solid}\FAIL{spotgac<23>}\trput{5}
     \psset{linestyle=solid}\FAIL{spotgac<24>}\trput{6}
     \psset{linestyle=solid}\FAIL{spotgac<25>}\trput{8}
    }
    \psset{linestyle=solid}\pstree{\CC{\hyperlink{spotgac<26>}{10}}\trput{9}}{
     \psset{linestyle=solid}\pstree{\CC{\hyperlink{spotgac<27>}{11}}\trput{5}}{
      \psset{linestyle=solid}\FAIL{spotgac<28>}\trput{6}
      \psset{linestyle=solid}\FAIL{spotgac<29>}\trput{8}
     }
     \psset{linestyle=solid}\FAIL{spotgac<30>}\trput{6}
     \psset{linestyle=solid}\FAIL{spotgac<31>}\trput{8}
    }
   }
  }
 }
 \psset{linestyle=solid}\pstree{\CC{\hyperlink{spotgac<32>}{5}}\trput{4}}{
  \psset{linestyle=solid}\pstree{\CC{\hyperlink{spotgac<33>}{6}}\trput{2}}{
   \psset{linestyle=solid}\FAIL{spotgac<34>}\trput{5}
   \psset{linestyle=solid}\FAIL{spotgac<35>}\trput{6}
  }
  \psset{linestyle=solid}\pstree{\CC{\hyperlink{spotgac<36>}{6}}\trput{6}}{
   \psset{linestyle=solid}\FAIL{spotgac<37>}\trput{2}
   \psset{linestyle=solid}\FAIL{spotgac<38>}\trput{5}
  }
 }
 \psset{linestyle=solid}\pstree{\CC{\hyperlink{spotgac<39>}{5}}\trput{5}}{
  \psset{linestyle=solid}\pstree{\CC{\hyperlink{spotgac<40>}{6}}\trput{2}}{
   \psset{linestyle=solid}\pstree{\CC{\hyperlink{spotgac<41>}{7}}\trput{4}}{
    \psset{linestyle=solid}\FAIL{spotgac<42>}\trput{6}
    \psset{linestyle=solid}\pstree{\CC{\hyperlink{spotgac<43>}{10}}\trput{9}}{
     \psset{linestyle=solid}\FAIL{spotgac<44>}\trput{6}
     \psset{linestyle=solid}\FAIL{spotgac<45>}\trput{8}
    }
   }
   \psset{linestyle=solid}\pstree{\CC{\hyperlink{spotgac<46>}{7}}\trput{6}}{
    \psset{linestyle=solid}\FAIL{spotgac<47>}\trput{4}
    \psset{linestyle=solid}\FAIL{spotgac<48>}\trput{9}
   }
  }
  \psset{linestyle=solid}\FAIL{spotgac<49>}\trput{4}
  \psset{linestyle=solid}\pstree{\CC{\hyperlink{spotgac<50>}{6}}\trput{6}}{
   \psset{linestyle=solid}\pstree{\CC{\hyperlink{spotgac<51>}{7}}\trput{2}}{
    \psset{linestyle=solid}\FAIL{spotgac<52>}\trput{4}
    \psset{linestyle=solid}\FAIL{spotgac<53>}\trput{9}
   }
   \psset{linestyle=solid}\pstree{\CC{\hyperlink{spotgac<54>}{7}}\trput{4}}{
    \psset{linestyle=solid}\pstree{\CC{\hyperlink{spotgac<55>}{10}}\trput{2}}{
     \psset{linestyle=solid}\pstree{\CC{\hyperlink{spotgac<56>}{12}}\trput{6}}{
      \psset{linestyle=solid}\FAIL{spotgac<57>}\trput{2}
      \psset{linestyle=solid}\FAIL{spotgac<58>}\trput{8}
     }
     \psset{linestyle=solid}\FAIL{spotgac<59>}\trput{8}
    }
    \psset{linestyle=solid}\pstree{\CC{\hyperlink{spotgac<60>}{10}}\trput{9}}{
     \psset{linestyle=solid}\pstree{\CC{\hyperlink{spotgac<61>}{12}}\trput{6}}{
      \psset{linestyle=solid}\pstree{\CC{\hyperlink{spotgac<62>}{22}}\trput{2}}{
       \psset{linestyle=solid}\pstree{\CC{\hyperlink{spotgac<63>}{24}}\trput{2}}{
        \psset{linestyle=solid}\pstree{\CC{\hyperlink{spotgac<64>}{31}}\trput{1}}{
         \psset{linestyle=solid}\FAIL{spotgac<65>}\trput{4}
         \psset{linestyle=solid}\SUCC{spotgac<66>}\trput{8}
        }
        \psset{linestyle=solid}\FAIL{spotgac<67>}\trput{7}
       }
       \psset{linestyle=solid}\FAIL{spotgac<68>}\trput{4}
      }
      \psset{linestyle=solid}\pstree{\CC{\hyperlink{spotgac<69>}{22}}\trput{8}}{
       \psset{linestyle=solid}\pstree{\CC{\hyperlink{spotgac<70>}{24}}\trput{2}}{
        \psset{linestyle=solid}\FAIL{spotgac<71>}\trput{1}
        \psset{linestyle=solid}\FAIL{spotgac<72>}\trput{7}
       }
       \psset{linestyle=solid}\FAIL{spotgac<73>}\trput{4}
      }
     }
     \psset{linestyle=solid}\FAIL{spotgac<74>}\trput{8}
    }
   }
  }
 }
 \psset{linestyle=solid}\pstree{\CC{\hyperlink{spotgac<75>}{5}}\trput{6}}{
  \psset{linestyle=solid}\pstree{\CC{\hyperlink{spotgac<76>}{6}}\trput{2}}{
   \psset{linestyle=solid}\FAIL{spotgac<77>}\trput{4}
   \psset{linestyle=solid}\pstree{\CC{\hyperlink{spotgac<78>}{7}}\trput{5}}{
    \psset{linestyle=solid}\FAIL{spotgac<79>}\trput{4}
    \psset{linestyle=solid}\FAIL{spotgac<80>}\trput{9}
   }
  }
  \psset{linestyle=solid}\pstree{\CC{\hyperlink{spotgac<81>}{6}}\trput{4}}{
  }
 }
}
}
}
\only<82>{
\scalebox{ 0.30}{\psset{nodesep=0pt,levelsep=36pt}
\pstree{\CC{\hyperlink{spotgac<1>}{4}}}{
 \psset{linestyle=solid}\pstree{\CC{\hyperlink{spotgac<2>}{5}}\trput{2}}{
  \psset{linestyle=solid}\pstree{\CC{\hyperlink{spotgac<3>}{6}}\trput{4}}{
   \psset{linestyle=solid}\pstree{\CC{\hyperlink{spotgac<4>}{7}}\trput{5}}{
    \psset{linestyle=solid}\FAIL{spotgac<5>}\trput{6}
    \psset{linestyle=solid}\pstree{\CC{\hyperlink{spotgac<6>}{12}}\trput{9}}{
     \psset{linestyle=solid}\pstree{\CC{\hyperlink{spotgac<7>}{19}}\trput{6}}{
      \psset{linestyle=solid}\FAIL{spotgac<8>}\trput{6}
      \psset{linestyle=solid}\pstree{\CC{\hyperlink{spotgac<9>}{20}}\trput{9}}{
       \psset{linestyle=solid}\pstree{\CC{\hyperlink{spotgac<10>}{23}}\trput{6}}{
        \psset{linestyle=solid}\FAIL{spotgac<11>}\trput{2}
        \psset{linestyle=solid}\FAIL{spotgac<12>}\trput{8}
       }
       \psset{linestyle=solid}\FAIL{spotgac<13>}\trput{8}
      }
     }
     \psset{linestyle=solid}\FAIL{spotgac<14>}\trput{8}
    }
   }
   \psset{linestyle=solid}\FAIL{spotgac<15>}\trput{6}
  }
  \psset{linestyle=solid}\pstree{\CC{\hyperlink{spotgac<16>}{6}}\trput{6}}{
   \psset{linestyle=solid}\pstree{\CC{\hyperlink{spotgac<17>}{10}}\trput{4}}{
    \psset{linestyle=solid}\FAIL{spotgac<18>}\trput{4}
    \psset{linestyle=solid}\FAIL{spotgac<19>}\trput{6}
    \psset{linestyle=solid}\FAIL{spotgac<20>}\trput{8}
   }
   \psset{linestyle=solid}\pstree{\CC{\hyperlink{spotgac<21>}{7}}\trput{5}}{
    \psset{linestyle=solid}\pstree{\CC{\hyperlink{spotgac<22>}{10}}\trput{4}}{
     \psset{linestyle=solid}\FAIL{spotgac<23>}\trput{5}
     \psset{linestyle=solid}\FAIL{spotgac<24>}\trput{6}
     \psset{linestyle=solid}\FAIL{spotgac<25>}\trput{8}
    }
    \psset{linestyle=solid}\pstree{\CC{\hyperlink{spotgac<26>}{10}}\trput{9}}{
     \psset{linestyle=solid}\pstree{\CC{\hyperlink{spotgac<27>}{11}}\trput{5}}{
      \psset{linestyle=solid}\FAIL{spotgac<28>}\trput{6}
      \psset{linestyle=solid}\FAIL{spotgac<29>}\trput{8}
     }
     \psset{linestyle=solid}\FAIL{spotgac<30>}\trput{6}
     \psset{linestyle=solid}\FAIL{spotgac<31>}\trput{8}
    }
   }
  }
 }
 \psset{linestyle=solid}\pstree{\CC{\hyperlink{spotgac<32>}{5}}\trput{4}}{
  \psset{linestyle=solid}\pstree{\CC{\hyperlink{spotgac<33>}{6}}\trput{2}}{
   \psset{linestyle=solid}\FAIL{spotgac<34>}\trput{5}
   \psset{linestyle=solid}\FAIL{spotgac<35>}\trput{6}
  }
  \psset{linestyle=solid}\pstree{\CC{\hyperlink{spotgac<36>}{6}}\trput{6}}{
   \psset{linestyle=solid}\FAIL{spotgac<37>}\trput{2}
   \psset{linestyle=solid}\FAIL{spotgac<38>}\trput{5}
  }
 }
 \psset{linestyle=solid}\pstree{\CC{\hyperlink{spotgac<39>}{5}}\trput{5}}{
  \psset{linestyle=solid}\pstree{\CC{\hyperlink{spotgac<40>}{6}}\trput{2}}{
   \psset{linestyle=solid}\pstree{\CC{\hyperlink{spotgac<41>}{7}}\trput{4}}{
    \psset{linestyle=solid}\FAIL{spotgac<42>}\trput{6}
    \psset{linestyle=solid}\pstree{\CC{\hyperlink{spotgac<43>}{10}}\trput{9}}{
     \psset{linestyle=solid}\FAIL{spotgac<44>}\trput{6}
     \psset{linestyle=solid}\FAIL{spotgac<45>}\trput{8}
    }
   }
   \psset{linestyle=solid}\pstree{\CC{\hyperlink{spotgac<46>}{7}}\trput{6}}{
    \psset{linestyle=solid}\FAIL{spotgac<47>}\trput{4}
    \psset{linestyle=solid}\FAIL{spotgac<48>}\trput{9}
   }
  }
  \psset{linestyle=solid}\FAIL{spotgac<49>}\trput{4}
  \psset{linestyle=solid}\pstree{\CC{\hyperlink{spotgac<50>}{6}}\trput{6}}{
   \psset{linestyle=solid}\pstree{\CC{\hyperlink{spotgac<51>}{7}}\trput{2}}{
    \psset{linestyle=solid}\FAIL{spotgac<52>}\trput{4}
    \psset{linestyle=solid}\FAIL{spotgac<53>}\trput{9}
   }
   \psset{linestyle=solid}\pstree{\CC{\hyperlink{spotgac<54>}{7}}\trput{4}}{
    \psset{linestyle=solid}\pstree{\CC{\hyperlink{spotgac<55>}{10}}\trput{2}}{
     \psset{linestyle=solid}\pstree{\CC{\hyperlink{spotgac<56>}{12}}\trput{6}}{
      \psset{linestyle=solid}\FAIL{spotgac<57>}\trput{2}
      \psset{linestyle=solid}\FAIL{spotgac<58>}\trput{8}
     }
     \psset{linestyle=solid}\FAIL{spotgac<59>}\trput{8}
    }
    \psset{linestyle=solid}\pstree{\CC{\hyperlink{spotgac<60>}{10}}\trput{9}}{
     \psset{linestyle=solid}\pstree{\CC{\hyperlink{spotgac<61>}{12}}\trput{6}}{
      \psset{linestyle=solid}\pstree{\CC{\hyperlink{spotgac<62>}{22}}\trput{2}}{
       \psset{linestyle=solid}\pstree{\CC{\hyperlink{spotgac<63>}{24}}\trput{2}}{
        \psset{linestyle=solid}\pstree{\CC{\hyperlink{spotgac<64>}{31}}\trput{1}}{
         \psset{linestyle=solid}\FAIL{spotgac<65>}\trput{4}
         \psset{linestyle=solid}\SUCC{spotgac<66>}\trput{8}
        }
        \psset{linestyle=solid}\FAIL{spotgac<67>}\trput{7}
       }
       \psset{linestyle=solid}\FAIL{spotgac<68>}\trput{4}
      }
      \psset{linestyle=solid}\pstree{\CC{\hyperlink{spotgac<69>}{22}}\trput{8}}{
       \psset{linestyle=solid}\pstree{\CC{\hyperlink{spotgac<70>}{24}}\trput{2}}{
        \psset{linestyle=solid}\FAIL{spotgac<71>}\trput{1}
        \psset{linestyle=solid}\FAIL{spotgac<72>}\trput{7}
       }
       \psset{linestyle=solid}\FAIL{spotgac<73>}\trput{4}
      }
     }
     \psset{linestyle=solid}\FAIL{spotgac<74>}\trput{8}
    }
   }
  }
 }
 \psset{linestyle=solid}\pstree{\CC{\hyperlink{spotgac<75>}{5}}\trput{6}}{
  \psset{linestyle=solid}\pstree{\CC{\hyperlink{spotgac<76>}{6}}\trput{2}}{
   \psset{linestyle=solid}\FAIL{spotgac<77>}\trput{4}
   \psset{linestyle=solid}\pstree{\CC{\hyperlink{spotgac<78>}{7}}\trput{5}}{
    \psset{linestyle=solid}\FAIL{spotgac<79>}\trput{4}
    \psset{linestyle=solid}\FAIL{spotgac<80>}\trput{9}
   }
  }
  \psset{linestyle=solid}\pstree{\CC{\hyperlink{spotgac<81>}{6}}\trput{4}}{
   \psset{linestyle=solid}\FAIL{spotgac<82>}\trput{2}
  }
 }
}
}
}
\only<83>{
\scalebox{ 0.30}{\psset{nodesep=0pt,levelsep=36pt}
\pstree{\CC{\hyperlink{spotgac<1>}{4}}}{
 \psset{linestyle=solid}\pstree{\CC{\hyperlink{spotgac<2>}{5}}\trput{2}}{
  \psset{linestyle=solid}\pstree{\CC{\hyperlink{spotgac<3>}{6}}\trput{4}}{
   \psset{linestyle=solid}\pstree{\CC{\hyperlink{spotgac<4>}{7}}\trput{5}}{
    \psset{linestyle=solid}\FAIL{spotgac<5>}\trput{6}
    \psset{linestyle=solid}\pstree{\CC{\hyperlink{spotgac<6>}{12}}\trput{9}}{
     \psset{linestyle=solid}\pstree{\CC{\hyperlink{spotgac<7>}{19}}\trput{6}}{
      \psset{linestyle=solid}\FAIL{spotgac<8>}\trput{6}
      \psset{linestyle=solid}\pstree{\CC{\hyperlink{spotgac<9>}{20}}\trput{9}}{
       \psset{linestyle=solid}\pstree{\CC{\hyperlink{spotgac<10>}{23}}\trput{6}}{
        \psset{linestyle=solid}\FAIL{spotgac<11>}\trput{2}
        \psset{linestyle=solid}\FAIL{spotgac<12>}\trput{8}
       }
       \psset{linestyle=solid}\FAIL{spotgac<13>}\trput{8}
      }
     }
     \psset{linestyle=solid}\FAIL{spotgac<14>}\trput{8}
    }
   }
   \psset{linestyle=solid}\FAIL{spotgac<15>}\trput{6}
  }
  \psset{linestyle=solid}\pstree{\CC{\hyperlink{spotgac<16>}{6}}\trput{6}}{
   \psset{linestyle=solid}\pstree{\CC{\hyperlink{spotgac<17>}{10}}\trput{4}}{
    \psset{linestyle=solid}\FAIL{spotgac<18>}\trput{4}
    \psset{linestyle=solid}\FAIL{spotgac<19>}\trput{6}
    \psset{linestyle=solid}\FAIL{spotgac<20>}\trput{8}
   }
   \psset{linestyle=solid}\pstree{\CC{\hyperlink{spotgac<21>}{7}}\trput{5}}{
    \psset{linestyle=solid}\pstree{\CC{\hyperlink{spotgac<22>}{10}}\trput{4}}{
     \psset{linestyle=solid}\FAIL{spotgac<23>}\trput{5}
     \psset{linestyle=solid}\FAIL{spotgac<24>}\trput{6}
     \psset{linestyle=solid}\FAIL{spotgac<25>}\trput{8}
    }
    \psset{linestyle=solid}\pstree{\CC{\hyperlink{spotgac<26>}{10}}\trput{9}}{
     \psset{linestyle=solid}\pstree{\CC{\hyperlink{spotgac<27>}{11}}\trput{5}}{
      \psset{linestyle=solid}\FAIL{spotgac<28>}\trput{6}
      \psset{linestyle=solid}\FAIL{spotgac<29>}\trput{8}
     }
     \psset{linestyle=solid}\FAIL{spotgac<30>}\trput{6}
     \psset{linestyle=solid}\FAIL{spotgac<31>}\trput{8}
    }
   }
  }
 }
 \psset{linestyle=solid}\pstree{\CC{\hyperlink{spotgac<32>}{5}}\trput{4}}{
  \psset{linestyle=solid}\pstree{\CC{\hyperlink{spotgac<33>}{6}}\trput{2}}{
   \psset{linestyle=solid}\FAIL{spotgac<34>}\trput{5}
   \psset{linestyle=solid}\FAIL{spotgac<35>}\trput{6}
  }
  \psset{linestyle=solid}\pstree{\CC{\hyperlink{spotgac<36>}{6}}\trput{6}}{
   \psset{linestyle=solid}\FAIL{spotgac<37>}\trput{2}
   \psset{linestyle=solid}\FAIL{spotgac<38>}\trput{5}
  }
 }
 \psset{linestyle=solid}\pstree{\CC{\hyperlink{spotgac<39>}{5}}\trput{5}}{
  \psset{linestyle=solid}\pstree{\CC{\hyperlink{spotgac<40>}{6}}\trput{2}}{
   \psset{linestyle=solid}\pstree{\CC{\hyperlink{spotgac<41>}{7}}\trput{4}}{
    \psset{linestyle=solid}\FAIL{spotgac<42>}\trput{6}
    \psset{linestyle=solid}\pstree{\CC{\hyperlink{spotgac<43>}{10}}\trput{9}}{
     \psset{linestyle=solid}\FAIL{spotgac<44>}\trput{6}
     \psset{linestyle=solid}\FAIL{spotgac<45>}\trput{8}
    }
   }
   \psset{linestyle=solid}\pstree{\CC{\hyperlink{spotgac<46>}{7}}\trput{6}}{
    \psset{linestyle=solid}\FAIL{spotgac<47>}\trput{4}
    \psset{linestyle=solid}\FAIL{spotgac<48>}\trput{9}
   }
  }
  \psset{linestyle=solid}\FAIL{spotgac<49>}\trput{4}
  \psset{linestyle=solid}\pstree{\CC{\hyperlink{spotgac<50>}{6}}\trput{6}}{
   \psset{linestyle=solid}\pstree{\CC{\hyperlink{spotgac<51>}{7}}\trput{2}}{
    \psset{linestyle=solid}\FAIL{spotgac<52>}\trput{4}
    \psset{linestyle=solid}\FAIL{spotgac<53>}\trput{9}
   }
   \psset{linestyle=solid}\pstree{\CC{\hyperlink{spotgac<54>}{7}}\trput{4}}{
    \psset{linestyle=solid}\pstree{\CC{\hyperlink{spotgac<55>}{10}}\trput{2}}{
     \psset{linestyle=solid}\pstree{\CC{\hyperlink{spotgac<56>}{12}}\trput{6}}{
      \psset{linestyle=solid}\FAIL{spotgac<57>}\trput{2}
      \psset{linestyle=solid}\FAIL{spotgac<58>}\trput{8}
     }
     \psset{linestyle=solid}\FAIL{spotgac<59>}\trput{8}
    }
    \psset{linestyle=solid}\pstree{\CC{\hyperlink{spotgac<60>}{10}}\trput{9}}{
     \psset{linestyle=solid}\pstree{\CC{\hyperlink{spotgac<61>}{12}}\trput{6}}{
      \psset{linestyle=solid}\pstree{\CC{\hyperlink{spotgac<62>}{22}}\trput{2}}{
       \psset{linestyle=solid}\pstree{\CC{\hyperlink{spotgac<63>}{24}}\trput{2}}{
        \psset{linestyle=solid}\pstree{\CC{\hyperlink{spotgac<64>}{31}}\trput{1}}{
         \psset{linestyle=solid}\FAIL{spotgac<65>}\trput{4}
         \psset{linestyle=solid}\SUCC{spotgac<66>}\trput{8}
        }
        \psset{linestyle=solid}\FAIL{spotgac<67>}\trput{7}
       }
       \psset{linestyle=solid}\FAIL{spotgac<68>}\trput{4}
      }
      \psset{linestyle=solid}\pstree{\CC{\hyperlink{spotgac<69>}{22}}\trput{8}}{
       \psset{linestyle=solid}\pstree{\CC{\hyperlink{spotgac<70>}{24}}\trput{2}}{
        \psset{linestyle=solid}\FAIL{spotgac<71>}\trput{1}
        \psset{linestyle=solid}\FAIL{spotgac<72>}\trput{7}
       }
       \psset{linestyle=solid}\FAIL{spotgac<73>}\trput{4}
      }
     }
     \psset{linestyle=solid}\FAIL{spotgac<74>}\trput{8}
    }
   }
  }
 }
 \psset{linestyle=solid}\pstree{\CC{\hyperlink{spotgac<75>}{5}}\trput{6}}{
  \psset{linestyle=solid}\pstree{\CC{\hyperlink{spotgac<76>}{6}}\trput{2}}{
   \psset{linestyle=solid}\FAIL{spotgac<77>}\trput{4}
   \psset{linestyle=solid}\pstree{\CC{\hyperlink{spotgac<78>}{7}}\trput{5}}{
    \psset{linestyle=solid}\FAIL{spotgac<79>}\trput{4}
    \psset{linestyle=solid}\FAIL{spotgac<80>}\trput{9}
   }
  }
  \psset{linestyle=solid}\pstree{\CC{\hyperlink{spotgac<81>}{6}}\trput{4}}{
   \psset{linestyle=solid}\FAIL{spotgac<82>}\trput{2}
   \psset{linestyle=solid}\pstree{\CC{\hyperlink{spotgac<83>}{7}}\trput{5}}{
   }
  }
 }
}
}
}
\only<84>{
\scalebox{ 0.30}{\psset{nodesep=0pt,levelsep=36pt}
\pstree{\CC{\hyperlink{spotgac<1>}{4}}}{
 \psset{linestyle=solid}\pstree{\CC{\hyperlink{spotgac<2>}{5}}\trput{2}}{
  \psset{linestyle=solid}\pstree{\CC{\hyperlink{spotgac<3>}{6}}\trput{4}}{
   \psset{linestyle=solid}\pstree{\CC{\hyperlink{spotgac<4>}{7}}\trput{5}}{
    \psset{linestyle=solid}\FAIL{spotgac<5>}\trput{6}
    \psset{linestyle=solid}\pstree{\CC{\hyperlink{spotgac<6>}{12}}\trput{9}}{
     \psset{linestyle=solid}\pstree{\CC{\hyperlink{spotgac<7>}{19}}\trput{6}}{
      \psset{linestyle=solid}\FAIL{spotgac<8>}\trput{6}
      \psset{linestyle=solid}\pstree{\CC{\hyperlink{spotgac<9>}{20}}\trput{9}}{
       \psset{linestyle=solid}\pstree{\CC{\hyperlink{spotgac<10>}{23}}\trput{6}}{
        \psset{linestyle=solid}\FAIL{spotgac<11>}\trput{2}
        \psset{linestyle=solid}\FAIL{spotgac<12>}\trput{8}
       }
       \psset{linestyle=solid}\FAIL{spotgac<13>}\trput{8}
      }
     }
     \psset{linestyle=solid}\FAIL{spotgac<14>}\trput{8}
    }
   }
   \psset{linestyle=solid}\FAIL{spotgac<15>}\trput{6}
  }
  \psset{linestyle=solid}\pstree{\CC{\hyperlink{spotgac<16>}{6}}\trput{6}}{
   \psset{linestyle=solid}\pstree{\CC{\hyperlink{spotgac<17>}{10}}\trput{4}}{
    \psset{linestyle=solid}\FAIL{spotgac<18>}\trput{4}
    \psset{linestyle=solid}\FAIL{spotgac<19>}\trput{6}
    \psset{linestyle=solid}\FAIL{spotgac<20>}\trput{8}
   }
   \psset{linestyle=solid}\pstree{\CC{\hyperlink{spotgac<21>}{7}}\trput{5}}{
    \psset{linestyle=solid}\pstree{\CC{\hyperlink{spotgac<22>}{10}}\trput{4}}{
     \psset{linestyle=solid}\FAIL{spotgac<23>}\trput{5}
     \psset{linestyle=solid}\FAIL{spotgac<24>}\trput{6}
     \psset{linestyle=solid}\FAIL{spotgac<25>}\trput{8}
    }
    \psset{linestyle=solid}\pstree{\CC{\hyperlink{spotgac<26>}{10}}\trput{9}}{
     \psset{linestyle=solid}\pstree{\CC{\hyperlink{spotgac<27>}{11}}\trput{5}}{
      \psset{linestyle=solid}\FAIL{spotgac<28>}\trput{6}
      \psset{linestyle=solid}\FAIL{spotgac<29>}\trput{8}
     }
     \psset{linestyle=solid}\FAIL{spotgac<30>}\trput{6}
     \psset{linestyle=solid}\FAIL{spotgac<31>}\trput{8}
    }
   }
  }
 }
 \psset{linestyle=solid}\pstree{\CC{\hyperlink{spotgac<32>}{5}}\trput{4}}{
  \psset{linestyle=solid}\pstree{\CC{\hyperlink{spotgac<33>}{6}}\trput{2}}{
   \psset{linestyle=solid}\FAIL{spotgac<34>}\trput{5}
   \psset{linestyle=solid}\FAIL{spotgac<35>}\trput{6}
  }
  \psset{linestyle=solid}\pstree{\CC{\hyperlink{spotgac<36>}{6}}\trput{6}}{
   \psset{linestyle=solid}\FAIL{spotgac<37>}\trput{2}
   \psset{linestyle=solid}\FAIL{spotgac<38>}\trput{5}
  }
 }
 \psset{linestyle=solid}\pstree{\CC{\hyperlink{spotgac<39>}{5}}\trput{5}}{
  \psset{linestyle=solid}\pstree{\CC{\hyperlink{spotgac<40>}{6}}\trput{2}}{
   \psset{linestyle=solid}\pstree{\CC{\hyperlink{spotgac<41>}{7}}\trput{4}}{
    \psset{linestyle=solid}\FAIL{spotgac<42>}\trput{6}
    \psset{linestyle=solid}\pstree{\CC{\hyperlink{spotgac<43>}{10}}\trput{9}}{
     \psset{linestyle=solid}\FAIL{spotgac<44>}\trput{6}
     \psset{linestyle=solid}\FAIL{spotgac<45>}\trput{8}
    }
   }
   \psset{linestyle=solid}\pstree{\CC{\hyperlink{spotgac<46>}{7}}\trput{6}}{
    \psset{linestyle=solid}\FAIL{spotgac<47>}\trput{4}
    \psset{linestyle=solid}\FAIL{spotgac<48>}\trput{9}
   }
  }
  \psset{linestyle=solid}\FAIL{spotgac<49>}\trput{4}
  \psset{linestyle=solid}\pstree{\CC{\hyperlink{spotgac<50>}{6}}\trput{6}}{
   \psset{linestyle=solid}\pstree{\CC{\hyperlink{spotgac<51>}{7}}\trput{2}}{
    \psset{linestyle=solid}\FAIL{spotgac<52>}\trput{4}
    \psset{linestyle=solid}\FAIL{spotgac<53>}\trput{9}
   }
   \psset{linestyle=solid}\pstree{\CC{\hyperlink{spotgac<54>}{7}}\trput{4}}{
    \psset{linestyle=solid}\pstree{\CC{\hyperlink{spotgac<55>}{10}}\trput{2}}{
     \psset{linestyle=solid}\pstree{\CC{\hyperlink{spotgac<56>}{12}}\trput{6}}{
      \psset{linestyle=solid}\FAIL{spotgac<57>}\trput{2}
      \psset{linestyle=solid}\FAIL{spotgac<58>}\trput{8}
     }
     \psset{linestyle=solid}\FAIL{spotgac<59>}\trput{8}
    }
    \psset{linestyle=solid}\pstree{\CC{\hyperlink{spotgac<60>}{10}}\trput{9}}{
     \psset{linestyle=solid}\pstree{\CC{\hyperlink{spotgac<61>}{12}}\trput{6}}{
      \psset{linestyle=solid}\pstree{\CC{\hyperlink{spotgac<62>}{22}}\trput{2}}{
       \psset{linestyle=solid}\pstree{\CC{\hyperlink{spotgac<63>}{24}}\trput{2}}{
        \psset{linestyle=solid}\pstree{\CC{\hyperlink{spotgac<64>}{31}}\trput{1}}{
         \psset{linestyle=solid}\FAIL{spotgac<65>}\trput{4}
         \psset{linestyle=solid}\SUCC{spotgac<66>}\trput{8}
        }
        \psset{linestyle=solid}\FAIL{spotgac<67>}\trput{7}
       }
       \psset{linestyle=solid}\FAIL{spotgac<68>}\trput{4}
      }
      \psset{linestyle=solid}\pstree{\CC{\hyperlink{spotgac<69>}{22}}\trput{8}}{
       \psset{linestyle=solid}\pstree{\CC{\hyperlink{spotgac<70>}{24}}\trput{2}}{
        \psset{linestyle=solid}\FAIL{spotgac<71>}\trput{1}
        \psset{linestyle=solid}\FAIL{spotgac<72>}\trput{7}
       }
       \psset{linestyle=solid}\FAIL{spotgac<73>}\trput{4}
      }
     }
     \psset{linestyle=solid}\FAIL{spotgac<74>}\trput{8}
    }
   }
  }
 }
 \psset{linestyle=solid}\pstree{\CC{\hyperlink{spotgac<75>}{5}}\trput{6}}{
  \psset{linestyle=solid}\pstree{\CC{\hyperlink{spotgac<76>}{6}}\trput{2}}{
   \psset{linestyle=solid}\FAIL{spotgac<77>}\trput{4}
   \psset{linestyle=solid}\pstree{\CC{\hyperlink{spotgac<78>}{7}}\trput{5}}{
    \psset{linestyle=solid}\FAIL{spotgac<79>}\trput{4}
    \psset{linestyle=solid}\FAIL{spotgac<80>}\trput{9}
   }
  }
  \psset{linestyle=solid}\pstree{\CC{\hyperlink{spotgac<81>}{6}}\trput{4}}{
   \psset{linestyle=solid}\FAIL{spotgac<82>}\trput{2}
   \psset{linestyle=solid}\pstree{\CC{\hyperlink{spotgac<83>}{7}}\trput{5}}{
    \psset{linestyle=solid}\FAIL{spotgac<84>}\trput{2}
   }
  }
 }
}
}
}
\only<85>{
\scalebox{ 0.30}{\psset{nodesep=0pt,levelsep=36pt}
\pstree{\CC{\hyperlink{spotgac<1>}{4}}}{
 \psset{linestyle=solid}\pstree{\CC{\hyperlink{spotgac<2>}{5}}\trput{2}}{
  \psset{linestyle=solid}\pstree{\CC{\hyperlink{spotgac<3>}{6}}\trput{4}}{
   \psset{linestyle=solid}\pstree{\CC{\hyperlink{spotgac<4>}{7}}\trput{5}}{
    \psset{linestyle=solid}\FAIL{spotgac<5>}\trput{6}
    \psset{linestyle=solid}\pstree{\CC{\hyperlink{spotgac<6>}{12}}\trput{9}}{
     \psset{linestyle=solid}\pstree{\CC{\hyperlink{spotgac<7>}{19}}\trput{6}}{
      \psset{linestyle=solid}\FAIL{spotgac<8>}\trput{6}
      \psset{linestyle=solid}\pstree{\CC{\hyperlink{spotgac<9>}{20}}\trput{9}}{
       \psset{linestyle=solid}\pstree{\CC{\hyperlink{spotgac<10>}{23}}\trput{6}}{
        \psset{linestyle=solid}\FAIL{spotgac<11>}\trput{2}
        \psset{linestyle=solid}\FAIL{spotgac<12>}\trput{8}
       }
       \psset{linestyle=solid}\FAIL{spotgac<13>}\trput{8}
      }
     }
     \psset{linestyle=solid}\FAIL{spotgac<14>}\trput{8}
    }
   }
   \psset{linestyle=solid}\FAIL{spotgac<15>}\trput{6}
  }
  \psset{linestyle=solid}\pstree{\CC{\hyperlink{spotgac<16>}{6}}\trput{6}}{
   \psset{linestyle=solid}\pstree{\CC{\hyperlink{spotgac<17>}{10}}\trput{4}}{
    \psset{linestyle=solid}\FAIL{spotgac<18>}\trput{4}
    \psset{linestyle=solid}\FAIL{spotgac<19>}\trput{6}
    \psset{linestyle=solid}\FAIL{spotgac<20>}\trput{8}
   }
   \psset{linestyle=solid}\pstree{\CC{\hyperlink{spotgac<21>}{7}}\trput{5}}{
    \psset{linestyle=solid}\pstree{\CC{\hyperlink{spotgac<22>}{10}}\trput{4}}{
     \psset{linestyle=solid}\FAIL{spotgac<23>}\trput{5}
     \psset{linestyle=solid}\FAIL{spotgac<24>}\trput{6}
     \psset{linestyle=solid}\FAIL{spotgac<25>}\trput{8}
    }
    \psset{linestyle=solid}\pstree{\CC{\hyperlink{spotgac<26>}{10}}\trput{9}}{
     \psset{linestyle=solid}\pstree{\CC{\hyperlink{spotgac<27>}{11}}\trput{5}}{
      \psset{linestyle=solid}\FAIL{spotgac<28>}\trput{6}
      \psset{linestyle=solid}\FAIL{spotgac<29>}\trput{8}
     }
     \psset{linestyle=solid}\FAIL{spotgac<30>}\trput{6}
     \psset{linestyle=solid}\FAIL{spotgac<31>}\trput{8}
    }
   }
  }
 }
 \psset{linestyle=solid}\pstree{\CC{\hyperlink{spotgac<32>}{5}}\trput{4}}{
  \psset{linestyle=solid}\pstree{\CC{\hyperlink{spotgac<33>}{6}}\trput{2}}{
   \psset{linestyle=solid}\FAIL{spotgac<34>}\trput{5}
   \psset{linestyle=solid}\FAIL{spotgac<35>}\trput{6}
  }
  \psset{linestyle=solid}\pstree{\CC{\hyperlink{spotgac<36>}{6}}\trput{6}}{
   \psset{linestyle=solid}\FAIL{spotgac<37>}\trput{2}
   \psset{linestyle=solid}\FAIL{spotgac<38>}\trput{5}
  }
 }
 \psset{linestyle=solid}\pstree{\CC{\hyperlink{spotgac<39>}{5}}\trput{5}}{
  \psset{linestyle=solid}\pstree{\CC{\hyperlink{spotgac<40>}{6}}\trput{2}}{
   \psset{linestyle=solid}\pstree{\CC{\hyperlink{spotgac<41>}{7}}\trput{4}}{
    \psset{linestyle=solid}\FAIL{spotgac<42>}\trput{6}
    \psset{linestyle=solid}\pstree{\CC{\hyperlink{spotgac<43>}{10}}\trput{9}}{
     \psset{linestyle=solid}\FAIL{spotgac<44>}\trput{6}
     \psset{linestyle=solid}\FAIL{spotgac<45>}\trput{8}
    }
   }
   \psset{linestyle=solid}\pstree{\CC{\hyperlink{spotgac<46>}{7}}\trput{6}}{
    \psset{linestyle=solid}\FAIL{spotgac<47>}\trput{4}
    \psset{linestyle=solid}\FAIL{spotgac<48>}\trput{9}
   }
  }
  \psset{linestyle=solid}\FAIL{spotgac<49>}\trput{4}
  \psset{linestyle=solid}\pstree{\CC{\hyperlink{spotgac<50>}{6}}\trput{6}}{
   \psset{linestyle=solid}\pstree{\CC{\hyperlink{spotgac<51>}{7}}\trput{2}}{
    \psset{linestyle=solid}\FAIL{spotgac<52>}\trput{4}
    \psset{linestyle=solid}\FAIL{spotgac<53>}\trput{9}
   }
   \psset{linestyle=solid}\pstree{\CC{\hyperlink{spotgac<54>}{7}}\trput{4}}{
    \psset{linestyle=solid}\pstree{\CC{\hyperlink{spotgac<55>}{10}}\trput{2}}{
     \psset{linestyle=solid}\pstree{\CC{\hyperlink{spotgac<56>}{12}}\trput{6}}{
      \psset{linestyle=solid}\FAIL{spotgac<57>}\trput{2}
      \psset{linestyle=solid}\FAIL{spotgac<58>}\trput{8}
     }
     \psset{linestyle=solid}\FAIL{spotgac<59>}\trput{8}
    }
    \psset{linestyle=solid}\pstree{\CC{\hyperlink{spotgac<60>}{10}}\trput{9}}{
     \psset{linestyle=solid}\pstree{\CC{\hyperlink{spotgac<61>}{12}}\trput{6}}{
      \psset{linestyle=solid}\pstree{\CC{\hyperlink{spotgac<62>}{22}}\trput{2}}{
       \psset{linestyle=solid}\pstree{\CC{\hyperlink{spotgac<63>}{24}}\trput{2}}{
        \psset{linestyle=solid}\pstree{\CC{\hyperlink{spotgac<64>}{31}}\trput{1}}{
         \psset{linestyle=solid}\FAIL{spotgac<65>}\trput{4}
         \psset{linestyle=solid}\SUCC{spotgac<66>}\trput{8}
        }
        \psset{linestyle=solid}\FAIL{spotgac<67>}\trput{7}
       }
       \psset{linestyle=solid}\FAIL{spotgac<68>}\trput{4}
      }
      \psset{linestyle=solid}\pstree{\CC{\hyperlink{spotgac<69>}{22}}\trput{8}}{
       \psset{linestyle=solid}\pstree{\CC{\hyperlink{spotgac<70>}{24}}\trput{2}}{
        \psset{linestyle=solid}\FAIL{spotgac<71>}\trput{1}
        \psset{linestyle=solid}\FAIL{spotgac<72>}\trput{7}
       }
       \psset{linestyle=solid}\FAIL{spotgac<73>}\trput{4}
      }
     }
     \psset{linestyle=solid}\FAIL{spotgac<74>}\trput{8}
    }
   }
  }
 }
 \psset{linestyle=solid}\pstree{\CC{\hyperlink{spotgac<75>}{5}}\trput{6}}{
  \psset{linestyle=solid}\pstree{\CC{\hyperlink{spotgac<76>}{6}}\trput{2}}{
   \psset{linestyle=solid}\FAIL{spotgac<77>}\trput{4}
   \psset{linestyle=solid}\pstree{\CC{\hyperlink{spotgac<78>}{7}}\trput{5}}{
    \psset{linestyle=solid}\FAIL{spotgac<79>}\trput{4}
    \psset{linestyle=solid}\FAIL{spotgac<80>}\trput{9}
   }
  }
  \psset{linestyle=solid}\pstree{\CC{\hyperlink{spotgac<81>}{6}}\trput{4}}{
   \psset{linestyle=solid}\FAIL{spotgac<82>}\trput{2}
   \psset{linestyle=solid}\pstree{\CC{\hyperlink{spotgac<83>}{7}}\trput{5}}{
    \psset{linestyle=solid}\FAIL{spotgac<84>}\trput{2}
    \psset{linestyle=solid}\FAIL{spotgac<85>}\trput{9}
   }
  }
 }
}
}
}
 %search tree
\vfill
\hyperlinkframestart<2-85>{\beamerreturnbutton{Back to Start}}
\hyperlinkframeend<1-84>{\beamerskipbutton{Skip Animation}}
\end{frame}

\begin{frame}[label=spotbc]
\frametitle{9Spot Search (Bounds Consistency)}
\scalebox{ 0.30}{\psset{nodesep=0pt,levelsep=36pt}
\pstree{\CC{\hyperlink{spotbc<1>}{1}}}{
 \psset{linestyle=solid}\pstree{\CC{\hyperlink{spotbc<2>}{2}}\trput{1}}{
  \psset{linestyle=solid}\pstree{\CC{\hyperlink{spotbc<3>}{4}}\trput{3}}{
   \psset{linestyle=solid}\FAILTREE{spotbc<383>}{149 231}\trput{2}
   \psset{linestyle=solid}\FAILTREE{spotbc<509>}{53 73}\trput{4}
   \psset{linestyle=solid}\pstree{\CC{\hyperlink{spotbc<510>}{5}}\trput{5}}{
    \psset{linestyle=solid}\FAILTREE{spotbc<719>}{82 127}\trput{2}
    \psset{linestyle=solid}\FAILTREE{spotbc<996>}{107 170}\trput{4}
    \psset{linestyle=solid}\pstree{\CC{\hyperlink{spotbc<997>}{6}}\trput{6}}{
     \psset{linestyle=solid}\FAILTREE{spotbc<1041>}{16 28}\trput{2}
     \psset{linestyle=solid}\pstree{\CC{\hyperlink{spotbc<1042>}{7}}\trput{4}}{
      \psset{linestyle=solid}\FAILTREE{spotbc<1220>}{70 108}\trput{2}
      \psset{linestyle=solid}\pstree{\CC{\hyperlink{spotbc<1221>}{10}}\trput{9}}{
       \psset{linestyle=solid}\FAILTREE{spotbc<1272>}{20 31}\trput{4}
       \psset{linestyle=solid}\pstree{\CC{\hyperlink{spotbc<1273>}{12}}\trput{6}}{
        \psset{linestyle=solid}\pstree{\CC{\hyperlink{spotbc<1274>}{22}}\trput{2}}{
         \psset{linestyle=solid}\FAIL{spotbc<1275>}\trput{1}
         \psset{linestyle=solid}\pstree{\CC{\hyperlink{spotbc<1276>}{23}}\trput{2}}{
          \psset{linestyle=solid}\FAIL{spotbc<1277>}\trput{1}
          \psset{linestyle=solid}\pstree{\CC{\hyperlink{spotbc<1278>}{24}}\trput{4}}{
           \psset{linestyle=solid}\pstree{\CC{\hyperlink{spotbc<1279>}{31}}\trput{1}}{
            \psset{linestyle=solid}\FAIL{spotbc<1280>}\trput{4}
            \psset{linestyle=solid}\SUCC{spotbc<1281>}\trput{8}
           }
           \psset{linestyle=solid}\FAILTREE{spotbc<1287>}{2 4}\trput{7}
          }
         }
         \psset{linestyle=solid}\FAILTREE{spotbc<1292>}{2 3}\trput{4}
         \psset{linestyle=solid}\FAIL{spotbc<1293>}\trput{6}
        }
        \psset{linestyle=solid}\FAILTREE{spotbc<1315>}{9 13}\trput{8}
       }
       \psset{linestyle=solid}\FAILTREE{spotbc<1332>}{5 12}\trput{8}
      }
     }
    }
   }
   \psset{linestyle=solid}\FAILTREE{spotbc<1619>}{121 166}\trput{6}
  }
  \psset{linestyle=solid}\FAILTREE{spotbc<1630>}{4 7}\trput{4}
  \psset{linestyle=solid}\FAILTREE{spotbc<1652>}{9 13}\trput{5}
  \psset{linestyle=solid}\FAILTREE{spotbc<1702>}{23 27}\trput{6}
 }
 \psset{linestyle=solid}\FAILTREE{spotbc<1780>}{36 42}\trput{4}
 \psset{linestyle=solid}\FAILTREE{spotbc<1817>}{17 20}\trput{5}
 \psset{linestyle=solid}\pstree{\CC{\hyperlink{spotbc<1818>}{2}}\trput{6}}{
  \psset{linestyle=solid}\FAILTREE{spotbc<1866>}{23 25}\trput{1}
  \psset{linestyle=solid}\pstree{\CC{\hyperlink{spotbc<1867>}{5}}\trput{3}}{
   \psset{linestyle=solid}\FAILTREE{spotbc<1896>}{13 16}\trput{2}
   \psset{linestyle=solid}\pstree{\CC{\hyperlink{spotbc<1897>}{6}}\trput{4}}{
    \psset{linestyle=solid}\FAILTREE{spotbc<1907>}{4 6}\trput{2}
    \psset{linestyle=solid}\pstree{\CC{\hyperlink{spotbc<1908>}{7}}\trput{5}}{
     \psset{linestyle=solid}\FAILTREE{spotbc<1920>}{5 7}\trput{2}
     \psset{linestyle=solid}\pstree{\CC{\hyperlink{spotbc<1921>}{10}}\trput{9}}{
      \psset{linestyle=solid}\FAILTREE{spotbc<1926>}{2 3}\trput{4}
      \psset{linestyle=solid}\FAILTREE{spotbc<1937>}{4 7}\trput{5}
      \psset{linestyle=solid}\FAIL{spotbc<1938>}\trput{8}
     }
    }
   }
  }
 }
}
}
 %search tree
\end{frame}

\begin{frame}[label=spotgac1]
\frametitle{9Spot Search (GAC, compact)}
\only<1>{
\scalebox{ 0.40}{\psset{nodesep=0pt,levelsep=36pt}
\pstree{\CC{\hyperlink{spotgac1<1>}{4}}}{
}
}
}
\only<2>{
\scalebox{ 0.40}{\psset{nodesep=0pt,levelsep=36pt}
\pstree{\CC{\hyperlink{spotgac1<1>}{4}}}{
 \psset{linestyle=solid}\pstree{\CC{\hyperlink{spotgac1<2>}{5}}\trput{2}}{
 }
}
}
}
\only<3>{
\scalebox{ 0.40}{\psset{nodesep=0pt,levelsep=36pt}
\pstree{\CC{\hyperlink{spotgac1<1>}{4}}}{
 \psset{linestyle=solid}\pstree{\CC{\hyperlink{spotgac1<2>}{5}}\trput{2}}{
  \psset{linestyle=solid}\pstree{\CC{\hyperlink{spotgac1<3>}{6}}\trput{4}}{
  }
 }
}
}
}
\only<4>{
\scalebox{ 0.40}{\psset{nodesep=0pt,levelsep=36pt}
\pstree{\CC{\hyperlink{spotgac1<1>}{4}}}{
 \psset{linestyle=solid}\pstree{\CC{\hyperlink{spotgac1<2>}{5}}\trput{2}}{
  \psset{linestyle=solid}\pstree{\CC{\hyperlink{spotgac1<3>}{6}}\trput{4}}{
   \psset{linestyle=solid}\pstree{\CC{\hyperlink{spotgac1<4>}{7}}\trput{5}}{
   }
  }
 }
}
}
}
\only<5>{
\scalebox{ 0.40}{\psset{nodesep=0pt,levelsep=36pt}
\pstree{\CC{\hyperlink{spotgac1<1>}{4}}}{
 \psset{linestyle=solid}\pstree{\CC{\hyperlink{spotgac1<2>}{5}}\trput{2}}{
  \psset{linestyle=solid}\pstree{\CC{\hyperlink{spotgac1<3>}{6}}\trput{4}}{
   \psset{linestyle=solid}\pstree{\CC{\hyperlink{spotgac1<4>}{7}}\trput{5}}{
    \psset{linestyle=solid}\FAIL{spotgac1<5>}\trput{6}
   }
  }
 }
}
}
}
\only<6>{
\scalebox{ 0.40}{\psset{nodesep=0pt,levelsep=36pt}
\pstree{\CC{\hyperlink{spotgac1<1>}{4}}}{
 \psset{linestyle=solid}\pstree{\CC{\hyperlink{spotgac1<2>}{5}}\trput{2}}{
  \psset{linestyle=solid}\pstree{\CC{\hyperlink{spotgac1<3>}{6}}\trput{4}}{
   \psset{linestyle=solid}\pstree{\CC{\hyperlink{spotgac1<4>}{7}}\trput{5}}{
    \psset{linestyle=solid}\FAIL{spotgac1<5>}\trput{6}
    \psset{linestyle=solid}\pstree{\CC{\hyperlink{spotgac1<6>}{12}}\trput{9}}{
    }
   }
  }
 }
}
}
}
\only<7>{
\scalebox{ 0.40}{\psset{nodesep=0pt,levelsep=36pt}
\pstree{\CC{\hyperlink{spotgac1<1>}{4}}}{
 \psset{linestyle=solid}\pstree{\CC{\hyperlink{spotgac1<2>}{5}}\trput{2}}{
  \psset{linestyle=solid}\pstree{\CC{\hyperlink{spotgac1<3>}{6}}\trput{4}}{
   \psset{linestyle=solid}\pstree{\CC{\hyperlink{spotgac1<4>}{7}}\trput{5}}{
    \psset{linestyle=solid}\FAIL{spotgac1<5>}\trput{6}
    \psset{linestyle=solid}\pstree{\CC{\hyperlink{spotgac1<6>}{12}}\trput{9}}{
     \psset{linestyle=solid}\pstree{\CC{\hyperlink{spotgac1<7>}{19}}\trput{6}}{
     }
    }
   }
  }
 }
}
}
}
\only<8>{
\scalebox{ 0.40}{\psset{nodesep=0pt,levelsep=36pt}
\pstree{\CC{\hyperlink{spotgac1<1>}{4}}}{
 \psset{linestyle=solid}\pstree{\CC{\hyperlink{spotgac1<2>}{5}}\trput{2}}{
  \psset{linestyle=solid}\pstree{\CC{\hyperlink{spotgac1<3>}{6}}\trput{4}}{
   \psset{linestyle=solid}\pstree{\CC{\hyperlink{spotgac1<4>}{7}}\trput{5}}{
    \psset{linestyle=solid}\FAIL{spotgac1<5>}\trput{6}
    \psset{linestyle=solid}\pstree{\CC{\hyperlink{spotgac1<6>}{12}}\trput{9}}{
     \psset{linestyle=solid}\pstree{\CC{\hyperlink{spotgac1<7>}{19}}\trput{6}}{
      \psset{linestyle=solid}\FAIL{spotgac1<8>}\trput{6}
     }
    }
   }
  }
 }
}
}
}
\only<9>{
\scalebox{ 0.40}{\psset{nodesep=0pt,levelsep=36pt}
\pstree{\CC{\hyperlink{spotgac1<1>}{4}}}{
 \psset{linestyle=solid}\pstree{\CC{\hyperlink{spotgac1<2>}{5}}\trput{2}}{
  \psset{linestyle=solid}\pstree{\CC{\hyperlink{spotgac1<3>}{6}}\trput{4}}{
   \psset{linestyle=solid}\pstree{\CC{\hyperlink{spotgac1<4>}{7}}\trput{5}}{
    \psset{linestyle=solid}\FAIL{spotgac1<5>}\trput{6}
    \psset{linestyle=solid}\pstree{\CC{\hyperlink{spotgac1<6>}{12}}\trput{9}}{
     \psset{linestyle=solid}\pstree{\CC{\hyperlink{spotgac1<7>}{19}}\trput{6}}{
      \psset{linestyle=solid}\FAIL{spotgac1<8>}\trput{6}
      \psset{linestyle=solid}\pstree{\CC{\hyperlink{spotgac1<9>}{20}}\trput{9}}{
      }
     }
    }
   }
  }
 }
}
}
}
\only<10>{
\scalebox{ 0.40}{\psset{nodesep=0pt,levelsep=36pt}
\pstree{\CC{\hyperlink{spotgac1<1>}{4}}}{
 \psset{linestyle=solid}\pstree{\CC{\hyperlink{spotgac1<2>}{5}}\trput{2}}{
  \psset{linestyle=solid}\pstree{\CC{\hyperlink{spotgac1<3>}{6}}\trput{4}}{
   \psset{linestyle=solid}\pstree{\CC{\hyperlink{spotgac1<4>}{7}}\trput{5}}{
    \psset{linestyle=solid}\FAIL{spotgac1<5>}\trput{6}
    \psset{linestyle=solid}\pstree{\CC{\hyperlink{spotgac1<6>}{12}}\trput{9}}{
     \psset{linestyle=solid}\pstree{\CC{\hyperlink{spotgac1<7>}{19}}\trput{6}}{
      \psset{linestyle=solid}\FAIL{spotgac1<8>}\trput{6}
      \psset{linestyle=solid}\pstree{\CC{\hyperlink{spotgac1<9>}{20}}\trput{9}}{
       \psset{linestyle=solid}\pstree{\CC{\hyperlink{spotgac1<10>}{23}}\trput{6}}{
       }
      }
     }
    }
   }
  }
 }
}
}
}
\only<11>{
\scalebox{ 0.40}{\psset{nodesep=0pt,levelsep=36pt}
\pstree{\CC{\hyperlink{spotgac1<1>}{4}}}{
 \psset{linestyle=solid}\pstree{\CC{\hyperlink{spotgac1<2>}{5}}\trput{2}}{
  \psset{linestyle=solid}\pstree{\CC{\hyperlink{spotgac1<3>}{6}}\trput{4}}{
   \psset{linestyle=solid}\pstree{\CC{\hyperlink{spotgac1<4>}{7}}\trput{5}}{
    \psset{linestyle=solid}\FAIL{spotgac1<5>}\trput{6}
    \psset{linestyle=solid}\pstree{\CC{\hyperlink{spotgac1<6>}{12}}\trput{9}}{
     \psset{linestyle=solid}\pstree{\CC{\hyperlink{spotgac1<7>}{19}}\trput{6}}{
      \psset{linestyle=solid}\FAIL{spotgac1<8>}\trput{6}
      \psset{linestyle=solid}\pstree{\CC{\hyperlink{spotgac1<9>}{20}}\trput{9}}{
       \psset{linestyle=solid}\pstree{\CC{\hyperlink{spotgac1<10>}{23}}\trput{6}}{
        \psset{linestyle=solid}\FAIL{spotgac1<11>}\trput{2}
       }
      }
     }
    }
   }
  }
 }
}
}
}
\only<12>{
\scalebox{ 0.40}{\psset{nodesep=0pt,levelsep=36pt}
\pstree{\CC{\hyperlink{spotgac1<1>}{4}}}{
 \psset{linestyle=solid}\pstree{\CC{\hyperlink{spotgac1<2>}{5}}\trput{2}}{
  \psset{linestyle=solid}\pstree{\CC{\hyperlink{spotgac1<3>}{6}}\trput{4}}{
   \psset{linestyle=solid}\pstree{\CC{\hyperlink{spotgac1<4>}{7}}\trput{5}}{
    \psset{linestyle=solid}\FAIL{spotgac1<5>}\trput{6}
    \psset{linestyle=solid}\pstree{\CC{\hyperlink{spotgac1<6>}{12}}\trput{9}}{
     \psset{linestyle=solid}\pstree{\CC{\hyperlink{spotgac1<7>}{19}}\trput{6}}{
      \psset{linestyle=solid}\FAIL{spotgac1<8>}\trput{6}
      \psset{linestyle=solid}\pstree{\CC{\hyperlink{spotgac1<9>}{20}}\trput{9}}{
       \psset{linestyle=solid}\pstree{\CC{\hyperlink{spotgac1<10>}{23}}\trput{6}}{
        \psset{linestyle=solid}\FAIL{spotgac1<11>}\trput{2}
        \psset{linestyle=solid}\FAIL{spotgac1<12>}\trput{8}
       }
      }
     }
    }
   }
  }
 }
}
}
}
\only<13>{
\scalebox{ 0.40}{\psset{nodesep=0pt,levelsep=36pt}
\pstree{\CC{\hyperlink{spotgac1<1>}{4}}}{
 \psset{linestyle=solid}\pstree{\CC{\hyperlink{spotgac1<2>}{5}}\trput{2}}{
  \psset{linestyle=solid}\pstree{\CC{\hyperlink{spotgac1<3>}{6}}\trput{4}}{
   \psset{linestyle=solid}\pstree{\CC{\hyperlink{spotgac1<4>}{7}}\trput{5}}{
    \psset{linestyle=solid}\FAIL{spotgac1<5>}\trput{6}
    \psset{linestyle=solid}\pstree{\CC{\hyperlink{spotgac1<6>}{12}}\trput{9}}{
     \psset{linestyle=solid}\pstree{\CC{\hyperlink{spotgac1<7>}{19}}\trput{6}}{
      \psset{linestyle=solid}\FAIL{spotgac1<8>}\trput{6}
      \psset{linestyle=solid}\pstree{\CC{\hyperlink{spotgac1<9>}{20}}\trput{9}}{
       \psset{linestyle=solid}\FAILTREE{spotgac1<12>}{1 2}\trput{6}
       \psset{linestyle=solid}\FAIL{spotgac1<13>}\trput{8}
      }
     }
    }
   }
  }
 }
}
}
}
\only<14>{
\scalebox{ 0.40}{\psset{nodesep=0pt,levelsep=36pt}
\pstree{\CC{\hyperlink{spotgac1<1>}{4}}}{
 \psset{linestyle=solid}\pstree{\CC{\hyperlink{spotgac1<2>}{5}}\trput{2}}{
  \psset{linestyle=solid}\pstree{\CC{\hyperlink{spotgac1<3>}{6}}\trput{4}}{
   \psset{linestyle=solid}\pstree{\CC{\hyperlink{spotgac1<4>}{7}}\trput{5}}{
    \psset{linestyle=solid}\FAIL{spotgac1<5>}\trput{6}
    \psset{linestyle=solid}\pstree{\CC{\hyperlink{spotgac1<6>}{12}}\trput{9}}{
     \psset{linestyle=solid}\FAILTREE{spotgac1<13>}{3 4}\trput{6}
     \psset{linestyle=solid}\FAIL{spotgac1<14>}\trput{8}
    }
   }
  }
 }
}
}
}
\only<15>{
\scalebox{ 0.40}{\psset{nodesep=0pt,levelsep=36pt}
\pstree{\CC{\hyperlink{spotgac1<1>}{4}}}{
 \psset{linestyle=solid}\pstree{\CC{\hyperlink{spotgac1<2>}{5}}\trput{2}}{
  \psset{linestyle=solid}\pstree{\CC{\hyperlink{spotgac1<3>}{6}}\trput{4}}{
   \psset{linestyle=solid}\FAILTREE{spotgac1<14>}{5 6}\trput{5}
   \psset{linestyle=solid}\FAIL{spotgac1<15>}\trput{6}
  }
 }
}
}
}
\only<16>{
\scalebox{ 0.40}{\psset{nodesep=0pt,levelsep=36pt}
\pstree{\CC{\hyperlink{spotgac1<1>}{4}}}{
 \psset{linestyle=solid}\pstree{\CC{\hyperlink{spotgac1<2>}{5}}\trput{2}}{
  \psset{linestyle=solid}\FAILTREE{spotgac1<15>}{6 7}\trput{4}
  \psset{linestyle=solid}\pstree{\CC{\hyperlink{spotgac1<16>}{6}}\trput{6}}{
  }
 }
}
}
}
\only<17>{
\scalebox{ 0.40}{\psset{nodesep=0pt,levelsep=36pt}
\pstree{\CC{\hyperlink{spotgac1<1>}{4}}}{
 \psset{linestyle=solid}\pstree{\CC{\hyperlink{spotgac1<2>}{5}}\trput{2}}{
  \psset{linestyle=solid}\FAILTREE{spotgac1<15>}{6 7}\trput{4}
  \psset{linestyle=solid}\pstree{\CC{\hyperlink{spotgac1<16>}{6}}\trput{6}}{
   \psset{linestyle=solid}\pstree{\CC{\hyperlink{spotgac1<17>}{10}}\trput{4}}{
   }
  }
 }
}
}
}
\only<18>{
\scalebox{ 0.40}{\psset{nodesep=0pt,levelsep=36pt}
\pstree{\CC{\hyperlink{spotgac1<1>}{4}}}{
 \psset{linestyle=solid}\pstree{\CC{\hyperlink{spotgac1<2>}{5}}\trput{2}}{
  \psset{linestyle=solid}\FAILTREE{spotgac1<15>}{6 7}\trput{4}
  \psset{linestyle=solid}\pstree{\CC{\hyperlink{spotgac1<16>}{6}}\trput{6}}{
   \psset{linestyle=solid}\pstree{\CC{\hyperlink{spotgac1<17>}{10}}\trput{4}}{
    \psset{linestyle=solid}\FAIL{spotgac1<18>}\trput{4}
   }
  }
 }
}
}
}
\only<19>{
\scalebox{ 0.40}{\psset{nodesep=0pt,levelsep=36pt}
\pstree{\CC{\hyperlink{spotgac1<1>}{4}}}{
 \psset{linestyle=solid}\pstree{\CC{\hyperlink{spotgac1<2>}{5}}\trput{2}}{
  \psset{linestyle=solid}\FAILTREE{spotgac1<15>}{6 7}\trput{4}
  \psset{linestyle=solid}\pstree{\CC{\hyperlink{spotgac1<16>}{6}}\trput{6}}{
   \psset{linestyle=solid}\pstree{\CC{\hyperlink{spotgac1<17>}{10}}\trput{4}}{
    \psset{linestyle=solid}\FAIL{spotgac1<18>}\trput{4}
    \psset{linestyle=solid}\FAIL{spotgac1<19>}\trput{6}
   }
  }
 }
}
}
}
\only<20>{
\scalebox{ 0.40}{\psset{nodesep=0pt,levelsep=36pt}
\pstree{\CC{\hyperlink{spotgac1<1>}{4}}}{
 \psset{linestyle=solid}\pstree{\CC{\hyperlink{spotgac1<2>}{5}}\trput{2}}{
  \psset{linestyle=solid}\FAILTREE{spotgac1<15>}{6 7}\trput{4}
  \psset{linestyle=solid}\pstree{\CC{\hyperlink{spotgac1<16>}{6}}\trput{6}}{
   \psset{linestyle=solid}\pstree{\CC{\hyperlink{spotgac1<17>}{10}}\trput{4}}{
    \psset{linestyle=solid}\FAIL{spotgac1<18>}\trput{4}
    \psset{linestyle=solid}\FAIL{spotgac1<19>}\trput{6}
    \psset{linestyle=solid}\FAIL{spotgac1<20>}\trput{8}
   }
  }
 }
}
}
}
\only<21>{
\scalebox{ 0.40}{\psset{nodesep=0pt,levelsep=36pt}
\pstree{\CC{\hyperlink{spotgac1<1>}{4}}}{
 \psset{linestyle=solid}\pstree{\CC{\hyperlink{spotgac1<2>}{5}}\trput{2}}{
  \psset{linestyle=solid}\FAILTREE{spotgac1<15>}{6 7}\trput{4}
  \psset{linestyle=solid}\pstree{\CC{\hyperlink{spotgac1<16>}{6}}\trput{6}}{
   \psset{linestyle=solid}\FAILTREE{spotgac1<20>}{1 3}\trput{4}
   \psset{linestyle=solid}\pstree{\CC{\hyperlink{spotgac1<21>}{7}}\trput{5}}{
   }
  }
 }
}
}
}
\only<22>{
\scalebox{ 0.40}{\psset{nodesep=0pt,levelsep=36pt}
\pstree{\CC{\hyperlink{spotgac1<1>}{4}}}{
 \psset{linestyle=solid}\pstree{\CC{\hyperlink{spotgac1<2>}{5}}\trput{2}}{
  \psset{linestyle=solid}\FAILTREE{spotgac1<15>}{6 7}\trput{4}
  \psset{linestyle=solid}\pstree{\CC{\hyperlink{spotgac1<16>}{6}}\trput{6}}{
   \psset{linestyle=solid}\FAILTREE{spotgac1<20>}{1 3}\trput{4}
   \psset{linestyle=solid}\pstree{\CC{\hyperlink{spotgac1<21>}{7}}\trput{5}}{
    \psset{linestyle=solid}\pstree{\CC{\hyperlink{spotgac1<22>}{10}}\trput{4}}{
    }
   }
  }
 }
}
}
}
\only<23>{
\scalebox{ 0.40}{\psset{nodesep=0pt,levelsep=36pt}
\pstree{\CC{\hyperlink{spotgac1<1>}{4}}}{
 \psset{linestyle=solid}\pstree{\CC{\hyperlink{spotgac1<2>}{5}}\trput{2}}{
  \psset{linestyle=solid}\FAILTREE{spotgac1<15>}{6 7}\trput{4}
  \psset{linestyle=solid}\pstree{\CC{\hyperlink{spotgac1<16>}{6}}\trput{6}}{
   \psset{linestyle=solid}\FAILTREE{spotgac1<20>}{1 3}\trput{4}
   \psset{linestyle=solid}\pstree{\CC{\hyperlink{spotgac1<21>}{7}}\trput{5}}{
    \psset{linestyle=solid}\pstree{\CC{\hyperlink{spotgac1<22>}{10}}\trput{4}}{
     \psset{linestyle=solid}\FAIL{spotgac1<23>}\trput{5}
    }
   }
  }
 }
}
}
}
\only<24>{
\scalebox{ 0.40}{\psset{nodesep=0pt,levelsep=36pt}
\pstree{\CC{\hyperlink{spotgac1<1>}{4}}}{
 \psset{linestyle=solid}\pstree{\CC{\hyperlink{spotgac1<2>}{5}}\trput{2}}{
  \psset{linestyle=solid}\FAILTREE{spotgac1<15>}{6 7}\trput{4}
  \psset{linestyle=solid}\pstree{\CC{\hyperlink{spotgac1<16>}{6}}\trput{6}}{
   \psset{linestyle=solid}\FAILTREE{spotgac1<20>}{1 3}\trput{4}
   \psset{linestyle=solid}\pstree{\CC{\hyperlink{spotgac1<21>}{7}}\trput{5}}{
    \psset{linestyle=solid}\pstree{\CC{\hyperlink{spotgac1<22>}{10}}\trput{4}}{
     \psset{linestyle=solid}\FAIL{spotgac1<23>}\trput{5}
     \psset{linestyle=solid}\FAIL{spotgac1<24>}\trput{6}
    }
   }
  }
 }
}
}
}
\only<25>{
\scalebox{ 0.40}{\psset{nodesep=0pt,levelsep=36pt}
\pstree{\CC{\hyperlink{spotgac1<1>}{4}}}{
 \psset{linestyle=solid}\pstree{\CC{\hyperlink{spotgac1<2>}{5}}\trput{2}}{
  \psset{linestyle=solid}\FAILTREE{spotgac1<15>}{6 7}\trput{4}
  \psset{linestyle=solid}\pstree{\CC{\hyperlink{spotgac1<16>}{6}}\trput{6}}{
   \psset{linestyle=solid}\FAILTREE{spotgac1<20>}{1 3}\trput{4}
   \psset{linestyle=solid}\pstree{\CC{\hyperlink{spotgac1<21>}{7}}\trput{5}}{
    \psset{linestyle=solid}\pstree{\CC{\hyperlink{spotgac1<22>}{10}}\trput{4}}{
     \psset{linestyle=solid}\FAIL{spotgac1<23>}\trput{5}
     \psset{linestyle=solid}\FAIL{spotgac1<24>}\trput{6}
     \psset{linestyle=solid}\FAIL{spotgac1<25>}\trput{8}
    }
   }
  }
 }
}
}
}
\only<26>{
\scalebox{ 0.40}{\psset{nodesep=0pt,levelsep=36pt}
\pstree{\CC{\hyperlink{spotgac1<1>}{4}}}{
 \psset{linestyle=solid}\pstree{\CC{\hyperlink{spotgac1<2>}{5}}\trput{2}}{
  \psset{linestyle=solid}\FAILTREE{spotgac1<15>}{6 7}\trput{4}
  \psset{linestyle=solid}\pstree{\CC{\hyperlink{spotgac1<16>}{6}}\trput{6}}{
   \psset{linestyle=solid}\FAILTREE{spotgac1<20>}{1 3}\trput{4}
   \psset{linestyle=solid}\pstree{\CC{\hyperlink{spotgac1<21>}{7}}\trput{5}}{
    \psset{linestyle=solid}\FAILTREE{spotgac1<25>}{1 3}\trput{4}
    \psset{linestyle=solid}\pstree{\CC{\hyperlink{spotgac1<26>}{10}}\trput{9}}{
    }
   }
  }
 }
}
}
}
\only<27>{
\scalebox{ 0.40}{\psset{nodesep=0pt,levelsep=36pt}
\pstree{\CC{\hyperlink{spotgac1<1>}{4}}}{
 \psset{linestyle=solid}\pstree{\CC{\hyperlink{spotgac1<2>}{5}}\trput{2}}{
  \psset{linestyle=solid}\FAILTREE{spotgac1<15>}{6 7}\trput{4}
  \psset{linestyle=solid}\pstree{\CC{\hyperlink{spotgac1<16>}{6}}\trput{6}}{
   \psset{linestyle=solid}\FAILTREE{spotgac1<20>}{1 3}\trput{4}
   \psset{linestyle=solid}\pstree{\CC{\hyperlink{spotgac1<21>}{7}}\trput{5}}{
    \psset{linestyle=solid}\FAILTREE{spotgac1<25>}{1 3}\trput{4}
    \psset{linestyle=solid}\pstree{\CC{\hyperlink{spotgac1<26>}{10}}\trput{9}}{
     \psset{linestyle=solid}\pstree{\CC{\hyperlink{spotgac1<27>}{11}}\trput{5}}{
     }
    }
   }
  }
 }
}
}
}
\only<28>{
\scalebox{ 0.40}{\psset{nodesep=0pt,levelsep=36pt}
\pstree{\CC{\hyperlink{spotgac1<1>}{4}}}{
 \psset{linestyle=solid}\pstree{\CC{\hyperlink{spotgac1<2>}{5}}\trput{2}}{
  \psset{linestyle=solid}\FAILTREE{spotgac1<15>}{6 7}\trput{4}
  \psset{linestyle=solid}\pstree{\CC{\hyperlink{spotgac1<16>}{6}}\trput{6}}{
   \psset{linestyle=solid}\FAILTREE{spotgac1<20>}{1 3}\trput{4}
   \psset{linestyle=solid}\pstree{\CC{\hyperlink{spotgac1<21>}{7}}\trput{5}}{
    \psset{linestyle=solid}\FAILTREE{spotgac1<25>}{1 3}\trput{4}
    \psset{linestyle=solid}\pstree{\CC{\hyperlink{spotgac1<26>}{10}}\trput{9}}{
     \psset{linestyle=solid}\pstree{\CC{\hyperlink{spotgac1<27>}{11}}\trput{5}}{
      \psset{linestyle=solid}\FAIL{spotgac1<28>}\trput{6}
     }
    }
   }
  }
 }
}
}
}
\only<29>{
\scalebox{ 0.40}{\psset{nodesep=0pt,levelsep=36pt}
\pstree{\CC{\hyperlink{spotgac1<1>}{4}}}{
 \psset{linestyle=solid}\pstree{\CC{\hyperlink{spotgac1<2>}{5}}\trput{2}}{
  \psset{linestyle=solid}\FAILTREE{spotgac1<15>}{6 7}\trput{4}
  \psset{linestyle=solid}\pstree{\CC{\hyperlink{spotgac1<16>}{6}}\trput{6}}{
   \psset{linestyle=solid}\FAILTREE{spotgac1<20>}{1 3}\trput{4}
   \psset{linestyle=solid}\pstree{\CC{\hyperlink{spotgac1<21>}{7}}\trput{5}}{
    \psset{linestyle=solid}\FAILTREE{spotgac1<25>}{1 3}\trput{4}
    \psset{linestyle=solid}\pstree{\CC{\hyperlink{spotgac1<26>}{10}}\trput{9}}{
     \psset{linestyle=solid}\pstree{\CC{\hyperlink{spotgac1<27>}{11}}\trput{5}}{
      \psset{linestyle=solid}\FAIL{spotgac1<28>}\trput{6}
      \psset{linestyle=solid}\FAIL{spotgac1<29>}\trput{8}
     }
    }
   }
  }
 }
}
}
}
\only<30>{
\scalebox{ 0.40}{\psset{nodesep=0pt,levelsep=36pt}
\pstree{\CC{\hyperlink{spotgac1<1>}{4}}}{
 \psset{linestyle=solid}\pstree{\CC{\hyperlink{spotgac1<2>}{5}}\trput{2}}{
  \psset{linestyle=solid}\FAILTREE{spotgac1<15>}{6 7}\trput{4}
  \psset{linestyle=solid}\pstree{\CC{\hyperlink{spotgac1<16>}{6}}\trput{6}}{
   \psset{linestyle=solid}\FAILTREE{spotgac1<20>}{1 3}\trput{4}
   \psset{linestyle=solid}\pstree{\CC{\hyperlink{spotgac1<21>}{7}}\trput{5}}{
    \psset{linestyle=solid}\FAILTREE{spotgac1<25>}{1 3}\trput{4}
    \psset{linestyle=solid}\pstree{\CC{\hyperlink{spotgac1<26>}{10}}\trput{9}}{
     \psset{linestyle=solid}\FAILTREE{spotgac1<29>}{1 2}\trput{5}
     \psset{linestyle=solid}\FAIL{spotgac1<30>}\trput{6}
    }
   }
  }
 }
}
}
}
\only<31>{
\scalebox{ 0.40}{\psset{nodesep=0pt,levelsep=36pt}
\pstree{\CC{\hyperlink{spotgac1<1>}{4}}}{
 \psset{linestyle=solid}\pstree{\CC{\hyperlink{spotgac1<2>}{5}}\trput{2}}{
  \psset{linestyle=solid}\FAILTREE{spotgac1<15>}{6 7}\trput{4}
  \psset{linestyle=solid}\pstree{\CC{\hyperlink{spotgac1<16>}{6}}\trput{6}}{
   \psset{linestyle=solid}\FAILTREE{spotgac1<20>}{1 3}\trput{4}
   \psset{linestyle=solid}\pstree{\CC{\hyperlink{spotgac1<21>}{7}}\trput{5}}{
    \psset{linestyle=solid}\FAILTREE{spotgac1<25>}{1 3}\trput{4}
    \psset{linestyle=solid}\pstree{\CC{\hyperlink{spotgac1<26>}{10}}\trput{9}}{
     \psset{linestyle=solid}\FAILTREE{spotgac1<29>}{1 2}\trput{5}
     \psset{linestyle=solid}\FAIL{spotgac1<30>}\trput{6}
     \psset{linestyle=solid}\FAIL{spotgac1<31>}\trput{8}
    }
   }
  }
 }
}
}
}
\only<32>{
\scalebox{ 0.40}{\psset{nodesep=0pt,levelsep=36pt}
\pstree{\CC{\hyperlink{spotgac1<1>}{4}}}{
 \psset{linestyle=solid}\FAILTREE{spotgac1<31>}{13 17}\trput{2}
 \psset{linestyle=solid}\pstree{\CC{\hyperlink{spotgac1<32>}{5}}\trput{4}}{
 }
}
}
}
\only<33>{
\scalebox{ 0.40}{\psset{nodesep=0pt,levelsep=36pt}
\pstree{\CC{\hyperlink{spotgac1<1>}{4}}}{
 \psset{linestyle=solid}\FAILTREE{spotgac1<31>}{13 17}\trput{2}
 \psset{linestyle=solid}\pstree{\CC{\hyperlink{spotgac1<32>}{5}}\trput{4}}{
  \psset{linestyle=solid}\pstree{\CC{\hyperlink{spotgac1<33>}{6}}\trput{2}}{
  }
 }
}
}
}
\only<34>{
\scalebox{ 0.40}{\psset{nodesep=0pt,levelsep=36pt}
\pstree{\CC{\hyperlink{spotgac1<1>}{4}}}{
 \psset{linestyle=solid}\FAILTREE{spotgac1<31>}{13 17}\trput{2}
 \psset{linestyle=solid}\pstree{\CC{\hyperlink{spotgac1<32>}{5}}\trput{4}}{
  \psset{linestyle=solid}\pstree{\CC{\hyperlink{spotgac1<33>}{6}}\trput{2}}{
   \psset{linestyle=solid}\FAIL{spotgac1<34>}\trput{5}
  }
 }
}
}
}
\only<35>{
\scalebox{ 0.40}{\psset{nodesep=0pt,levelsep=36pt}
\pstree{\CC{\hyperlink{spotgac1<1>}{4}}}{
 \psset{linestyle=solid}\FAILTREE{spotgac1<31>}{13 17}\trput{2}
 \psset{linestyle=solid}\pstree{\CC{\hyperlink{spotgac1<32>}{5}}\trput{4}}{
  \psset{linestyle=solid}\pstree{\CC{\hyperlink{spotgac1<33>}{6}}\trput{2}}{
   \psset{linestyle=solid}\FAIL{spotgac1<34>}\trput{5}
   \psset{linestyle=solid}\FAIL{spotgac1<35>}\trput{6}
  }
 }
}
}
}
\only<36>{
\scalebox{ 0.40}{\psset{nodesep=0pt,levelsep=36pt}
\pstree{\CC{\hyperlink{spotgac1<1>}{4}}}{
 \psset{linestyle=solid}\FAILTREE{spotgac1<31>}{13 17}\trput{2}
 \psset{linestyle=solid}\pstree{\CC{\hyperlink{spotgac1<32>}{5}}\trput{4}}{
  \psset{linestyle=solid}\FAILTREE{spotgac1<35>}{1 2}\trput{2}
  \psset{linestyle=solid}\pstree{\CC{\hyperlink{spotgac1<36>}{6}}\trput{6}}{
  }
 }
}
}
}
\only<37>{
\scalebox{ 0.40}{\psset{nodesep=0pt,levelsep=36pt}
\pstree{\CC{\hyperlink{spotgac1<1>}{4}}}{
 \psset{linestyle=solid}\FAILTREE{spotgac1<31>}{13 17}\trput{2}
 \psset{linestyle=solid}\pstree{\CC{\hyperlink{spotgac1<32>}{5}}\trput{4}}{
  \psset{linestyle=solid}\FAILTREE{spotgac1<35>}{1 2}\trput{2}
  \psset{linestyle=solid}\pstree{\CC{\hyperlink{spotgac1<36>}{6}}\trput{6}}{
   \psset{linestyle=solid}\FAIL{spotgac1<37>}\trput{2}
  }
 }
}
}
}
\only<38>{
\scalebox{ 0.40}{\psset{nodesep=0pt,levelsep=36pt}
\pstree{\CC{\hyperlink{spotgac1<1>}{4}}}{
 \psset{linestyle=solid}\FAILTREE{spotgac1<31>}{13 17}\trput{2}
 \psset{linestyle=solid}\pstree{\CC{\hyperlink{spotgac1<32>}{5}}\trput{4}}{
  \psset{linestyle=solid}\FAILTREE{spotgac1<35>}{1 2}\trput{2}
  \psset{linestyle=solid}\pstree{\CC{\hyperlink{spotgac1<36>}{6}}\trput{6}}{
   \psset{linestyle=solid}\FAIL{spotgac1<37>}\trput{2}
   \psset{linestyle=solid}\FAIL{spotgac1<38>}\trput{5}
  }
 }
}
}
}
\only<39>{
\scalebox{ 0.40}{\psset{nodesep=0pt,levelsep=36pt}
\pstree{\CC{\hyperlink{spotgac1<1>}{4}}}{
 \psset{linestyle=solid}\FAILTREE{spotgac1<31>}{13 17}\trput{2}
 \psset{linestyle=solid}\FAILTREE{spotgac1<38>}{3 4}\trput{4}
 \psset{linestyle=solid}\pstree{\CC{\hyperlink{spotgac1<39>}{5}}\trput{5}}{
 }
}
}
}
\only<40>{
\scalebox{ 0.40}{\psset{nodesep=0pt,levelsep=36pt}
\pstree{\CC{\hyperlink{spotgac1<1>}{4}}}{
 \psset{linestyle=solid}\FAILTREE{spotgac1<31>}{13 17}\trput{2}
 \psset{linestyle=solid}\FAILTREE{spotgac1<38>}{3 4}\trput{4}
 \psset{linestyle=solid}\pstree{\CC{\hyperlink{spotgac1<39>}{5}}\trput{5}}{
  \psset{linestyle=solid}\pstree{\CC{\hyperlink{spotgac1<40>}{6}}\trput{2}}{
  }
 }
}
}
}
\only<41>{
\scalebox{ 0.40}{\psset{nodesep=0pt,levelsep=36pt}
\pstree{\CC{\hyperlink{spotgac1<1>}{4}}}{
 \psset{linestyle=solid}\FAILTREE{spotgac1<31>}{13 17}\trput{2}
 \psset{linestyle=solid}\FAILTREE{spotgac1<38>}{3 4}\trput{4}
 \psset{linestyle=solid}\pstree{\CC{\hyperlink{spotgac1<39>}{5}}\trput{5}}{
  \psset{linestyle=solid}\pstree{\CC{\hyperlink{spotgac1<40>}{6}}\trput{2}}{
   \psset{linestyle=solid}\pstree{\CC{\hyperlink{spotgac1<41>}{7}}\trput{4}}{
   }
  }
 }
}
}
}
\only<42>{
\scalebox{ 0.40}{\psset{nodesep=0pt,levelsep=36pt}
\pstree{\CC{\hyperlink{spotgac1<1>}{4}}}{
 \psset{linestyle=solid}\FAILTREE{spotgac1<31>}{13 17}\trput{2}
 \psset{linestyle=solid}\FAILTREE{spotgac1<38>}{3 4}\trput{4}
 \psset{linestyle=solid}\pstree{\CC{\hyperlink{spotgac1<39>}{5}}\trput{5}}{
  \psset{linestyle=solid}\pstree{\CC{\hyperlink{spotgac1<40>}{6}}\trput{2}}{
   \psset{linestyle=solid}\pstree{\CC{\hyperlink{spotgac1<41>}{7}}\trput{4}}{
    \psset{linestyle=solid}\FAIL{spotgac1<42>}\trput{6}
   }
  }
 }
}
}
}
\only<43>{
\scalebox{ 0.40}{\psset{nodesep=0pt,levelsep=36pt}
\pstree{\CC{\hyperlink{spotgac1<1>}{4}}}{
 \psset{linestyle=solid}\FAILTREE{spotgac1<31>}{13 17}\trput{2}
 \psset{linestyle=solid}\FAILTREE{spotgac1<38>}{3 4}\trput{4}
 \psset{linestyle=solid}\pstree{\CC{\hyperlink{spotgac1<39>}{5}}\trput{5}}{
  \psset{linestyle=solid}\pstree{\CC{\hyperlink{spotgac1<40>}{6}}\trput{2}}{
   \psset{linestyle=solid}\pstree{\CC{\hyperlink{spotgac1<41>}{7}}\trput{4}}{
    \psset{linestyle=solid}\FAIL{spotgac1<42>}\trput{6}
    \psset{linestyle=solid}\pstree{\CC{\hyperlink{spotgac1<43>}{10}}\trput{9}}{
    }
   }
  }
 }
}
}
}
\only<44>{
\scalebox{ 0.40}{\psset{nodesep=0pt,levelsep=36pt}
\pstree{\CC{\hyperlink{spotgac1<1>}{4}}}{
 \psset{linestyle=solid}\FAILTREE{spotgac1<31>}{13 17}\trput{2}
 \psset{linestyle=solid}\FAILTREE{spotgac1<38>}{3 4}\trput{4}
 \psset{linestyle=solid}\pstree{\CC{\hyperlink{spotgac1<39>}{5}}\trput{5}}{
  \psset{linestyle=solid}\pstree{\CC{\hyperlink{spotgac1<40>}{6}}\trput{2}}{
   \psset{linestyle=solid}\pstree{\CC{\hyperlink{spotgac1<41>}{7}}\trput{4}}{
    \psset{linestyle=solid}\FAIL{spotgac1<42>}\trput{6}
    \psset{linestyle=solid}\pstree{\CC{\hyperlink{spotgac1<43>}{10}}\trput{9}}{
     \psset{linestyle=solid}\FAIL{spotgac1<44>}\trput{6}
    }
   }
  }
 }
}
}
}
\only<45>{
\scalebox{ 0.40}{\psset{nodesep=0pt,levelsep=36pt}
\pstree{\CC{\hyperlink{spotgac1<1>}{4}}}{
 \psset{linestyle=solid}\FAILTREE{spotgac1<31>}{13 17}\trput{2}
 \psset{linestyle=solid}\FAILTREE{spotgac1<38>}{3 4}\trput{4}
 \psset{linestyle=solid}\pstree{\CC{\hyperlink{spotgac1<39>}{5}}\trput{5}}{
  \psset{linestyle=solid}\pstree{\CC{\hyperlink{spotgac1<40>}{6}}\trput{2}}{
   \psset{linestyle=solid}\pstree{\CC{\hyperlink{spotgac1<41>}{7}}\trput{4}}{
    \psset{linestyle=solid}\FAIL{spotgac1<42>}\trput{6}
    \psset{linestyle=solid}\pstree{\CC{\hyperlink{spotgac1<43>}{10}}\trput{9}}{
     \psset{linestyle=solid}\FAIL{spotgac1<44>}\trput{6}
     \psset{linestyle=solid}\FAIL{spotgac1<45>}\trput{8}
    }
   }
  }
 }
}
}
}
\only<46>{
\scalebox{ 0.40}{\psset{nodesep=0pt,levelsep=36pt}
\pstree{\CC{\hyperlink{spotgac1<1>}{4}}}{
 \psset{linestyle=solid}\FAILTREE{spotgac1<31>}{13 17}\trput{2}
 \psset{linestyle=solid}\FAILTREE{spotgac1<38>}{3 4}\trput{4}
 \psset{linestyle=solid}\pstree{\CC{\hyperlink{spotgac1<39>}{5}}\trput{5}}{
  \psset{linestyle=solid}\pstree{\CC{\hyperlink{spotgac1<40>}{6}}\trput{2}}{
   \psset{linestyle=solid}\FAILTREE{spotgac1<45>}{2 3}\trput{4}
   \psset{linestyle=solid}\pstree{\CC{\hyperlink{spotgac1<46>}{7}}\trput{6}}{
   }
  }
 }
}
}
}
\only<47>{
\scalebox{ 0.40}{\psset{nodesep=0pt,levelsep=36pt}
\pstree{\CC{\hyperlink{spotgac1<1>}{4}}}{
 \psset{linestyle=solid}\FAILTREE{spotgac1<31>}{13 17}\trput{2}
 \psset{linestyle=solid}\FAILTREE{spotgac1<38>}{3 4}\trput{4}
 \psset{linestyle=solid}\pstree{\CC{\hyperlink{spotgac1<39>}{5}}\trput{5}}{
  \psset{linestyle=solid}\pstree{\CC{\hyperlink{spotgac1<40>}{6}}\trput{2}}{
   \psset{linestyle=solid}\FAILTREE{spotgac1<45>}{2 3}\trput{4}
   \psset{linestyle=solid}\pstree{\CC{\hyperlink{spotgac1<46>}{7}}\trput{6}}{
    \psset{linestyle=solid}\FAIL{spotgac1<47>}\trput{4}
   }
  }
 }
}
}
}
\only<48>{
\scalebox{ 0.40}{\psset{nodesep=0pt,levelsep=36pt}
\pstree{\CC{\hyperlink{spotgac1<1>}{4}}}{
 \psset{linestyle=solid}\FAILTREE{spotgac1<31>}{13 17}\trput{2}
 \psset{linestyle=solid}\FAILTREE{spotgac1<38>}{3 4}\trput{4}
 \psset{linestyle=solid}\pstree{\CC{\hyperlink{spotgac1<39>}{5}}\trput{5}}{
  \psset{linestyle=solid}\pstree{\CC{\hyperlink{spotgac1<40>}{6}}\trput{2}}{
   \psset{linestyle=solid}\FAILTREE{spotgac1<45>}{2 3}\trput{4}
   \psset{linestyle=solid}\pstree{\CC{\hyperlink{spotgac1<46>}{7}}\trput{6}}{
    \psset{linestyle=solid}\FAIL{spotgac1<47>}\trput{4}
    \psset{linestyle=solid}\FAIL{spotgac1<48>}\trput{9}
   }
  }
 }
}
}
}
\only<49>{
\scalebox{ 0.40}{\psset{nodesep=0pt,levelsep=36pt}
\pstree{\CC{\hyperlink{spotgac1<1>}{4}}}{
 \psset{linestyle=solid}\FAILTREE{spotgac1<31>}{13 17}\trput{2}
 \psset{linestyle=solid}\FAILTREE{spotgac1<38>}{3 4}\trput{4}
 \psset{linestyle=solid}\pstree{\CC{\hyperlink{spotgac1<39>}{5}}\trput{5}}{
  \psset{linestyle=solid}\FAILTREE{spotgac1<48>}{4 5}\trput{2}
  \psset{linestyle=solid}\FAIL{spotgac1<49>}\trput{4}
 }
}
}
}
\only<50>{
\scalebox{ 0.40}{\psset{nodesep=0pt,levelsep=36pt}
\pstree{\CC{\hyperlink{spotgac1<1>}{4}}}{
 \psset{linestyle=solid}\FAILTREE{spotgac1<31>}{13 17}\trput{2}
 \psset{linestyle=solid}\FAILTREE{spotgac1<38>}{3 4}\trput{4}
 \psset{linestyle=solid}\pstree{\CC{\hyperlink{spotgac1<39>}{5}}\trput{5}}{
  \psset{linestyle=solid}\FAILTREE{spotgac1<48>}{4 5}\trput{2}
  \psset{linestyle=solid}\FAIL{spotgac1<49>}\trput{4}
  \psset{linestyle=solid}\pstree{\CC{\hyperlink{spotgac1<50>}{6}}\trput{6}}{
  }
 }
}
}
}
\only<51>{
\scalebox{ 0.40}{\psset{nodesep=0pt,levelsep=36pt}
\pstree{\CC{\hyperlink{spotgac1<1>}{4}}}{
 \psset{linestyle=solid}\FAILTREE{spotgac1<31>}{13 17}\trput{2}
 \psset{linestyle=solid}\FAILTREE{spotgac1<38>}{3 4}\trput{4}
 \psset{linestyle=solid}\pstree{\CC{\hyperlink{spotgac1<39>}{5}}\trput{5}}{
  \psset{linestyle=solid}\FAILTREE{spotgac1<48>}{4 5}\trput{2}
  \psset{linestyle=solid}\FAIL{spotgac1<49>}\trput{4}
  \psset{linestyle=solid}\pstree{\CC{\hyperlink{spotgac1<50>}{6}}\trput{6}}{
   \psset{linestyle=solid}\pstree{\CC{\hyperlink{spotgac1<51>}{7}}\trput{2}}{
   }
  }
 }
}
}
}
\only<52>{
\scalebox{ 0.40}{\psset{nodesep=0pt,levelsep=36pt}
\pstree{\CC{\hyperlink{spotgac1<1>}{4}}}{
 \psset{linestyle=solid}\FAILTREE{spotgac1<31>}{13 17}\trput{2}
 \psset{linestyle=solid}\FAILTREE{spotgac1<38>}{3 4}\trput{4}
 \psset{linestyle=solid}\pstree{\CC{\hyperlink{spotgac1<39>}{5}}\trput{5}}{
  \psset{linestyle=solid}\FAILTREE{spotgac1<48>}{4 5}\trput{2}
  \psset{linestyle=solid}\FAIL{spotgac1<49>}\trput{4}
  \psset{linestyle=solid}\pstree{\CC{\hyperlink{spotgac1<50>}{6}}\trput{6}}{
   \psset{linestyle=solid}\pstree{\CC{\hyperlink{spotgac1<51>}{7}}\trput{2}}{
    \psset{linestyle=solid}\FAIL{spotgac1<52>}\trput{4}
   }
  }
 }
}
}
}
\only<53>{
\scalebox{ 0.40}{\psset{nodesep=0pt,levelsep=36pt}
\pstree{\CC{\hyperlink{spotgac1<1>}{4}}}{
 \psset{linestyle=solid}\FAILTREE{spotgac1<31>}{13 17}\trput{2}
 \psset{linestyle=solid}\FAILTREE{spotgac1<38>}{3 4}\trput{4}
 \psset{linestyle=solid}\pstree{\CC{\hyperlink{spotgac1<39>}{5}}\trput{5}}{
  \psset{linestyle=solid}\FAILTREE{spotgac1<48>}{4 5}\trput{2}
  \psset{linestyle=solid}\FAIL{spotgac1<49>}\trput{4}
  \psset{linestyle=solid}\pstree{\CC{\hyperlink{spotgac1<50>}{6}}\trput{6}}{
   \psset{linestyle=solid}\pstree{\CC{\hyperlink{spotgac1<51>}{7}}\trput{2}}{
    \psset{linestyle=solid}\FAIL{spotgac1<52>}\trput{4}
    \psset{linestyle=solid}\FAIL{spotgac1<53>}\trput{9}
   }
  }
 }
}
}
}
\only<54>{
\scalebox{ 0.40}{\psset{nodesep=0pt,levelsep=36pt}
\pstree{\CC{\hyperlink{spotgac1<1>}{4}}}{
 \psset{linestyle=solid}\FAILTREE{spotgac1<31>}{13 17}\trput{2}
 \psset{linestyle=solid}\FAILTREE{spotgac1<38>}{3 4}\trput{4}
 \psset{linestyle=solid}\pstree{\CC{\hyperlink{spotgac1<39>}{5}}\trput{5}}{
  \psset{linestyle=solid}\FAILTREE{spotgac1<48>}{4 5}\trput{2}
  \psset{linestyle=solid}\FAIL{spotgac1<49>}\trput{4}
  \psset{linestyle=solid}\pstree{\CC{\hyperlink{spotgac1<50>}{6}}\trput{6}}{
   \psset{linestyle=solid}\FAILTREE{spotgac1<53>}{1 2}\trput{2}
   \psset{linestyle=solid}\pstree{\CC{\hyperlink{spotgac1<54>}{7}}\trput{4}}{
   }
  }
 }
}
}
}
\only<55>{
\scalebox{ 0.40}{\psset{nodesep=0pt,levelsep=36pt}
\pstree{\CC{\hyperlink{spotgac1<1>}{4}}}{
 \psset{linestyle=solid}\FAILTREE{spotgac1<31>}{13 17}\trput{2}
 \psset{linestyle=solid}\FAILTREE{spotgac1<38>}{3 4}\trput{4}
 \psset{linestyle=solid}\pstree{\CC{\hyperlink{spotgac1<39>}{5}}\trput{5}}{
  \psset{linestyle=solid}\FAILTREE{spotgac1<48>}{4 5}\trput{2}
  \psset{linestyle=solid}\FAIL{spotgac1<49>}\trput{4}
  \psset{linestyle=solid}\pstree{\CC{\hyperlink{spotgac1<50>}{6}}\trput{6}}{
   \psset{linestyle=solid}\FAILTREE{spotgac1<53>}{1 2}\trput{2}
   \psset{linestyle=solid}\pstree{\CC{\hyperlink{spotgac1<54>}{7}}\trput{4}}{
    \psset{linestyle=solid}\pstree{\CC{\hyperlink{spotgac1<55>}{10}}\trput{2}}{
    }
   }
  }
 }
}
}
}
\only<56>{
\scalebox{ 0.40}{\psset{nodesep=0pt,levelsep=36pt}
\pstree{\CC{\hyperlink{spotgac1<1>}{4}}}{
 \psset{linestyle=solid}\FAILTREE{spotgac1<31>}{13 17}\trput{2}
 \psset{linestyle=solid}\FAILTREE{spotgac1<38>}{3 4}\trput{4}
 \psset{linestyle=solid}\pstree{\CC{\hyperlink{spotgac1<39>}{5}}\trput{5}}{
  \psset{linestyle=solid}\FAILTREE{spotgac1<48>}{4 5}\trput{2}
  \psset{linestyle=solid}\FAIL{spotgac1<49>}\trput{4}
  \psset{linestyle=solid}\pstree{\CC{\hyperlink{spotgac1<50>}{6}}\trput{6}}{
   \psset{linestyle=solid}\FAILTREE{spotgac1<53>}{1 2}\trput{2}
   \psset{linestyle=solid}\pstree{\CC{\hyperlink{spotgac1<54>}{7}}\trput{4}}{
    \psset{linestyle=solid}\pstree{\CC{\hyperlink{spotgac1<55>}{10}}\trput{2}}{
     \psset{linestyle=solid}\pstree{\CC{\hyperlink{spotgac1<56>}{12}}\trput{6}}{
     }
    }
   }
  }
 }
}
}
}
\only<57>{
\scalebox{ 0.40}{\psset{nodesep=0pt,levelsep=36pt}
\pstree{\CC{\hyperlink{spotgac1<1>}{4}}}{
 \psset{linestyle=solid}\FAILTREE{spotgac1<31>}{13 17}\trput{2}
 \psset{linestyle=solid}\FAILTREE{spotgac1<38>}{3 4}\trput{4}
 \psset{linestyle=solid}\pstree{\CC{\hyperlink{spotgac1<39>}{5}}\trput{5}}{
  \psset{linestyle=solid}\FAILTREE{spotgac1<48>}{4 5}\trput{2}
  \psset{linestyle=solid}\FAIL{spotgac1<49>}\trput{4}
  \psset{linestyle=solid}\pstree{\CC{\hyperlink{spotgac1<50>}{6}}\trput{6}}{
   \psset{linestyle=solid}\FAILTREE{spotgac1<53>}{1 2}\trput{2}
   \psset{linestyle=solid}\pstree{\CC{\hyperlink{spotgac1<54>}{7}}\trput{4}}{
    \psset{linestyle=solid}\pstree{\CC{\hyperlink{spotgac1<55>}{10}}\trput{2}}{
     \psset{linestyle=solid}\pstree{\CC{\hyperlink{spotgac1<56>}{12}}\trput{6}}{
      \psset{linestyle=solid}\FAIL{spotgac1<57>}\trput{2}
     }
    }
   }
  }
 }
}
}
}
\only<58>{
\scalebox{ 0.40}{\psset{nodesep=0pt,levelsep=36pt}
\pstree{\CC{\hyperlink{spotgac1<1>}{4}}}{
 \psset{linestyle=solid}\FAILTREE{spotgac1<31>}{13 17}\trput{2}
 \psset{linestyle=solid}\FAILTREE{spotgac1<38>}{3 4}\trput{4}
 \psset{linestyle=solid}\pstree{\CC{\hyperlink{spotgac1<39>}{5}}\trput{5}}{
  \psset{linestyle=solid}\FAILTREE{spotgac1<48>}{4 5}\trput{2}
  \psset{linestyle=solid}\FAIL{spotgac1<49>}\trput{4}
  \psset{linestyle=solid}\pstree{\CC{\hyperlink{spotgac1<50>}{6}}\trput{6}}{
   \psset{linestyle=solid}\FAILTREE{spotgac1<53>}{1 2}\trput{2}
   \psset{linestyle=solid}\pstree{\CC{\hyperlink{spotgac1<54>}{7}}\trput{4}}{
    \psset{linestyle=solid}\pstree{\CC{\hyperlink{spotgac1<55>}{10}}\trput{2}}{
     \psset{linestyle=solid}\pstree{\CC{\hyperlink{spotgac1<56>}{12}}\trput{6}}{
      \psset{linestyle=solid}\FAIL{spotgac1<57>}\trput{2}
      \psset{linestyle=solid}\FAIL{spotgac1<58>}\trput{8}
     }
    }
   }
  }
 }
}
}
}
\only<59>{
\scalebox{ 0.40}{\psset{nodesep=0pt,levelsep=36pt}
\pstree{\CC{\hyperlink{spotgac1<1>}{4}}}{
 \psset{linestyle=solid}\FAILTREE{spotgac1<31>}{13 17}\trput{2}
 \psset{linestyle=solid}\FAILTREE{spotgac1<38>}{3 4}\trput{4}
 \psset{linestyle=solid}\pstree{\CC{\hyperlink{spotgac1<39>}{5}}\trput{5}}{
  \psset{linestyle=solid}\FAILTREE{spotgac1<48>}{4 5}\trput{2}
  \psset{linestyle=solid}\FAIL{spotgac1<49>}\trput{4}
  \psset{linestyle=solid}\pstree{\CC{\hyperlink{spotgac1<50>}{6}}\trput{6}}{
   \psset{linestyle=solid}\FAILTREE{spotgac1<53>}{1 2}\trput{2}
   \psset{linestyle=solid}\pstree{\CC{\hyperlink{spotgac1<54>}{7}}\trput{4}}{
    \psset{linestyle=solid}\pstree{\CC{\hyperlink{spotgac1<55>}{10}}\trput{2}}{
     \psset{linestyle=solid}\FAILTREE{spotgac1<58>}{1 2}\trput{6}
     \psset{linestyle=solid}\FAIL{spotgac1<59>}\trput{8}
    }
   }
  }
 }
}
}
}
\only<60>{
\scalebox{ 0.40}{\psset{nodesep=0pt,levelsep=36pt}
\pstree{\CC{\hyperlink{spotgac1<1>}{4}}}{
 \psset{linestyle=solid}\FAILTREE{spotgac1<31>}{13 17}\trput{2}
 \psset{linestyle=solid}\FAILTREE{spotgac1<38>}{3 4}\trput{4}
 \psset{linestyle=solid}\pstree{\CC{\hyperlink{spotgac1<39>}{5}}\trput{5}}{
  \psset{linestyle=solid}\FAILTREE{spotgac1<48>}{4 5}\trput{2}
  \psset{linestyle=solid}\FAIL{spotgac1<49>}\trput{4}
  \psset{linestyle=solid}\pstree{\CC{\hyperlink{spotgac1<50>}{6}}\trput{6}}{
   \psset{linestyle=solid}\FAILTREE{spotgac1<53>}{1 2}\trput{2}
   \psset{linestyle=solid}\pstree{\CC{\hyperlink{spotgac1<54>}{7}}\trput{4}}{
    \psset{linestyle=solid}\FAILTREE{spotgac1<59>}{2 3}\trput{2}
    \psset{linestyle=solid}\pstree{\CC{\hyperlink{spotgac1<60>}{10}}\trput{9}}{
    }
   }
  }
 }
}
}
}
\only<61>{
\scalebox{ 0.40}{\psset{nodesep=0pt,levelsep=36pt}
\pstree{\CC{\hyperlink{spotgac1<1>}{4}}}{
 \psset{linestyle=solid}\FAILTREE{spotgac1<31>}{13 17}\trput{2}
 \psset{linestyle=solid}\FAILTREE{spotgac1<38>}{3 4}\trput{4}
 \psset{linestyle=solid}\pstree{\CC{\hyperlink{spotgac1<39>}{5}}\trput{5}}{
  \psset{linestyle=solid}\FAILTREE{spotgac1<48>}{4 5}\trput{2}
  \psset{linestyle=solid}\FAIL{spotgac1<49>}\trput{4}
  \psset{linestyle=solid}\pstree{\CC{\hyperlink{spotgac1<50>}{6}}\trput{6}}{
   \psset{linestyle=solid}\FAILTREE{spotgac1<53>}{1 2}\trput{2}
   \psset{linestyle=solid}\pstree{\CC{\hyperlink{spotgac1<54>}{7}}\trput{4}}{
    \psset{linestyle=solid}\FAILTREE{spotgac1<59>}{2 3}\trput{2}
    \psset{linestyle=solid}\pstree{\CC{\hyperlink{spotgac1<60>}{10}}\trput{9}}{
     \psset{linestyle=solid}\pstree{\CC{\hyperlink{spotgac1<61>}{12}}\trput{6}}{
     }
    }
   }
  }
 }
}
}
}
\only<62>{
\scalebox{ 0.40}{\psset{nodesep=0pt,levelsep=36pt}
\pstree{\CC{\hyperlink{spotgac1<1>}{4}}}{
 \psset{linestyle=solid}\FAILTREE{spotgac1<31>}{13 17}\trput{2}
 \psset{linestyle=solid}\FAILTREE{spotgac1<38>}{3 4}\trput{4}
 \psset{linestyle=solid}\pstree{\CC{\hyperlink{spotgac1<39>}{5}}\trput{5}}{
  \psset{linestyle=solid}\FAILTREE{spotgac1<48>}{4 5}\trput{2}
  \psset{linestyle=solid}\FAIL{spotgac1<49>}\trput{4}
  \psset{linestyle=solid}\pstree{\CC{\hyperlink{spotgac1<50>}{6}}\trput{6}}{
   \psset{linestyle=solid}\FAILTREE{spotgac1<53>}{1 2}\trput{2}
   \psset{linestyle=solid}\pstree{\CC{\hyperlink{spotgac1<54>}{7}}\trput{4}}{
    \psset{linestyle=solid}\FAILTREE{spotgac1<59>}{2 3}\trput{2}
    \psset{linestyle=solid}\pstree{\CC{\hyperlink{spotgac1<60>}{10}}\trput{9}}{
     \psset{linestyle=solid}\pstree{\CC{\hyperlink{spotgac1<61>}{12}}\trput{6}}{
      \psset{linestyle=solid}\pstree{\CC{\hyperlink{spotgac1<62>}{22}}\trput{2}}{
      }
     }
    }
   }
  }
 }
}
}
}
\only<63>{
\scalebox{ 0.40}{\psset{nodesep=0pt,levelsep=36pt}
\pstree{\CC{\hyperlink{spotgac1<1>}{4}}}{
 \psset{linestyle=solid}\FAILTREE{spotgac1<31>}{13 17}\trput{2}
 \psset{linestyle=solid}\FAILTREE{spotgac1<38>}{3 4}\trput{4}
 \psset{linestyle=solid}\pstree{\CC{\hyperlink{spotgac1<39>}{5}}\trput{5}}{
  \psset{linestyle=solid}\FAILTREE{spotgac1<48>}{4 5}\trput{2}
  \psset{linestyle=solid}\FAIL{spotgac1<49>}\trput{4}
  \psset{linestyle=solid}\pstree{\CC{\hyperlink{spotgac1<50>}{6}}\trput{6}}{
   \psset{linestyle=solid}\FAILTREE{spotgac1<53>}{1 2}\trput{2}
   \psset{linestyle=solid}\pstree{\CC{\hyperlink{spotgac1<54>}{7}}\trput{4}}{
    \psset{linestyle=solid}\FAILTREE{spotgac1<59>}{2 3}\trput{2}
    \psset{linestyle=solid}\pstree{\CC{\hyperlink{spotgac1<60>}{10}}\trput{9}}{
     \psset{linestyle=solid}\pstree{\CC{\hyperlink{spotgac1<61>}{12}}\trput{6}}{
      \psset{linestyle=solid}\pstree{\CC{\hyperlink{spotgac1<62>}{22}}\trput{2}}{
       \psset{linestyle=solid}\pstree{\CC{\hyperlink{spotgac1<63>}{24}}\trput{2}}{
       }
      }
     }
    }
   }
  }
 }
}
}
}
\only<64>{
\scalebox{ 0.40}{\psset{nodesep=0pt,levelsep=36pt}
\pstree{\CC{\hyperlink{spotgac1<1>}{4}}}{
 \psset{linestyle=solid}\FAILTREE{spotgac1<31>}{13 17}\trput{2}
 \psset{linestyle=solid}\FAILTREE{spotgac1<38>}{3 4}\trput{4}
 \psset{linestyle=solid}\pstree{\CC{\hyperlink{spotgac1<39>}{5}}\trput{5}}{
  \psset{linestyle=solid}\FAILTREE{spotgac1<48>}{4 5}\trput{2}
  \psset{linestyle=solid}\FAIL{spotgac1<49>}\trput{4}
  \psset{linestyle=solid}\pstree{\CC{\hyperlink{spotgac1<50>}{6}}\trput{6}}{
   \psset{linestyle=solid}\FAILTREE{spotgac1<53>}{1 2}\trput{2}
   \psset{linestyle=solid}\pstree{\CC{\hyperlink{spotgac1<54>}{7}}\trput{4}}{
    \psset{linestyle=solid}\FAILTREE{spotgac1<59>}{2 3}\trput{2}
    \psset{linestyle=solid}\pstree{\CC{\hyperlink{spotgac1<60>}{10}}\trput{9}}{
     \psset{linestyle=solid}\pstree{\CC{\hyperlink{spotgac1<61>}{12}}\trput{6}}{
      \psset{linestyle=solid}\pstree{\CC{\hyperlink{spotgac1<62>}{22}}\trput{2}}{
       \psset{linestyle=solid}\pstree{\CC{\hyperlink{spotgac1<63>}{24}}\trput{2}}{
        \psset{linestyle=solid}\pstree{\CC{\hyperlink{spotgac1<64>}{31}}\trput{1}}{
        }
       }
      }
     }
    }
   }
  }
 }
}
}
}
\only<65>{
\scalebox{ 0.40}{\psset{nodesep=0pt,levelsep=36pt}
\pstree{\CC{\hyperlink{spotgac1<1>}{4}}}{
 \psset{linestyle=solid}\FAILTREE{spotgac1<31>}{13 17}\trput{2}
 \psset{linestyle=solid}\FAILTREE{spotgac1<38>}{3 4}\trput{4}
 \psset{linestyle=solid}\pstree{\CC{\hyperlink{spotgac1<39>}{5}}\trput{5}}{
  \psset{linestyle=solid}\FAILTREE{spotgac1<48>}{4 5}\trput{2}
  \psset{linestyle=solid}\FAIL{spotgac1<49>}\trput{4}
  \psset{linestyle=solid}\pstree{\CC{\hyperlink{spotgac1<50>}{6}}\trput{6}}{
   \psset{linestyle=solid}\FAILTREE{spotgac1<53>}{1 2}\trput{2}
   \psset{linestyle=solid}\pstree{\CC{\hyperlink{spotgac1<54>}{7}}\trput{4}}{
    \psset{linestyle=solid}\FAILTREE{spotgac1<59>}{2 3}\trput{2}
    \psset{linestyle=solid}\pstree{\CC{\hyperlink{spotgac1<60>}{10}}\trput{9}}{
     \psset{linestyle=solid}\pstree{\CC{\hyperlink{spotgac1<61>}{12}}\trput{6}}{
      \psset{linestyle=solid}\pstree{\CC{\hyperlink{spotgac1<62>}{22}}\trput{2}}{
       \psset{linestyle=solid}\pstree{\CC{\hyperlink{spotgac1<63>}{24}}\trput{2}}{
        \psset{linestyle=solid}\pstree{\CC{\hyperlink{spotgac1<64>}{31}}\trput{1}}{
         \psset{linestyle=solid}\FAIL{spotgac1<65>}\trput{4}
        }
       }
      }
     }
    }
   }
  }
 }
}
}
}
\only<66>{
\scalebox{ 0.40}{\psset{nodesep=0pt,levelsep=36pt}
\pstree{\CC{\hyperlink{spotgac1<1>}{4}}}{
 \psset{linestyle=solid}\FAILTREE{spotgac1<31>}{13 17}\trput{2}
 \psset{linestyle=solid}\FAILTREE{spotgac1<38>}{3 4}\trput{4}
 \psset{linestyle=solid}\pstree{\CC{\hyperlink{spotgac1<39>}{5}}\trput{5}}{
  \psset{linestyle=solid}\FAILTREE{spotgac1<48>}{4 5}\trput{2}
  \psset{linestyle=solid}\FAIL{spotgac1<49>}\trput{4}
  \psset{linestyle=solid}\pstree{\CC{\hyperlink{spotgac1<50>}{6}}\trput{6}}{
   \psset{linestyle=solid}\FAILTREE{spotgac1<53>}{1 2}\trput{2}
   \psset{linestyle=solid}\pstree{\CC{\hyperlink{spotgac1<54>}{7}}\trput{4}}{
    \psset{linestyle=solid}\FAILTREE{spotgac1<59>}{2 3}\trput{2}
    \psset{linestyle=solid}\pstree{\CC{\hyperlink{spotgac1<60>}{10}}\trput{9}}{
     \psset{linestyle=solid}\pstree{\CC{\hyperlink{spotgac1<61>}{12}}\trput{6}}{
      \psset{linestyle=solid}\pstree{\CC{\hyperlink{spotgac1<62>}{22}}\trput{2}}{
       \psset{linestyle=solid}\pstree{\CC{\hyperlink{spotgac1<63>}{24}}\trput{2}}{
        \psset{linestyle=solid}\pstree{\CC{\hyperlink{spotgac1<64>}{31}}\trput{1}}{
         \psset{linestyle=solid}\FAIL{spotgac1<65>}\trput{4}
         \psset{linestyle=solid}\SUCC{spotgac1<66>}\trput{8}
        }
       }
      }
     }
    }
   }
  }
 }
}
}
}
\only<67>{
\scalebox{ 0.40}{\psset{nodesep=0pt,levelsep=36pt}
\pstree{\CC{\hyperlink{spotgac1<1>}{4}}}{
 \psset{linestyle=solid}\FAILTREE{spotgac1<31>}{13 17}\trput{2}
 \psset{linestyle=solid}\FAILTREE{spotgac1<38>}{3 4}\trput{4}
 \psset{linestyle=solid}\pstree{\CC{\hyperlink{spotgac1<39>}{5}}\trput{5}}{
  \psset{linestyle=solid}\FAILTREE{spotgac1<48>}{4 5}\trput{2}
  \psset{linestyle=solid}\FAIL{spotgac1<49>}\trput{4}
  \psset{linestyle=solid}\pstree{\CC{\hyperlink{spotgac1<50>}{6}}\trput{6}}{
   \psset{linestyle=solid}\FAILTREE{spotgac1<53>}{1 2}\trput{2}
   \psset{linestyle=solid}\pstree{\CC{\hyperlink{spotgac1<54>}{7}}\trput{4}}{
    \psset{linestyle=solid}\FAILTREE{spotgac1<59>}{2 3}\trput{2}
    \psset{linestyle=solid}\pstree{\CC{\hyperlink{spotgac1<60>}{10}}\trput{9}}{
     \psset{linestyle=solid}\pstree{\CC{\hyperlink{spotgac1<61>}{12}}\trput{6}}{
      \psset{linestyle=solid}\pstree{\CC{\hyperlink{spotgac1<62>}{22}}\trput{2}}{
       \psset{linestyle=solid}\pstree{\CC{\hyperlink{spotgac1<63>}{24}}\trput{2}}{
        \psset{linestyle=solid}\pstree{\CC{\hyperlink{spotgac1<64>}{31}}\trput{1}}{
         \psset{linestyle=solid}\FAIL{spotgac1<65>}\trput{4}
         \psset{linestyle=solid}\SUCC{spotgac1<66>}\trput{8}
        }
        \psset{linestyle=solid}\FAIL{spotgac1<67>}\trput{7}
       }
      }
     }
    }
   }
  }
 }
}
}
}
\only<68>{
\scalebox{ 0.40}{\psset{nodesep=0pt,levelsep=36pt}
\pstree{\CC{\hyperlink{spotgac1<1>}{4}}}{
 \psset{linestyle=solid}\FAILTREE{spotgac1<31>}{13 17}\trput{2}
 \psset{linestyle=solid}\FAILTREE{spotgac1<38>}{3 4}\trput{4}
 \psset{linestyle=solid}\pstree{\CC{\hyperlink{spotgac1<39>}{5}}\trput{5}}{
  \psset{linestyle=solid}\FAILTREE{spotgac1<48>}{4 5}\trput{2}
  \psset{linestyle=solid}\FAIL{spotgac1<49>}\trput{4}
  \psset{linestyle=solid}\pstree{\CC{\hyperlink{spotgac1<50>}{6}}\trput{6}}{
   \psset{linestyle=solid}\FAILTREE{spotgac1<53>}{1 2}\trput{2}
   \psset{linestyle=solid}\pstree{\CC{\hyperlink{spotgac1<54>}{7}}\trput{4}}{
    \psset{linestyle=solid}\FAILTREE{spotgac1<59>}{2 3}\trput{2}
    \psset{linestyle=solid}\pstree{\CC{\hyperlink{spotgac1<60>}{10}}\trput{9}}{
     \psset{linestyle=solid}\pstree{\CC{\hyperlink{spotgac1<61>}{12}}\trput{6}}{
      \psset{linestyle=solid}\pstree{\CC{\hyperlink{spotgac1<62>}{22}}\trput{2}}{
       \psset{linestyle=solid}\pstree{\CC{\hyperlink{spotgac1<63>}{24}}\trput{2}}{
        \psset{linestyle=solid}\pstree{\CC{\hyperlink{spotgac1<64>}{31}}\trput{1}}{
         \psset{linestyle=solid}\FAIL{spotgac1<65>}\trput{4}
         \psset{linestyle=solid}\SUCC{spotgac1<66>}\trput{8}
        }
        \psset{linestyle=solid}\FAIL{spotgac1<67>}\trput{7}
       }
       \psset{linestyle=solid}\FAIL{spotgac1<68>}\trput{4}
      }
     }
    }
   }
  }
 }
}
}
}
\only<69>{
\scalebox{ 0.40}{\psset{nodesep=0pt,levelsep=36pt}
\pstree{\CC{\hyperlink{spotgac1<1>}{4}}}{
 \psset{linestyle=solid}\FAILTREE{spotgac1<31>}{13 17}\trput{2}
 \psset{linestyle=solid}\FAILTREE{spotgac1<38>}{3 4}\trput{4}
 \psset{linestyle=solid}\pstree{\CC{\hyperlink{spotgac1<39>}{5}}\trput{5}}{
  \psset{linestyle=solid}\FAILTREE{spotgac1<48>}{4 5}\trput{2}
  \psset{linestyle=solid}\FAIL{spotgac1<49>}\trput{4}
  \psset{linestyle=solid}\pstree{\CC{\hyperlink{spotgac1<50>}{6}}\trput{6}}{
   \psset{linestyle=solid}\FAILTREE{spotgac1<53>}{1 2}\trput{2}
   \psset{linestyle=solid}\pstree{\CC{\hyperlink{spotgac1<54>}{7}}\trput{4}}{
    \psset{linestyle=solid}\FAILTREE{spotgac1<59>}{2 3}\trput{2}
    \psset{linestyle=solid}\pstree{\CC{\hyperlink{spotgac1<60>}{10}}\trput{9}}{
     \psset{linestyle=solid}\pstree{\CC{\hyperlink{spotgac1<61>}{12}}\trput{6}}{
      \psset{linestyle=solid}\pstree{\CC{\hyperlink{spotgac1<62>}{22}}\trput{2}}{
       \psset{linestyle=solid}\pstree{\CC{\hyperlink{spotgac1<63>}{24}}\trput{2}}{
        \psset{linestyle=solid}\pstree{\CC{\hyperlink{spotgac1<64>}{31}}\trput{1}}{
         \psset{linestyle=solid}\FAIL{spotgac1<65>}\trput{4}
         \psset{linestyle=solid}\SUCC{spotgac1<66>}\trput{8}
        }
        \psset{linestyle=solid}\FAIL{spotgac1<67>}\trput{7}
       }
       \psset{linestyle=solid}\FAIL{spotgac1<68>}\trput{4}
      }
      \psset{linestyle=solid}\pstree{\CC{\hyperlink{spotgac1<69>}{22}}\trput{8}}{
      }
     }
    }
   }
  }
 }
}
}
}
\only<70>{
\scalebox{ 0.40}{\psset{nodesep=0pt,levelsep=36pt}
\pstree{\CC{\hyperlink{spotgac1<1>}{4}}}{
 \psset{linestyle=solid}\FAILTREE{spotgac1<31>}{13 17}\trput{2}
 \psset{linestyle=solid}\FAILTREE{spotgac1<38>}{3 4}\trput{4}
 \psset{linestyle=solid}\pstree{\CC{\hyperlink{spotgac1<39>}{5}}\trput{5}}{
  \psset{linestyle=solid}\FAILTREE{spotgac1<48>}{4 5}\trput{2}
  \psset{linestyle=solid}\FAIL{spotgac1<49>}\trput{4}
  \psset{linestyle=solid}\pstree{\CC{\hyperlink{spotgac1<50>}{6}}\trput{6}}{
   \psset{linestyle=solid}\FAILTREE{spotgac1<53>}{1 2}\trput{2}
   \psset{linestyle=solid}\pstree{\CC{\hyperlink{spotgac1<54>}{7}}\trput{4}}{
    \psset{linestyle=solid}\FAILTREE{spotgac1<59>}{2 3}\trput{2}
    \psset{linestyle=solid}\pstree{\CC{\hyperlink{spotgac1<60>}{10}}\trput{9}}{
     \psset{linestyle=solid}\pstree{\CC{\hyperlink{spotgac1<61>}{12}}\trput{6}}{
      \psset{linestyle=solid}\pstree{\CC{\hyperlink{spotgac1<62>}{22}}\trput{2}}{
       \psset{linestyle=solid}\pstree{\CC{\hyperlink{spotgac1<63>}{24}}\trput{2}}{
        \psset{linestyle=solid}\pstree{\CC{\hyperlink{spotgac1<64>}{31}}\trput{1}}{
         \psset{linestyle=solid}\FAIL{spotgac1<65>}\trput{4}
         \psset{linestyle=solid}\SUCC{spotgac1<66>}\trput{8}
        }
        \psset{linestyle=solid}\FAIL{spotgac1<67>}\trput{7}
       }
       \psset{linestyle=solid}\FAIL{spotgac1<68>}\trput{4}
      }
      \psset{linestyle=solid}\pstree{\CC{\hyperlink{spotgac1<69>}{22}}\trput{8}}{
       \psset{linestyle=solid}\pstree{\CC{\hyperlink{spotgac1<70>}{24}}\trput{2}}{
       }
      }
     }
    }
   }
  }
 }
}
}
}
\only<71>{
\scalebox{ 0.40}{\psset{nodesep=0pt,levelsep=36pt}
\pstree{\CC{\hyperlink{spotgac1<1>}{4}}}{
 \psset{linestyle=solid}\FAILTREE{spotgac1<31>}{13 17}\trput{2}
 \psset{linestyle=solid}\FAILTREE{spotgac1<38>}{3 4}\trput{4}
 \psset{linestyle=solid}\pstree{\CC{\hyperlink{spotgac1<39>}{5}}\trput{5}}{
  \psset{linestyle=solid}\FAILTREE{spotgac1<48>}{4 5}\trput{2}
  \psset{linestyle=solid}\FAIL{spotgac1<49>}\trput{4}
  \psset{linestyle=solid}\pstree{\CC{\hyperlink{spotgac1<50>}{6}}\trput{6}}{
   \psset{linestyle=solid}\FAILTREE{spotgac1<53>}{1 2}\trput{2}
   \psset{linestyle=solid}\pstree{\CC{\hyperlink{spotgac1<54>}{7}}\trput{4}}{
    \psset{linestyle=solid}\FAILTREE{spotgac1<59>}{2 3}\trput{2}
    \psset{linestyle=solid}\pstree{\CC{\hyperlink{spotgac1<60>}{10}}\trput{9}}{
     \psset{linestyle=solid}\pstree{\CC{\hyperlink{spotgac1<61>}{12}}\trput{6}}{
      \psset{linestyle=solid}\pstree{\CC{\hyperlink{spotgac1<62>}{22}}\trput{2}}{
       \psset{linestyle=solid}\pstree{\CC{\hyperlink{spotgac1<63>}{24}}\trput{2}}{
        \psset{linestyle=solid}\pstree{\CC{\hyperlink{spotgac1<64>}{31}}\trput{1}}{
         \psset{linestyle=solid}\FAIL{spotgac1<65>}\trput{4}
         \psset{linestyle=solid}\SUCC{spotgac1<66>}\trput{8}
        }
        \psset{linestyle=solid}\FAIL{spotgac1<67>}\trput{7}
       }
       \psset{linestyle=solid}\FAIL{spotgac1<68>}\trput{4}
      }
      \psset{linestyle=solid}\pstree{\CC{\hyperlink{spotgac1<69>}{22}}\trput{8}}{
       \psset{linestyle=solid}\pstree{\CC{\hyperlink{spotgac1<70>}{24}}\trput{2}}{
        \psset{linestyle=solid}\FAIL{spotgac1<71>}\trput{1}
       }
      }
     }
    }
   }
  }
 }
}
}
}
\only<72>{
\scalebox{ 0.40}{\psset{nodesep=0pt,levelsep=36pt}
\pstree{\CC{\hyperlink{spotgac1<1>}{4}}}{
 \psset{linestyle=solid}\FAILTREE{spotgac1<31>}{13 17}\trput{2}
 \psset{linestyle=solid}\FAILTREE{spotgac1<38>}{3 4}\trput{4}
 \psset{linestyle=solid}\pstree{\CC{\hyperlink{spotgac1<39>}{5}}\trput{5}}{
  \psset{linestyle=solid}\FAILTREE{spotgac1<48>}{4 5}\trput{2}
  \psset{linestyle=solid}\FAIL{spotgac1<49>}\trput{4}
  \psset{linestyle=solid}\pstree{\CC{\hyperlink{spotgac1<50>}{6}}\trput{6}}{
   \psset{linestyle=solid}\FAILTREE{spotgac1<53>}{1 2}\trput{2}
   \psset{linestyle=solid}\pstree{\CC{\hyperlink{spotgac1<54>}{7}}\trput{4}}{
    \psset{linestyle=solid}\FAILTREE{spotgac1<59>}{2 3}\trput{2}
    \psset{linestyle=solid}\pstree{\CC{\hyperlink{spotgac1<60>}{10}}\trput{9}}{
     \psset{linestyle=solid}\pstree{\CC{\hyperlink{spotgac1<61>}{12}}\trput{6}}{
      \psset{linestyle=solid}\pstree{\CC{\hyperlink{spotgac1<62>}{22}}\trput{2}}{
       \psset{linestyle=solid}\pstree{\CC{\hyperlink{spotgac1<63>}{24}}\trput{2}}{
        \psset{linestyle=solid}\pstree{\CC{\hyperlink{spotgac1<64>}{31}}\trput{1}}{
         \psset{linestyle=solid}\FAIL{spotgac1<65>}\trput{4}
         \psset{linestyle=solid}\SUCC{spotgac1<66>}\trput{8}
        }
        \psset{linestyle=solid}\FAIL{spotgac1<67>}\trput{7}
       }
       \psset{linestyle=solid}\FAIL{spotgac1<68>}\trput{4}
      }
      \psset{linestyle=solid}\pstree{\CC{\hyperlink{spotgac1<69>}{22}}\trput{8}}{
       \psset{linestyle=solid}\pstree{\CC{\hyperlink{spotgac1<70>}{24}}\trput{2}}{
        \psset{linestyle=solid}\FAIL{spotgac1<71>}\trput{1}
        \psset{linestyle=solid}\FAIL{spotgac1<72>}\trput{7}
       }
      }
     }
    }
   }
  }
 }
}
}
}
\only<73>{
\scalebox{ 0.40}{\psset{nodesep=0pt,levelsep=36pt}
\pstree{\CC{\hyperlink{spotgac1<1>}{4}}}{
 \psset{linestyle=solid}\FAILTREE{spotgac1<31>}{13 17}\trput{2}
 \psset{linestyle=solid}\FAILTREE{spotgac1<38>}{3 4}\trput{4}
 \psset{linestyle=solid}\pstree{\CC{\hyperlink{spotgac1<39>}{5}}\trput{5}}{
  \psset{linestyle=solid}\FAILTREE{spotgac1<48>}{4 5}\trput{2}
  \psset{linestyle=solid}\FAIL{spotgac1<49>}\trput{4}
  \psset{linestyle=solid}\pstree{\CC{\hyperlink{spotgac1<50>}{6}}\trput{6}}{
   \psset{linestyle=solid}\FAILTREE{spotgac1<53>}{1 2}\trput{2}
   \psset{linestyle=solid}\pstree{\CC{\hyperlink{spotgac1<54>}{7}}\trput{4}}{
    \psset{linestyle=solid}\FAILTREE{spotgac1<59>}{2 3}\trput{2}
    \psset{linestyle=solid}\pstree{\CC{\hyperlink{spotgac1<60>}{10}}\trput{9}}{
     \psset{linestyle=solid}\pstree{\CC{\hyperlink{spotgac1<61>}{12}}\trput{6}}{
      \psset{linestyle=solid}\pstree{\CC{\hyperlink{spotgac1<62>}{22}}\trput{2}}{
       \psset{linestyle=solid}\pstree{\CC{\hyperlink{spotgac1<63>}{24}}\trput{2}}{
        \psset{linestyle=solid}\pstree{\CC{\hyperlink{spotgac1<64>}{31}}\trput{1}}{
         \psset{linestyle=solid}\FAIL{spotgac1<65>}\trput{4}
         \psset{linestyle=solid}\SUCC{spotgac1<66>}\trput{8}
        }
        \psset{linestyle=solid}\FAIL{spotgac1<67>}\trput{7}
       }
       \psset{linestyle=solid}\FAIL{spotgac1<68>}\trput{4}
      }
      \psset{linestyle=solid}\pstree{\CC{\hyperlink{spotgac1<69>}{22}}\trput{8}}{
       \psset{linestyle=solid}\FAILTREE{spotgac1<72>}{1 2}\trput{2}
       \psset{linestyle=solid}\FAIL{spotgac1<73>}\trput{4}
      }
     }
    }
   }
  }
 }
}
}
}
\only<74>{
\scalebox{ 0.40}{\psset{nodesep=0pt,levelsep=36pt}
\pstree{\CC{\hyperlink{spotgac1<1>}{4}}}{
 \psset{linestyle=solid}\FAILTREE{spotgac1<31>}{13 17}\trput{2}
 \psset{linestyle=solid}\FAILTREE{spotgac1<38>}{3 4}\trput{4}
 \psset{linestyle=solid}\pstree{\CC{\hyperlink{spotgac1<39>}{5}}\trput{5}}{
  \psset{linestyle=solid}\FAILTREE{spotgac1<48>}{4 5}\trput{2}
  \psset{linestyle=solid}\FAIL{spotgac1<49>}\trput{4}
  \psset{linestyle=solid}\pstree{\CC{\hyperlink{spotgac1<50>}{6}}\trput{6}}{
   \psset{linestyle=solid}\FAILTREE{spotgac1<53>}{1 2}\trput{2}
   \psset{linestyle=solid}\pstree{\CC{\hyperlink{spotgac1<54>}{7}}\trput{4}}{
    \psset{linestyle=solid}\FAILTREE{spotgac1<59>}{2 3}\trput{2}
    \psset{linestyle=solid}\pstree{\CC{\hyperlink{spotgac1<60>}{10}}\trput{9}}{
     \psset{linestyle=solid}\pstree{\CC{\hyperlink{spotgac1<61>}{12}}\trput{6}}{
      \psset{linestyle=solid}\pstree{\CC{\hyperlink{spotgac1<62>}{22}}\trput{2}}{
       \psset{linestyle=solid}\pstree{\CC{\hyperlink{spotgac1<63>}{24}}\trput{2}}{
        \psset{linestyle=solid}\pstree{\CC{\hyperlink{spotgac1<64>}{31}}\trput{1}}{
         \psset{linestyle=solid}\FAIL{spotgac1<65>}\trput{4}
         \psset{linestyle=solid}\SUCC{spotgac1<66>}\trput{8}
        }
        \psset{linestyle=solid}\FAIL{spotgac1<67>}\trput{7}
       }
       \psset{linestyle=solid}\FAIL{spotgac1<68>}\trput{4}
      }
      \psset{linestyle=solid}\FAILTREE{spotgac1<73>}{2 3}\trput{8}
     }
     \psset{linestyle=solid}\FAIL{spotgac1<74>}\trput{8}
    }
   }
  }
 }
}
}
}
\only<75>{
\scalebox{ 0.40}{\psset{nodesep=0pt,levelsep=36pt}
\pstree{\CC{\hyperlink{spotgac1<1>}{4}}}{
 \psset{linestyle=solid}\FAILTREE{spotgac1<31>}{13 17}\trput{2}
 \psset{linestyle=solid}\FAILTREE{spotgac1<38>}{3 4}\trput{4}
 \psset{linestyle=solid}\pstree{\CC{\hyperlink{spotgac1<39>}{5}}\trput{5}}{
  \psset{linestyle=solid}\FAILTREE{spotgac1<48>}{4 5}\trput{2}
  \psset{linestyle=solid}\FAIL{spotgac1<49>}\trput{4}
  \psset{linestyle=solid}\pstree{\CC{\hyperlink{spotgac1<50>}{6}}\trput{6}}{
   \psset{linestyle=solid}\FAILTREE{spotgac1<53>}{1 2}\trput{2}
   \psset{linestyle=solid}\pstree{\CC{\hyperlink{spotgac1<54>}{7}}\trput{4}}{
    \psset{linestyle=solid}\FAILTREE{spotgac1<59>}{2 3}\trput{2}
    \psset{linestyle=solid}\pstree{\CC{\hyperlink{spotgac1<60>}{10}}\trput{9}}{
     \psset{linestyle=solid}\pstree{\CC{\hyperlink{spotgac1<61>}{12}}\trput{6}}{
      \psset{linestyle=solid}\pstree{\CC{\hyperlink{spotgac1<62>}{22}}\trput{2}}{
       \psset{linestyle=solid}\pstree{\CC{\hyperlink{spotgac1<63>}{24}}\trput{2}}{
        \psset{linestyle=solid}\pstree{\CC{\hyperlink{spotgac1<64>}{31}}\trput{1}}{
         \psset{linestyle=solid}\FAIL{spotgac1<65>}\trput{4}
         \psset{linestyle=solid}\SUCC{spotgac1<66>}\trput{8}
        }
        \psset{linestyle=solid}\FAIL{spotgac1<67>}\trput{7}
       }
       \psset{linestyle=solid}\FAIL{spotgac1<68>}\trput{4}
      }
      \psset{linestyle=solid}\FAILTREE{spotgac1<73>}{2 3}\trput{8}
     }
     \psset{linestyle=solid}\FAIL{spotgac1<74>}\trput{8}
    }
   }
  }
 }
 \psset{linestyle=solid}\pstree{\CC{\hyperlink{spotgac1<75>}{5}}\trput{6}}{
 }
}
}
}
\only<76>{
\scalebox{ 0.40}{\psset{nodesep=0pt,levelsep=36pt}
\pstree{\CC{\hyperlink{spotgac1<1>}{4}}}{
 \psset{linestyle=solid}\FAILTREE{spotgac1<31>}{13 17}\trput{2}
 \psset{linestyle=solid}\FAILTREE{spotgac1<38>}{3 4}\trput{4}
 \psset{linestyle=solid}\pstree{\CC{\hyperlink{spotgac1<39>}{5}}\trput{5}}{
  \psset{linestyle=solid}\FAILTREE{spotgac1<48>}{4 5}\trput{2}
  \psset{linestyle=solid}\FAIL{spotgac1<49>}\trput{4}
  \psset{linestyle=solid}\pstree{\CC{\hyperlink{spotgac1<50>}{6}}\trput{6}}{
   \psset{linestyle=solid}\FAILTREE{spotgac1<53>}{1 2}\trput{2}
   \psset{linestyle=solid}\pstree{\CC{\hyperlink{spotgac1<54>}{7}}\trput{4}}{
    \psset{linestyle=solid}\FAILTREE{spotgac1<59>}{2 3}\trput{2}
    \psset{linestyle=solid}\pstree{\CC{\hyperlink{spotgac1<60>}{10}}\trput{9}}{
     \psset{linestyle=solid}\pstree{\CC{\hyperlink{spotgac1<61>}{12}}\trput{6}}{
      \psset{linestyle=solid}\pstree{\CC{\hyperlink{spotgac1<62>}{22}}\trput{2}}{
       \psset{linestyle=solid}\pstree{\CC{\hyperlink{spotgac1<63>}{24}}\trput{2}}{
        \psset{linestyle=solid}\pstree{\CC{\hyperlink{spotgac1<64>}{31}}\trput{1}}{
         \psset{linestyle=solid}\FAIL{spotgac1<65>}\trput{4}
         \psset{linestyle=solid}\SUCC{spotgac1<66>}\trput{8}
        }
        \psset{linestyle=solid}\FAIL{spotgac1<67>}\trput{7}
       }
       \psset{linestyle=solid}\FAIL{spotgac1<68>}\trput{4}
      }
      \psset{linestyle=solid}\FAILTREE{spotgac1<73>}{2 3}\trput{8}
     }
     \psset{linestyle=solid}\FAIL{spotgac1<74>}\trput{8}
    }
   }
  }
 }
 \psset{linestyle=solid}\pstree{\CC{\hyperlink{spotgac1<75>}{5}}\trput{6}}{
  \psset{linestyle=solid}\pstree{\CC{\hyperlink{spotgac1<76>}{6}}\trput{2}}{
  }
 }
}
}
}
\only<77>{
\scalebox{ 0.40}{\psset{nodesep=0pt,levelsep=36pt}
\pstree{\CC{\hyperlink{spotgac1<1>}{4}}}{
 \psset{linestyle=solid}\FAILTREE{spotgac1<31>}{13 17}\trput{2}
 \psset{linestyle=solid}\FAILTREE{spotgac1<38>}{3 4}\trput{4}
 \psset{linestyle=solid}\pstree{\CC{\hyperlink{spotgac1<39>}{5}}\trput{5}}{
  \psset{linestyle=solid}\FAILTREE{spotgac1<48>}{4 5}\trput{2}
  \psset{linestyle=solid}\FAIL{spotgac1<49>}\trput{4}
  \psset{linestyle=solid}\pstree{\CC{\hyperlink{spotgac1<50>}{6}}\trput{6}}{
   \psset{linestyle=solid}\FAILTREE{spotgac1<53>}{1 2}\trput{2}
   \psset{linestyle=solid}\pstree{\CC{\hyperlink{spotgac1<54>}{7}}\trput{4}}{
    \psset{linestyle=solid}\FAILTREE{spotgac1<59>}{2 3}\trput{2}
    \psset{linestyle=solid}\pstree{\CC{\hyperlink{spotgac1<60>}{10}}\trput{9}}{
     \psset{linestyle=solid}\pstree{\CC{\hyperlink{spotgac1<61>}{12}}\trput{6}}{
      \psset{linestyle=solid}\pstree{\CC{\hyperlink{spotgac1<62>}{22}}\trput{2}}{
       \psset{linestyle=solid}\pstree{\CC{\hyperlink{spotgac1<63>}{24}}\trput{2}}{
        \psset{linestyle=solid}\pstree{\CC{\hyperlink{spotgac1<64>}{31}}\trput{1}}{
         \psset{linestyle=solid}\FAIL{spotgac1<65>}\trput{4}
         \psset{linestyle=solid}\SUCC{spotgac1<66>}\trput{8}
        }
        \psset{linestyle=solid}\FAIL{spotgac1<67>}\trput{7}
       }
       \psset{linestyle=solid}\FAIL{spotgac1<68>}\trput{4}
      }
      \psset{linestyle=solid}\FAILTREE{spotgac1<73>}{2 3}\trput{8}
     }
     \psset{linestyle=solid}\FAIL{spotgac1<74>}\trput{8}
    }
   }
  }
 }
 \psset{linestyle=solid}\pstree{\CC{\hyperlink{spotgac1<75>}{5}}\trput{6}}{
  \psset{linestyle=solid}\pstree{\CC{\hyperlink{spotgac1<76>}{6}}\trput{2}}{
   \psset{linestyle=solid}\FAIL{spotgac1<77>}\trput{4}
  }
 }
}
}
}
\only<78>{
\scalebox{ 0.40}{\psset{nodesep=0pt,levelsep=36pt}
\pstree{\CC{\hyperlink{spotgac1<1>}{4}}}{
 \psset{linestyle=solid}\FAILTREE{spotgac1<31>}{13 17}\trput{2}
 \psset{linestyle=solid}\FAILTREE{spotgac1<38>}{3 4}\trput{4}
 \psset{linestyle=solid}\pstree{\CC{\hyperlink{spotgac1<39>}{5}}\trput{5}}{
  \psset{linestyle=solid}\FAILTREE{spotgac1<48>}{4 5}\trput{2}
  \psset{linestyle=solid}\FAIL{spotgac1<49>}\trput{4}
  \psset{linestyle=solid}\pstree{\CC{\hyperlink{spotgac1<50>}{6}}\trput{6}}{
   \psset{linestyle=solid}\FAILTREE{spotgac1<53>}{1 2}\trput{2}
   \psset{linestyle=solid}\pstree{\CC{\hyperlink{spotgac1<54>}{7}}\trput{4}}{
    \psset{linestyle=solid}\FAILTREE{spotgac1<59>}{2 3}\trput{2}
    \psset{linestyle=solid}\pstree{\CC{\hyperlink{spotgac1<60>}{10}}\trput{9}}{
     \psset{linestyle=solid}\pstree{\CC{\hyperlink{spotgac1<61>}{12}}\trput{6}}{
      \psset{linestyle=solid}\pstree{\CC{\hyperlink{spotgac1<62>}{22}}\trput{2}}{
       \psset{linestyle=solid}\pstree{\CC{\hyperlink{spotgac1<63>}{24}}\trput{2}}{
        \psset{linestyle=solid}\pstree{\CC{\hyperlink{spotgac1<64>}{31}}\trput{1}}{
         \psset{linestyle=solid}\FAIL{spotgac1<65>}\trput{4}
         \psset{linestyle=solid}\SUCC{spotgac1<66>}\trput{8}
        }
        \psset{linestyle=solid}\FAIL{spotgac1<67>}\trput{7}
       }
       \psset{linestyle=solid}\FAIL{spotgac1<68>}\trput{4}
      }
      \psset{linestyle=solid}\FAILTREE{spotgac1<73>}{2 3}\trput{8}
     }
     \psset{linestyle=solid}\FAIL{spotgac1<74>}\trput{8}
    }
   }
  }
 }
 \psset{linestyle=solid}\pstree{\CC{\hyperlink{spotgac1<75>}{5}}\trput{6}}{
  \psset{linestyle=solid}\pstree{\CC{\hyperlink{spotgac1<76>}{6}}\trput{2}}{
   \psset{linestyle=solid}\FAIL{spotgac1<77>}\trput{4}
   \psset{linestyle=solid}\pstree{\CC{\hyperlink{spotgac1<78>}{7}}\trput{5}}{
   }
  }
 }
}
}
}
\only<79>{
\scalebox{ 0.40}{\psset{nodesep=0pt,levelsep=36pt}
\pstree{\CC{\hyperlink{spotgac1<1>}{4}}}{
 \psset{linestyle=solid}\FAILTREE{spotgac1<31>}{13 17}\trput{2}
 \psset{linestyle=solid}\FAILTREE{spotgac1<38>}{3 4}\trput{4}
 \psset{linestyle=solid}\pstree{\CC{\hyperlink{spotgac1<39>}{5}}\trput{5}}{
  \psset{linestyle=solid}\FAILTREE{spotgac1<48>}{4 5}\trput{2}
  \psset{linestyle=solid}\FAIL{spotgac1<49>}\trput{4}
  \psset{linestyle=solid}\pstree{\CC{\hyperlink{spotgac1<50>}{6}}\trput{6}}{
   \psset{linestyle=solid}\FAILTREE{spotgac1<53>}{1 2}\trput{2}
   \psset{linestyle=solid}\pstree{\CC{\hyperlink{spotgac1<54>}{7}}\trput{4}}{
    \psset{linestyle=solid}\FAILTREE{spotgac1<59>}{2 3}\trput{2}
    \psset{linestyle=solid}\pstree{\CC{\hyperlink{spotgac1<60>}{10}}\trput{9}}{
     \psset{linestyle=solid}\pstree{\CC{\hyperlink{spotgac1<61>}{12}}\trput{6}}{
      \psset{linestyle=solid}\pstree{\CC{\hyperlink{spotgac1<62>}{22}}\trput{2}}{
       \psset{linestyle=solid}\pstree{\CC{\hyperlink{spotgac1<63>}{24}}\trput{2}}{
        \psset{linestyle=solid}\pstree{\CC{\hyperlink{spotgac1<64>}{31}}\trput{1}}{
         \psset{linestyle=solid}\FAIL{spotgac1<65>}\trput{4}
         \psset{linestyle=solid}\SUCC{spotgac1<66>}\trput{8}
        }
        \psset{linestyle=solid}\FAIL{spotgac1<67>}\trput{7}
       }
       \psset{linestyle=solid}\FAIL{spotgac1<68>}\trput{4}
      }
      \psset{linestyle=solid}\FAILTREE{spotgac1<73>}{2 3}\trput{8}
     }
     \psset{linestyle=solid}\FAIL{spotgac1<74>}\trput{8}
    }
   }
  }
 }
 \psset{linestyle=solid}\pstree{\CC{\hyperlink{spotgac1<75>}{5}}\trput{6}}{
  \psset{linestyle=solid}\pstree{\CC{\hyperlink{spotgac1<76>}{6}}\trput{2}}{
   \psset{linestyle=solid}\FAIL{spotgac1<77>}\trput{4}
   \psset{linestyle=solid}\pstree{\CC{\hyperlink{spotgac1<78>}{7}}\trput{5}}{
    \psset{linestyle=solid}\FAIL{spotgac1<79>}\trput{4}
   }
  }
 }
}
}
}
\only<80>{
\scalebox{ 0.40}{\psset{nodesep=0pt,levelsep=36pt}
\pstree{\CC{\hyperlink{spotgac1<1>}{4}}}{
 \psset{linestyle=solid}\FAILTREE{spotgac1<31>}{13 17}\trput{2}
 \psset{linestyle=solid}\FAILTREE{spotgac1<38>}{3 4}\trput{4}
 \psset{linestyle=solid}\pstree{\CC{\hyperlink{spotgac1<39>}{5}}\trput{5}}{
  \psset{linestyle=solid}\FAILTREE{spotgac1<48>}{4 5}\trput{2}
  \psset{linestyle=solid}\FAIL{spotgac1<49>}\trput{4}
  \psset{linestyle=solid}\pstree{\CC{\hyperlink{spotgac1<50>}{6}}\trput{6}}{
   \psset{linestyle=solid}\FAILTREE{spotgac1<53>}{1 2}\trput{2}
   \psset{linestyle=solid}\pstree{\CC{\hyperlink{spotgac1<54>}{7}}\trput{4}}{
    \psset{linestyle=solid}\FAILTREE{spotgac1<59>}{2 3}\trput{2}
    \psset{linestyle=solid}\pstree{\CC{\hyperlink{spotgac1<60>}{10}}\trput{9}}{
     \psset{linestyle=solid}\pstree{\CC{\hyperlink{spotgac1<61>}{12}}\trput{6}}{
      \psset{linestyle=solid}\pstree{\CC{\hyperlink{spotgac1<62>}{22}}\trput{2}}{
       \psset{linestyle=solid}\pstree{\CC{\hyperlink{spotgac1<63>}{24}}\trput{2}}{
        \psset{linestyle=solid}\pstree{\CC{\hyperlink{spotgac1<64>}{31}}\trput{1}}{
         \psset{linestyle=solid}\FAIL{spotgac1<65>}\trput{4}
         \psset{linestyle=solid}\SUCC{spotgac1<66>}\trput{8}
        }
        \psset{linestyle=solid}\FAIL{spotgac1<67>}\trput{7}
       }
       \psset{linestyle=solid}\FAIL{spotgac1<68>}\trput{4}
      }
      \psset{linestyle=solid}\FAILTREE{spotgac1<73>}{2 3}\trput{8}
     }
     \psset{linestyle=solid}\FAIL{spotgac1<74>}\trput{8}
    }
   }
  }
 }
 \psset{linestyle=solid}\pstree{\CC{\hyperlink{spotgac1<75>}{5}}\trput{6}}{
  \psset{linestyle=solid}\pstree{\CC{\hyperlink{spotgac1<76>}{6}}\trput{2}}{
   \psset{linestyle=solid}\FAIL{spotgac1<77>}\trput{4}
   \psset{linestyle=solid}\pstree{\CC{\hyperlink{spotgac1<78>}{7}}\trput{5}}{
    \psset{linestyle=solid}\FAIL{spotgac1<79>}\trput{4}
    \psset{linestyle=solid}\FAIL{spotgac1<80>}\trput{9}
   }
  }
 }
}
}
}
\only<81>{
\scalebox{ 0.40}{\psset{nodesep=0pt,levelsep=36pt}
\pstree{\CC{\hyperlink{spotgac1<1>}{4}}}{
 \psset{linestyle=solid}\FAILTREE{spotgac1<31>}{13 17}\trput{2}
 \psset{linestyle=solid}\FAILTREE{spotgac1<38>}{3 4}\trput{4}
 \psset{linestyle=solid}\pstree{\CC{\hyperlink{spotgac1<39>}{5}}\trput{5}}{
  \psset{linestyle=solid}\FAILTREE{spotgac1<48>}{4 5}\trput{2}
  \psset{linestyle=solid}\FAIL{spotgac1<49>}\trput{4}
  \psset{linestyle=solid}\pstree{\CC{\hyperlink{spotgac1<50>}{6}}\trput{6}}{
   \psset{linestyle=solid}\FAILTREE{spotgac1<53>}{1 2}\trput{2}
   \psset{linestyle=solid}\pstree{\CC{\hyperlink{spotgac1<54>}{7}}\trput{4}}{
    \psset{linestyle=solid}\FAILTREE{spotgac1<59>}{2 3}\trput{2}
    \psset{linestyle=solid}\pstree{\CC{\hyperlink{spotgac1<60>}{10}}\trput{9}}{
     \psset{linestyle=solid}\pstree{\CC{\hyperlink{spotgac1<61>}{12}}\trput{6}}{
      \psset{linestyle=solid}\pstree{\CC{\hyperlink{spotgac1<62>}{22}}\trput{2}}{
       \psset{linestyle=solid}\pstree{\CC{\hyperlink{spotgac1<63>}{24}}\trput{2}}{
        \psset{linestyle=solid}\pstree{\CC{\hyperlink{spotgac1<64>}{31}}\trput{1}}{
         \psset{linestyle=solid}\FAIL{spotgac1<65>}\trput{4}
         \psset{linestyle=solid}\SUCC{spotgac1<66>}\trput{8}
        }
        \psset{linestyle=solid}\FAIL{spotgac1<67>}\trput{7}
       }
       \psset{linestyle=solid}\FAIL{spotgac1<68>}\trput{4}
      }
      \psset{linestyle=solid}\FAILTREE{spotgac1<73>}{2 3}\trput{8}
     }
     \psset{linestyle=solid}\FAIL{spotgac1<74>}\trput{8}
    }
   }
  }
 }
 \psset{linestyle=solid}\pstree{\CC{\hyperlink{spotgac1<75>}{5}}\trput{6}}{
  \psset{linestyle=solid}\FAILTREE{spotgac1<80>}{2 3}\trput{2}
  \psset{linestyle=solid}\pstree{\CC{\hyperlink{spotgac1<81>}{6}}\trput{4}}{
  }
 }
}
}
}
\only<82>{
\scalebox{ 0.40}{\psset{nodesep=0pt,levelsep=36pt}
\pstree{\CC{\hyperlink{spotgac1<1>}{4}}}{
 \psset{linestyle=solid}\FAILTREE{spotgac1<31>}{13 17}\trput{2}
 \psset{linestyle=solid}\FAILTREE{spotgac1<38>}{3 4}\trput{4}
 \psset{linestyle=solid}\pstree{\CC{\hyperlink{spotgac1<39>}{5}}\trput{5}}{
  \psset{linestyle=solid}\FAILTREE{spotgac1<48>}{4 5}\trput{2}
  \psset{linestyle=solid}\FAIL{spotgac1<49>}\trput{4}
  \psset{linestyle=solid}\pstree{\CC{\hyperlink{spotgac1<50>}{6}}\trput{6}}{
   \psset{linestyle=solid}\FAILTREE{spotgac1<53>}{1 2}\trput{2}
   \psset{linestyle=solid}\pstree{\CC{\hyperlink{spotgac1<54>}{7}}\trput{4}}{
    \psset{linestyle=solid}\FAILTREE{spotgac1<59>}{2 3}\trput{2}
    \psset{linestyle=solid}\pstree{\CC{\hyperlink{spotgac1<60>}{10}}\trput{9}}{
     \psset{linestyle=solid}\pstree{\CC{\hyperlink{spotgac1<61>}{12}}\trput{6}}{
      \psset{linestyle=solid}\pstree{\CC{\hyperlink{spotgac1<62>}{22}}\trput{2}}{
       \psset{linestyle=solid}\pstree{\CC{\hyperlink{spotgac1<63>}{24}}\trput{2}}{
        \psset{linestyle=solid}\pstree{\CC{\hyperlink{spotgac1<64>}{31}}\trput{1}}{
         \psset{linestyle=solid}\FAIL{spotgac1<65>}\trput{4}
         \psset{linestyle=solid}\SUCC{spotgac1<66>}\trput{8}
        }
        \psset{linestyle=solid}\FAIL{spotgac1<67>}\trput{7}
       }
       \psset{linestyle=solid}\FAIL{spotgac1<68>}\trput{4}
      }
      \psset{linestyle=solid}\FAILTREE{spotgac1<73>}{2 3}\trput{8}
     }
     \psset{linestyle=solid}\FAIL{spotgac1<74>}\trput{8}
    }
   }
  }
 }
 \psset{linestyle=solid}\pstree{\CC{\hyperlink{spotgac1<75>}{5}}\trput{6}}{
  \psset{linestyle=solid}\FAILTREE{spotgac1<80>}{2 3}\trput{2}
  \psset{linestyle=solid}\pstree{\CC{\hyperlink{spotgac1<81>}{6}}\trput{4}}{
   \psset{linestyle=solid}\FAIL{spotgac1<82>}\trput{2}
  }
 }
}
}
}
\only<83>{
\scalebox{ 0.40}{\psset{nodesep=0pt,levelsep=36pt}
\pstree{\CC{\hyperlink{spotgac1<1>}{4}}}{
 \psset{linestyle=solid}\FAILTREE{spotgac1<31>}{13 17}\trput{2}
 \psset{linestyle=solid}\FAILTREE{spotgac1<38>}{3 4}\trput{4}
 \psset{linestyle=solid}\pstree{\CC{\hyperlink{spotgac1<39>}{5}}\trput{5}}{
  \psset{linestyle=solid}\FAILTREE{spotgac1<48>}{4 5}\trput{2}
  \psset{linestyle=solid}\FAIL{spotgac1<49>}\trput{4}
  \psset{linestyle=solid}\pstree{\CC{\hyperlink{spotgac1<50>}{6}}\trput{6}}{
   \psset{linestyle=solid}\FAILTREE{spotgac1<53>}{1 2}\trput{2}
   \psset{linestyle=solid}\pstree{\CC{\hyperlink{spotgac1<54>}{7}}\trput{4}}{
    \psset{linestyle=solid}\FAILTREE{spotgac1<59>}{2 3}\trput{2}
    \psset{linestyle=solid}\pstree{\CC{\hyperlink{spotgac1<60>}{10}}\trput{9}}{
     \psset{linestyle=solid}\pstree{\CC{\hyperlink{spotgac1<61>}{12}}\trput{6}}{
      \psset{linestyle=solid}\pstree{\CC{\hyperlink{spotgac1<62>}{22}}\trput{2}}{
       \psset{linestyle=solid}\pstree{\CC{\hyperlink{spotgac1<63>}{24}}\trput{2}}{
        \psset{linestyle=solid}\pstree{\CC{\hyperlink{spotgac1<64>}{31}}\trput{1}}{
         \psset{linestyle=solid}\FAIL{spotgac1<65>}\trput{4}
         \psset{linestyle=solid}\SUCC{spotgac1<66>}\trput{8}
        }
        \psset{linestyle=solid}\FAIL{spotgac1<67>}\trput{7}
       }
       \psset{linestyle=solid}\FAIL{spotgac1<68>}\trput{4}
      }
      \psset{linestyle=solid}\FAILTREE{spotgac1<73>}{2 3}\trput{8}
     }
     \psset{linestyle=solid}\FAIL{spotgac1<74>}\trput{8}
    }
   }
  }
 }
 \psset{linestyle=solid}\pstree{\CC{\hyperlink{spotgac1<75>}{5}}\trput{6}}{
  \psset{linestyle=solid}\FAILTREE{spotgac1<80>}{2 3}\trput{2}
  \psset{linestyle=solid}\pstree{\CC{\hyperlink{spotgac1<81>}{6}}\trput{4}}{
   \psset{linestyle=solid}\FAIL{spotgac1<82>}\trput{2}
   \psset{linestyle=solid}\pstree{\CC{\hyperlink{spotgac1<83>}{7}}\trput{5}}{
   }
  }
 }
}
}
}
\only<84>{
\scalebox{ 0.40}{\psset{nodesep=0pt,levelsep=36pt}
\pstree{\CC{\hyperlink{spotgac1<1>}{4}}}{
 \psset{linestyle=solid}\FAILTREE{spotgac1<31>}{13 17}\trput{2}
 \psset{linestyle=solid}\FAILTREE{spotgac1<38>}{3 4}\trput{4}
 \psset{linestyle=solid}\pstree{\CC{\hyperlink{spotgac1<39>}{5}}\trput{5}}{
  \psset{linestyle=solid}\FAILTREE{spotgac1<48>}{4 5}\trput{2}
  \psset{linestyle=solid}\FAIL{spotgac1<49>}\trput{4}
  \psset{linestyle=solid}\pstree{\CC{\hyperlink{spotgac1<50>}{6}}\trput{6}}{
   \psset{linestyle=solid}\FAILTREE{spotgac1<53>}{1 2}\trput{2}
   \psset{linestyle=solid}\pstree{\CC{\hyperlink{spotgac1<54>}{7}}\trput{4}}{
    \psset{linestyle=solid}\FAILTREE{spotgac1<59>}{2 3}\trput{2}
    \psset{linestyle=solid}\pstree{\CC{\hyperlink{spotgac1<60>}{10}}\trput{9}}{
     \psset{linestyle=solid}\pstree{\CC{\hyperlink{spotgac1<61>}{12}}\trput{6}}{
      \psset{linestyle=solid}\pstree{\CC{\hyperlink{spotgac1<62>}{22}}\trput{2}}{
       \psset{linestyle=solid}\pstree{\CC{\hyperlink{spotgac1<63>}{24}}\trput{2}}{
        \psset{linestyle=solid}\pstree{\CC{\hyperlink{spotgac1<64>}{31}}\trput{1}}{
         \psset{linestyle=solid}\FAIL{spotgac1<65>}\trput{4}
         \psset{linestyle=solid}\SUCC{spotgac1<66>}\trput{8}
        }
        \psset{linestyle=solid}\FAIL{spotgac1<67>}\trput{7}
       }
       \psset{linestyle=solid}\FAIL{spotgac1<68>}\trput{4}
      }
      \psset{linestyle=solid}\FAILTREE{spotgac1<73>}{2 3}\trput{8}
     }
     \psset{linestyle=solid}\FAIL{spotgac1<74>}\trput{8}
    }
   }
  }
 }
 \psset{linestyle=solid}\pstree{\CC{\hyperlink{spotgac1<75>}{5}}\trput{6}}{
  \psset{linestyle=solid}\FAILTREE{spotgac1<80>}{2 3}\trput{2}
  \psset{linestyle=solid}\pstree{\CC{\hyperlink{spotgac1<81>}{6}}\trput{4}}{
   \psset{linestyle=solid}\FAIL{spotgac1<82>}\trput{2}
   \psset{linestyle=solid}\pstree{\CC{\hyperlink{spotgac1<83>}{7}}\trput{5}}{
    \psset{linestyle=solid}\FAIL{spotgac1<84>}\trput{2}
   }
  }
 }
}
}
}
\only<85>{
\scalebox{ 0.40}{\psset{nodesep=0pt,levelsep=36pt}
\pstree{\CC{\hyperlink{spotgac1<1>}{4}}}{
 \psset{linestyle=solid}\FAILTREE{spotgac1<31>}{13 17}\trput{2}
 \psset{linestyle=solid}\FAILTREE{spotgac1<38>}{3 4}\trput{4}
 \psset{linestyle=solid}\pstree{\CC{\hyperlink{spotgac1<39>}{5}}\trput{5}}{
  \psset{linestyle=solid}\FAILTREE{spotgac1<48>}{4 5}\trput{2}
  \psset{linestyle=solid}\FAIL{spotgac1<49>}\trput{4}
  \psset{linestyle=solid}\pstree{\CC{\hyperlink{spotgac1<50>}{6}}\trput{6}}{
   \psset{linestyle=solid}\FAILTREE{spotgac1<53>}{1 2}\trput{2}
   \psset{linestyle=solid}\pstree{\CC{\hyperlink{spotgac1<54>}{7}}\trput{4}}{
    \psset{linestyle=solid}\FAILTREE{spotgac1<59>}{2 3}\trput{2}
    \psset{linestyle=solid}\pstree{\CC{\hyperlink{spotgac1<60>}{10}}\trput{9}}{
     \psset{linestyle=solid}\pstree{\CC{\hyperlink{spotgac1<61>}{12}}\trput{6}}{
      \psset{linestyle=solid}\pstree{\CC{\hyperlink{spotgac1<62>}{22}}\trput{2}}{
       \psset{linestyle=solid}\pstree{\CC{\hyperlink{spotgac1<63>}{24}}\trput{2}}{
        \psset{linestyle=solid}\pstree{\CC{\hyperlink{spotgac1<64>}{31}}\trput{1}}{
         \psset{linestyle=solid}\FAIL{spotgac1<65>}\trput{4}
         \psset{linestyle=solid}\SUCC{spotgac1<66>}\trput{8}
        }
        \psset{linestyle=solid}\FAIL{spotgac1<67>}\trput{7}
       }
       \psset{linestyle=solid}\FAIL{spotgac1<68>}\trput{4}
      }
      \psset{linestyle=solid}\FAILTREE{spotgac1<73>}{2 3}\trput{8}
     }
     \psset{linestyle=solid}\FAIL{spotgac1<74>}\trput{8}
    }
   }
  }
 }
 \psset{linestyle=solid}\pstree{\CC{\hyperlink{spotgac1<75>}{5}}\trput{6}}{
  \psset{linestyle=solid}\FAILTREE{spotgac1<80>}{2 3}\trput{2}
  \psset{linestyle=solid}\pstree{\CC{\hyperlink{spotgac1<81>}{6}}\trput{4}}{
   \psset{linestyle=solid}\FAIL{spotgac1<82>}\trput{2}
   \psset{linestyle=solid}\pstree{\CC{\hyperlink{spotgac1<83>}{7}}\trput{5}}{
    \psset{linestyle=solid}\FAIL{spotgac1<84>}\trput{2}
    \psset{linestyle=solid}\FAIL{spotgac1<85>}\trput{9}
   }
  }
 }
}
}
}
\vfill
\hyperlinkframestart<2-85>{\beamerreturnbutton{Back to Start}}
\hyperlinkframeend<1-84>{\beamerskipbutton{Skip Animation}}
\end{frame}
 %search tree

\begin{frame}
\frametitle{Solution}

\end{frame}

-------------------------------------------------------------------------
